% concatenated from qgzzgq1.tex
\documentclass[10pt,reqno]{amsart}
\usepackage[T1]{fontenc}
\usepackage[utf8]{inputenc}
\usepackage{lmodern}
\usepackage[a4paper,top=0.82in,bottom=0.84in,left=0.92in,right=0.92in]{geometry}
\usepackage{amsmath,amssymb,amsfonts,amsthm,mathtools,mathrsfs,amscd}
\usepackage{enumitem,microtype,booktabs,array,xcolor,hyperref}
\usepackage[nameinlink,capitalize,noabbrev]{cleveref}
\allowdisplaybreaks[2]
\setcounter{secnumdepth}{3}
\emergencystretch=2.5em
\raggedbottom

\hypersetup{
 bookmarksdepth=2,
 colorlinks=true,
 linkcolor=blue!45!black,
 citecolor=blue!45!black,
 urlcolor=blue!45!black,
 pdftitle={Infinite-Dimensional Levy Area and Spectral Ancestry},
 pdfauthor={Guangqian Zhao},
 pdfsubject={Primitive Rees geometry, critical weak-Schatten response, and infinite-dimensional Levy area},
 pdfkeywords={infinite-dimensional Levy area, causal Rees response kernel, critical spectral ancestry, weak Schatten ideals, Hall commutator, bidirectional recovery}
}

\theoremstyle{plain}
\newtheorem{theorem}{Theorem}[section]
\newtheorem{proposition}[theorem]{Proposition}
\newtheorem{lemma}[theorem]{Lemma}
\newtheorem{corollary}[theorem]{Corollary}
\theoremstyle{definition}
\newtheorem{definition}[theorem]{Definition}
\newtheorem{assumption}[theorem]{Assumption}
\newtheorem{example}[theorem]{Example}
\theoremstyle{remark}
\newtheorem{remark}[theorem]{Remark}
\begin{document}

\title[L\'evy area and spectral ancestry]{Infinite-Dimensional L\'evy Area and Spectral Ancestry}
\author{Guangqian Zhao}
\address{School of Mathematical Sciences, University of Science and Technology of China, Hefei, Anhui 230026, China}
\email{zhaoguangqian@mail.ustc.edu.cn}
\date{}
\subjclass[2020]{16W50, 17B01, 46L52, 47B10, 47B35, 60G15, 60G44, 60G51, 60L20}
\keywords{causal characters; primitive Rees geometry; response kernel;
infinite-dimensional L\'evy area; weak Schatten ideals; Hall commutator; spectral ancestry; stochastic convolution}

\begin{abstract}
We relate primitive ancestry in free-Lie expansions of infinite-dimensional
L\'evy area to critical weak-Schatten spectra.  For
$L=\mathfrak L(V)$, $D=[L,L]$, and $J=\ker(T(V)\to S(V))$, we prove
$L\cap J^r=\gamma_r(D)$ and identify the exact normal layers with
$\mathfrak L_r(D/[D,D])$.  Continuous Hilbert-martingale area has the sharp
split endpoint
$H\times\mathcal S^0_{1,\infty}(H)_{\rm skew}\times\mathcal S_1(H)_{\rm sa}$.
The critical grade $r$ has singular profile
$(\log(e+N))^{r-1}/N$.  A Rees response kernel characterizes strictness,
equality of ancestry and spectral depth, and recovery of the finite labelled
ancestry flag.  Pure-area Hall experiments give full-support kernels and exact
critical laws.  Revealed brackets and evolution families provide conditional
covariance and second-order L\'evy/SPDE interfaces.
\end{abstract}

\maketitle

\section{Introduction}
\label{sec:introduction}

The paper studies the chain
\begin{equation}
 \text{causal source}
 \longrightarrow
 \text{critical second-order geometry}
 \longrightarrow
 \text{primitive Rees flag}
 \longrightarrow
 \text{native response kernel}.
 \label{eq:main-architecture}
\end{equation}
The three arrows are, respectively, probabilistic, algebraic, and analytic.
Primitive depth by itself does not determine a compact spectrum; the relevant
invariant is the associated-graded response carried by the declared readers.

The main results are as follows.

\paragraph{Sharp stochastic entrance.}
For a continuous square-integrable martingale with bounded terminal trace
energy, the antisymmetric area has a dimension-free weak-trace bound and every
fixed sample belongs to the little weak ideal, whereas the collision
coordinate is trace class.  The exponent one is optimal in three different
ways: no $\mathcal S_{p,\infty}$ with $p<1$ is universal, no finite critical
Lorentz ideal is universal, and no prescribed dimension-free little modulus
is available.  The resulting Polish split state is
\begin{equation}
 \mathscr G_{\rm crit}(H)
 =H\times\mathcal S^0_{1,\infty}(H)_{\rm skew}
   \times\mathcal S_1(H)_{\rm sa}.
 \label{eq:intro-critical-state}
\end{equation}
It is the exact degree-two quotient of the typed It\^o quasi-shuffle state,
and its anchor provides canonical increments on bounded stopping intervals.

\paragraph{Primitive Rees geometry.}
Let
\[
 L=\mathfrak L(V),\qquad D=[L,L],\qquad
 J=\ker(T(V)\to S(V)).
\]
The primitive normal-cone theorem identifies
\begin{equation}
 F^rL:=L\cap J^r=\gamma_r(D),\qquad
 F^rL/F^{r+1}L\cong\mathfrak L_r(D/[D,D]).
 \label{eq:intro-normal-cone}
\end{equation}
Normal order and derived ancestry are therefore the same canonical valuation.
At fixed tensor degree the filtration is finite and lives on the finite
multiplicity spaces in the Schur decomposition; no splitting of an exact
layer is intrinsic.

\paragraph{One theorem for forward and inverse response.}
For a finite decreasing ancestry filtration $F^\bullet M$, the correct
increasing filtered object is the dual flag
\[
 \mathscr D_rM^*:=(F^{r+1}M)^\perp.
\]
A filtered response $\mathcal R:M^*\to\mathcal Y$ has the canonical kernel
\begin{equation}
 \mathfrak K_r(\mathcal R)
 \in (F^rM/F^{r+1}M)\otimes\partial_r\mathcal Y.
 \label{eq:intro-kernel}
\end{equation}
Full left support is equivalent to strictness and valuation exchange.  Hence
for every nonzero labelled probe $q$,
\begin{equation}
 d_{\rm sp}(q)=a_F(q),
 \qquad
 F^rM=\bigcap_{d_{\rm sp}(q)\le r-1}\ker q.
 \label{eq:intro-bidirectional}
\end{equation}
The first equality is the forward statement; the second reconstructs the
finite labelled ancestry flag.  Both are contractions of the same response kernel.

\paragraph{Native Hall support and the finite-order classification.}
The analytic filtration is generated by the grade-one harmonic atom
$D_\infty=\operatorname{diag}(1,1/2,\ldots)$, whose tensor powers satisfy
\begin{equation}
 s_N(D_\infty^{\otimes r})
 \sim \frac1{\Gamma(r)}\frac{(\log N)^{r-1}}{N}.
 \label{eq:intro-harmonic-profile}
\end{equation}
Every compact Hilbert--Schmidt operator is an off-diagonal corner of a skew
pure-area direction.  Recursive Hall commutators place a prescribed primitive
word in its exact ordinary degree, and signed tensor corners expose the
corresponding tensor product without a lower-cost remainder.  These readers
have full support on every finite Rees layer.  In particular,
\begin{equation}
 \mathbf S_{(3,1)}:\ \frac1N,
 \qquad
 \mathbf S_{(2,1,1)}:\ \frac{\log(e+N)}N,
 \label{eq:intro-degree-four}
\end{equation}
the long hook has depth one in every degree, and exact finite-order depths are
unbounded.  The abstract lower theorem is therefore deterministic; Brownian
blocks certify that the grade-one boundary also occurs naturally in a
stochastic source.

Two additional interfaces are retained because they clarify the meaning and
reach of the main theorem.  If the predictable bracket is revealed at the
initial sigma-field, the exterior lift has a conditional operator-valued
Dol\'eans propagator.  Its particle-two block measures the loss caused by
atomic bracket delivery, and its determinant profile transfers bracket
ideals and Mellin growth to area envelopes.  A representation-theoretic
obstruction shows that higher exterior degree is not itself higher free-Lie
ancestry.  For a deterministic evolution family,
terminal-frame stochastic convolutions satisfy a transported Chen law, and
operational-energy convergence yields Galerkin stability of the critical
lift.  For a specific SPDE, this entrance is combined with the equation's regularity and fixed-point estimates.

Brownian motion supplies the decisive type check.  The raw Chen-local area,
the deterministic second-chaos coefficient root, and a nonlinear
covariance-shell Gramian root live on different spaces and retain different
multiplicities.  Radonification relates the first two only in $L^2$, while the
third is a global auxiliary sample observable.  In the polynomial covariance
regime, the raw area has zero asymptotic spectral density relative to the
annealed root.  Together with the point-evaluation obstruction, a registration
trilemma rules out a bounded spectrum-preserving identification of all three
branches.

\paragraph{Context.}
Infinite-dimensional L\'evy area and rough-path liftability depend on the
chosen tensor topology
\cite{LedouxLyonsQian02,Dereich10,GarridoLuSchmalfuss15,GrongNilssenSchmeding22,Kalinichenko25}.
Semimartingale and jump lifts are developed in
\cite{Chevyrev18,FrizShekhar17,ChevyrevFriz19,ChevyrevFerrucci26}, and
stochastic-convolution regularity in
\cite{DaPratoKwapienZabczyk87,Brzezniak97,DaPratoZabczyk14,GubinelliTindel10}.
The present problem is complementary: the area coordinate selects a critical
operator ideal, primitive normal order propagates its logarithmic cost, and
labelled native responses determine which exact layers remain visible.  The
algebraic input uses free-Lie and commutator-filtration theory, the typed It\^o
interface uses quasi-shuffle/Hoffman calculus, and the analytic side uses
regular variation and tensor-product spectral asymptotics
\cite{Reutenauer93,Kapranov98,Hoffman00,CurryEtAl14,BGT87,KarolNazarovNikitin08,Rastegaev18}.
The finite-dimensional Brownian-area spectrum of
\cite{DelaCruzOberhauser26} supplies a natural-source certificate, while the
conditional covariance results are related to semimartingale signature and
cumulant methods \cite{FrizHagerTapia24}.

\paragraph{Organization.}
Section~\ref{cc:sec-causal-completions} develops causal completion;
Section~\ref{cc:sec-algebraic-ancestry} proves the primitive Rees theorem;
Sections~\ref{cr:sec-analytic}--\ref{rk:sec-kernel} establish the analytic
filtration and response-kernel theorem; Section~\ref{cc:sec-critical-geometry}
proves the sharp stochastic entrance; Section~\ref{hall:sec-native} gives
native Hall support and the finite-order classification.  The final sections
contain conditional covariance, L\'evy/SPDE interfaces, Brownian type
separation, and the infinite-tower consequences.
\section{Causal interval geometry and the two completions}
\label{cc:sec-causal-completions}

The primitive object is a transported multiplicative section over directed
time intervals.  Its endpoint anchor determines the two-parameter field.
Optional completion enlarges the domain to stopping intervals, whereas
singular extension enlarges the target group; their interchange is governed
by convergence of lifted anchors.

\subsection{Causal characters and endpoint anchors}

Fix a time horizon \(T<\infty\).  A \emph{transported Polish coefficient
system} is a Polish group \((G,\star,{\mathbf 1})\) together with continuous group
endomorphisms
\begin{equation}
 \alpha_{t,s}:G\longrightarrow G,
 \qquad
 \alpha_{t,t}={\mathrm{Id}},
 \qquad
 \alpha_{t,u}\alpha_{u,s}=\alpha_{t,s},
 \qquad 0\le s\le u\le t\le T,
 \label{cc:eq-transport-system}
\end{equation}
whose joint evaluation map \((t,s,g)\mapsto\alpha_{t,s}(g)\) is continuous.
We write \(G_t:=G\) and regard \(\alpha_{t,s}:G_s\to G_t\).  This is a fixed
trivialization; no measurable-field assertion for varying fibres is implicit.

\begin{definition}[Causal evolution character]
\label{cc:def-causal-character}
A causal evolution character is a normalized family

\[
 \Phi_{s,t}\in G_t\;(=G),
 \qquad \Phi_{t,t}={\mathbf 1},
\]

satisfying the transported Chen identity
\begin{equation}
 \Phi_{s,t}
 =\alpha_{t,u}(\Phi_{s,u})\star\Phi_{u,t},
 \qquad s\le u\le t.
 \label{cc:eq-transported-Chen}
\end{equation}
On a filtered probability space we additionally require
\(\Phi_{s,t}\) to be \({\mathcal F}_t\)-measurable for deterministic \(s\le t\),
and the normalization and Chen identities are understood almost surely for
each fixed deterministic pair and triple.  No simultaneous exceptional-set
claim is made until the anchor is perfected below.
The triple \(\mathfrak C=(G,\alpha,\Phi)\), or equivalently the anchored
triple introduced below, is the underlying causal object.
\end{definition}

Geometric ancestry requires additional registration data.  Choose
topological alphabet spaces \(V_t\) and continuous linear transports
\begin{equation}
 b_{t,s}:V_s\longrightarrow V_t,
 \qquad b_{t,t}={\mathrm{Id}},
 \qquad b_{t,u}b_{u,s}=b_{t,s},
 \label{cc:eq-alphabet-transport}
\end{equation}
and let
\[
 \beta_{t,s}:=(b_{t,s})_\#:
 \mathsf G_{\mathrm{geo}}^{\le N}(V_s)
 \longrightarrow
 \mathsf G_{\mathrm{geo}}^{\le N}(V_t)
\]
be the induced homomorphism of truncated geometric character groups and of
their free Lie algebras.  We assume that finite products, logarithms, and
the induced maps are continuous.  A registered chart is a family of
continuous homomorphisms
\begin{equation}
 \chi_t:G_t\longrightarrow\mathsf G_{\mathrm{geo}}^{\le N}(V_t),
 \qquad
 \chi_t\alpha_{t,s}=\beta_{t,s}\chi_s.
 \label{cc:eq-geometric-chart}
\end{equation}
whose joint evaluation is Borel.  Whenever charted anchors must be càdlàg,
joint continuity, or direct preservation of the selected paths, is assumed.
The charted logarithm
\begin{equation}
 X_{s,t}:=\log_\star\chi_t(\Phi_{s,t})
 \label{cc:eq-charted-logarithm}
\end{equation}
lies in the truncated free Lie algebra and satisfies
\[
 X_{s,t}
 ={\operatorname{BCH}}\!\left(\beta_{t,u}X_{s,u},X_{u,t}\right).
\]
Naturality of commutativization gives
\[
 (b_{t,s})_\#(J_s^r)\subseteq J_t^r,
 \qquad
 J_t:=\ker\!\left(T(V_t)\longrightarrow S(V_t)\right),
\]
and total tensor degree is preserved.  Exact valuation is preserved when
the induced map on the relevant associated-graded layer is injective; in
general transport may annihilate an initial symbol and raise valuation.

A chart may instead be defined on an admissible transported subgroup of a
lifted target:
\[
 \widetilde\chi_t:\widetilde G_t^{\mathrm{adm}}
 \longrightarrow\mathsf G_{\mathrm{geo}}^{\le N}(V_t),
 \qquad
 \widetilde\chi_t\widetilde\alpha_{t,s}
 =\beta_{t,s}\widetilde\chi_s.
\]
For a selected lift \(L\), put
\(X^L_{s,t}:=\log_\star\widetilde\chi_t(L_{s,t})\).  It descends to the
base when \(\widetilde\chi_t\) is trivial on the admissible kernel and is
otherwise gauge dependent.  Random charts are interpreted samplewise.
An \(L^0\)-valued chart supplies simultaneous representatives for a
countable registered family; an uncountable family requires a jointly
measurable representative as an additional hypothesis.  Hence the chart,
normal filtration, and spectral readers are realization data, not
components of \(\mathfrak C\).

The interval category has the initial object \(0\).  Consequently its
multiplicative sections have no independent two-parameter freedom.

\begin{theorem}[Anchor representation]
\label{cc:thm-anchor-representation}
Causal evolution characters are in natural bijection with normalized
endpoint sections \(A_t\in G_t\), \(A_0={\mathbf 1}\), equivalently paths
\(A:[0,T]\to G\) in the fixed trivialization.  The correspondence is
\begin{equation}
 A_t=\Phi_{0,t},
 \qquad
 \Phi_{s,t}
 =\bigl(\alpha_{t,s}(A_s)\bigr)^{-1}\star A_t.
 \label{cc:eq-anchor-reconstruction}
\end{equation}
In the stochastic category the correspondence is between deterministic-time
characters modulo almost-sure equality and adapted anchors modulo
modification; the identities are initially asserted for each fixed time
tuple.  It becomes a simultaneous pathwise correspondence after the
optional perfection below.
It commutes with every continuous group homomorphism intertwining the
transports.  If \(G\) is a finite-step exponential group and the transports
intertwine the exponential maps, write
\[
 \alpha^{\mathrm{Lie}}_{t,s}
 :=\log_\star\circ\alpha_{t,s}\circ\exp_\star
\]
for the induced Lie-algebra homomorphism.  Then, with
\(a_t=\log_\star A_t\) and \(\Psi_{s,t}=\log_\star\Phi_{s,t}\),
\begin{equation}
 \Psi_{s,t}
 ={\operatorname{BCH}}\!\left(-\alpha^{\mathrm{Lie}}_{t,s}(a_s),a_t\right)
 ={\operatorname{BCH}}\!\left(
   \alpha^{\mathrm{Lie}}_{t,u}(\Psi_{s,u}),\Psi_{u,t}\right).
 \label{cc:eq-transported-BCH}
\end{equation}
Thus causal composition and the logarithmic coordinate are two descriptions
of the same object.  The chart \(\chi\), when present, turns the latter into
the geometric Lie coordinate \(X\); no \(J\)-adic structure is asserted
without that extra datum.
\end{theorem}

\begin{proof}
Applying \eqref{cc:eq-transported-Chen} to \(0\le s\le t\) gives

\[
 A_t=\alpha_{t,s}(A_s)\star\Phi_{s,t},
\]

and therefore forces \eqref{cc:eq-anchor-reconstruction}.  Conversely,
for \(s\le u\le t\),
\begin{align*}
 &\alpha_{t,u}\!\left(
   \bigl(\alpha_{u,s}(A_s)\bigr)^{-1}\star A_u\right)
 \star\bigl(\alpha_{t,u}(A_u)\bigr)^{-1}\star A_t\\
 &\qquad=\bigl(\alpha_{t,s}(A_s)\bigr)^{-1}\star A_t,
\end{align*}
which is \eqref{cc:eq-transported-Chen}.  Naturality follows by applying
the intertwiner to \eqref{cc:eq-anchor-reconstruction}.  In the
stochastic case, adaptation of \(A\) implies
\({\mathcal F}_t\)-measurability of the reconstructed increment because
\({\mathcal F}_s\subseteq{\mathcal F}_t\); the converse follows from
\(A_t=\Phi_{0,t}\).  Finally use
\(\exp_\star({\operatorname{BCH}}(x,y))=\exp_\star(x)\star\exp_\star(y)\) and exponential
equivariance
\(\alpha_{t,s}\exp_\star
 =\exp_\star\alpha^{\mathrm{Lie}}_{t,s}\).
\end{proof}

Thus exceptional sets and both extensions can be treated at the anchor
level and then propagated by \eqref{cc:eq-anchor-reconstruction}.

\subsection{Optional completion: extension of the interval domain}

Let \((\Omega,{\mathcal F},({\mathcal F}_t)_{t\ge0},{\mathbb P})\) satisfy the usual conditions.
Suppose that a deterministic-time causal character is given only up to an
exceptional null set which may depend on \((s,t)\), and put
\(A_t=\Phi_{0,t}\).  We use the standard stopped sigma-fields and right
approximation convention; see \cite[Ch.~I]{JacodShiryaev03}.  The following
theorem removes all parameter-dependent exceptional sets by perfecting one
process.

\begin{theorem}[Anchored optional completion]
\label{cc:thm-optional-completion}
Assume that the adapted anchor \(A\) has a \(G\)-valued càdlàg modification
\(\bar A\).  For bounded stopping times \(0\le\sigma\le\tau\le T\), set
\begin{equation}
 \bar\Phi_{\sigma,\tau}
 :=\bigl(\alpha_{\tau,\sigma}(\bar A_\sigma)\bigr)^{-1}
    \star\bar A_\tau.
 \label{cc:eq-optional-anchor}
\end{equation}
Then:

\begin{enumerate}[label=\textup{(\roman*)},leftmargin=2.8em]
\item \(\bar\Phi_{\sigma,\tau}\) is \({\mathcal F}_\tau\)-measurable.  For each
fixed deterministic pair \((s,t)\), it agrees almost surely with
\(\Phi_{s,t}\).

\item On one event, simultaneously for all bounded stopping triples
\(\sigma\le\rho\le\tau\),
\begin{equation}
 \bar\Phi_{\sigma,\tau}
 =\alpha_{\tau,\rho}(\bar\Phi_{\sigma,\rho})
  \star\bar\Phi_{\rho,\tau}.
 \label{cc:eq-optional-Chen}
\end{equation}

\item If bounded stopping times
\(\sigma_n\downarrow\sigma\) and \(\tau_n\downarrow\tau\), with
\(\sigma_n\le\tau_n\), then
\begin{equation}
 \bar\Phi_{\sigma_n,\tau_n}\longrightarrow
 \bar\Phi_{\sigma,\tau}
 \qquad\text{almost surely}.
 \label{cc:eq-right-simple-limit}
\end{equation}

\item This extension is unique among assignments \(Y_{\sigma,\tau}\) with
the following two properties: for every finite-valued stopping pair,
\(Y_{\sigma,\tau}\) is the pointwise pasting of the perfected deterministic
increments \(\bar\Phi_{s,t}\); and for the right endpoints
\(\sigma^{(n)},\tau^{(n)}\) of the nested dyadic partitions of \([0,T]\),
\(Y_{\sigma^{(n)},\tau^{(n)}}\to Y_{\sigma,\tau}\) in probability.  It is
natural under continuous transport-intertwining homomorphisms.
\end{enumerate}
\end{theorem}

\begin{proof}
The deterministic anchor identity gives
\(A_t=\alpha_{t,s}(A_s)\star\Phi_{s,t}\) almost surely for every fixed
pair.  Replacing \(A\) by its modification proves deterministic agreement.
By completeness, we may redefine \(\bar A\) on the single null event on
which its path is not càdlàg; hence all subsequent pathwise identities may
be read literally on one common event.  Moreover
\(\bar A_t=A_t\) almost surely and completeness of each \({\mathcal F}_t\) make
\(\bar A\) adapted.
The calculation
\begin{align*}
 &\alpha_{\tau,\rho}\!\left(
  \bigl(\alpha_{\rho,\sigma}(\bar A_\sigma)\bigr)^{-1}
  \star\bar A_\rho\right)
  \star
  \bigl(\alpha_{\tau,\rho}(\bar A_\rho)\bigr)^{-1}
  \star\bar A_\tau\\
 &\qquad=
 \bigl(\alpha_{\tau,\sigma}(\bar A_\sigma)\bigr)^{-1}
 \star\bar A_\tau
\end{align*}
is pathwise.  It holds for every numerical triple and hence after
substitution of arbitrary stopping times; no uncountable intersection of
null events is involved.

An adapted càdlàg process is optional, so \(\bar A_\sigma\) and
\(\bar A_\tau\) are respectively \({\mathcal F}_\sigma\)- and
\({\mathcal F}_\tau\)-measurable.  Since \(\sigma\le\tau\), one has
\({\mathcal F}_\sigma\subseteq{\mathcal F}_\tau\), and the random variables \(\sigma\) and
\(\tau\) are also \({\mathcal F}_\tau\)-measurable.  Joint continuity of the
transport evaluation, followed by continuity of inversion and
multiplication, therefore proves \({\mathcal F}_\tau\)-measurability of
\eqref{cc:eq-optional-anchor}.  Pathwise right continuity gives
\(\bar A_{\sigma_n}\to\bar A_\sigma\) and
\(\bar A_{\tau_n}\to\bar A_\tau\); joint continuity of the same three maps
proves \eqref{cc:eq-right-simple-limit}.

For uniqueness, the case \(T=0\) is immediate.  If \(T>0\), let
\[
 d_n(r):=T2^{-n}\Bigl\lceil 2^n r/T\Bigr\rceil,
 \qquad d_n(0):=0,
\]
and set \(\sigma^{(n)}=d_n(\sigma)\),
\(\tau^{(n)}=d_n(\tau)\).  These are finite-valued stopping times,
decrease to the original endpoints, and preserve
\(\sigma^{(n)}\le\tau^{(n)}\).  Finite-valued pasting identifies every
competing assignment with \(\bar\Phi\) on this countable approximating
family.  Its assumed convergence and part~\textup{(iii)} give two limits
in probability of the same sequence.  Uniqueness of limits in probability
therefore gives equality almost surely at \((\sigma,\tau)\).  Naturality
follows by applying the intertwiner directly to
\eqref{cc:eq-optional-anchor}.
\end{proof}

\begin{corollary}[Countable projective optional completion]
\label{cc:cor-projective-optional-completion}
Suppose
\[
 G=\varprojlim_{j\in\mathbb N}G_j
\]
is realized as a closed subgroup of a countable product of Polish groups,
with continuous bonding homomorphisms, and suppose that the transports and
anchors commute with those maps.  If every coordinate anchor has an adapted
càdlàg modification, then these modifications can be chosen
indistinguishably compatible.  They assemble to a càdlàg \(G\)-valued
anchor, and \eqref{cc:eq-optional-anchor} gives the unique projective
optional completion.
\end{corollary}

\begin{proof}
For every bonding map, compatibility holds almost surely at every
deterministic time.  Intersect over the countable family of bonding maps
and rational times, then use right continuity to obtain compatibility at
all times on one event.  The coordinate path therefore takes values in the
closed inverse-limit subgroup.  For a standard bounded product metric,
coordinatewise right continuity and coordinatewise left limits imply
product càdlàgness: control finitely many leading coordinates and then the
metric tail.  Apply \Cref{cc:thm-optional-completion} in every coordinate.
\end{proof}

For an uncountable projective system, coordinatewise perfection alone does
not supply one common full-probability event.  A countable cofinal
realization, a countable separating family with a reconstruction theorem,
or an explicit path-separability assumption is then required.

The right approximation fixes the half-open convention \((s,t]\).  At a
jump, left approximations generally recover \(\bar A_{s-}\) or
\(\bar A_{t-}\) and define a different interval convention.  The càdlàg
hypothesis is also substantive: optional measurability and deterministic
agreement alone do not determine values on the graph of an atomless
stopping time.  Indeed, let \(\theta\) be an atomless bounded stopping time.
The optional process \(Y_t={\mathbf 1}_{\{t=\theta\}}\) satisfies
\(Y_t=0\) almost surely for every deterministic \(t\), whereas
\(Y_\theta=1\) surely.  Thus deterministic sections cannot distinguish
\(Y\) from the zero process; càdlàg perfection rules out precisely this
graph modification.

For a deterministic terminal cap \(q\), the transported anchor
\(A^{[q]}_r:=\alpha_{q,r}(\bar A_r)\) satisfies
\begin{equation}
 (A^{[q]}_\sigma)^{-1}\star A^{[q]}_\tau
 =\alpha_{q,\tau}(\bar\Phi_{\sigma,\tau}),
 \qquad \sigma\le\tau\le q.
 \label{cc:eq-terminal-fibre}
\end{equation}
If a first coordinate is realized as
\(A^{(1)}_t=\int_{(0,t]}U(t,a)\,{\mathop{}\!\mathrm d} X_a\), with
\(\alpha^{(1)}_{t,s}=U(t,s)\), this algebraic identity does not by itself
define \(\int_{(\sigma,\tau]}U(\tau,a)\,{\mathop{}\!\mathrm d} X_a\) as an Itô integral: the
random-terminal integrand need not be predictable.  For example, take
\(\sigma=0\), Brownian motion, \(U(t,s)=e^{t-s}\), and
\[
 \tau=t_0{\mathbf 1}_{\{B_{t_0}>0\}}+T{\mathbf 1}_{\{B_{t_0}\le0\}}.
\]
For every \(a\in(0,t_0)\), the formal integrand distinguishes
\(\{B_{t_0}>0\}\), which is not \({\mathcal F}_a\)-measurable.  The identity does
give that identification when \(\tau\) is \({\mathcal F}_\sigma\)-measurable, the
integrand has a predictable stochastically integrable version, and
finite-valued \({\mathcal F}_\sigma\)-measurable approximations converge in both the
integrand and endpoint topologies, by locality and stability of stochastic
integration; see \cite[Ch.~4]{Applebaum09}.

\subsection{Singular extension: lifting the target}

We now keep the interval domain fixed and enlarge the value group, in the
spirit of character-valued renormalization and enlarged model spaces
\cite{ConnesKreimer00,Hairer14}.

Let
\begin{equation}
 1\longrightarrow K\longrightarrow\widetilde G
 \xrightarrow{\,p\,}G\longrightarrow1
 \label{cc:eq-target-extension}
\end{equation}
be an exact sequence of groups.  Assume that \(\widetilde G\) has transports
\(\widetilde\alpha_{t,s}\) satisfying the analogue of
\eqref{cc:eq-transport-system}, preserving \(K\), and
\begin{equation}
 p\widetilde\alpha_{t,s}=\alpha_{t,s}p.
 \label{cc:eq-lifted-transport}
\end{equation}
For the algebraic statements no topology is needed.  Whenever Borel,
càdlàg, or optional assertions are made, we additionally assume that
\(\widetilde G\) is a Polish group, that \(K\) is closed, and that \(p\) and
the joint lifted-transport evaluation are continuous.
In the fixed trivialization we write
\(\widetilde G_t:=\widetilde G\) and \(K_t:=K\).
Choose a normalized set-theoretic section \(q:G\to\widetilde G\), and for
a base character \(\Phi\) define the pointwise lift and its interval defect
\begin{equation}
 Q_{s,t}:=q(\Phi_{s,t}),
 \qquad
 \Omega_{s,u,t}
 :=\widetilde\alpha_{t,u}(Q_{s,u})\star Q_{u,t}
   \star Q_{s,t}^{-1}\in K.
 \label{cc:eq-lifting-defect}
\end{equation}
The section \(q\) and the anchor selections are purely algebraic unless
measurability, continuity, or admissibility is explicitly imposed.

\begin{theorem}[Singular lifting, counterterms, and gauge]
\label{cc:thm-singular-lifting}
The defects in \eqref{cc:eq-lifting-defect} obey the nonabelian identity
\begin{equation}
 \widetilde\alpha_{t,v}(\Omega_{s,u,v})\star\Omega_{s,v,t}
 =\operatorname{Ad}_{\widetilde\alpha_{t,u}(Q_{s,u})}
      (\Omega_{u,v,t})\star\Omega_{s,u,t}
 \label{cc:eq-nonabelian-two-cocycle}
\end{equation}
for \(s\le u\le v\le t\).  If
\(Q'_{s,t}=r_{s,t}\star Q_{s,t}\), \(r_{s,t}\in K\), then
\begin{equation}
 \Omega'_{s,u,t}
 =\widetilde\alpha_{t,u}(r_{s,u})\star
  \operatorname{Ad}_{\widetilde\alpha_{t,u}(Q_{s,u})}(r_{u,t})
  \star\Omega_{s,u,t}\star r_{s,t}^{-1}.
 \label{cc:eq-defect-gauge-change}
\end{equation}
Consequently, corrected lifts
\(\widetilde\Phi_{s,t}=k_{s,t}\star Q_{s,t}\) are causal exactly when
\begin{equation}
 k_{s,t}=\widetilde\alpha_{t,u}(k_{s,u})\star
 \operatorname{Ad}_{\widetilde\alpha_{t,u}(Q_{s,u})}(k_{u,t})
 \star\Omega_{s,u,t}.
 \label{cc:eq-counterterm-equation}
\end{equation}

Nevertheless, the unrestricted interval problem has no algebraic existence
obstruction.  If \(H_t\in p^{-1}(A_t)\), \(H_0={\mathbf 1}\), then
\begin{equation}
 L^H_{s,t}:=
 \bigl(\widetilde\alpha_{t,s}(H_s)\bigr)^{-1}\star H_t
 \label{cc:eq-rooted-lift}
\end{equation}
is a causal lift, and every causal lift has this form with
\(H_t=L_{0,t}\).  The based vertex gauges \(b_t\in K\), \(b_0={\mathbf 1}\), act by
\begin{equation}
 (L^b)_{s,t}
 :=\widetilde\alpha_{t,s}(b_s)^{-1}\star L_{s,t}\star b_t,
 \label{cc:eq-vertex-gauge}
\end{equation}
freely and transitively on all unrestricted lifts.  Thus causality
shows that \(\Omega\) is the defect of the chosen pointwise section rather
than an obstruction to an unrestricted rooted lift; it selects no
renormalization gauge.
\end{theorem}

\begin{proof}
Compose the pointwise lifts over \((s,u]\), \((u,v]\), and \((v,t]\).
Left association equals
\[
 \widetilde\alpha_{t,v}(\Omega_{s,u,v})\star
 \Omega_{s,v,t}\star Q_{s,t},
\]
whereas right association equals
\[
 \operatorname{Ad}_{\widetilde\alpha_{t,u}(Q_{s,u})}
 (\Omega_{u,v,t})\star\Omega_{s,u,t}\star Q_{s,t}.
\]
Associativity and cancellation prove
\eqref{cc:eq-nonabelian-two-cocycle}.  Substitution of
\(Q'=r\star Q\) into \eqref{cc:eq-lifting-defect} gives
\eqref{cc:eq-defect-gauge-change}; setting \(r=k\) and requiring the new
defect to be the identity gives \eqref{cc:eq-counterterm-equation}.

Since \(p(H_t)=A_t\), the anchor identity in \(G\) gives

\[
 p(L^H_{s,t})
 =\alpha_{t,s}(A_s)^{-1}\star A_t=\Phi_{s,t}.
\]

Cancellation of the middle \(H_u\)-factors proves causal composition of
\(L^H\).  Conversely, applying causality of a lift \(L\) to \(0\le s\le t\)
gives \(L_{0,t}=\widetilde\alpha_{t,s}(L_{0,s})\star L_{s,t}\), and hence
the rooted formula.  If \(L,L'\) are two lifts, then
\(b_t=L_{0,t}^{-1}\star L'_{0,t}\in K\) is the unique based vertex gauge
with \(L^b=L'\).  This proves freeness and transitivity.
\end{proof}

Relative to a causal lift \(L\), every lift has the unique form
\(L'_{s,t}=a_{s,t}\star L_{s,t}\), \(a_{s,t}\in K\), and is causal
exactly when
\begin{equation}
 a_{s,t}=\widetilde\alpha_{t,u}(a_{s,u})\star
 \operatorname{Ad}_{\widetilde\alpha_{t,u}(L_{s,u})}(a_{u,t}).
 \label{cc:eq-relative-one-cocycle}
\end{equation}
A vertex gauge transforms this relative cocycle by
\begin{equation}
 a^b_{s,t}=\widetilde\alpha_{t,s}(b_s)^{-1}\star a_{s,t}\star
 \operatorname{Ad}_{L_{s,t}}(b_t).
 \label{cc:eq-relative-gauge}
\end{equation}
These formulas follow by direct substitution.  For nonabelian \(K\) they
define a cocycle--gauge groupoid rather than a group under pointwise
multiplication.  For abelian \(K\), the causal lifts form a torsor under
the twisted cocycle group \(Z^1\); quotienting by the allowed vertex
gauges gives the corresponding relative \(H^1\).

An admissible singular problem restricts
the pointwise lifts and gauges by continuity, adaptedness, locality,
optional compatibility, spatial naturality, or a completed topology.  The
rooted construction solves that restricted problem only when the \(H_t\)
can be selected in the admissible class.

The set section \(q\) records the extension by its Schreier data
\begin{align}
 \rho(g)(k)&:=q(g)\star k\star q(g)^{-1},\notag\\
 f(g,h)&:=q(g)\star q(h)\star q(g\star h)^{-1},\notag\\
 d_{t,s}(g)&:=\widetilde\alpha_{t,s}(q(g))
       \star q(\alpha_{t,s}g)^{-1}.
 \label{cc:eq-Schreier-data}
\end{align}
Here \(f(g,h),d_{t,s}(g)\in K\), and \(\rho(g)\) is an automorphism of
\(K\).  These quantities depend on \(q\), while the extension and its
outer action do not.

\begin{theorem}[Canonical splittings and optional naturality]
\label{cc:prop-canonicality-criterion}
A source-independent pointwise lifting rule \(j:G\to\widetilde G\) is
multiplicative and transport-natural if and only if
\begin{equation}
 pj={\mathrm{Id}},\qquad
 j(g\star h)=j(g)\star j(h),\qquad
 \widetilde\alpha_{t,s}j=j\alpha_{t,s}.
 \label{cc:eq-strict-splitting}
\end{equation}
Writing \(j(g)=r(g)\star q(g)\), these conditions are equivalent to
\begin{align}
 r(g\star h)
 &=r(g)\star\rho(g)(r(h))\star f(g,h),
 \label{cc:eq-splitting-factor-equation}\\
 r(\alpha_{t,s}g)
 &=\widetilde\alpha_{t,s}(r(g))\star d_{t,s}(g),
 \label{cc:eq-splitting-transport-equation}
\end{align}
where \(r:G\to K\) and \(r({\mathbf 1})={\mathbf 1}\).  Every such splitting sends
causal characters to causal lifts.  If \(j\) is continuous and \(\bar A\)
is a càdlàg base anchor, then \(j(\bar A)\) is a càdlàg lifted anchor and
\begin{equation}
 \widetilde\Phi_{\sigma,\tau}=j(\bar\Phi_{\sigma,\tau})
 \label{cc:eq-optional-splitting-naturality}
\end{equation}
for all bounded stopping pairs.

More generally, if \(H\) is an adapted càdlàg lifted anchor and \(b\) is
an adapted càdlàg \(K\)-valued gauge with \(b_0={\mathbf 1}\), then
\(H^b_t=H_t\star b_t\) has stopped increments
\begin{equation}
 (L^b)_{\sigma,\tau}
 =\widetilde\alpha_{\tau,\sigma}(b_\sigma)^{-1}
  \star L_{\sigma,\tau}\star b_\tau.
 \label{cc:eq-stopped-vertex-gauge}
\end{equation}
Thus optional completion preserves gauge arrows but selects no gauge.
Existence of a splitting does not imply uniqueness; within this
source-independent class, a canonical rule requires the admissibility
conditions to select a unique splitting.  For random characters the
unstopped rule requires \(j\) to be Borel, and the stopped assertion
requires continuity.
\end{theorem}

\begin{proof}
Using the Schreier data, direct multiplication gives
\[
 j(g)\star j(h)
 =r(g)\star\rho(g)(r(h))\star f(g,h)\star q(g\star h),
\]
and transport gives
\[
 \widetilde\alpha_{t,s}(j(g))
 =\widetilde\alpha_{t,s}(r(g))\star d_{t,s}(g)
  \star q(\alpha_{t,s}g).
\]
Cancellation proves the equivalence with the two equations for \(r\).
Applying a strict splitting to the transported Chen identity proves the
lifting assertion.  Continuity gives càdlàgness of \(j(\bar A)\), and
the anchor formula, multiplicativity, and transport equivariance give
\[
 \bigl(\widetilde\alpha_{\tau,\sigma}(j(\bar A_\sigma))\bigr)^{-1}
 \star j(\bar A_\tau)
 =j\!\left(
  \bigl(\alpha_{\tau,\sigma}(\bar A_\sigma)\bigr)^{-1}
  \star\bar A_\tau\right).
\]
Finally, insertion of \(H^b=H\star b\) into the anchor formula yields
\eqref{cc:eq-stopped-vertex-gauge}.
\end{proof}

\subsection{Interchange of optional completion and singular extension}

The preceding axes commute only after specifying a topology strong enough
to evaluate at all stopping times.  We state the criterion in the finite
nilpotent targets used by truncated character theories.  Let
\(\mathbb K\in\{\mathbb R,\mathbb C\}\), and let
\(\widetilde{\mathcal A}=\mathbb K{\mathbf 1}\oplus\mathfrak n\) be a separable unital Banach
algebra with \(\mathfrak n^{\star(N+1)}=0\), and put
\(\widetilde G={\mathbf 1}+\mathfrak n\).  Inversion is the finite polynomial
\begin{equation}
 ({\mathbf 1}+x)^{-1}=\sum_{j=0}^{N}(-x)^{\star j},
 \label{cc:eq-polynomial-inverse}
\end{equation}
so multiplication and inversion are locally Lipschitz.  Assume also that
the extension projection \(p\) is continuous and that the lifted
transports are uniformly locally Lipschitz on compact time
triangles: for each \(R\) there is \(L_{T,R}\) such that
\begin{equation}
 \|\widetilde\alpha_{t,s}(x)-\widetilde\alpha_{t,s}(y)\|
 \le L_{T,R}\|x-y\|
 \label{cc:eq-uniform-transport-Lipschitz}
\end{equation}
whenever \(s\le t\le T\) and \(\|x\|,\|y\|\le R\).

\begin{lemma}[The stopping-time meaning of ucp]
\label{cc:lem-ucp-characterization}
For adapted càdlàg processes \(X^{(m)},X\) in a separable metric space,
the following are equivalent on every compact time interval:
\begin{enumerate}[label=\textup{(\roman*)},leftmargin=2.8em]
\item \(\sup_{t\le T}d(X^{(m)}_t,X_t)\to0\) in probability;
\item for every sequence of possibly \(m\)-dependent stopping times
\(\tau_m\le T\), one has
\(d(X^{(m)}_{\tau_m},X_{\tau_m})\to0\) in probability.
\end{enumerate}
\end{lemma}

\begin{proof}
The first condition immediately dominates every stopped evaluation.  If it
fails, pass to and relabel a subsequence for which

\[
 {\mathbb P}\!\left\{\sup_{t\le T}d(X^{(m)}_t,X_t)>2\varepsilon\right\}
 >\delta
\]

for some \(\varepsilon,\delta>0\).  With
\(Y^{(m)}_t=d(X^{(m)}_t,X_t)\), the first entrance time

\[
 \tau_m=\inf\{t\le T:d(X^{(m)}_t,X_t)>\varepsilon\}\wedge T
\]

is a stopping time because \(Y^{(m)}\) is adapted and càdlàg and the
filtration is right-continuous.  On
\(\{\sup_{t\le T}Y^{(m)}_t>2\varepsilon\}\), the open set
\((\varepsilon,\infty)\) is entered before or at \(T\); right continuity at
the entrance gives \(Y^{(m)}_{\tau_m}\ge\varepsilon\).  Consequently
\({\mathbb P}\{d(X^{(m)}_{\tau_m},X_{\tau_m})>\varepsilon/2\}>\delta\),
contradicting part~\textup{(ii)}.
\end{proof}

\medskip
\noindent\emph{Fixed times do not control stopping times.}
Assume \(\mathfrak n\ne0\), choose \(0\ne z\in\mathfrak n\) with
\(z\star z=0\), and take trivial transport on \([0,1]\).  For
\[
 t_m=1-m^{-1},\qquad \varepsilon_m=(4m^2)^{-1},\qquad
 h_m(t)=\left(1-\frac{|t-t_m|}{\varepsilon_m}\right)_+,
\]
the continuous anchors \(A^{(m)}_t={\mathbf 1}+h_m(t)z\) induce exact characters
\[
 \Phi^{(m)}_{s,t}={\mathbf 1}+(h_m(t)-h_m(s))z.
\]
The supports escape every fixed \(t<1\), while \(h_m(1)=0\); hence
\(\Phi^{(m)}_{s,t}\to{\mathbf 1}\) for each deterministic pair.  Nevertheless,
at \(\tau_m=t_m\),
\(\Phi^{(m)}_{0,\tau_m}={\mathbf 1}+z\).  Such a \(z\) always exists in a
nonzero finite nilpotent algebra: take a nonzero element of its last
nonzero ideal power.  Thus fixed-time convergence cannot support optional
interchange.
\medskip

\begin{theorem}[Optional--singular interchange]
\label{cc:thm-optional-singular-interchange}
Let \(\widetilde A^{(m)},\widetilde A\) be normalized adapted càdlàg
\(\widetilde G\)-valued anchors, write
\(A^{(m)}:=p(\widetilde A^{(m)})\) and
\(A:=p(\widetilde A)\) for their base anchors, and assume
\begin{equation}
 \sup_{t\le T}\|\widetilde A^{(m)}_t-\widetilde A_t\|
 \longrightarrow0
 \qquad\text{in probability}.
 \label{cc:eq-lifted-anchor-ucp}
\end{equation}
Define their perfected lifted characters by
\begin{equation}
 \widetilde\Phi^{(m)}_{s,t}
 :=\bigl(\widetilde\alpha_{t,s}(\widetilde A^{(m)}_s)\bigr)^{-1}
   \star\widetilde A^{(m)}_t,
 \qquad
 \widetilde\Phi_{s,t}
 :=\bigl(\widetilde\alpha_{t,s}(\widetilde A_s)\bigr)^{-1}
   \star\widetilde A_t.
 \label{cc:eq-lifted-perfect-characters}
\end{equation}
Then
\begin{equation}
 \sup_{0\le s\le t\le T}
 \|\widetilde\Phi^{(m)}_{s,t}-\widetilde\Phi_{s,t}\|
 \longrightarrow0
 \qquad\text{in probability}.
 \label{cc:eq-uniform-triangle-convergence}
\end{equation}
In particular, for every possibly \(m\)-dependent stopping pair
\(\sigma_m\le\tau_m\le T\),
\begin{equation}
 \widetilde\Phi^{(m)}_{\sigma_m,\tau_m}
 \longrightarrow\widetilde\Phi_{\sigma_m,\tau_m}
 \qquad\text{in probability}.
 \label{cc:eq-stopped-lift-convergence}
\end{equation}
For a fixed stopping pair, the two routes agree in the following
commuting-limit sense:
\begin{equation}
 \bigl(\widetilde\alpha_{\tau,\sigma}
       (\widetilde A_\sigma)\bigr)^{-1}
       \star\widetilde A_\tau
 =
 \lim_{m\to\infty}
 \bigl(\widetilde\alpha_{\tau,\sigma}
       (\widetilde A^{(m)}_\sigma)\bigr)^{-1}
       \star\widetilde A^{(m)}_\tau
 \quad\text{in probability}.
 \label{cc:eq-optional-limit-interchange}
\end{equation}
It is compatible with the projection \(p:\widetilde G\to G\).

More generally, suppose raw counterterm-corrected intervals
\(\widehat\Phi^{(m)}_{s,t}\) have measurable stopped evaluations and
measurable displayed suprema (as holds, for example, for jointly measurable
fields with separable sample paths), have anchors
\(\widetilde A^{(m)}_t=\widehat\Phi^{(m)}_{0,t}\), satisfy
\eqref{cc:eq-lifted-anchor-ucp}, and have uniform anchored defect
\begin{equation}
 \sup_{s\le t\le T}
 \left\|
 \widetilde A^{(m)}_t-
 \widetilde\alpha_{t,s}(\widetilde A^{(m)}_s)
   \star\widehat\Phi^{(m)}_{s,t}
 \right\|
 \longrightarrow0
 \quad\text{in probability}.
 \label{cc:eq-uniform-anchored-defect}
\end{equation}
Then \(\widehat\Phi^{(m)}\) has the same uniform triangle limit
\(\widetilde\Phi\), in the explicit sense that
\[
 \sup_{s\le t\le T}
 \|\widehat\Phi^{(m)}_{s,t}-\widetilde\Phi_{s,t}\|
 \longrightarrow0
 \qquad\text{in probability},
\]
and hence the same stopped limits.  Thus
\eqref{cc:eq-lifted-anchor-ucp} together with
\eqref{cc:eq-uniform-anchored-defect} is the robust upgrade of a
fixed-interval renormalization certificate.
\end{theorem}

\begin{proof}
Right approximation of both endpoints shows that each perfected increment
field is right-continuous on the time triangle in the product order.
Consequently every supremum below equals the supremum over the countable
union of dyadic time triangles and is measurable.
On each norm ball, \eqref{cc:eq-polynomial-inverse}, the
submultiplicativity of the Banach-algebra norm, and the telescoping identity
for noncommutative powers give a constant \(K_R\) such that
\[
 \|x^{-1}\star x'-y^{-1}\star y'\|
 \le K_R(\|x-y\|+\|x'-y'\|).
\]
A càdlàg path is bounded on \([0,T]\).  Given \(\eta>0\), first choose
\(R_0\) so that
\({\mathbb P}\{\sup_{t\le T}\|\widetilde A_t\|>R_0\}<\eta\); then
\eqref{cc:eq-lifted-anchor-ucp} shows that, for all sufficiently large
\(m\), both anchors lie in the ball of radius \(R_0+1\) outside an event of
probability at most \(2\eta\).  Since every group endomorphism fixes
\({\mathbf 1}\), \eqref{cc:eq-uniform-transport-Lipschitz} also bounds all
transported anchors there, uniformly over the time triangle.  The preceding
polynomial estimate and \eqref{cc:eq-uniform-transport-Lipschitz} now
yield
\[
 \sup_{s\le t}
 \|\widetilde\Phi^{(m)}_{s,t}-\widetilde\Phi_{s,t}\|
 \le K_{T,R}\sup_{r\le T}
 \|\widetilde A^{(m)}_r-\widetilde A_r\|.
\]
For fixed error threshold, \eqref{cc:eq-lifted-anchor-ucp} makes the
right-hand side negligible in probability; then letting \(\eta\downarrow0\)
removes the localization event.  This proves
\eqref{cc:eq-uniform-triangle-convergence}; evaluation gives
\eqref{cc:eq-stopped-lift-convergence}.  Projection compatibility follows
from \eqref{cc:eq-lifted-transport}: explicitly,
\[
 p(\widetilde\Phi^{(m)}_{s,t})
 =\bigl(\alpha_{t,s}(A^{(m)}_s)\bigr)^{-1}\star A^{(m)}_t,
 \qquad
 p(\widetilde\Phi_{s,t})
 =\bigl(\alpha_{t,s}(A_s)\bigr)^{-1}\star A_t.
\]
Continuity of \(p\) passes these identities to every asserted limit.

For the raw family, write
\(a^{(m)}_{s,t}=\widetilde\alpha_{t,s}(\widetilde A^{(m)}_s)\).  Then
\begin{align*}
 \widehat\Phi^{(m)}_{s,t}
 -(a^{(m)}_{s,t})^{-1}\star\widetilde A^{(m)}_t
 &=(a^{(m)}_{s,t})^{-1}\star\bigl(
   a^{(m)}_{s,t}\star\widehat\Phi^{(m)}_{s,t}
   -\widetilde A^{(m)}_t\bigr).
\end{align*}
The inverse suprema are tight, so
\eqref{cc:eq-uniform-anchored-defect} makes the raw family uniformly
close to the perfected one.  The first part completes the proof.
\end{proof}

By \Cref{cc:lem-ucp-characterization}, ucp is exactly the convergence of
anchors detected by all approximation-dependent stopping times; the bump
construction shows that fixed-time convergence is insufficient.  Optional
completion is canonical after perfection of the anchor, whereas singular
extension remains relative to an admissible lift.  Once that lift is chosen,
ucp convergence gives the asserted interchange.

\begin{lemma}[Borel closure of probabilistic limits]
\label{cc:lem-borel-probability-limit}
Let \(S,Y\) be Polish spaces, let \(\mu\) be a probability measure on
\(S\), and let \(\Gamma_m:S\to Y\) be Borel maps converging in
\(\mu\)-probability to a unique \(Y\)-valued limit.  Then a deterministic
subsequence converges pointwise on a Borel set of full \(\mu\)-measure;
extending that limit arbitrarily off the set gives a Borel version
\(\Gamma:S\to Y\).
\end{lemma}

\begin{proof}
Choose a subsequence whose distances to the limit are summable in
probability and apply Borel--Cantelli.  Its Cauchy set is Borel, and the
pointwise limit of Borel maps into a Polish space is Borel there.
\end{proof}

\subsection{The causal normal object}
\label{cc:subsec-causal-normal-object}

\begin{definition}[Causal normal geometry]
\label{cc:def-causal-normal-geometry}
A causal normal geometry of step \(N\) consists of a causal evolution
character \(\mathfrak C=(G,\alpha,\Phi)\), a chosen admissible causal lift
\(\widetilde\Phi\) through a transported target extension (the identity
extension is allowed), and a transport-compatible geometric chart
\[
 \widetilde\chi_t:\widetilde G_t^{\rm adm}
 \longrightarrow\mathsf G_{\rm geo}^{\le N}(V_t).
\]
We write
\[
 \mathfrak X=(\mathfrak C;\widetilde\Phi,\widetilde\chi),
 \qquad
 X^{\mathfrak X}_{s,t}
 :=\log_\star\widetilde\chi_t(\widetilde\Phi_{s,t}).
\]
The lift and chart are realization data; the normal filtration itself is the
canonical commutativization filtration of the charted free Lie algebra.  The
geometry is \emph{base-canonical} when the chart is trivial on the
admissible kernel, in which case it descends to \(\Phi\).
\end{definition}

\begin{proposition}[Causal-normal completion principle]
\label{cc:prop-causal-normal-completion}
Assume the hypotheses of
\Cref{cc:thm-optional-completion,cc:thm-optional-singular-interchange} for a
chosen admissible lifted anchor, and assume the chart preserves the selected
c\`adl\`ag paths.  Then the causal normal geometry extends uniquely to
bounded stopping intervals.  Optional completion, the selected singular
limit, and application of the chart commute in probability; adapted
c\`adl\`ag gauges act by the stopped vertex formula
\eqref{cc:eq-stopped-vertex-gauge}.  Thus stopping changes the domain,
singular lifting changes the target, and normal geometry is read only after
both operations have been fixed.
\end{proposition}

\begin{proof}
The stopped lifted character and its uniqueness are given by
\Cref{cc:thm-optional-completion}.  The commuting limit is
\Cref{cc:thm-optional-singular-interchange}; continuity of the chart permits
its application after the limit.  Gauge naturality is
\eqref{cc:eq-stopped-vertex-gauge}.
\end{proof}
\section{Primitive Rees geometry}
\label{cc:sec-algebraic-ancestry}

Let \(\mathbf k\) be a field of characteristic zero and let \(V\) be an
arbitrary \(\mathbf k\)-vector space.  Put
\[
 L:={\mathfrak L}(V),\qquad D:=[L,L],\qquad U(L)=T(V),
\]
where the last identification is the canonical enveloping realization of a
free Lie algebra \cite[Thm.~0.5]{Reutenauer93}.  Let
\[
 \mathfrak a:T(V)\longrightarrow S(V),\qquad
 J:=\ker\mathfrak a
\]
be associative abelianization and the two-sided commutator ideal defining
commutativization.  This is the standard commutator thickening of \(S(V)\)
\cite{Kapranov98,FeiginShoikhet07,JordanOrem13}.  Its primitive normal
cone will identify the exact bracket ancestry of a causal logarithm.

Indeed, after a transport-compatible geometric chart has been fixed, the
charted logarithm of a causal character satisfies
\[
 X_{s,t}={\operatorname{BCH}}\!\left(\beta_{t,u}X_{s,u},X_{u,t}\right).
\]
For \(0\ne x\in L\), put
\[
 \operatorname{ord}_{\mathrm{anc}}(x)
 :=\max\{r\geq0:x\in L\cap J^r\}.
\]
The theorem below identifies this normal order with lower-central bracket
length inside the derived ideal.

Define
\[
 \gamma_1(D):=D,\qquad \gamma_{r+1}(D):=[D,\gamma_r(D)],
 \qquad Q:=D/[D,D].
\]
For \(v\in V\) and \(d\in D\), set
\begin{equation}
 \rho_v(\bar d):=\overline{[v,d]}.
 \label{cc:eq-rho}
\end{equation}
This is well defined.  Jacobi gives
\([\rho_v,\rho_w](\bar d)=\overline{[[v,w],d]}=0\), so \(Q\) is
an \(S(V)\)-module.  Each \(\rho_v\) extends uniquely to a derivation of
the free Lie algebra \({\mathfrak L}(Q)\), and the extensions still commute because
they commute on its free generators.  We write
\(T(Q)\#_\rho S(V)\) for the smash product generated by \(V\) and \(Q\)
with relations
\begin{equation}
 vw=wv,\qquad vq-qv=\rho_vq.
 \label{cc:eq-smash-relations}
\end{equation}
The normal degrees of \(V\) and \(Q\) are zero and one, respectively.  Its
Hopf structure is the enveloping Hopf structure for which every element of
\(V\oplus Q\) is primitive.

\begin{theorem}[Primitive normal-cone theorem]
\label{cc:thm-normal-cone}
There is a natural isomorphism of normal-graded Hopf algebras
\begin{equation}
 {\operatorname{gr}}_JT(V)\cong T(Q)\#_\rho S(V)
 \cong U\!\left(V_{\mathrm{ab}}\ltimes_\rho{\mathfrak L}(Q)\right).
 \label{cc:eq-normal-cone}
\end{equation}
The isomorphism is induced by the identity on the degree-zero primitives
\(L/D\cong V\) and the degree-one primitives \(D/[D,D]=Q\), and it
respects ordinary tensor degree.
Consequently,
\begin{align}
 J^r/J^{r+1}&\cong Q^{\otimes r}\otimes S(V),
 \label{cc:eq-normal-layer}\\
 {\operatorname{Prim}}({\operatorname{gr}}_JT(V))_0&\cong V,
 \qquad
 {\operatorname{Prim}}({\operatorname{gr}}_JT(V))_r
 \cong{\mathfrak L}_r(Q)\cong\gamma_r(D)/\gamma_{r+1}(D),
 \quad r\geq1,
 \label{cc:eq-primitive-layers}
\end{align}
and, most importantly,
\begin{equation}
 L\cap J^r=\gamma_r(D),\qquad r\geq1.
 \label{cc:eq-intersection}
\end{equation}
Thus normal order and derived-ideal ancestry agree exactly.
\end{theorem}

\begin{proof}
Give \(L\) the filtration
\[
 H_0L:=L,\qquad H_rL:=\gamma_r(D)\quad(r\geq1).
\]
Since \(D\) is an ideal, Jacobi and the standard lower-central induction
give
\begin{equation}
 [H_aL,H_bL]\subseteq H_{a+b}L
 \quad(a,b\geq0).
 \label{cc:eq-strong-filtration}
\end{equation}
The action of \(H_0L/H_1L=L/D\cong V_{\mathrm{ab}}\) on
\(H_1L/H_2L=Q\) is \(\rho\).

The derived algebra \(D\) is free by the Shirshov--Witt theorem
\cite[\S2.2]{Reutenauer93}.  For any free Lie algebra \(\mathfrak f\),
bracketing representatives gives a canonical surjection
\[
 {\mathfrak L}\!\left(\mathfrak f/[\mathfrak f,\mathfrak f]\right)
 \longrightarrow {\operatorname{gr}}_\gamma\mathfrak f
\]
of graded Lie algebras.  After choosing a free generating space \(W\),
both sides identify with the free Lie algebra on \(W\), graded by bracket
length.  The map is therefore an isomorphism; because it was defined
without \(W\), the isomorphism is canonical.  Applied to \(D\), it
preserves ordinary tensor degree because the lower central series of
\(D\) is homogeneous.
Moreover,
\({\operatorname{ad}}_x\) preserves
\(\gamma_r(D)\), and its action on \({\operatorname{gr}}_\gamma D\) depends only on
\(x+D\), because
\([D,\gamma_r(D)]\subseteq\gamma_{r+1}(D)\).  Hence the preceding
isomorphism is \(L/D\)-equivariant when \(\mathfrak f=D\).  The
commuting operators \(\rho_v\) extend by derivations from \(Q\) to
\({\mathfrak L}(Q)\), and therefore
\begin{equation}
 {\operatorname{gr}}_HL\cong V_{\mathrm{ab}}\ltimes_\rho{\mathfrak L}(Q).
 \label{cc:eq-graded-lie}
\end{equation}

We next identify the induced enveloping filtration.
Let \(\mathcal U_0:=U(L)\), and for \(r\geq1\) let \(\mathcal U_r\) be
the span of products
\(x_1\cdots x_m\), where \(x_i\in H_{a_i}L\) and
\(\sum_i a_i\geq r\).  Strong filteredness makes
\(\mathcal U_\bullet\) multiplicative and every \(\mathcal U_r\) a
two-sided ideal.  Expanding a bracket in
\(H_aL=\gamma_a(D)\) produces words with at least \(a\) factors from
\(D\); hence \(H_aL\subseteq J^a\) and
\(\mathcal U_r\subseteq J^r\).  Conversely, \(\mathcal U_1\) is the
two-sided ideal generated by \(H_1L=D\), so \(\mathcal U_1=J\), and
multiplicativity gives \(J^r=\mathcal U_1^r\subseteq\mathcal U_r\).
Therefore
\begin{equation}
 \mathcal U_r=J^r.
 \label{cc:eq-induced-is-j}
\end{equation}

We record the filtered PBW argument because it also proves the crucial
intersection statement; compare \cite{Quillen68}.  In ordinary
tensor degree \(n\), one has
\(H_rL_n=0\) for \(2r>n\).  Choose degreewise complements of
\(H_{a+1}L_n\) in \(H_aL_n\), assign weight \(a\) to their basis
vectors, and order their union.  The PBW theorem
\cite[\S0.1]{Reutenauer93} says that the ordered monomials in this basis
form a basis of \(U(L)\).  Interchanging vectors of weights \(a,b\)
creates a correction bracket of weight at least \(a+b\), so reordering
never lowers total weight.  It follows that \(\mathcal U_r\) is exactly
the span of ordered PBW monomials of total weight at least \(r\).
Consequently the exact-weight PBW
monomials form bases on both sides of
\[
 U({\operatorname{gr}}_HL)\longrightarrow{\operatorname{gr}}_{\mathcal U}U(L),
\]
so this canonical map is an isomorphism.  Since \(L\subset U(L)\) is
exactly the span of the length-one PBW monomials, the same basis
description also gives
\begin{equation}
 L\cap\mathcal U_r=H_rL.
 \label{cc:eq-pbw-intersection}
\end{equation}
Equations \eqref{cc:eq-induced-is-j} and
\eqref{cc:eq-pbw-intersection} prove
\eqref{cc:eq-intersection}, while
\eqref{cc:eq-graded-lie} gives
\[
 {\operatorname{gr}}_JT(V)\cong U({\operatorname{gr}}_HL)
 \cong U\!\left(V_{\mathrm{ab}}\ltimes_\rho{\mathfrak L}(Q)\right).
\]
Since \(U({\mathfrak L}(Q))=T(Q)\), PBW for the semidirect decomposition identifies
the last enveloping algebra with the smash product in
\eqref{cc:eq-normal-cone}; multiplication gives the vector-space
isomorphism \(T(Q)\otimes S(V)\to T(Q)\#_\rho S(V)\).  Its normal
degree-\(r\) vector space is therefore
\(Q^{\otimes r}\otimes S(V)\), proving
\eqref{cc:eq-normal-layer}.

The identification is also Hopf.  Indeed \(J\) is generated by the
primitive Lie ideal \(D\), so
\[
 \Delta(J)\subseteq J\otimes U(L)+U(L)\otimes J,
 \qquad S(J)\subseteq J,
\]
and multiplicativity of \(\Delta\) gives
\begin{equation}
 \Delta(J^r)\subseteq\sum_{a+b=r}J^a\otimes J^b,
 \qquad S(J^r)\subseteq J^r.
 \label{cc:eq-hopf-filtration}
\end{equation}
The filtered PBW map carries every primitive generator to its own normal
class.  Since these classes generate, it intertwines coproduct, counit,
and antipode.

In characteristic zero, \({\operatorname{Prim}} U(\mathfrak g)=\mathfrak g\): the top
PBW symbol of a primitive element is primitive in \(S(\mathfrak g)\),
whose only primitives have symmetric degree one, and the counit removes
the scalar term.  This is the Friedrichs characterization in the free
case \cite[Thm.~1.4]{Reutenauer93}.  Applying it to
\(\mathfrak g={\operatorname{gr}}_HL\) and using \eqref{cc:eq-graded-lie} gives the
degree-zero primitive space \(V\) and the degree-\(r\) primitive space
\({\mathfrak L}_r(Q)\).  The latter is
\(\gamma_r(D)/\gamma_{r+1}(D)\) by construction, proving
\eqref{cc:eq-primitive-layers} and completing the proof.
\end{proof}

Put \(F_0L(V):=L(V)\) and
\(F_rL(V):=L(V)\cap J(V)^r=\gamma_r(D(V))\) for \(r\geq1\), and write
\(F_rL_n(V):=F_rL(V)\cap L_n(V)\).  The 
\emph{full} primal layer and its algebraic dual are
\begin{equation}
 G^{\mathrm{full}}_{n,r}(V):=
 \frac{F_rL_n(V)}{F_{r+1}L_n(V)}
 \cong\bigl({\mathfrak L}_r(Q(V))\bigr)_n,
 \qquad
 E^{\mathrm{full}}_{n,r}(V):=
 \bigl(G^{\mathrm{full}}_{n,r}(V)\bigr)^{\vee}.
 \label{cc:eq-exact-layer}
\end{equation}
We retain the superscript ``\(\mathrm{full}\)'' throughout, so these
algebraic layers cannot be confused with the finite registered layers used
by the spectral reader.

For a time-dependent alphabet \(V_t\), abbreviate
\begin{equation}
 \begin{gathered}
 L_t:=L(V_t),\qquad D_t:=[L_t,L_t],\qquad
 J_t:=\ker\bigl(T(V_t)\to S(V_t)\bigr),\\
 F_0L_t:=L_t,\quad
 F_rL_t:=L_t\cap J_t^r=\gamma_r(D_t)\ (r\geq1),\qquad
 G^{\mathrm{full}}_{n,r}(t):=
 \dfrac{F_rL_t\cap(L_t)_n}{F_{r+1}L_t\cap(L_t)_n}.
 \end{gathered}
 \label{cc:eq-fibrewise-full-layer}
\end{equation}
The canonical map at order \(r\) is only the quotient
\begin{equation}
 \varpi^{(t)}_{n,r}:
 F_rL_t\cap(L_t)_n\longrightarrow G^{\mathrm{full}}_{n,r}(t).
 \label{cc:eq-normal-symbol-quotient}
\end{equation}
This quotient map is defined on the filtered subspace, not on all of
\((L_t)_n\).  For \(0\ne x\in(L_t)_n\), set
\(\nu_{F,t}(x):=\max\{r:x\in F_rL_t\}\).  Its only canonical homogeneous
component is its \emph{full initial symbol}
\begin{equation}
 \operatorname{in}^{\mathrm{full}}_{F,t}(x)
 :=\varpi^{(t)}_{n,\nu_{F,t}(x)}(x)
 \in G^{\mathrm{full}}_{n,\nu_{F,t}(x)}(t).
 \label{cc:eq-canonical-initial-symbol}
\end{equation}
All further homogeneous components require a splitting or a Rees lift.

Suppose that a chart
\(\chi_t:G_t\to\mathsf G_{\mathrm{geo}}^{\le N}(V_t)\) intertwines
\(\alpha_{t,s}\) with the geometric transport \(\beta_{t,s}\), which
preserves tensor degree and the \(J\)-adic filtration.  For
\(1\leq n\leq N\), put
\begin{equation}
 x_{s,t}^{(n)}:=[\log_\star\chi_t(\Phi_{s,t})]_n,
 \label{cc:eq-charted-degree-n-logarithm}
\end{equation}
If \(x_{s,t}^{(n)}\ne0\) and
\(r_{s,t}^{(n)}:=\nu_{F,t}(x_{s,t}^{(n)})\geq1\), its full causal normal
symbol is
\begin{equation}
 \sigma_n^{\chi,\mathrm{full}}(\Phi_{s,t})
 :=\operatorname{in}^{\mathrm{full}}_{F,t}(x_{s,t}^{(n)})
 \in G^{\mathrm{full}}_{n,r_{s,t}^{(n)}}(t).
 \label{cc:eq-registered-normal-symbol}
\end{equation}

Relative to the chart, this initial symbol is canonical.  Valuation zero
records the abelianized component rather than positive ancestry.  For a
causal normal geometry \(\mathfrak X\) we abbreviate the same object by
\begin{equation}
 \sigma_{n,r}^{\mathfrak X}(s,t)
 :=\operatorname{in}^{\mathrm{full}}_{F,t}
   \bigl([X^{\mathfrak X}_{s,t}]_n\bigr),
 \qquad r=\nu_{F,t}([X^{\mathfrak X}_{s,t}]_n).
 \label{cc:eq-causal-normal-symbol}
\end{equation}

\begin{corollary}[Finite jets and naturality]
\label{cc:cor-faithful-jets}
\label{cc:prop-naturality}
For every \(M\geq0\), the associative realization induces an injection
\[
 L/\gamma_{M+1}(D)\lhook\joinrel\longrightarrow T(V)/J^{M+1}.
\]
Its successive normal quotients are \(L/D\) in order zero and
\(\gamma_r(D)/\gamma_{r+1}(D)\) in orders \(1\leq r\leq M\).
Every linear map \(A:V\to W\) induces compatible maps
\begin{equation}
 Q(A):Q(V)\longrightarrow Q(W),
 \qquad
 Q(A)(\rho_vq)=\rho_{Av}Q(A)(q),
 \label{cc:eq-Q-naturality}
\end{equation}
and hence a graded Hopf homomorphism
\begin{equation}
 \mathfrak N(A):
 T(Q(V))\#_\rho S(V)\longrightarrow T(Q(W))\#_\rho S(W).
 \label{cc:eq-normal-functor}
\end{equation}
Under \eqref{cc:eq-normal-cone}, this is the normal-cone map induced
by \(T(A)\).  Hence \(G^{\mathrm{full}}_{n,r}\) is a homogeneous
degree-\(n\) polynomial functor, and \(E^{\mathrm{full}}_{n,r}\) is its
contravariant algebraic dual.  These assertions remain algebraically
valid for infinite-dimensional \(V\).
\end{corollary}

\begin{proof}
The kernel of the displayed associative realization is
\(L\cap J^{M+1}=\gamma_{M+1}(D)\), proving the first assertion.
Functoriality of the free Lie algebra sends \(D(V)\) into \(D(W)\) and
\([D(V),D(V)]\) into \([D(W),D(W)]\), so it descends to \(Q(A)\).
The identity
\[
 L(A)[v,d]=[Av,L(A)d]
\]
gives \eqref{cc:eq-Q-naturality}.  Therefore the images of \(V\) and
\(Q(V)\) satisfy the two relations \eqref{cc:eq-smash-relations} in the
target, producing \eqref{cc:eq-normal-functor}.  Both this map and the
map induced on \({\operatorname{gr}}_JT(V)\) agree on the primitive generators \(V\) and
\(Q(V)\), hence everywhere.  Taking normal degree \(r\) and ordinary
degree \(n\) gives the induced natural transformation.  Scalar dilation
by \(c\) acts in tensor degree \(n\) by \(c^n\), and the construction
takes subquotients of \(V^{\otimes n}\).
\end{proof}

\begin{remark}[Algebraic scope]
\label{cc:rem-algebraic-scope}
The theorem is purely algebraic.  It allows arbitrary \(V\) because
\[
 (L\cap J^r)_n=0\qquad(2r>n).
\]
Thus the filtration has finite length in every tensor degree.  Completed
tensor algebras and continuity of \(\rho\) require separate topological
hypotheses.
\end{remark}

For the rest of the section we use the extended valuation
\begin{equation}
 \nu_F(0):=+\infty,
 \qquad
 \nu_F(x):=\max\{r:x\in F_rL\}\quad(x\ne0).
 \label{cc:eq-extended-F-valuation}
\end{equation}

\begin{proposition}[Filtered BCH and the initial ancestry symbol]
\label{cc:prop-filtered-bch}
\label{cc:cor-causal-special-fibre}
\label{cc:cor-bch-cobracket}
Let \(A:V\to V\) be linear and write \(L(A)\) for its functorial action
on \(L\).  Then
\[
 L(A)\gamma_r(D)\subseteq\gamma_r(D),
 \qquad
 [\gamma_r(D),\gamma_s(D)]\subseteq\gamma_{r+s}(D).
\]
Consequently every truncation in ordinary tensor degree is nilpotent and
filtered for
\(F_0L=L\), \(F_rL=L\cap J^r=\gamma_r(D)\).  It therefore defines BCH
laws on the ordinary-degree completions of \({\operatorname{gr}}_FL\) and of the extended
Rees Lie algebra
\begin{equation}
 \mathcal R_F(L):=
 \sum_{r\geq0}F_rL\,z^{-r}\mathbf k[z]
 \subseteq L[z,z^{-1}],
 \qquad
 \mathcal R_F(L)/z\mathcal R_F(L)\cong{\operatorname{gr}}_FL,
 \label{cc:eq-rees-lie}
\end{equation}
Reduction modulo \(z\) records the associated-graded deformation.  If
\(x\in F_rL\setminus F_{r+1}L\),
its full initial symbol is the nonzero class
\begin{equation}
 \operatorname{in}^{\mathrm{full}}_F(x)
 :=x+F_{r+1}L\in{\mathfrak L}_r(Q).
 \label{cc:eq-initial-symbol}
\end{equation}
Thus filtration order cannot decrease under linear transport (with the
value \(+\infty\) if the transported element vanishes), and
\begin{equation}
 \nu_F([x,y])\geq\nu_F(x)+\nu_F(y)
 \label{cc:eq-valuation-superadditive}
\end{equation}
for nonzero \(x,y\), with the convention \(\nu_F(0)=+\infty\).  Exact
valuation is preserved
by transport precisely when the induced associated-graded map does not kill
the initial symbol; injectivity of that map on the containing graded layer
is sufficient.  Likewise,
equality holds in \eqref{cc:eq-valuation-superadditive} precisely when
the associated-graded bracket of the two initial symbols is nonzero.
Otherwise transport,
bracketing, cancellation in a sum, or a BCH component may raise the order
or annihilate the element.  A full Rees lift of a causal logarithm requires
a Rees jet or a splitting and is not supplied canonically by the filtration.

In the ordinary-degree completion, exponentiation identifies the special
fibre with
\begin{equation}
 \exp\!\left(V_{\mathrm{ab}}\ltimes_\rho\widehat{{\mathfrak L}}(Q)\right).
 \label{cc:eq-special-fibre-group}
\end{equation}
Here
\(\widehat{{\mathfrak L}}(Q)=\prod_{n\geq1}{\mathfrak L}(Q)_n\) denotes completion in
ordinary tensor degree.
Moreover, the associated-graded bracket induces
\begin{equation}
 \mathfrak b_{a,b}:G^{\mathrm{full}}_{n,a}(V)\otimes
 G^{\mathrm{full}}_{m,b}(V)
 \longrightarrow G^{\mathrm{full}}_{n+m,a+b}(V),
 \qquad \mathfrak b_{a,b}([x],[y])=[x,y],
 \label{cc:eq-ancestry-bracket}
\end{equation}
and this is the mixed quadratic symbol of BCH:
\begin{equation}
 \left.\partial_s\partial_t\right|_{s=t=0}
 {\operatorname{BCH}}(sx,ty)=\tfrac12[x,y].
 \label{cc:eq-bch-mixed-hessian}
\end{equation}
If the three layers in \eqref{cc:eq-ancestry-bracket} are finite
dimensional, duality gives
\begin{equation}
 \mathfrak b_{a,b}^{\vee}:E^{\mathrm{full}}_{n+m,a+b}(V)
 \longrightarrow E^{\mathrm{full}}_{n,a}(V)\otimes
 E^{\mathrm{full}}_{m,b}(V).
 \label{cc:eq-ancestry-cobracket}
\end{equation}
\end{proposition}

\begin{proof}
Functoriality sends \(D\) into \(D\), and induction sends each term of its
lower central series into itself.  The bracket inclusion is
\eqref{cc:eq-strong-filtration}.  Every homogeneous BCH term is an
iterated Lie bracket, so its filtration degree is at least the sum of the
positive degrees of its entries.  This proves closure of
\eqref{cc:eq-rees-lie} and filteredness of BCH\@.  The coefficient of
\(z^{-r}\) in \(\mathcal R_F(L)\) is \(F_rL\), whereas its coefficient
in \(z\mathcal R_F(L)\) is \(F_{r+1}L\); hence the stated special-fibre
isomorphism follows degree by degree.  The description of the
initial symbol follows from \eqref{cc:eq-primitive-layers}.  Naturality
of the normal-cone isomorphism says that its order-\(r\) class is sent to
the induced associated-graded image.  If that image is nonzero, the
valuation remains \(r\); if it vanishes, the image lies in
\(F_{r+1}L\) (or is zero), so the valuation increases rather than
decreases.  Finally, if \(a=\nu_F(x)\) and \(b=\nu_F(y)\), then
\([x,y]\in F_{a+b}L\), which proves
\eqref{cc:eq-valuation-superadditive}.  Its class in
\(F_{a+b}L/F_{a+b+1}L\) is exactly
\([\operatorname{in}_F^{\mathrm{full}}(x),
  \operatorname{in}_F^{\mathrm{full}}(y)]\).  This class is nonzero if
and only if the valuation of \([x,y]\) is exactly \(a+b\); if the bracket
itself vanishes, both statements are read using \(\nu_F(0)=+\infty\).
Applying the degreewise exponential to
\(\mathcal R_F(L)/z\mathcal R_F(L)\cong{\operatorname{gr}}_FL\) and using
\eqref{cc:eq-graded-lie} gives \eqref{cc:eq-special-fibre-group}.
Well-definedness of \(\mathfrak b_{a,b}\) is exactly
\([F_aL,F_bL]\subseteq F_{a+b}L\), together with
\([F_{a+1}L,F_bL]+[F_aL,F_{b+1}L]\subseteq F_{a+b+1}L\).
The quadratic term of BCH is \(\frac12[x,y]\), which proves
\eqref{cc:eq-bch-mixed-hessian}; the canonical identification
\((G\otimes H)^\vee\cong G^\vee\otimes H^\vee\) in finite dimension gives
\eqref{cc:eq-ancestry-cobracket}.
\end{proof}

Thus no completion enters the normal-cone theorem; ordinary-degree
completion is used only to sum BCH.

The first normal layer \(Q\) admits an intrinsic presentation that contains
both the curvature and its Bianchi relation.  It is the Koszul model for
the infinitesimal Alexander invariant of a free Lie algebra; compare
\cite[Thm.~6.2]{PapadimaSuciu04}.  Set \(\mathsf S:=S(V)\), give
both polynomial and exterior generators degree one, and define
\begin{align}
 d_3(f\otimes u\wedge v\wedge w)
 &:=fu\otimes v\wedge w-fv\otimes u\wedge w
      +fw\otimes u\wedge v,                                    
 \label{cc:eq-d3}\\
 d_2(f\otimes u\wedge v)&:=fu\otimes v-fv\otimes u.
 \label{cc:eq-d2}
\end{align}

\begin{theorem}[Universal Bianchi--Koszul module]
\label{cc:thm-bianchi-module}
The map
\[
 \kappa:\mathsf S\otimes\Lambda^2V\longrightarrow Q,
 \qquad
 \kappa(f\otimes u\wedge v)=f\cdot\overline{[u,v]},
\]
induces natural graded \(\mathsf S\)-module isomorphisms
\begin{equation}
 Q\cong\operatorname{coker}d_3
 \cong\ker\bigl(\mathsf S\otimes V\xrightarrow{\mu}\mathsf S\bigr),
 \qquad \mu(f\otimes v)=fv.
 \label{cc:eq-koszul-model}
\end{equation}
Equivalently, there are exact sequences
\begin{equation}
 \mathsf S\otimes\Lambda^3V\xrightarrow{d_3}
 \mathsf S\otimes\Lambda^2V\xrightarrow{\kappa}Q\longrightarrow0
 \label{cc:eq-bianchi-presentation}
\end{equation}
Writing \(\bar d_2\) for the map induced by \(d_2\) on
\(\operatorname{coker}d_3\cong Q\), one also has
\begin{equation}
 0\longrightarrow Q\xrightarrow{\bar d_2}\mathsf S\otimes V
 \xrightarrow{\mu}\mathsf S\xrightarrow{\varepsilon}\mathbf k
 \longrightarrow0.
 \label{cc:eq-first-syzygy}
\end{equation}
For every \(\mathsf S\)-module \(E\), viewed as a
\(V_{\mathrm{ab}}\)-module through its commuting multiplication
operators, there is a natural isomorphism
\begin{equation}
 {\operatorname{Hom}}_{\mathsf S}(Q,E)\cong Z^2_{\mathrm{CE}}(V_{\mathrm{ab}};E).
 \label{cc:eq-represents-cocycles}
\end{equation}
If \(V\) is finite dimensional, then, as a polynomial functor over
\(\mathbf k\), for \(n\geq2\),
\begin{equation}
 Q_n\cong\mathbf S_{(n-1,1)}V.
 \label{cc:eq-hook}
\end{equation}
\end{theorem}

\begin{proof}
Let \(B:=\operatorname{coker}d_3\), and denote by \(c(u,v)\) the class
of \(1\otimes u\wedge v\).  Its defining relation is
\begin{equation}
 u\cdot c(v,w)-v\cdot c(u,w)+w\cdot c(u,v)=0.
 \label{cc:eq-bianchi-relation}
\end{equation}
Define on \(V\oplus B\)
\[
 [(u,a),(v,b)]:=(0,c(u,v)+u\cdot b-v\cdot a).
\]
Jacobi for three vectors in \(V\) is
\eqref{cc:eq-bianchi-relation}; Jacobi for two vectors and one element
of \(B\) is commutativity of the \(\mathsf S\)-action, and all remaining
cases are immediate.  If \(f:V\to\mathfrak g\) and
\(\mathfrak g''=0\), then
\([{\operatorname{ad}}_{f(u)},{\operatorname{ad}}_{f(v)}]={\operatorname{ad}}_{[f(u),f(v)]}\) vanishes on
\(\mathfrak g'\).  Hence \(\mathfrak g'\) is an \(\mathsf S\)-module,
and \(c_f(u,v)=[f(u),f(v)]\) satisfies
\eqref{cc:eq-bianchi-relation}.  It induces a unique
\(\mathsf S\)-linear map \(\bar f:B\to\mathfrak g'\), and
\((u,b)\mapsto f(u)+\bar f(b)\) is the unique Lie map extending \(f\).
Thus \(V\oplus B\) is the free metabelian Lie algebra, and
\(L/[D,D]\cong V\oplus B\).  The derived ideal of \(V\oplus B\) is all
of \(B\), because \(B\) is generated as an \(\mathsf S\)-module by the
elements \(c(u,v)\); comparison of derived ideals yields \(Q\cong B\).

The Koszul complex
\[
 \mathsf S\otimes\Lambda^3V\xrightarrow{d_3}
 \mathsf S\otimes\Lambda^2V\xrightarrow{d_2}
 \mathsf S\otimes V\xrightarrow{\mu}\mathsf S
 \xrightarrow{\varepsilon}\mathbf k\longrightarrow0
\]
is exact.  In infinite dimension this follows from the finite-dimensional
case because every cycle belongs to the complex of a finite-dimensional
subspace.  This proves \eqref{cc:eq-koszul-model} and
\eqref{cc:eq-first-syzygy}.
An \(\mathsf S\)-linear
map from \(\mathsf S\otimes\Lambda^2V\) to \(E\) is an alternating map
\(b:\Lambda^2V\to E\), and it kills \(\operatorname{im}d_3\) exactly
when \(b\) satisfies the abelian Chevalley--Eilenberg cocycle equation.
This proves \eqref{cc:eq-represents-cocycles}.  Finally, in degree
\(n\), \eqref{cc:eq-first-syzygy} reads
\[
 Q_n=\ker\bigl({\operatorname{Sym}}^{n-1}V\otimes V\to{\operatorname{Sym}}^nV\bigr).
\]
Pieri's rule \cite[\S6.1 and Ch.~8]{Fulton97} decomposes the source as
\(\mathbf S_{(n)}V\oplus\mathbf S_{(n-1,1)}V\), with Schur functors of
excess row length understood to vanish.  Multiplication is the projection
onto \(\mathbf S_{(n)}V\), so its kernel is
\(\mathbf S_{(n-1,1)}V\).  The construction and the corresponding Young
symmetrizer are defined over every characteristic-zero field, proving
\eqref{cc:eq-hook} over \(\mathbf k\).
\end{proof}

Let \(s:V\hookrightarrow L\) be the generator section and define
\begin{equation}
 \nabla_ud:=[s(u),d],\qquad
 \Omega(u,v):=[s(u),s(v)]\in D.
 \label{cc:eq-connection-curvature}
\end{equation}

\begin{corollary}[Curvature raises ancestry]
\label{cc:prop-curvature}
The exact identities
\begin{align}
 [\nabla_u,\nabla_v]&={\operatorname{ad}}_{\Omega(u,v)}|_D,
 \label{cc:eq-curvature-commutator}\\
 \nabla_u\Omega(v,w)-\nabla_v\Omega(u,w)+\nabla_w\Omega(u,v)&=0
 \label{cc:eq-curvature-bianchi}
\end{align}
hold.  Moreover,
\(D\) is the smallest Lie subalgebra of \(L\) which contains
\(\Omega(\Lambda^2V)\) and is stable under every \(\nabla_u\), and
\begin{equation}
 \nabla_u\gamma_r(D)\subseteq\gamma_r(D),\qquad
 [\nabla_u,\nabla_v]\gamma_r(D)\subseteq\gamma_{r+1}(D),
 \label{cc:eq-curvature-raises}
\end{equation}
and curvature has the well-defined transverse symbol
\begin{equation}
 \mathcal R^{(r)}_{u,v}:
 \gamma_r(D)/\gamma_{r+1}(D)\longrightarrow
 \gamma_{r+1}(D)/\gamma_{r+2}(D),
 \qquad[d]\longmapsto[\Omega(u,v),d].
 \label{cc:eq-curvature-symbol}
\end{equation}
The induced \(V\)-actions on each fixed ancestry layer commute, and the
degree-one curvature class is
\(\kappa(1\otimes u\wedge v)\in Q\).
\end{corollary}

\begin{proof}
Equations \eqref{cc:eq-curvature-commutator} and
\eqref{cc:eq-curvature-bianchi} are the two corresponding forms of
Jacobi.  Let \(K\) be the Lie subalgebra generated from
\(\Omega(\Lambda^2V)\) by the operators \(\nabla_u\).  Stability under
all \({\operatorname{ad}}_{s(u)}\), together with
\([{\operatorname{ad}}_x,{\operatorname{ad}}_y]={\operatorname{ad}}_{[x,y]}\), implies stability under \({\operatorname{ad}}_x\) for
every \(x\in L\), because \(s(V)\) generates \(L\).  Thus \(K\) is a
Lie ideal of \(L\).  It contains every \([s(u),s(v)]\), so it contains
the derived ideal \(D\); conversely its generators and all closure
operations remain inside the ideal \(D\).  Hence \(K=D\), proving the
minimality assertion.

Since every \(\gamma_r(D)\) is stable under \({\operatorname{ad}}_L\), the first
inclusion in \eqref{cc:eq-curvature-raises} holds.  Explicitly, the
case \(r=1\) uses that \(D\) is an ideal, and if it holds at level \(r\),
then Jacobi gives
\[
 [L,[D,\gamma_r(D)]]
 \subseteq[[L,D],\gamma_r(D)]+[D,[L,\gamma_r(D)]]
 \subseteq[D,\gamma_r(D)].
\]
This proves the induction.  The second inclusion follows
from \eqref{cc:eq-curvature-commutator}, because
\(\Omega(u,v)\in D\).  Finally,
\([D,\gamma_{r+1}(D)]\subseteq\gamma_{r+2}(D)\), so
\eqref{cc:eq-curvature-symbol} is independent of the representative.
The remaining assertions follow on passing to the associated graded.
\end{proof}

Thus each fixed ancestry layer is flat, while curvature survives as a
transverse map to the next layer.  The whole positive tower is the free
Lie algebra generated by the Bianchi module \(Q\).  Notice that
\eqref{cc:eq-represents-cocycles} represents actual cocycles
\(Z^2_{\mathrm{CE}}\), not cohomology classes: the canonical generator
section fixes the curvature tensor before any change-of-section quotient.

\medskip
\noindent\textit{Finite-dimensional characters.}

For the remainder of this section, assume that \(\mathbf k=\mathbb C\) and
\(1\leq d:=\dim V<\infty\).  All character identities are taken in the
\((t,u)\)-adic completion of the \(\lambda\)-ring of polynomial
\(\mathrm{GL}(V)\)-representations.  Let \(\psi_k\) denote its \(k\)-th
Adams operation, extended by \(\psi_k(t)=t^k\), and put
\begin{equation}
 \mathcal Q(t):=\sum_{n\geq2}\operatorname{ch}(Q_n)t^n
 =\sum_{n\geq2}s_{(n-1,1)}t^n.
 \label{cc:eq-Q-character}
\end{equation}

\begin{corollary}[Witt ancestry character]
\label{cc:cor-witt-character}
For the primal full layers \(G^{\mathrm{full}}_{n,r}(V)\) of
\eqref{cc:eq-exact-layer}, one has
\begin{align}
 \mathcal A(t,u)
 &:=\sum_{r\geq1}\sum_{n\geq2}
 \operatorname{ch}\!\left(G^{\mathrm{full}}_{n,r}(V)\right)t^nu^r
 \notag\\
 &=\sum_{k\geq1}\frac{\mu(k)}{k}
 \log\!\left(\frac{1}{1-u^k\psi_k(\mathcal Q(t))}\right),
 \label{cc:eq-plethystic-witt}\\
 \sum_n\operatorname{ch}\!\left(G^{\mathrm{full}}_{n,r}(V)\right)t^n
 &=\frac1r\sum_{k\mid r}\mu(k)
 \psi_k(\mathcal Q(t))^{r/k}.
 \label{cc:eq-layer-witt}
\end{align}
Moreover,
\begin{equation}
 q_d(t):=\sum_{n\geq2}(\dim Q_n)t^n
 =1-\frac{1-dt}{(1-t)^d},
 \label{cc:eq-Q-hilbert}
\end{equation}
and
\begin{equation}
 \sum_{r\geq1,n\geq2}b_{n,r}(d)t^nu^r
 =\sum_{k\geq1}\frac{\mu(k)}{k}
 \log\!\left(\frac{1}{1-u^kq_d(t^k)}\right),
 \label{cc:eq-dimension-witt}
\end{equation}
where \(b_{n,r}(d)=\dim G^{\mathrm{full}}_{n,r}(V)\).
Equivalently, for each fixed ancestry depth,
\begin{equation}
 \sum_n b_{n,r}(d)t^n
 =\frac1r\sum_{k\mid r}\mu(k)
 q_d(t^k)^{r/k}.
 \label{cc:eq-fixed-depth-dimension}
\end{equation}
Moreover,
\begin{equation}
 \dim Q_n
 =\frac{(d-1)(n-1)}{n}\binom{d+n-2}{n-1},
 \qquad n\geq2,
 \label{cc:eq-Q-dimension}
\end{equation}
and \(G^{\mathrm{full}}_{n,r}(V)=0\) whenever \(2r>n\).
\end{corollary}

\begin{proof}
By \eqref{cc:eq-primitive-layers}, exact ancestry \(r\) is the
length-\(r\) component of the free Lie algebra on the graded module \(Q\).
The multigraded Witt--Brandt formula
\cite[\S0.4.2 and \S8.2]{Reutenauer93} states in a characteristic-zero
\(\lambda\)-ring that
\[
 \operatorname{ch}({\mathfrak L}_r(Q))
 =\frac1r\sum_{k\mid r}\mu(k)
   \psi_k(\operatorname{ch}Q)^{r/k}.
\]
Applying it to the internally graded object \(Q=\bigoplus_{n\ge2}Q_n\)
gives \eqref{cc:eq-layer-witt}.  With \(r=km\), summing over \(r\)
gives
\[
 \sum_{k\ge1}\frac{\mu(k)}{k}
 \sum_{m\ge1}\frac{(u^k\psi_k(\mathcal Q(t)))^m}{m},
\]
which is \eqref{cc:eq-plethystic-witt}.  Taking dimensions yields
\eqref{cc:eq-dimension-witt} and
\eqref{cc:eq-fixed-depth-dimension}.  Finally,
\eqref{cc:eq-first-syzygy} gives, with generators placed in internal
degree one,
\[
 q_d(t)=\frac{dt}{(1-t)^d}-\frac1{(1-t)^d}+1,
\]
which is \eqref{cc:eq-Q-hilbert}.  Its degree-\(n\) coefficient is
\[
 d\binom{d+n-2}{n-1}-\binom{d+n-1}{n}
 =\frac{(d-1)(n-1)}{n}\binom{d+n-2}{n-1},
\]
proving \eqref{cc:eq-Q-dimension}.  Since every
generator of \(Q\) has internal degree at least two, a length-\(r\) Lie
word has internal degree at least \(2r\), proving the last assertion.
\end{proof}

For a partition \(\lambda\vdash n\) with \(\ell(\lambda)\leq d\), define
the sector ancestry polynomial
\begin{equation}
 A_{\lambda,n}(u):=
 \sum_{r\geq1}\dim{\operatorname{Hom}}_{\mathrm{GL}(V)}
 \bigl(\mathbf S_\lambda V,G^{\mathrm{full}}_{n,r}(V)\bigr)u^r.
 \label{cc:eq-ancestry-polynomial}
\end{equation}
Its degree records maximal visible ancestry, while its coefficients retain
the multiplicities at every depth.  They stabilize once
\(d\geq\ell(\lambda)\); the uniform condition \(d\geq n\) makes every
partition of \(n\) visible.

\begin{example}[Multiplicity and mixed ancestry]
\label{cc:ex-mixed-ancestry}
In total degree six, the sector \(\mathbf S_{(3,2,1)}V\) occurs twice in
ancestry two and once in ancestry three.  Equivalently,
\begin{equation}
 A_{(3,2,1),6}(u)=2u^2+u^3.
 \label{cc:eq-mixed-polynomial}
\end{equation}
Indeed, \(Q_2=\mathbf S_{(1,1)}V\),
\(Q_3=\mathbf S_{(2,1)}V\), and
\(Q_4=\mathbf S_{(3,1)}V\), so
\[
 ({\mathfrak L}_2(Q))_6=(Q_2\otimes Q_4)\oplus\Lambda^2Q_3,
 \qquad
 ({\mathfrak L}_3(Q))_6={\mathfrak L}_3(Q_2),
\]
and hence, at the character level,
\begin{align*}
 \operatorname{ch}G^{\mathrm{full}}_{6,2}
 &=s_{(1,1)}s_{(3,1)}
   +\tfrac12\bigl(s_{(2,1)}^2-\psi_2(s_{(2,1)})\bigr),\\
 \operatorname{ch}G^{\mathrm{full}}_{6,3}
 &=\tfrac13\bigl(s_{(1,1)}^3-\psi_3(s_{(1,1)})\bigr).
\end{align*}
The Littlewood--Richardson rule \cite[Chs.~5--6 and 8]{Fulton97} and the two
displayed Witt operations give
\begin{align*}
 G^{\mathrm{full}}_{6,2}(V)
 &\cong \mathbf S_{(4,2)}V
 \oplus2\mathbf S_{(4,1,1)}V
 \oplus\mathbf S_{(3,3)}V
 \oplus2\mathbf S_{(3,2,1)}V\\
 &\quad\oplus\mathbf S_{(3,1,1,1)}V
 \oplus\mathbf S_{(2,2,1,1)}V,\\
 G^{\mathrm{full}}_{6,3}(V)
 &\cong\mathbf S_{(3,2,1)}V
 \oplus\mathbf S_{(2,2,1,1)}V
 \oplus\mathbf S_{(2,1,1,1,1)}V.
\end{align*}
These are identities of polynomial functors (partitions with too many rows
evaluate to zero).  Thus the relevant multiplicities are two and one
whenever \(d\geq3\); all displayed summands are nonzero when \(d\geq5\).
Since every generator
of \(Q\) has internal degree at least two, no ancestry greater than three
occurs in degree six.
Thus even inside one irreducible symmetry sector, maximal depth is not a
substitute for the full filtered multiplicity space.
\end{example}

For each tensor degree \(n\), define the finite ancestry package
\begin{equation}
 \mathcal G_n^{\mathrm{full}}(V):=
 \bigoplus_{1\leq r\leq\lfloor n/2\rfloor}
 G^{\mathrm{full}}_{n,r}(V),
 \qquad
 \mathcal N_n|_{G^{\mathrm{full}}_{n,r}(V)}:=r\,\mathrm{Id},
 \label{cc:eq-finite-ancestry-package}
\end{equation}
and its algebraic dual
\begin{equation}
 \mathcal E_n^{\mathrm{full}}(V):=
 \bigl(\mathcal G_n^{\mathrm{full}}(V)\bigr)^{\vee}
 =\bigoplus_{r=1}^{\lfloor n/2\rfloor}
 E^{\mathrm{full}}_{n,r}(V).
 \label{cc:eq-finite-readout-package}
\end{equation}
The vanishing for \(2r>n\) makes these sums finite.  In the
\(\mathbf S_\lambda\)-sector, put
\begin{equation}
 \mathcal M_{\lambda,n}:=
 \bigoplus_{r=1}^{\lfloor n/2\rfloor}
 {\operatorname{Hom}}_{\mathrm{GL}(V)}
 \bigl(\mathbf S_\lambda V,G^{\mathrm{full}}_{n,r}(V)\bigr),
 \label{cc:eq-sector-multiplicity-space}
\end{equation}
and the ancestry polynomial is its marked trace:
\begin{equation}
 A_{\lambda,n}(u)
 =\operatorname{Tr}_{\mathcal M_{\lambda,n}}(u^{\mathcal N_n}),
 \label{cc:eq-ancestry-marked-trace}
\end{equation}
where \(\mathcal N_n\) acts by \(r\) on the \(r\)-th summand.  Thus
\eqref{cc:eq-ancestry-marked-trace} follows directly from
\eqref{cc:eq-ancestry-polynomial}.  The normal grading makes
\((\mathcal G_n^{\mathrm{full}},\mathcal N_n)\) intrinsic, while any
finite analytic registration is an additional quotient of these full
layers.
\begin{theorem}[Real Schur exact-layer projectors]
\label{cc:thm-schur-exact-projectors}
Fix \(n\geq2\).  There is a system of stable natural transformations
\[
 P_{\lambda,r}:L_n(H)\longrightarrow L_n(H),
 \qquad
 \lambda\vdash n,\quad
 1\leq r\leq\lfloor n/2\rfloor,
\]
defined for real vector spaces \(H\), with the following properties.
\begin{align}
 P_{\lambda,r}P_{\mu,s}
 &=\mathbf1_{\{(\lambda,r)=(\mu,s)\}}P_{\lambda,r},
 &
 \sum_{\lambda,r}P_{\lambda,r}
 &=I_{L_n},
 \label{cc:eq-exact-projector-orthogonality}\\
 \operatorname{Ran}P_{\lambda,r}
 &\subseteq F_rL_n,
 &
 P_{\lambda,r}(F_{r+1}L_n)&=0.
 \label{cc:eq-exact-projector-filtration}
\end{align}
The quotient \(F_rL_n\to G^{\mathrm{full}}_{n,r}\) restricts to an
isomorphism from \(\operatorname{Ran}P_{\lambda,r}\) onto the
\(\lambda\)-isotypic part of the exact layer.  If \(e_n\) is a normalized
real Dynkin--Specht--Wever idempotent, each projector has the form
\begin{equation}
 a_{\lambda,r}
 :=\sum_{\sigma\in\mathfrak S_n}c_{\lambda,r}(\sigma)\sigma
 \in e_n\mathbb R[\mathfrak S_n]e_n,
 \qquad
 P_{\lambda,r}
 =\left.\rho(a_{\lambda,r})\right|_{L_n(H)},
 \qquad c_{\lambda,r}(\sigma)\in\mathbb R,
 \label{cc:eq-real-permutation-projector}
\end{equation}
where \(\rho(\sigma)\) permutes tensor legs.  Thus it belongs to a finite
permutation algebra and uses at most \(n!\) permutation terms.
For every bounded linear map \(A:H\to K\),
\begin{equation}
 L_n(A)P_{\lambda,r}^{H}
 =P_{\lambda,r}^{K}L_n(A).
 \label{cc:eq-projector-bounded-naturality}
\end{equation}
Moreover, on every fixed-degree all-cut multi-flattening grade with
leg-permutation-invariant norm,
\begin{equation}
 \|P_{\lambda,r}X\|_{\mathrm{mf},g}
 \leq C_{\lambda,r}^{(n)}\|X\|_{\mathrm{mf},g},
 \qquad
 C_{\lambda,r}^{(n)}
 :=\sum_{\sigma\in\mathfrak S_n}|c_{\lambda,r}(\sigma)|.
 \label{cc:eq-projector-multiflat-bound}
\end{equation}
The constant is finite and independent of \(\dim H\).
For a leg-permutation-invariant quasi-norm, the same conclusion holds
with \(C_{\lambda,r}^{(n)}\) multiplied by the finite quasi-triangle
constant for at most \(n!\) summands.

After complexification,
\begin{equation}
 P_{\lambda,r}^{\mathbb C}(\overline z)
 =\overline{P_{\lambda,r}^{\mathbb C}(z)},
 \label{cc:eq-projector-real-form}
\end{equation}
and, for \(d=\dim_{\mathbb R}H\geq n\), one has the stable decomposition
\begin{equation}
 L_n(H_{\mathbb C})
 \cong
 \bigoplus_{\lambda\vdash n}
 \mathbf S_\lambda H_{\mathbb C}\otimes M_{\lambda,n},
 \qquad
 F_rL_n(H_{\mathbb C})
 \cong
 \bigoplus_{\lambda\vdash n}
 \mathbf S_\lambda H_{\mathbb C}\otimes M^{\geq r}_{\lambda,n},
 \label{cc:eq-schur-ancestry-filtration}
\end{equation}
with real-form-compatible splittings
\begin{equation}
 M^{\geq r}_{\lambda,n}
 =M_{\lambda,n,r}\oplus M^{\geq r+1}_{\lambda,n},
 \qquad
 G^{\mathrm{full}}_{n,r}(H_{\mathbb C})
 \cong\bigoplus_{\lambda\vdash n}
 \mathbf S_\lambda H_{\mathbb C}\otimes M_{\lambda,n,r}.
 \label{cc:eq-schur-exact-layers}
\end{equation}
Here \(P_{\lambda,r}^{\mathbb C}\) is the identity on the displayed
\(\mathbf S_\lambda H_{\mathbb C}\otimes M_{\lambda,n,r}\) summand and
zero on the others.  For a fixed partition \(\lambda\), the same
description is already valid whenever \(d\geq\ell(\lambda)\), with
Schur functors having more than \(d\) rows interpreted as zero.
\end{theorem}

\begin{proof}
Multilinearization identifies the degree-\(n\) free-Lie functor with its
finite coefficient module \(\operatorname{Lie}(n)\), a real
\(\mathfrak S_n\)-module.  Because the filtration \(F_rL_n\) is natural
and homogeneous, its multilinearization is a decreasing filtration
\[
 \operatorname{Lie}(n)=\operatorname{Lie}(n)_{\geq1}
 \supseteq\operatorname{Lie}(n)_{\geq2}
 \supseteq\cdots
\]
by real \(\mathfrak S_n\)-submodules.  Maschke's theorem provides
\(\mathfrak S_n\)-invariant complements
\[
 \operatorname{Lie}(n)_{\geq r}
 =C_{n,r}\oplus\operatorname{Lie}(n)_{\geq r+1}.
\]
The real central idempotents of \(\mathbb R[\mathfrak S_n]\) split each
\(C_{n,r}\) into its Specht isotypic pieces \(C_{\lambda,n,r}\).
Projection onto \(C_{\lambda,n,r}\), along all the other chosen
summands, gives pairwise orthogonal equivariant idempotents whose sum is
the identity of \(\operatorname{Lie}(n)\).
Characteristic-zero polarization and depolarization carry these
coefficient-module idempotents back to homogeneous natural
transformations of \(L_n\).  They may be checked in the faithful range
\(\dim H\geq n\), and naturality then evaluates the same transformations
in every smaller dimension.

For completeness, these coefficient-module idempotents are finite
permutation polynomials.  A normalized real Dynkin--Specht--Wever
idempotent \(e_n\) realizes
\[
 L_n(H)=H^{\otimes n}e_n
\]
and \(\operatorname{Lie}(n)\) as a summand of the regular
\(\mathbb R[\mathfrak S_n]\)-module
\cite[Chs.~0--1]{Reutenauer93}.  The stable double-centralizer
identification represents every equivariant endomorphism of this summand
by the opposite action of the finite corner
\(e_n\mathbb R[\mathfrak S_n]e_n\).  The chosen idempotents therefore
give \eqref{cc:eq-real-permutation-projector}.  The range and kernel
relations in \eqref{cc:eq-exact-projector-filtration} are exactly the
chosen splitting of the coefficient filtration.

All complements and idempotents were chosen over \(\mathbb R\), so
complexification gives \eqref{cc:eq-projector-real-form}.  Schur--Weyl
duality identifies the Specht isotypic decomposition of the coefficient
module with
\eqref{cc:eq-schur-ancestry-filtration}--\eqref{cc:eq-schur-exact-layers};
see \cite[Chs.~4--6]{Fulton97}.  The group-algebra element in
\eqref{cc:eq-real-permutation-projector} is independent of \(H\), hence
commutes with every map \(A^{\otimes n}\).  It therefore defines the same
natural projectors in every dimension and proves
\eqref{cc:eq-projector-bounded-naturality}, including for bounded maps
between Hilbert completions.  Each leg permutation is an isometry of an
all-cut, permutation-invariant multi-flattening grade.  The triangle
inequality applied to \eqref{cc:eq-real-permutation-projector} proves
\eqref{cc:eq-projector-multiflat-bound}.  The range statements descend
to smaller dimensions; only the Schur summands with too many rows
vanish.  The faithful uniform stable range \(d\geq n\) displays every
partition of \(n\), while \(d\geq\ell(\lambda)\) displays the individual
\(\lambda\)-sector.
\end{proof}

\begin{remark}[Canonical data versus chosen lifts]
\label{cc:rem-projectors-noncanonical}
The spaces
\[
 M^{\geq r}_{\lambda,n}
 \quad\text{and}\quad
 M^{\geq r}_{\lambda,n}/M^{\geq r+1}_{\lambda,n}
\]
and hence their dimensions are intrinsic.  The complements
\(M_{\lambda,n,r}\) and the operators \(P_{\lambda,r}\) are auxiliary.
Changing them changes an exact-layer lift inside \(L_n(H)\), but not its
class in \(G^{\mathrm{full}}_{n,r}(H)\).  This distinction prevents a
chosen projector from being mistaken for a canonical homogeneous
component of a filtered logarithm.
\end{remark}

For a partition \(\lambda\vdash n\), let
\begin{equation}
 m_{\lambda,n,r}:=
 \dim\operatorname{Hom}_{\mathrm{GL}(H_{\mathbb C})}
 \!\left(\mathbf S_\lambda H_{\mathbb C},
 G^{\mathrm{full}}_{n,r}(H_{\mathbb C})\right),
 \qquad
 a_n(\lambda):=\max\{r:m_{\lambda,n,r}\neq0\},
 \label{cc:eq-schur-ancestry-multiplicity}
\end{equation}
where \(a_n(\lambda)=-\infty\) if the sector is absent.  The dimensions
are stable once \(d\geq\ell(\lambda)\).
Equivalently, for every nonzero sector,
\(a_n(\lambda)=\deg A_{\lambda,n}\), where
\(A_{\lambda,n}\) is the ancestry polynomial defined above.

\begin{corollary}[First-row ancestry bound]
\label{cc:cor-first-row-ancestry}
For every \(\lambda\vdash n\),
\begin{equation}
 m_{\lambda,n,r}=0
 \quad\text{if}\quad
 r>n-\lambda_1
 \quad\text{or}\quad
 2r>n.
 \label{cc:eq-first-row-layer-bound}
\end{equation}
Consequently,
\begin{equation}
 a_n(\lambda)\leq
 \min\!\left\{n-\lambda_1,\left\lfloor\frac n2\right\rfloor\right\}.
 \label{cc:eq-first-row-ancestry-bound}
\end{equation}
\end{corollary}

\begin{proof}
By \Cref{cc:thm-normal-cone,cc:thm-bianchi-module},
\[
 G^{\mathrm{full}}_{n,r}(H_{\mathbb C})
 \cong\bigl({\mathfrak L}_r(Q(H_{\mathbb C}))\bigr)_n,
 \qquad
 Q_m(H_{\mathbb C})\cong
 \mathbf S_{(m-1,1)}H_{\mathbb C}.
\]
The length-\(r\) free-Lie piece is a direct summand of
\(Q^{\otimes r}\) in characteristic zero.  Its internal degree-\(n\)
part is therefore a subrepresentation of the sum, over
\(m_1+\cdots+m_r=n\) and \(m_i\geq2\), of
\[
 \mathbf S_{(m_1-1,1)}H_{\mathbb C}\otimes\cdots\otimes
 \mathbf S_{(m_r-1,1)}H_{\mathbb C}.
\]
The Littlewood--Richardson first-row inequality gives
\[
 \lambda_1\leq\sum_{i=1}^r(m_i-1)=n-r
\]
for every Schur functor occurring in such a tensor product
\cite[Chs.~5--6]{Fulton97}.  Thus \(r\leq n-\lambda_1\).
The inequalities \(m_i\geq2\) also imply \(2r\leq n\).
\end{proof}

\begin{corollary}[Low degrees, long hooks, and unbounded depth]
\label{cc:cor-low-degree-long-hook}
Over \(\mathbb C\), with excess-row Schur functors understood to vanish,
the first exact layers are
\begin{align}
 G^{\mathrm{full}}_{3,1}(H_{\mathbb C})
 &\cong\mathbf S_{(2,1)}H_{\mathbb C},
 &
 G^{\mathrm{full}}_{3,r}(H_{\mathbb C})&=0\quad(r\geq2),
 \label{cc:eq-degree-three-ancestry}\\
 G^{\mathrm{full}}_{4,1}(H_{\mathbb C})
 &\cong\mathbf S_{(3,1)}H_{\mathbb C},
 &
 G^{\mathrm{full}}_{4,2}(H_{\mathbb C})
 &\cong\mathbf S_{(2,1,1)}H_{\mathbb C}.
 \label{cc:eq-degree-four-ancestry}
\end{align}
In the stable range \(\dim H\geq4\), the chosen real-form-compatible
projectors give the direct splitting
\begin{equation}
 L_4(H_{\mathbb C})
 =
 \operatorname{Ran}P_{(3,1),1}^{\mathbb C}
 \oplus
 \operatorname{Ran}P_{(2,1,1),2}^{\mathbb C},
 \label{cc:eq-degree-four-sector-split}
\end{equation}
so the \((3,1)\)-sector has exact ancestry one and the
\((2,1,1)\)-sector has exact ancestry two.

More generally, for every \(n\geq2\),
\begin{equation}
 G^{\mathrm{full}}_{n,1}(H_{\mathbb C})
 \cong\mathbf S_{(n-1,1)}H_{\mathbb C},
 \qquad
 m_{(n-1,1),n,r}=\mathbf1_{\{r=1\}}.
 \label{cc:eq-long-hook-depth-one}
\end{equation}
At the opposite, maximal-depth edge,
\begin{equation}
 G^{\mathrm{full}}_{2r,r}(H)
 \cong{\mathfrak L}_r(\Lambda^2H).
 \label{cc:eq-maximal-depth-layer}
\end{equation}
For every \(r\geq2\) this space is nonzero when \(\dim H\geq3\).
In the uniform stable range \(\dim H\geq2r\), there is consequently at
least one partition \(\lambda_r\vdash2r\) such that
\begin{equation}
 m_{\lambda_r,2r,r}>0,
 \qquad
 a_{2r}(\lambda_r)=r.
 \label{cc:eq-arbitrarily-deep-schur-sector}
\end{equation}
Thus the purely algebraic area genealogy has arbitrarily large finite
depth.
\end{corollary}

\begin{proof}
Exact ancestry one is \(Q_n\), so \Cref{cc:eq-hook} gives
\[
 G^{\mathrm{full}}_{3,1}\cong\mathbf S_{(2,1)},\qquad
 G^{\mathrm{full}}_{4,1}\cong\mathbf S_{(3,1)}.
\]
The bound \(2r\leq n\) eliminates higher depth in degree three.  In
degree four, the only remaining depth is two, and
\[
 G^{\mathrm{full}}_{4,2}
 =\bigl({\mathfrak L}_2(Q)\bigr)_4
 =\Lambda^2Q_2
 =\Lambda^2(\Lambda^2H_{\mathbb C})
 \cong\mathbf S_{(2,1,1)}H_{\mathbb C}.
\]
The last is the standard Littlewood--Richardson decomposition; it is an
identity of polynomial functors, so it also covers dimensions below the
stable range by the excess-row convention.  This proves
\eqref{cc:eq-degree-three-ancestry}--\eqref{cc:eq-degree-four-ancestry},
and \Cref{cc:thm-schur-exact-projectors} gives
\eqref{cc:eq-degree-four-sector-split}.

For general \(n\), exact ancestry one is again
\(Q_n\cong\mathbf S_{(n-1,1)}\).  The first-row bound
\eqref{cc:eq-first-row-layer-bound} excludes this same long-hook sector
from every depth \(r\geq2\).  It is visibly nonzero: for independent
\(x,y\in H\), the element
\[
 (\operatorname{ad}_x)^{n-2}[x,y]
\]
is a nonzero weight-\((n-1,1)\) vector and has exactly one area leaf.
This proves
\eqref{cc:eq-long-hook-depth-one}.

Finally, a length-\(r\) word in \(Q\) has internal degree at least
\(2r\).  Equality forces every free generator to lie in
\(Q_2=\Lambda^2H\), whence
\eqref{cc:eq-maximal-depth-layer}.  If \(\dim H\geq3\), then
\(\dim\Lambda^2H\geq3\), and the degree-\(r\) piece of the free Lie
algebra on \(\Lambda^2H\) is nonzero for every \(r\); for example, if
\(a,b\in\Lambda^2H\) are linearly independent, the Hall word
\((\operatorname{ad}_a)^{r-1}b\) is nonzero.  When
\(\dim H\geq2r\), Schur--Weyl decomposition displays every partition of
\(2r\).  Some \(\lambda_r\) must therefore occur in the nonzero module
\eqref{cc:eq-maximal-depth-layer}.  Since no degree-\(2r\) layer can
have ancestry greater than \(r\), its maximal ancestry is exactly \(r\).
\end{proof}

In summary, the normal-cone identity
\(L\cap J^r=\gamma_r(D)\) identifies ideal order, lower-central
ancestry, and the tagged area-leaf filtration.  The Bianchi quotient
\(Q=\operatorname{coker}d_3\) supplies the first intrinsic atom, while
\((\mathfrak L_r(Q))_n\) records its exact depth-\(r\),
degree-\(n\) descendants.  The real Schur projectors split these exact
layers functorially, the first-row estimate controls every symmetry
sector, and the low-degree and long-hook formulas exhibit both sharp
finite-depth sectors and descendants of arbitrarily large depth.  Thus
each fixed tensor degree has finite ancestry depth; completion is needed
for infinite causal logarithms, not for the algebraic normal cone.
\section{Critical analytic filtration and packet arithmetic}
\label{cr:sec-analytic}

\subsection{Spectral measures, asymptotics, and critical weak scales}
\label{cc:subsec:atom-rig}

For a compact operator \(T\) and an integer \(q\geq1\), put
\begin{equation}
 \|T\|_{\mathfrak W^{[q]}}
 :=\sup_{N\geq1}
 \frac{N s_N(T)}{(\log(e+N))^{q-1}},
 \qquad
 \mathfrak W^{[q]}_0
 :=\left\{T:\frac{Ns_N(T)}{(\log(e+N))^{q-1}}\to0\right\}.
 \label{cc:eq:critical-ideals}
\end{equation}
In particular,
\(\mathfrak W^{[1]}=\mathcal S_{1,\infty}\) and
\(\mathfrak W^{[1]}_0=\mathcal S^0_{1,\infty}\).
These are the weak Schatten scale and its separable part, with a logarithmic
weight when \(q>1\); see
\cite{Pietsch80,LordSukochevZanin12} for the underlying operator-ideal
framework.

\begin{definition}[Positive log-spectral data]
\label{cc:def:log-spectral-rig}
For compact \(T\), define
\begin{equation}
 \Lambda_T:=\sum_{j:s_j(T)>0}\delta_{-\log s_j(T)},
 \qquad
 \mathcal N_T(R):=\Lambda_T(({-}\infty,R]),
 \label{cc:eq:log-measure-counting}
\end{equation}
and, on their domains of finiteness,
\begin{align}
 \zeta_T(s)&:=\operatorname{Tr}|T|^s
 =\int e^{-sx}\,d\Lambda_T(x),                                
 \label{cc:eq:zeta-reader}\\
 \mathfrak F_T(a)&:=\frac12\sum_j\log(1+a^2s_j(T)^2)
 =\frac12\log\det(I+a^2T^*T).
 \label{cc:eq:fredholm-reader}
\end{align}
The determinant identity in \eqref{cc:eq:fredholm-reader} applies when
\(T\in\mathcal S_2\), because then \(T^*T\in\mathcal S_1\) and the
Fredholm determinant is defined; see \cite{Simon05}.  The series is otherwise understood in
\([0,\infty]\).
\end{definition}

\begin{proposition}[Exact positive spectralization]
\label{cc:prop:positive-spectralization}
For compact \(A,B\),
\begin{equation}
 \Lambda_{A\oplus B}=\Lambda_A+\Lambda_B,
 \qquad
 \Lambda_{A\otimes B}=\Lambda_A*\Lambda_B.
 \label{cc:eq:rig-law}
\end{equation}
Consequently
\begin{equation}
 \zeta_{A\oplus B}=\zeta_A+\zeta_B,
 \qquad
 \zeta_{A\otimes B}=\zeta_A\zeta_B,
 \qquad
 \mathfrak F_{A\oplus B}=\mathfrak F_A+\mathfrak F_B,
 \label{cc:eq:reader-rig-laws}
\end{equation}
whenever the indicated quantities are finite.  Moreover, with
\(K(u)=\frac12\log(1+e^{2u})\),
\begin{equation}
 \mathfrak F_T(e^t)=\int K(t-x)\,d\Lambda_T(x).
 \label{cc:eq:soft-threshold-reader}
\end{equation}
\end{proposition}

\begin{proof}
The singular multiset of \(A\oplus B\) is the disjoint union of the two
singular multisets, while that of \(A\otimes B\) is
\(\{s_j(A)s_k(B)\}_{j,k}\), by the Schmidt expansions of the two compact
arrows; see \cite{Simon05}.  The map \(s\mapsto-\log s\) turns the latter
products into sums.  Taking the Laplace transform and the soft-threshold
convolution gives \eqref{cc:eq:reader-rig-laws} and
\eqref{cc:eq:soft-threshold-reader}.
\end{proof}

We use the Laplace and integral forms of Karamata's theorem, Potter bounds,
and monotone asymptotic inversion in the standard form of
\cite[Chs.~1 and 4]{BGT87}.

\begin{theorem}[Counting, Mellin, and Fredholm dictionary]
\label{cc:thm:regular-dictionary}
Let \(p,r,C>0\), let \(\ell\) be eventually positive and slowly varying,
and suppose
\begin{equation}
 \mathcal N_T(R)
 \sim C e^{pR}R^{r-1}\ell(R)
 \qquad(R\to\infty).
 \label{cc:eq:regular-counting}
\end{equation}
Then, as \(\varepsilon\downarrow0\),
\begin{equation}
 \zeta_T(p+\varepsilon)
 \sim pC\Gamma(r)\varepsilon^{-r}\ell(1/\varepsilon).
 \label{cc:eq:regular-zeta}
\end{equation}
If \(0<p<2\), then \(T\in\mathcal S_2\) and
\begin{equation}
 \mathfrak F_T(a)
 \sim
 \frac{\pi C}{2\sin(\pi p/2)}
 a^p(\log a)^{r-1}\ell(\log a)
 \qquad(a\to\infty).
 \label{cc:eq:regular-fredholm}
\end{equation}
If \(p\geq2\), then \(\mathfrak F_T(a)=+\infty\) for every \(a>0\);
thus the restriction \(0<p<2\) is sharp.
The critical threshold mass satisfies
\begin{equation}
 \operatorname{Tr}\!\left(
 |T|^p\mathbf 1_{\{|T|\geq e^{-R}\}}
 \right)
 \sim\frac{pC}{r}R^r\ell(R).
 \label{cc:eq:threshold-mass}
\end{equation}
Finally,
\begin{equation}
 s_N(T)
 \sim
 \bigl(Cp^{1-r}\bigr)^{1/p}
 N^{-1/p}(\log N)^{(r-1)/p}\ell(\log N)^{1/p}.
 \label{cc:eq:regular-quantile}
\end{equation}
Thus \(T\in\mathcal S_q\) for every \(q>p\), but
\(T\notin\mathcal S_p\).  If, in addition, \(\zeta_T\) admits a
meromorphic continuation to a complex neighborhood of \(p\), then
\(r\in\mathbb N\), \(\ell(R)\) converges to a positive constant, and the
pole at \(p\) has order exactly \(r\).
\end{theorem}

\begin{proof}
The part of \(\Lambda_T\) in \(({-}\infty,0)\) is finite and contributes
only bounded terms to the limits below.  On \([0,\infty)\), Stieltjes
integration by parts gives
\begin{equation}
 \zeta_T(p+\varepsilon)
 =(p+\varepsilon)\int_0^\infty
 e^{-(p+\varepsilon)R}\mathcal N_T(R)\,dR+O(1).
 \label{cc:eq:zeta-layer-cake}
\end{equation}
Karamata's Laplace theorem applied to
\eqref{cc:eq:regular-counting} proves
\eqref{cc:eq:regular-zeta}.  If \(p<2\), the same counting bound makes
\(\int_0^\infty e^{-2R}\mathcal N_T(R)\,dR\) finite, so
\(T\in\mathcal S_2\).  Tonelli's theorem and scalar integration then give
the exact layer-cake identity
\begin{equation}
 \mathfrak F_T(e^t)
 =\int_{\mathbb R}
 \frac{\mathcal N_T(R)}{1+e^{2(R-t)}}\,dR
 \label{cc:eq:fredholm-layer-cake}
\end{equation}
and the change \(R=t+y\) reduces the leading constant to
\[
 \int_{-\infty}^{\infty}\frac{e^{py}}{1+e^{2y}}\,dy
 =\frac{\pi}{2\sin(\pi p/2)}.
\]
Potter bounds give an integrable majorant of the form
\(e^{(p+\eta)y}\mathbf1_{\{y<0\}}+
e^{-(2-p-\eta)y}\mathbf1_{\{y\geq0\}}\), for
\(0<\eta<\min\{p,2-p\}\).  Dominated convergence therefore applies
precisely for \(0<p<2\), proving
\eqref{cc:eq:regular-fredholm}.
If \(p\geq2\), the counting law gives
\(\sum_js_j(T)^2=\infty\).  Since \(s_j(T)\to0\) and
\(\log(1+a^2u^2)\geq \tfrac12a^2u^2\) for all sufficiently small \(u\),
the defining series for \(\mathfrak F_T(a)\) diverges for every \(a>0\).

For the threshold mass, put
\(U_p(R)=\int_{(-\infty,R]}e^{-px}\,d\Lambda_T(x)\).  Another Stieltjes
integration by parts gives
\[
 U_p(R)=e^{-pR}\mathcal N_T(R)
 +p\int_0^R e^{-px}\mathcal N_T(x)\,dx+O(1).
\]
The boundary term is of order \(R^{r-1}\ell(R)\), whereas Karamata's
integral theorem makes the integral asymptotic to
\(C R^r\ell(R)/r\).  This proves
\eqref{cc:eq:threshold-mass}.

Finally set
\[
 b_N:=\bigl(Cp^{1-r}\bigr)^{1/p}N^{-1/p}
 (\log N)^{(r-1)/p}\ell(\log N)^{1/p}.
\]
Slow variation and \(-\log b_N\sim p^{-1}\log N\) imply, for each fixed
\(\delta\in(0,1)\),
\[
 \frac{\mathcal N_T(-\log((1\pm\delta)b_N))}{N}
 \longrightarrow (1\pm\delta)^{-p}.
\]
Monotonicity, with the usual one-index adjustment at a multiplicity jump,
brackets \(s_N(T)/b_N\) between \(1-\delta\) and \(1+\delta\).
Letting \(\delta\downarrow0\) proves
\eqref{cc:eq:regular-quantile} with the displayed constant.  The same
counting law shows \(\sum_Ns_N(T)^q<\infty\) for \(q>p\), while
\eqref{cc:eq:threshold-mass} diverges for \(q=p\); hence the final ideal
claims follow as well.

For the last assertion, \eqref{cc:eq:regular-zeta} tends to infinity, so
the meromorphic singularity is a pole, say of order \(k\).  Comparison with
its leading Laurent term gives
\(\ell(x)\sim c x^{k-r}\) for some \(c>0\).  Since \(\ell\) is slowly
varying, necessarily \(k=r\).  Thus \(r\in\mathbb N\),
\(\ell(x)\to c\), and the pole order is \(r\).
\end{proof}

For \(p,r>0\), when the following limit belongs to \((0,\infty)\), set
\begin{equation}
 \mathscr R_p^{(r)}(T)
 :=\lim_{\varepsilon\downarrow0}
 \varepsilon^r\zeta_T(p+\varepsilon).
 \label{cc:eq:Mellin-coefficient}
\end{equation}
This is a real-axis boundary coefficient.  It is a Laurent coefficient
only when a meromorphic continuation is known and \(r\in\mathbb N\).

The corresponding tensor-product asymptotics are known in greater
generality; see \cite{KarolNazarovNikitin08,Rastegaev18}.  We record the
regularly varying case with its constants.

\begin{theorem}[Critical arithmetic under tensor product]
\label{cc:thm:tensor-critical-arithmetic}
Suppose \(p_A,p_B,r_A,r_B,C_A,C_B>0\), the functions
\(\ell_A,\ell_B\) are eventually positive and slowly varying, and
\(A,B\) satisfy
\[
 \mathcal N_A(R)\sim C_Ae^{p_AR}R^{r_A-1}\ell_A(R),
 \qquad
 \mathcal N_B(R)\sim C_Be^{p_BR}R^{r_B-1}\ell_B(R).
\]
If \(p_A>p_B\), then
\begin{equation}
 \mathcal N_{A\otimes B}(R)
 \sim C_A\zeta_B(p_A)e^{p_AR}R^{r_A-1}\ell_A(R).
 \label{cc:eq:nonresonant-product}
\end{equation}
The symmetric formula holds when \(p_B>p_A\).  If
\(p_A=p_B=p\), then
\begin{equation}
 \mathcal N_{A\otimes B}(R)
 \sim p\,\mathrm B(r_A,r_B)C_AC_B
 e^{pR}R^{r_A+r_B-1}\ell_A(R)\ell_B(R),
 \label{cc:eq:resonant-product}
\end{equation}
where
\(\mathrm B(a,b)=\Gamma(a)\Gamma(b)/\Gamma(a+b)\).
Consequently the power exponent with larger value dominates, while equal
power exponents resonate and their Mellin orders add.

In the pure logarithmic case, if \(p,r,\rho,a,b>0\) and
\[
 s_N(A)\sim aN^{-1/p}(\log N)^{(r-1)/p},\qquad
 s_N(B)\sim bN^{-1/p}(\log N)^{(\rho-1)/p},
\]
then
\begin{equation}
 s_N(A\otimes B)
 \sim ab\left(
 \frac{\Gamma(r)\Gamma(\rho)}{\Gamma(r+\rho)}
 \right)^{1/p}
 N^{-1/p}(\log N)^{(r+\rho-1)/p}.
 \label{cc:eq:amplitude-convolution}
\end{equation}
In particular, if \(s_N(A_j)\sim a_jN^{-1/p}\), then
\begin{equation}
 s_N(A_1\otimes\cdots\otimes A_k)
 \sim\frac{a_1\cdots a_k}{\Gamma(k)^{1/p}}
 N^{-1/p}(\log N)^{(k-1)/p}.
 \label{cc:eq:kfold-simple-atoms}
\end{equation}
Finally, for any \(p,r,\rho>0\), existence of the coefficients on the
right-hand sides below implies existence of those on the left, and
\begin{align}
 \mathscr R_p^{(r+\rho)}(A\otimes B)
 &=\mathscr R_p^{(r)}(A)\mathscr R_p^{(\rho)}(B),
 \label{cc:eq:coefficient-product}\\
 \mathscr R_p^{(r)}(A\otimes C)
 &=\mathscr R_p^{(r)}(A)\operatorname{Tr}|C|^p,
 \qquad 0\ne C\in\mathcal S_p,
 \label{cc:eq:coefficient-decoration}\\
 \mathscr R_p^{(r)}(cA)
 &=|c|^p\mathscr R_p^{(r)}(A),\qquad c\ne0.
 \label{cc:eq:coefficient-scalar}
\end{align}
\end{theorem}

\begin{proof}
The positive rig law gives the exact Stieltjes convolution
\[
 \mathcal N_{A\otimes B}(R)
 =\int_{\mathbb R}\mathcal N_A(R-y)\,d\Lambda_B(y).
\]
If \(p_A>p_B\), divide by the proposed right-hand side.
The integrand converges pointwise to \(e^{-p_Ay}\); Potter bounds and
\(\zeta_B(p_A-\eta)<\infty\) for a small
\(\eta\in(0,p_A-p_B)\) give an integrable majorant.  Dominated convergence
proves \eqref{cc:eq:nonresonant-product}.

At resonance, set
\[
 f_A(x):=e^{-px}\mathcal N_A(x),\qquad
 d\mu_B(y):=e^{-py}\,d\Lambda_B(y),\qquad
 U_B(R):=\mu_B(({-}\infty,R]).
\]
The threshold-mass formula \eqref{cc:eq:threshold-mass} gives
\[
 f_A(x)\sim C_Ax^{r_A-1}\ell_A(x),\qquad
 U_B(R)\sim\frac{pC_B}{r_B}R^{r_B}\ell_B(R).
\]
The portion of the convolution with \(y<0\) is a finite weighted sum and
has order at most
\(e^{pR}R^{r_A-1}\ell_A(R)\).  Since the log-spectrum of \(A\) is bounded
below, the portion with \(y>R\) is confined to an interval of fixed length;
the local estimate below makes it
\(O(e^{pR}R^{r_B-1}\ell_B(R))\).  Both are lower order than the claimed
main term because \(r_A,r_B>0\).  It therefore remains to analyze
\(0\leq y\leq R\).
The regular counting bound also gives, uniformly for
\(0\leq h\leq R/2\),
\[
 U_B(R)-U_B(R-h)
 \leq K(1+h)R^{r_B-1}\ell_B(R).
\]
Indeed, integration by parts on \((R-h,R]\), followed by the upper half of
the counting asymptotic and Potter's bound on \([R/2,R]\), gives
\begin{align*}
 U_B(R)-U_B(R-h)
 &\leq e^{-pR}\mathcal N_B(R)
   +p\int_{R-h}^R e^{-py}\mathcal N_B(y)\,dy\\
 &\leq K(1+h)R^{r_B-1}\ell_B(R).
\end{align*}
Fix \(0<\delta<1/2\).  On
\(\delta R\leq y\leq(1-\delta)R\), uniform convergence for regularly
varying functions and the vague convergence
\[
 \frac{dU_B(Rt)}{U_B(R)}
 \Longrightarrow r_Bt^{r_B-1}\,dt
\]
give the limiting integral over \([\delta,1-\delta]\).  Potter bounds
control the endpoint \(0\leq y\leq\delta R\) by
\(K\delta^{r_B-\eta}\) times the normalization.  Applying the displayed
increment estimate on unit intervals and summing the regular bound for
\(f_A\) controls \((1-\delta)R\leq y\leq R\) by
\(K\delta^{r_A-\eta}\), for any
\(0<\eta<\min\{r_A,r_B\}\).  Letting \(\delta\downarrow0\) therefore gives
\[
 \int_0^R f_A(R-y)\,dU_B(y)
 \sim pC_AC_B\,\mathrm B(r_A,r_B)
 R^{r_A+r_B-1}\ell_A(R)\ell_B(R).
\]
Multiplication by \(e^{pR}\) proves
\eqref{cc:eq:resonant-product}.  Its constant is consistent with zeta
multiplication because
\[
 p\,[p\,\mathrm B(r_A,r_B)C_AC_B]\Gamma(r_A+r_B)
 =(pC_A\Gamma(r_A))(pC_B\Gamma(r_B)).
\]

The counting constants corresponding to the two hypotheses are
\(a^pp^{r-1}\) and \(b^pp^{\rho-1}\).  Insert them into
\eqref{cc:eq:resonant-product} and use
\eqref{cc:eq:regular-quantile}.  Iteration with
\(\Gamma(1)=1\) proves \eqref{cc:eq:kfold-simple-atoms}.

The coefficient identities follow from
\(\zeta_{A\otimes B}=\zeta_A\zeta_B\),
\(\zeta_{cA}=|c|^s\zeta_A\), and
\(\zeta_C(p+\varepsilon)\to\operatorname{Tr}|C|^p\).
\end{proof}

At the ancestry face \(p=1\), if
\(s_N(T)\sim c(\log N)^{r-1}/N\), the dictionary becomes
\begin{equation}
 \zeta_T(1+\varepsilon)
 \sim c\Gamma(r)\varepsilon^{-r},
 \qquad
 \mathfrak F_T(a)
 \sim\frac{\pi c}{2}a(\log a)^{r-1}.
 \label{cc:eq:p-one-dictionary}
\end{equation}
These are real-axis asymptotics.  A literal complex pole follows only when
\(\zeta_T\) is meromorphic at \(p\) and \(r\in\mathbb N\); the harmonic atom
below has such a continuation.  If \(r\in\mathbb N\), the singular-value
asymptotic also gives
\(T\in\mathfrak W^{[r]}\setminus\mathfrak W^{[r]}_0\).

\begin{proposition}[Pointwise detection of the weak logarithmic scale]
\label{cc:prop:weak-detection}
Fix \(0<p<2\), \(r\geq1\), and set
\[
 \varphi_{p,r}(N)
 :=N^{-1/p}\bigl(\log(e+N)\bigr)^{(r-1)/p}.
\]
For a compact \(T\), the following are equivalent:
\begin{align}
 s_N(T)&=O(\varphi_{p,r}(N)),\label{cc:eq:weak-singular}\\
 \mathcal N_T(R)&=O(e^{pR}(1+R)^{r-1}),\label{cc:eq:weak-counting}\\
 \mathfrak F_T(a)&=O(a^p(\log(e+a))^{r-1}).
 \label{cc:eq:weak-fredholm}
\end{align}
The same equivalence holds with every \(O\) replaced by \(o\).
Thus the positive determinant detects both the weak and little
logarithmic ideals, even when no asymptotic equivalent exists.
\end{proposition}

\begin{proof}
Because \(N\mapsto\varphi_{p,r}(N)\) is eventually decreasing and regularly
varying with index \(-1/p\), monotone inversion gives the equivalence of
\eqref{cc:eq:weak-singular} and \eqref{cc:eq:weak-counting}, including
the little versions.  Put \(t=\log a\).  Inserting the counting bound into
\eqref{cc:eq:fredholm-layer-cake}, split at \(R=t\).  The part \(R\leq t\)
is controlled at its upper endpoint by \(p>0\); on \(R\geq t\), use
\((1+e^{2(R-t)})^{-1}\leq e^{-2(R-t)}\), which is integrable exactly when
\(p<2\).  This proves \eqref{cc:eq:weak-fredholm}.  Conversely, every singular value at least
\(e^{-R}\) contributes at least \(\frac12\log2\) to
\(\mathfrak F_T(e^R)\), so
\[
 \mathcal N_T(R)\leq\frac{2}{\log2}\mathfrak F_T(e^R).
\]
For the little version, first make the normalized counting bound smaller
than an arbitrary \(\delta>0\) beyond a fixed radius.  The finite initial
spectral part contributes only \(O(\log(e+a))\), which is
\(o(a^p(\log a)^{r-1})\); the same split makes the normalized Fredholm
limsup at most a constant times \(\delta\).  Letting \(\delta\downarrow0\)
and using the displayed threshold inequality proves both implications.
\end{proof}

\begin{proposition}[Uniform logarithmic envelopes]
\label{cc:prop-uniform-logarithmic-envelope}
Let \(r\ge1\), let \(T\) be compact, and suppose
\begin{equation}
 \|T\|_{\mathfrak W^{[r]}}\le b.
 \label{cc:eq-uniform-logarithmic-input}
\end{equation}
Then \(T\in\mathcal S_2\), and there is a constant \(C_r\), depending
only on \(r\), such that for \(a\ge0\) and \(0<\varepsilon\le1\),
\begin{align}
 \mathfrak F_T(a)
 &\le C_r(ab)\bigl(\log(e+ab)\bigr)^{r-1},
 \label{cc:eq-uniform-fredholm-envelope}\\
 \zeta_T(1+\varepsilon)
 &\le C_r b^{1+\varepsilon}\varepsilon^{-r}.
 \label{cc:eq-uniform-mellin-envelope}
\end{align}
The same bounds hold for every flattening of a tensor whose
multi-flattening \(\mathfrak W^{[r]}\)-norm is at most \(b\).  Thus, once a
Rees coefficient has been placed in grade \(r\), all upper Mellin and
Fredholm consequences are automatic.
\end{proposition}

\begin{proof}
The case \(b=0\) is immediate.  After replacing \(T\) by \(T/b\),
\eqref{cc:eq-uniform-logarithmic-input} gives
\[
 s_N(T/b)\le \frac{(\log(e+N))^{r-1}}{N}.
\]
The quantitative implication from the singular-value bound to the Fredholm
bound in \Cref{cc:prop:weak-detection}, followed by
\(\mathfrak F_T(a)=\mathfrak F_{T/b}(ab)\), proves
\eqref{cc:eq-uniform-fredholm-envelope}.  The same singular-value estimate
gives
\[
 \zeta_T(1+\varepsilon)
 \le b^{1+\varepsilon}
 \sum_{N\ge1}N^{-1-\varepsilon}
       (\log(e+N))^{(r-1)(1+\varepsilon)}.
\]
Integral comparison and the substitution \(u=\varepsilon\log x\) bound the
sum by
\[
 C_r\varepsilon^{-r-(r-1)\varepsilon}
 \le C'_r\varepsilon^{-r},
\]
because \(\sup_{0<\varepsilon\le1}
\varepsilon^{-(r-1)\varepsilon}<\infty\).  Membership in
\(\mathcal S_2\) follows from the same singular-value bound.  The
multi-flattening assertion is obtained by applying the scalar statement to
each cut.
\end{proof}
\subsection{Critical packets and ancestry amplification}
\label{cc:subsec-critical-packets}

The preceding dictionary converts a realized singular profile into scalar
spectral laws.  We now isolate the single analytic mechanism shared by all
finite-depth realizations.  Put
\begin{equation}
 \phi_r(N):=\frac{(\log(e+N))^{r-1}}{N},\qquad r\ge1.
 \label{cc:eq-critical-profile}
\end{equation}
A compact operator is a \emph{critical atom of amplitude \(a>0\)} if
\(s_N(B)\sim a/N\).  A family \((B_L)\) is a
\emph{window-critical atom family} if
\(\sup_L\|B_L\|_{\mathfrak W^{[1]}}<\infty\) and, for some \(a>0\),
\[
 s_k(B_L)\ge \frac ak,\qquad 1\le k\le L.
\]

For \(r\ge1\), call
\begin{equation}
 \mathcal P=
 \bigoplus_{\nu=1}^m c_\nu U_\nu
 \bigl(B_{\nu,1}\otimes\cdots\otimes B_{\nu,r}\otimes C_\nu\bigr)
 V_\nu^*+R
 \label{cc:eq-exact-critical-packet}
\end{equation}
an \emph{exact critical \(r\)-packet} when the outer input and output
packets are mutually orthogonal, every \(B_{\nu,j}\) is a critical atom,
\(0\ne C_\nu\in\mathcal S_1\), and
\begin{equation}
 s_N(R)=
 \begin{cases}
  o(N^{-1}),&r=1,\\
  O(\phi_{r-1}(N)),&r\ge2.
 \end{cases}
 \label{cc:eq-packet-lower-remainder}
\end{equation}
Zero coefficients are omitted.  A compact operator \(Y\) has a
\emph{principal \(r\)-exposure} if there are contractions \(Q_\pm\) for
which \(Q_+YQ_-\) is one nonzero \(r\)-fold packet plus a remainder of the
form \eqref{cc:eq-packet-lower-remainder}.  Exact packets produce full-scale
asymptotics; principal exposures are the weaker datum needed to certify a
sharp grade.

For \(\kappa>0\), write
\begin{equation}
 T\models\operatorname{Crit}_r(\kappa)
 \label{cc:eq-exact-critical-law-notation}
\end{equation}
when the following three equivalent critical readouts hold:
\begin{align}
 s_N(T)
 &\sim\frac{\kappa}{\Gamma(r)}\phi_r(N),
 \label{cc:eq-packet-singular}\\
 \operatorname{Tr}|T|^{1+\varepsilon}
 &\sim\kappa\varepsilon^{-r},
 \label{cc:eq-packet-mellin}\\
 \frac12\log\det(I+x^2T^*T)
 &\sim\frac{\pi\kappa}{2\Gamma(r)}
 x(\log x)^{r-1}.
 \label{cc:eq-packet-fredholm}
\end{align}
The limits are \(N\to\infty\), \(\varepsilon\downarrow0\), and
\(x\to\infty\).  The notation records real-axis data; meromorphic
continuation is an additional assertion.

\begin{lemma}[Critical product lower divisor law]
\label{cc:gen-lemma-divisor-lower}
Fix \(r\geq1\).  Suppose that, for every integer \(L\), compact
operators \(A_{j,L}\), \(1\leq j\leq r\), satisfy
\[
 s_k(A_{j,L})\geq\frac{c}{k},
 \qquad 1\leq k\leq L.
\]
For \(T_L=A_{1,L}\otimes\cdots\otimes A_{r,L}\), there are constants
\(b_r,B_r>0\), independent of \(L\), such that
\begin{equation}
 s_N(T_L)\geq b_r\frac{(\log(e+N))^{r-1}}{N}
 \label{cc:gen-divisor-lower}
\end{equation}
whenever
\[
 B_r\leq N\leq b_rL(\log(e+L))^{r-1}.
\]
\end{lemma}

\begin{proof}
The singular values of \(T_L\) are the decreasing rearrangement of
\(\prod_{j=1}^rs_{n_j}(A_{j,L})\).  For \(2\leq R\leq L\), every tuple
with \(n_1\cdots n_r\leq R\) contributes a singular value at least
\(c^r/R\).  The elementary divisor count
\[
 D_r(R):=
 \#\{(n_1,\ldots,n_r)\in\mathbb N^r:n_1\cdots n_r\leq R\}
 \geq c_rR(\log R)^{r-1}
\]
follows by induction from
\(D_r(R)=\sum_{n\leq R}D_{r-1}(R/n)\) and comparison of the resulting
harmonic sums with integrals.  Hence
\[
 s_{\lfloor c_rR(\log R)^{r-1}\rfloor}(T_L)\geq c^r/R.
\]
Monotonicity and inversion of
\(N\asymp_rR(\log R)^{r-1}\) prove
\eqref{cc:gen-divisor-lower}.
\end{proof}

\begin{lemma}[One-grade noncancellation]
\label{cc:gen-lemma-one-grade-noncancellation}
Suppose that
\(s_N(T_L)\geq c\phi_r(N)\) on arbitrarily long windows.
If \(r\geq2\) and
\[
 \sup_L\|R_L\|_{\mathfrak W^{[r-1]}}\leq C,
\]
then, above a fixed lower cutoff and on a fixed fraction of each common
window,
\[
 s_N(T_L+R_L)\geq c'\phi_r(N).
\]
For \(r=1\), the same holds after a fixed index shift if
\(\sup_L\operatorname{rank}R_L<\infty\).
\end{lemma}

\begin{proof}
For \(r\geq2\), Weyl's inequality applied to
\(T_L=(T_L+R_L)-R_L\) gives
\[
 s_N(T_L+R_L)\geq s_{2N-1}(T_L)-s_N(R_L).
\]
On the half-window on which \(2N-1\) remains admissible, the first term
is at least \(c_1\phi_r(N)\), whereas
\[
 s_N(R_L)\leq C\phi_{r-1}(N)
 \leq\frac{C_rC}{\log(e+N)}\phi_r(N).
\]
A cutoff independent of \(L\) makes the last coefficient at most
\(c_1/2\).  If \(r=1\) and the rank bound is \(R_0\), the
rank-perturbation form of Weyl's inequality gives
\[
 s_N(T_L+R_L)\geq s_{N+R_0}(T_L).
\]
\end{proof}



\begin{lemma}[Trace-class decoration of a logarithmic critical profile]
\label{cc:lem-trace-decoration-tauberian}
Let $r\ge1$, let $A$ be compact, and suppose
\begin{equation}
 s_N(A)\sim c_A\phi_r(N),\qquad c_A>0.
 \label{cc:eq-decoration-base-asymptotic}
\end{equation}
If $C\in\mathcal S_1$ is nonzero, then
\begin{equation}
 s_N(A\otimes C)\sim c_A\operatorname{Tr}|C|\,\phi_r(N).
 \label{cc:eq-trace-decoration-asymptotic}
\end{equation}
The same statement holds for a finite orthogonal direct sum, with the
constants added.  Moreover, trace-class decoration preserves the little grade
$\mathfrak W^{[r]}_0$.  For a finite-window lower bound it is enough to retain
one nonzero singular direction of $C$; if $s_k(A_L)\ge c\phi_r(k)$ on a
window, then
\[
 s_k(A_L\otimes C)\ge c\,s_1(C)\phi_r(k)
\]
on the same window.
\end{lemma}

\begin{proof}
Write $c_j=s_j(C)$.  The tensor-product counting identity is
\begin{equation}
 N_{A\otimes C}(u)=\sum_{j:c_j>0}N_A(u/c_j).
 \label{cc:eq-decoration-counting-identity}
\end{equation}
Monotone inversion of \eqref{cc:eq-decoration-base-asymptotic} gives
\[
 N_A(u)\sim c_Au^{-1}\bigl(\log(1/u)\bigr)^{r-1}
 \qquad (u\downarrow0),
\]
with the logarithmic factor interpreted as $1$ when $r=1$.  The same
inversion also gives a global envelope
\[
 N_A(v)\leq C_Av^{-1}
 \bigl(\log(e+1/v)\bigr)^{r-1},
 \qquad v>0.
\]
After a scalar rescaling of $C$, all but finitely many $c_j$ lie in $(0,1]$.
For those indices,
\[
 \frac{N_A(u/c_j)}
 {u^{-1}(\log(e+1/u))^{r-1}}
 \leq C_Ac_j
 \left(\frac{\log(e+c_j/u)}{\log(e+1/u)}\right)^{r-1}
 \leq C_Ac_j.
\]
For each fixed $j$ the same ratio converges to $c_Ac_j$.  The finitely many
remaining $c_j>1$ are handled separately.  Dominated convergence in
\eqref{cc:eq-decoration-counting-identity}, using $\sum_jc_j<\infty$, yields
\[
 N_{A\otimes C}(u)\sim
 c_A\operatorname{Tr}|C|\,u^{-1}(\log(1/u))^{r-1}.
\]
A second monotone inversion proves
\eqref{cc:eq-trace-decoration-asymptotic}.  Orthogonal sums add counting
functions.  The little assertion follows by approximating $C$ in trace norm
by finite-rank operators and using the continuous trace-decoration estimate in
$\mathfrak W^{[r]}$.  The window bound follows from the tensor slice
corresponding to a singular vector of $C$ for $s_1(C)$.
\end{proof}

\begin{theorem}[Critical ancestry amplification principle]
\label{cc:thm-critical-packet-principle}
Let \(r\ge1\).

\begin{enumerate}[label=\textup{(\roman*)},leftmargin=2.8em]
\item Suppose \(\mathcal P\) is an exact critical \(r\)-packet and
\[
 s_N(B_{\nu,j})\sim \frac{a_{\nu,j}}N,
 \qquad a_{\nu,j}>0.
\]
Set
\begin{equation}
 \kappa(\mathcal P)
 :=\sum_{\nu=1}^m |c_\nu|\,\operatorname{Tr}|C_\nu|
       \prod_{j=1}^r a_{\nu,j}.
 \label{cc:eq-packet-amplitude}
\end{equation}
Then
\[
 \mathcal P\models\operatorname{Crit}_r(\kappa(\mathcal P)).
\]

\item Let \(Y_L\) be uniformly bounded in \(\mathfrak W^{[r]}\).  Suppose
that, after fixed two-sided contractions, it has a principal exposure
\[
 Q_+Y_LQ_-=c_LU
 (B_{1,L}\otimes\cdots\otimes B_{r,L})V+R_L,
\]
where \(\inf_L|c_L|>0\), the \(B_{j,L}\) are window-critical with uniform
constants, and the remainders are uniformly bounded in
\(\mathfrak W^{[r-1]}\) when \(r\ge2\), or have uniformly bounded rank and
operator norm when \(r=1\).  Then there are constants
\(N_r,c_r,C_r>0\), independent of \(L\), such that
\begin{equation}
 s_N(Y_L)\ge c_r\phi_r(N)
 \label{cc:eq-window-exposure-lower}
\end{equation}
whenever
\[
 N_r\le N\le C_rL(\log(e+L))^{r-1}.
\]
Consequently no grade \(q<r\) can bound such a family uniformly on
arbitrarily long windows.
\end{enumerate}
\end{theorem}

\begin{proof}
For part~\textup{(i)}, repeated use of
\Cref{cc:thm:tensor-critical-arithmetic}, followed by
\Cref{cc:lem-trace-decoration-tauberian}, gives
\[
 s_N(B_{\nu,1}\otimes\cdots\otimes B_{\nu,r}\otimes C_\nu)
 \sim\frac{\operatorname{Tr}|C_\nu|\prod_ja_{\nu,j}}{\Gamma(r)}
 \phi_r(N).
\]
The orthogonal direct sum adds the corresponding log-counting constants,
and scalar multiplication contributes \(|c_\nu|\).  Thus the packet without
\(R\) has the equivalent \eqref{cc:eq-packet-singular}.  The remainder is
strictly lower grade.  Ky Fan inequalities, with an index split
\(\lfloor\delta N\rfloor\), show that it does not change a regularly varying
singular sequence of index \(-1\); first let \(N\to\infty\), then
\(\delta\downarrow0\).  The Mellin and Fredholm laws follow from
\Cref{cc:thm:regular-dictionary} at \(p=1\).

For part~\textup{(ii)},
\Cref{cc:gen-lemma-divisor-lower} gives the \(\phi_r\)-lower profile for the
principal tensor product on the stated divisor window.  The uniform
lower-grade bound cannot cancel it by
\Cref{cc:gen-lemma-one-grade-noncancellation}.  Singular values decrease
under two-sided contractions, so the same lower bound holds for \(Y_L\).
The ratio \(\phi_r(N)/\phi_q(N)=(\log(e+N))^{r-q}\) diverges for \(q<r\).
\end{proof}

The packet theorem is source-independent.  Its source-level
formulation will be given only after the canonical critical entrance and the
filtered realization and response-kernel structures have been constructed;
at that point exact, window, and bounded response certificates become three
modes of one theorem.
\section{Rees response kernels and valuation exchange}
\label{rk:sec-kernel}

The previous sections provide two filtrations with opposite orientations.  A
finite ancestry flag is decreasing, while analytic critical cost is
increasing.  The correct filtered object is therefore the dual ancestry
flag.  This observation makes the primal and dual directions contractions of
one associated-graded kernel.

We use superscripts $F^r$ for an abstract decreasing filtration and retain
subscripts $F_r$ for the concrete free-Lie and Schur filtrations introduced
above.  The concrete algebraic and analytic Rees objects are displayed with
$z^{-r}$ so that equal ancestry and critical costs occupy the same formal
weight.  The abstract saturation theorem uses the equivalent positive-degree
parameter $u^r$; formally $u=z^{-1}$.  These are bookkeeping conventions for
one filtration comparison, not two independent deformations.

\subsection{Strict filtered responses}
\label{rk:subsec-strict}

Let $M$ be a finite-dimensional vector space over a field $\mathbf k$ of
characteristic zero, with a finite decreasing filtration
\begin{equation}
 M=F^0M\supseteq F^1M\supseteq\cdots\supseteq F^{d+1}M=0,
 \qquad G_r:=F^rM/F^{r+1}M.
 \label{rk:eq-source-filtration}
\end{equation}
Its dual filtration is increasing:
\begin{equation}
 \mathscr D_rM^*:=(F^{r+1}M)^\perp,
 \qquad \mathscr D_{-1}M^*:=0.
 \label{rk:eq-dual-filtration}
\end{equation}
Restriction identifies
$\mathscr D_rM^*/\mathscr D_{r-1}M^*\cong G_r^*$.  For $q\ne0$ put
\begin{equation}
 a_F(q):=\min\{r:q\in\mathscr D_rM^*\}
       =\max\{r:q(F^rM)\ne0\},
 \qquad a_F(0):=-\infty.
 \label{rk:eq-algebraic-depth}
\end{equation}

Let $\mathcal Y$ be a vector space with an increasing filtration
\[
 0=K_{-1}\mathcal Y\subseteq K_0\mathcal Y\subseteq\cdots
 \subseteq K_d\mathcal Y=\mathcal Y,
 \qquad
 \partial_r\mathcal Y:=K_r\mathcal Y/K_{r-1}\mathcal Y.
\]
For $0\ne y\in\mathcal Y$, define
$\nu_K(y):=\min\{r:y\in K_r\mathcal Y\}$ and put
$\nu_K(0):=-\infty$.

\begin{definition}[Filtered and strict response]
\label{rk:def-filtered-strict}
A linear map $\mathcal R:M^*\to\mathcal Y$ is \emph{filtered} if
\[
 \mathcal R(\mathscr D_rM^*)\subseteq K_r\mathcal Y
 \quad\text{for every }r.
\]
It is \emph{strict} if
\begin{equation}
 \mathcal R^{-1}(K_r\mathcal Y)=\mathscr D_rM^*,
 \qquad -1\le r\le d.
 \label{rk:eq-strictness}
\end{equation}
A filtered map induces
\begin{equation}
 \operatorname{gr}_r\mathcal R:
 G_r^*\longrightarrow\partial_r\mathcal Y.
 \label{rk:eq-graded-response}
\end{equation}
\end{definition}

Extend both filtrations constantly above degree \(d\).  The induced Rees map is
\begin{equation}
 \operatorname{Rees}(\mathcal R):
 \bigoplus_{r\ge0}\mathscr D_rM^*u^r
 \longrightarrow
 \bigoplus_{r\ge0}K_r\mathcal Yu^r,
 \qquad qu^r\longmapsto\mathcal R(q)u^r.
 \label{rk:eq-Rees-map}
\end{equation}
A graded submodule \(I\) of the target is \emph{\(u\)-saturated} when
\(uy\in I\) implies \(y\in I\).  The clause \(r=-1\) in
\eqref{rk:eq-strictness} says exactly that \(\mathcal R\) is injective.

\begin{theorem}[Valuation exchange and Rees saturation]
\label{rk:thm-valuation-exchange}
For a filtered response $\mathcal R:M^*\to\mathcal Y$, the following are
equivalent.
\begin{enumerate}[label=\textup{(\roman*)},leftmargin=2.8em]
\item $\mathcal R$ is strict;
\item every map $\operatorname{gr}_r\mathcal R$ is injective;
\item $\nu_K(\mathcal Rq)=a_F(q)$ for every $q\in M^*$;
\item the induced Rees map is injective and has $u$-saturated image;
\item the Rees map is injective and its cokernel has no $u$-torsion.
\end{enumerate}
Under these equivalent conditions,
\begin{equation}
 \boxed{
 F^rM=
 \bigcap_{\substack{q\in M^*\\
              \nu_K(\mathcal Rq)\le r-1}}
 \ker q.}
 \label{rk:eq-flag-reconstruction}
\end{equation}
Thus forward depth and reconstruction of the primal flag are the same
strictness statement.
\end{theorem}

\begin{proof}
Filteredness gives $\nu_K(\mathcal Rq)\le a_F(q)$.  If all graded maps are
injective and $a_F(q)=r$, the nonzero class of $q$ in
$\mathscr D_r/\mathscr D_{r-1}$ cannot map into $K_{r-1}$, so equality
holds.  Conversely, a nonzero class in the kernel of
$\operatorname{gr}_r\mathcal R$ contradicts valuation equality.  This is
equivalent to \eqref{rk:eq-strictness}.

For the Rees clauses, use the map \eqref{rk:eq-Rees-map}.  If
$\mathcal R$ is strict and
$u(yu^r)=\mathcal R(q)u^{r+1}$ lies in its image, then
$y=\mathcal R(q)\in K_r$, hence $q\in\mathscr D_r$ and $yu^r$ already lies
in the image.  Thus the image is $u$-saturated.  Conversely, if
$\mathcal R(q)\in K_r$ while $q$ first belongs to $\mathscr D_s$ with
$s>r$, then $\mathcal R(q)u^s$ lies in the image and repeated saturation
would force $\mathcal R(q)u^r$ to lie in the image, hence produce
$q'\in\mathscr D_r$ with $\mathcal R(q')=\mathcal R(q)$.  Injectivity of the
Rees map gives $q'=q$, a contradiction.  Therefore saturation gives
strictness.  A submodule is $u$-saturated exactly when its quotient has no
$u$-torsion.

Finally, the probes of analytic depth at most $r-1$ are exactly
$\mathscr D_{r-1}=(F^rM)^\perp$.  Finite-dimensional annihilator duality
gives \eqref{rk:eq-flag-reconstruction}.
\end{proof}

\begin{proposition}[All filtered lifts of one special fibre]
\label{rk:prop-filtered-lifts}
Let $\operatorname{Hom}_{\rm fil}(M^*,\mathcal Y)$ be the filtered maps and
let $\operatorname{Hom}_{-1}(M^*,\mathcal Y)$ consist of maps satisfying
$T(\mathscr D_r)\subseteq K_{r-1}\mathcal Y$.  There is a natural exact
sequence
\begin{equation}
 0\longrightarrow\operatorname{Hom}_{-1}(M^*,\mathcal Y)
 \longrightarrow\operatorname{Hom}_{\rm fil}(M^*,\mathcal Y)
 \longrightarrow
 \bigoplus_r\operatorname{Hom}(G_r^*,\partial_r\mathcal Y)
 \longrightarrow0.
 \label{rk:eq-filtered-lift-exact}
\end{equation}
Consequently, two filtered responses with the same associated-graded map
differ only by a one-step-lowering response and have the same valuation
theorem.  A filtered splitting selects a lift but creates no new invariant.
\end{proposition}

\begin{proof}
The kernel of passage to the associated graded is precisely
$\operatorname{Hom}_{-1}$.  Surjectivity follows by choosing complements of
successive filtration steps in source and target and lifting each graded map
between those complements.
\end{proof}

\subsection{The canonical response kernel}
\label{rk:subsec-kernel}

For a finite-dimensional space $G$, let
$\operatorname{coev}_G(1)\in G\otimes G^*$ be the canonical coevaluation
tensor.  In any pair of dual bases it is $\sum_je_j\otimes e_j^*$.

\begin{definition}[Rees response kernel]
\label{rk:def-response-kernel}
For a filtered response $\mathcal R$, its order-$r$ Rees response kernel is
\begin{equation}
 \mathfrak K_r(\mathcal R)
 :=\bigl(\operatorname{Id}_{G_r}\otimes
          \operatorname{gr}_r\mathcal R\bigr)
    \operatorname{coev}_{G_r}(1)
 \in G_r\otimes\partial_r\mathcal Y.
 \label{rk:eq-response-kernel}
\end{equation}
For $K\in G\otimes Z$, define its left support by
\begin{equation}
 \operatorname{supp}_L(K)
 :=\operatorname{span}\bigl\{
 (\operatorname{Id}_G\otimes\lambda)K:\lambda\in Z^*\bigr\}
 \subseteq G.
 \label{rk:eq-left-support}
\end{equation}
Equivalently, this is the smallest subspace $S\subseteq G$ for which
$K\in S\otimes Z$.
\end{definition}

\begin{theorem}[Kernel support theorem]
\label{rk:thm-kernel-support}
For every $r$,
\begin{equation}
 \ker(\operatorname{gr}_r\mathcal R)
 =\operatorname{supp}_L\bigl(\mathfrak K_r(\mathcal R)\bigr)^\perp
 \subseteq G_r^*.
 \label{rk:eq-kernel-annihilator}
\end{equation}
Consequently, $\mathcal R$ is strict if and only if
\begin{equation}
 \operatorname{supp}_L\bigl(\mathfrak K_r(\mathcal R)\bigr)=G_r
 \quad\text{for every }r.
 \label{rk:eq-full-support}
\end{equation}
Primal symbol exposure and dual depth detection are the two tensor legs of
one canonical kernel.
\end{theorem}

\begin{proof}
For $q\in G_r^*$, contraction in the first tensor leg gives
\[
 (q\otimes\operatorname{Id})\mathfrak K_r(\mathcal R)
 =\operatorname{gr}_r\mathcal R(q).
\]
This vanishes exactly when $q$ annihilates every contraction in the second
leg, proving \eqref{rk:eq-kernel-annihilator}.  The final assertion follows
from \Cref{rk:thm-valuation-exchange}.
\end{proof}

Every declared native reader merely pushes this kernel forward.

\begin{theorem}[Reader symbols and Yoneda factorization]
\label{rk:thm-reader-symbols}
Let $\mathscr A_r$ be a class of allowed boundary readers
$A:\partial_r\mathcal Y\to Z_A$, and put
\[
 K_{r,A}:=(\operatorname{Id}\otimes A)\mathfrak K_r(\mathcal R)
 \in G_r\otimes Z_A,
 \qquad
 S_r(\mathscr A):=
 \sum_{A\in\mathscr A_r}\operatorname{supp}_L(K_{r,A}).
\]
The total allowed response on $G_r^*$ is injective if and only if
\begin{equation}
 S_r(\mathscr A)=G_r.
 \label{rk:eq-reader-support}
\end{equation}
If $Z_A$ contains a fixed nonzero line $\mathbf k\tau_{r,A}$ and
$K_{r,A}=\sigma_A\otimes\tau_{r,A}$, then
\begin{equation}
 A\,\operatorname{gr}_r\mathcal R(q)
 =q(\sigma_A)\tau_{r,A}.
 \label{rk:eq-rank-one-reader}
\end{equation}
Thus rank-one symbol factorization is not an additional axiom; it is a
coordinate form of the canonical response kernel.  The reader symbols detect
order $r$ exactly when they span $G_r$.
\end{theorem}

\begin{proof}
By \Cref{rk:thm-kernel-support}, the kernel of the product of the allowed
graded maps is
$\bigcap_A\operatorname{supp}_L(K_{r,A})^\perp
=S_r(\mathscr A)^\perp$.  Formula \eqref{rk:eq-rank-one-reader} is contraction
in the first tensor leg.
\end{proof}

\begin{proposition}[Naturality and base change of response kernels]
\label{rk:prop-kernel-naturality}
Let $f:(M,F^\bullet)\to(M',F^\bullet)$ preserve the decreasing
filtrations, let
$B:(\mathcal Y,K_\bullet)\to(\mathcal Y',K_\bullet)$ preserve the
increasing filtrations, and let
$\mathcal R:M^*\to\mathcal Y$ and
$\mathcal R':(M')^*\to\mathcal Y'$ be filtered responses satisfying
\begin{equation}
 \mathcal R'=B\circ\mathcal R\circ f^*.
 \label{rk:eq-response-base-change}
\end{equation}
Write $f_r:G_r(M)\to G_r(M')$ and
$B_r:\partial_r\mathcal Y\to\partial_r\mathcal Y'$ for the induced
maps.  Then
\begin{equation}
 \mathfrak K_r(\mathcal R')
 =(f_r\otimes B_r)\mathfrak K_r(\mathcal R).
 \label{rk:eq-kernel-base-change}
\end{equation}
Consequently
\[
 \operatorname{supp}_L\mathfrak K_r(\mathcal R')
 \subseteq f_r\bigl(\operatorname{supp}_L\mathfrak K_r(\mathcal R)\bigr).
\]
If $\mathcal R$ has full support, $f_r$ is surjective, and $B_r$ is
injective on the image of
$\operatorname{gr}_r\mathcal R\circ f_r^*$, then $\mathcal R'$ has full
support on $G_r(M')$.  Thus a quotient observation recovers exactly its
observable ancestry layer.  In particular, spatial pushforwards,
exact-layer identifications, and native boundary readers act on one canonical
kernel rather than creating new detector data.
\end{proposition}

\begin{proof}
Under the canonical identification
$\operatorname{Hom}(G_r^*,Z)\cong G_r\otimes Z$, the tensor corresponding
to the composite
$B_r\circ\operatorname{gr}_r\mathcal R\circ f_r^*$ is
$(f_r\otimes B_r)\mathfrak K_r(\mathcal R)$.  By
\eqref{rk:eq-response-base-change}, that composite is
$\operatorname{gr}_r\mathcal R'$, proving
\eqref{rk:eq-kernel-base-change}.  The support inclusion follows by
contracting the second tensor leg.  Under the final hypotheses,
$f_r^*$, $\operatorname{gr}_r\mathcal R$, and the restriction of $B_r$
are all injective.  Their composite is therefore injective, and
\Cref{rk:thm-kernel-support} gives full support for $\mathcal R'$.
\end{proof}

\begin{proposition}[Finite certificates and the Rees observability Gramian]
\label{rk:prop-robust-support}
Let $G$ be a finite-dimensional exact layer and let a labelled native reader
family have full support on $G$.  Then at most $\dim G$ scalar contractions of
allowed reader coordinates suffice to detect $G$.  Thus there are
$\sigma_1,\ldots,\sigma_m\in G$, $m\leq\dim G$, spanning $G$, and a scalar
analysis map
\[
 \mathsf C_G:G^*\longrightarrow\mathbf k^m,
 \qquad q\longmapsto(q(\sigma_1),\ldots,q(\sigma_m)).
\]
After fixing Hilbert norms, write
\[
 \kappa(A):=\inf_{\|x\|=1}\|Ax\|
\]
for the minimum modulus of a linear map $A$, and define the finite Rees
observability Gramian
\begin{equation}
 \Gamma_G:=\mathsf C_G^*\mathsf C_G,
 \qquad
 \kappa_G:=\kappa(\mathsf C_G)
 =\lambda_{\min}(\Gamma_G)^{1/2}.
 \label{rk:eq-rees-observability-gramian}
\end{equation}
Then full support is equivalent to $\Gamma_G\succ0$, equivalently
$\kappa_G>0$.  For another analysis map $\widetilde{\mathsf C}_G$,
\begin{equation}
 |\kappa(\widetilde{\mathsf C}_G)-\kappa(\mathsf C_G)|
 \leq\|\widetilde{\mathsf C}_G-\mathsf C_G\|.
 \label{rk:eq-observability-gap-stability}
\end{equation}
Consequently full support is stable under perturbations smaller than
$\kappa_G$, and operator-norm convergence of finite reader maps transports a
uniform positive Gramian gap between Galerkin models and their limit.
\end{proposition}

\begin{proof}
Full support means that contractions of the labelled kernels span $G$.
Choose a basis from these contractions; this uses at most $\dim G$ scalar
readers and gives the displayed analysis map.  Its kernel is the annihilator
of $\operatorname{span}\{\sigma_j\}$, proving the Gramian criterion.  For a
unit vector $q$,
\[
 \|\widetilde{\mathsf C}_Gq\|
 \geq\|\mathsf C_Gq\|-\|\widetilde{\mathsf C}_G-\mathsf C_G\|.
\]
Taking infima and then interchanging the two maps proves
\eqref{rk:eq-observability-gap-stability}.
\end{proof}

\begin{remark}[The genuinely native question]
If every algebraic boundary functional is admitted as a reader, full support
is formal.  A source theorem is nonformal because its allowed readers are
restricted to operations intrinsic to the observable: tensor flattenings,
spatial compressions, conditional readouts, or another declared class.  The
only native issue is whether these restricted pushforwards retain full left
support.
\end{remark}


\subsection{Native reachability and finite Rees observability}
\label{rk:subsec-reachability}
Fix an exact layer $G$.  A native experiment is described by finite-dimensional
parameter data $E_e$, a source-symbol map $i_e:E_e\to G$, and a boundary
transfer $T_e:E_e^*\to B_e$.  The corresponding labelled graded readout is
\begin{equation}
 \mathcal O_e=T_e i_e^*:G^*\longrightarrow B_e.
 \label{rk:eq-native-experiment-factorization}
\end{equation}
For a subspace $W\subseteq E_e^*$, write
$W^\circ:=\{x\in E_e:\ell(x)=0\text{ for every }\ell\in W\}$.
Set
\[
 \mathsf R:=\sum_e i_e(E_e),\qquad
 \mathsf V:=\sum_e i_e\bigl((\ker T_e)^\circ\bigr)\subseteq G.
\]
The spaces $\mathsf R$ and $\mathsf V$ are respectively the reachable and
visible symbol spaces.

\begin{proposition}[Reachability--visibility decomposition]
\label{rk:prop-reachability-visibility}
For the product readout $\mathcal O=(\mathcal O_e)_e$ one has
\begin{equation}
 \ker\mathcal O=\mathsf V^\circ,
 \qquad
 \mathsf V\subseteq\mathsf R\subseteq G.
 \label{rk:eq-reachability-visibility}
\end{equation}
If every $T_e$ is injective, then $\mathsf V=\mathsf R$; hence full Rees
support is equivalent to $\mathsf R=G$.
\end{proposition}

\begin{proof}
The kernel tensor of $\mathcal O_e=T_ei_e^*$ is
$(i_e\otimes T_e)\operatorname{coev}_{E_e}$.  Its left support is
$i_e(\operatorname{Ran}T_e^*)=i_e((\ker T_e)^\circ)$.  Summing over the
labelled experiments and applying the kernel--annihilator identity proves
\eqref{rk:eq-reachability-visibility}.  Injectivity of $T_e$ gives
$(\ker T_e)^\circ=E_e$.
\end{proof}

Here finite Rees observability means separation of a finite-dimensional
labelled exact layer; it is distinct from state or generator observability for
an evolution equation.

\subsection{Critical boundary consequences}
\label{rk:subsec-critical-consequences}

For the logarithmic weak filtration introduced in
\Cref{cr:sec-analytic}, write
\[
 \tau_r:=\bigl[D_\infty^{\otimes r}\bigr],
 \qquad
 D_\infty=\operatorname{diag}(1,1/2,1/3,\ldots).
\]
By \Cref{cc:thm:tensor-critical-arithmetic}, $\tau_r$ is a nonzero class of
critical cost $r$, $\tau_a\otimes\tau_b=\tau_{a+b}$, and a representative
satisfies
\begin{equation}
 s_N(D_\infty^{\otimes r})
 \sim\frac1{\Gamma(r)}\phi_r(N).
 \label{rk:eq-critical-boundary-line}
\end{equation}

\begin{definition}[Full-support and coefficient-pure response]
\label{rk:def-full-pure}
A Rees response kernel is \emph{full-support} in order $r$ if its left
support is $G_r$.  Relative to declared native readers, it is
\emph{coefficient-pure} if the total pushed-forward kernel lies in
\[
 G_r\otimes E_r\otimes\mathbf k\tau_r
\]
for a finite-dimensional label space $E_r$.  A window-pure family is defined
by replacing $D_\infty$ by finite truncations which reproduce the same
singular values on every prescribed finite output window.
\end{definition}

\begin{theorem}[Qualitative and quantitative response consequences]
\label{rk:thm-response-consequences}
Let $\mathcal R$ be a filtered response.
\begin{enumerate}[label=\textup{(\roman*)},leftmargin=2.8em]
\item Full support in every order is equivalent to strictness.  It gives
      valuation exchange and the flag reconstruction formula
      \eqref{rk:eq-flag-reconstruction}.
\item Coefficient-pure full support gives, for every nonzero probe of top
      depth $r$, at least one labelled native coordinate with a nonzero exact
      critical $r$-packet.  Its singular, Mellin, and Fredholm laws are those
      of \Cref{cc:thm-critical-packet-principle}.
\item Window-pure full support gives the corresponding family-valued
      finite-window lower profile.  It preserves the valuation but does not
      assert one full-tail extremizing realization.
\end{enumerate}
Full support alone need not imply a regularly varying equivalent: it records
nonvanishing in the analytic quotient, while coefficient-purity identifies a
specific critical boundary line.
\end{theorem}

\begin{proof}
Part~\textup{(i)} is
\Cref{rk:thm-valuation-exchange,rk:thm-kernel-support}.  In the
coefficient-pure case, contraction by a nonzero top dual class leaves a
nonzero coefficient of $\tau_r$ in at least one labelled coordinate; apply
the exact part of \Cref{cc:thm-critical-packet-principle}.  The same argument
with truncations and the window clause of that theorem proves part~\textup{(iii)}.
\end{proof}

\begin{proposition}[Independent failure mechanisms]
\label{rk:prop-minimality}
Each part of the architecture has an independent role.
\begin{enumerate}[label=\textup{(\roman*)},leftmargin=2.8em]
\item Without filteredness, an algebraic order-one vector may be sent to
      singular values $N^{-1/2}$, so no finite critical upper grade exists.
\item With filteredness but proper kernel support, every nonzero probe in the
      annihilator of that support is invisible at the correct grade.
\item A full-support kernel may carry an oscillatory boundary representative;
      valuation equality then holds without an asymptotic equivalent.
\item If the grade-one boundary class is nilpotent under tensor product, one
      grade-one certificate cannot propagate to every depth.
\item If labels are discarded, equal singular-value profiles can represent
      distinct coefficient directions, so annihilator reconstruction fails.
\item A universal coefficient-pure kernel says nothing about an unrelated raw
      observable until its native readers are shown to push forward that
      kernel with full support.
\end{enumerate}
\end{proposition}

\begin{proof}
The first two assertions follow from one-dimensional and finite-dimensional
examples.  For~\textup{(iii)}, choose an element of the critical quotient
whose normalized singular values oscillate.  Part~\textup{(iv)} is obtained
by declaring a higher product zero.  For~\textup{(v)}, quotient a
two-dimensional labelled response by the transposition of its coordinates.
For~\textup{(vi)}, prescribe arbitrary compact operators on the raw source
space independently of the registered response.
\end{proof}

\begin{theorem}[Sharp depth-only no-go and determinant blindness]
\label{rk:thm-sharp-depth-no-go}
Fix $r\ge1$.  For every nonzero compact operator $K$ on a separable
Hilbert space $H$, there is a block-lower-triangular operator model of a
nonzero depth-$r$ primitive word whose only nonzero endpoint block is
unitarily equivalent to $K$.  In particular,
\begin{equation}
 s_j(W_r)=s_j(K),\qquad j\ge1.
 \label{rk:eq-arbitrary-depth-spectrum}
\end{equation}
Consequently primitive depth alone forces no Schatten membership, decay
law, regular variation, Mellin order, or Fredholm growth.

More explicitly, on
$\mathcal H=H_0\oplus\cdots\oplus H_r$, with $H_j\cong H$, choose
adjacent compact blocks $X_j:H_{j-1}\to H_j$ and set
$W_1=X_1$, $W_j=[X_j,W_{j-1}]$.  Then $W_r$ has only its $(r,0)$ block and
realizes the free normal word
\[
 \operatorname{ad}_{\mathsf q_r}\operatorname{ad}_{\mathsf q_{r-1}}
 \cdots\operatorname{ad}_{\mathsf q_2}(\mathsf q_1)
 \in\mathfrak L_r(Q^\sharp).
\]
Moreover $W_r^2=0$.  If $K\in\mathcal S_1$, then
\begin{equation}
 \det(I+zW_r)=1,
 \qquad
 \frac12\log\det(I+z^2W_r^*W_r)
 =\frac12\log\det(I+z^2K^*K).
 \label{rk:eq-two-determinants}
\end{equation}
More generally, if
$D=\operatorname{diag}(D_0,\ldots,D_r)$ and $R_\triangle$ are trace
class, with $R_\triangle$ strictly block triangular, then
\begin{equation}
 \det(I+z(D+R_\triangle))
 =\prod_{j=0}^r\det(I+zD_j).
 \label{rk:eq-triangular-determinant-blindness}
\end{equation}
Thus an ordinary Fredholm determinant can erase an entire strictly
triangular ancestry layer, while hermitized and labelled responses retain it.
\end{theorem}

\begin{proof}
Write $K=U|K|$ and take
\[
 A_1=\cdots=A_{r-1}=|K|^{1/r},
 \qquad A_r=U|K|^{1/r},
\]
with $A_1=K$ when $r=1$.  Let $X_j$ have $A_j$ as its only nonzero
adjacent block.  The block orientation gives inductively
$W_j=X_jX_{j-1}\cdots X_1$, because the reversed product vanishes.  Hence
the only nonzero block of $W_r$ is $A_r\cdots A_1=K$, proving
\eqref{rk:eq-arbitrary-depth-spectrum}; the universal Lie map
$\mathsf q_j\mapsto X_j$ realizes the displayed nonzero depth-$r$ word.
The same block support gives $W_r^2=0$.  For trace-class $K$, the nilpotent
trace-class operator $W_r$ has Fredholm determinant one, while the nonzero
spectrum of $W_r^*W_r$ is that of $K^*K$, proving
\eqref{rk:eq-two-determinants}.

Wherever $I+zD$ is invertible, factor
\[
 I+z(D+R_\triangle)
 =(I+zD)\bigl[I+z(I+zD)^{-1}R_\triangle\bigr].
\]
The second factor is identity plus a trace-class strictly triangular
nilpotent, hence has determinant one.  Multiplicativity proves
\eqref{rk:eq-triangular-determinant-blindness} away from the discrete
noninvertibility set of $I+zD$; analyticity extends the identity to every
$z$.
\end{proof}

\section{Canonical stochastic entrance and the critical endpoint}
\label{cc:sec-critical-geometry}

\subsection{The exterior realization of critical area}

The logarithmic normal filtration records primitive ancestry.  The exterior
construction below has a different purpose: it is a multiplicative reader of
one degree--two primitive and, after conditioning, a generating function for
all repeated uses of its covariance.  Exterior particle number is therefore
not a second ancestry filtration.  This distinction lets the same reader
connect causal composition, predictable brackets, Fredholm determinants, and
operator ideals without identifying algebraically different gradings.

Throughout this subsection \(H\) is a real separable Hilbert space.  For vectors
\(u,v\) in a Hilbert space, \(u\otimes v\) denotes the rank--one operator
\(w\mapsto\langle w,v\rangle u\).  We use the exterior convention

\[
 \langle u_1\wedge\cdots\wedge u_j,
         v_1\wedge\cdots\wedge v_j\rangle
 =\det(\langle u_a,v_b\rangle)_{a,b=1}^j.
\]

If \(A\in\mathcal S_2(H)\) is skew-adjoint, let
\(\alpha_A\in\Lambda^2H\) be determined by
\begin{equation}
 \langle Av,u\rangle
 =\langle\alpha_A,u\wedge v\rangle.
 \label{cc:eq-two-form-operator}
\end{equation}

\subsubsection{The multiplicative exterior reader}

Consider the sector-resolved degree--two space
\[
 \mathcal G_2(H)
 :=H\times\mathcal S_2(H)_{\mathrm{skew}}
       \times\mathcal S_1(H)_{\mathrm{sa}}
\]
with product
\begin{equation}
 (x,A,Q)\circ(y,B,R)
 :=\left(x+y,
 A+B+\frac12(x\otimes y-y\otimes x),Q+R\right).
 \label{cc:eq-sector-resolved-product}
\end{equation}
It is a topological group, with
\[
 e=(0,0,0),\qquad (x,A,Q)^{-1}=(-x,-A,-Q).
\]
The third coordinate records the collision sector, while the second
coordinate records chronological antisymmetry.

Put
\[
 \mathcal F_{\wedge}(H):=\bigoplus_{j\geq0}\Lambda^jH
\]
and write
\(\exp_{\wedge}(\alpha)=\sum_{j\geq0}\alpha^{\wedge j}/j!\).
For \(\lambda\in\mathbb R\), define
\begin{equation}
 \Sigma_{\wedge}^{\lambda}(x,A,Q)
 :=(1+\lambda x)\wedge
   \exp_{\wedge}(2\lambda^2\alpha_A).
 \label{cc:eq-exterior-reader}
\end{equation}

\begin{theorem}[Multiplicative exterior area reader]
\label{cc:thm-multiplicative-exterior-reader}
For every \(g,h\in\mathcal G_2(H)\),
\begin{equation}
 \Sigma_{\wedge}^{\lambda}(g\circ h)
 =\Sigma_{\wedge}^{\lambda}(g)\wedge
  \Sigma_{\wedge}^{\lambda}(h).
 \label{cc:eq-exterior-multiplicativity}
\end{equation}
If \(a_1(A)\geq a_2(A)\geq\cdots\) are the positive two--plane
singular values of \(A\), then
\begin{align}
 \mathfrak P_A(z)
 &:=\log\|\exp_{\wedge}(z\alpha_A)\|^2
   =\sum_{j\geq1}\log(1+z^2a_j(A)^2) \notag\\
 &=\frac12\log\det(I+z^2A^*A)
   =\mathfrak F_A(z).
 \label{cc:eq-exterior-profile}
\end{align}
Consequently, for every \(0<p<2\), the weak and little
\(\mathcal S_{p,\infty}\) scales are read directly from \(\mathfrak P_A\)
by \Cref{cc:prop:weak-detection}.  The full vector
\(\Sigma_{\wedge}^{\lambda}\) is faithful in \(A\) for
\(\lambda\ne0\), whereas the scalar profile remembers precisely the nonzero
singular spectrum, with multiplicity, and forgets the singular subspaces.
\end{theorem}

\begin{proof}
For \(x,y\in H\),
\[
 2\alpha_{\frac12(x\otimes y-y\otimes x)}=x\wedge y.
\]
The decomposable two--form \(x\wedge y\) squares to zero and
\((x+y)\wedge x\wedge y=0\).  Since even exterior forms commute, the
right-hand side of \eqref{cc:eq-exterior-multiplicativity} is
\[
 (1+\lambda(x+y))\wedge
 \exp_{\wedge}\!\left(
  2\lambda^2(\alpha_A+\alpha_B)+\lambda^2x\wedge y
 \right),
\]
which is the left-hand side.

An orthogonal two--plane decomposition of a compact skew operator gives
\(
 \alpha_A=\sum_j a_j(A)e_{2j-1}\wedge e_{2j}
\)
up to signs.  The mutually commuting decomposable summands yield
\[
 \|\exp_{\wedge}(z\alpha_A)\|^2
 =\prod_j(1+z^2a_j(A)^2).
\]
Every \(a_j(A)\) occurs twice in the singular-value list of \(A\), proving
\eqref{cc:eq-exterior-profile}.

Finally, \eqref{cc:eq-exterior-profile} identifies
\(\mathfrak P_A\) with the positive Fredholm reader of
\Cref{cc:def:log-spectral-rig}, so the asserted weak and little detection is
exactly the case \(r=1\) of \Cref{cc:prop:weak-detection}.
\end{proof}

\begin{remark}[Reader degree versus ancestry degree]
\label{cc:rem-reader-versus-ancestry}
The degree--two component of
\(\exp_{\wedge}(2\lambda^2\alpha_A)\) is
\(2\lambda^2\alpha_A\), so the exterior vector determines \(A\).
Its higher even components are exterior powers of that same two--form.  They
are not the higher primitive Lie layers of the normal cone; those are
measured by the lower-central filtration of the convolution logarithm.
\end{remark}

\subsubsection{Fermionic superoperators}

Let \(\Omega_{\wedge}=1\in\Lambda^0H\), let \(P_j\) be the projection onto
\(\Lambda^jH\), and identify
\(P_0=\Omega_{\wedge}\otimes\Omega_{\wedge}\) with the vacuum projection.
Right creation and annihilation are
\[
 c^{\dagger}(h)\xi:=\xi\wedge h,
 \qquad c(h):=c^{\dagger}(h)^*.
\]
They satisfy the canonical anticommutation relations.  For
\(Q=\sum_i u_i\otimes v_i\in\mathcal S_1(H)\), define
\begin{equation}
 \mathcal B_Q(T):=\sum_i c^{\dagger}(u_i)T c(v_i),
 \qquad
 \mathrm d\Gamma(Q):=\sum_i c^{\dagger}(u_i)c(v_i).
 \label{cc:eq-fermionic-superoperators}
\end{equation}
For \(Q\succeq0\), put
\begin{equation}
 \Gamma_{\wedge}(Q):=\bigoplus_{j\geq0}\Lambda^jQ.
 \label{cc:eq-fermionic-second-quantization}
\end{equation}
This is the standard CAR calculus used here only as a reader of classical
martingale integrals; see
\cite{BarnettStreaterWilde82,ApplebaumHudson84} for its stochastic-analysis
background.

\begin{lemma}[Fermionic trace-class calculus]
\label{cc:lem-fermionic-calculus}
The operators in \eqref{cc:eq-fermionic-superoperators} do not depend on
the chosen nuclear representation of \(Q\), and
\begin{equation}
 \|\mathrm d\Gamma(Q)\|\leq\|Q\|_1,
 \qquad
 \|\mathcal B_Q(T)\|_1\leq\|Q\|_1\|T\|_1.
 \label{cc:eq-fermionic-trace-bound}
\end{equation}
For \(P,Q\in\mathcal S_1(H)\),
\begin{equation}
 \mathcal B_P\mathcal B_Q=\mathcal B_Q\mathcal B_P.
 \label{cc:eq-fermionic-commutation}
\end{equation}
If \(Q,T\succeq0\), then \(\mathcal B_Q(T)\succeq0\), and
\begin{equation}
 \operatorname{Tr}\mathcal B_Q(T)
 =(\operatorname{Tr}Q)(\operatorname{Tr}T)
  -\operatorname{Tr}(\mathrm d\Gamma(Q)T)
 \leq(\operatorname{Tr}Q)(\operatorname{Tr}T).
 \label{cc:eq-particle-hole-trace}
\end{equation}
Moreover,
\begin{equation}
 \frac1{j!}\mathcal B_Q^j(P_0)=\Lambda^jQ,
 \qquad
 \|\Lambda^jQ\|_1\leq\frac{\|Q\|_1^j}{j!}.
 \label{cc:eq-fermionic-powers}
\end{equation}
Consequently, for \(Q\succeq0\) and \(x\geq0\),
\begin{equation}
 e^{x\mathcal B_Q}(P_0)=\Gamma_{\wedge}(xQ),
 \qquad
 \operatorname{Tr}_{\mathcal F_{\wedge}(H)}\Gamma_{\wedge}(xQ)
 =\det(I+xQ).
 \label{cc:eq-second-quantization-determinant}
\end{equation}
\end{lemma}

\begin{proof}
For rank--one input,
\[
 \|c^{\dagger}(u)c(v)\|\leq\|u\|\|v\|,
 \qquad
 \|c^{\dagger}(u)Tc(v)\|_1
 \leq\|u\|\|v\|\|T\|_1.
\]
The projective tensor description of nuclear operators gives the unique
extensions and \eqref{cc:eq-fermionic-trace-bound}.  For
\(P=a\otimes b\) and \(Q=u\otimes v\), the two CAR interchanges give two
minus signs:
\[
 c^{\dagger}(a)c^{\dagger}(u)T c(v)c(b)
 =c^{\dagger}(u)c^{\dagger}(a)T c(b)c(v).
\]
This proves \eqref{cc:eq-fermionic-commutation} by approximation.

If \(Q=\sum_iq_i e_i\otimes e_i\succeq0\), then
\(
 \mathcal B_Q(T)=\sum_iq_i c(e_i)^*Tc(e_i)\succeq0
\).
Cyclicity of the trace and
\(c(e_i)c^{\dagger}(e_i)=I-c^{\dagger}(e_i)c(e_i)\) give
\eqref{cc:eq-particle-hole-trace}.  Expanding
\(\mathcal B_Q^j(P_0)\), repeated indices vanish.  Every unordered
antisymmetrized \(j\)-tuple occurs \(j!\) times because the signs from the
creation and annihilation strings cancel.  Thus
\(\mathcal B_Q^j(P_0)=j!\Lambda^jQ\).  Trace-norm approximation proves
this for arbitrary trace-class \(Q\), and the norm estimate follows from
\eqref{cc:eq-fermionic-trace-bound}.  Summing the powers proves the first
identity in \eqref{cc:eq-second-quantization-determinant}.  If
\((q_i)\) are the eigenvalues of \(Q\succeq0\), then
\[
 \operatorname{Tr}\Gamma_{\wedge}(xQ)
 =\sum_{j\geq0}\operatorname{Tr}\Lambda^j(xQ)
 =\prod_i(1+xq_i)=\det(I+xQ),
\]
with absolute convergence because \(Q\) is trace class.
\end{proof}

\begin{corollary}[Free capacity and the Fredholm derivative]
\label{cc:cor-free-capacity-Fredholm}
For positive trace-class operators \(Q\) on \(H\) and \(T\) on
\(\mathcal F_{\wedge}(H)\), define the free capacity of \(T\) in the
direction \(Q\) by
\begin{equation}
 \operatorname{Cap}^{\mathrm{free}}_Q(T)
 :=(\operatorname{Tr}Q)(\operatorname{Tr}T)
   -\operatorname{Tr}(\mathrm d\Gamma(Q)T)
 =\operatorname{Tr}\mathcal B_Q(T).
 \label{cc:eq-free-capacity-definition}
\end{equation}
Then
\[
 0\leq\operatorname{Cap}^{\mathrm{free}}_Q(T)
 \leq(\operatorname{Tr}Q)(\operatorname{Tr}T).
\]
Moreover, for positive trace-class \(P,Q\) on \(H\),
\begin{equation}
 \operatorname{Cap}^{\mathrm{free}}_Q(\Gamma_{\wedge}(P))
 =\det(I+P)\operatorname{Tr}((I+P)^{-1}Q).
 \label{cc:eq-free-capacity-Fredholm}
\end{equation}
Consequently,
\begin{equation}
 \frac{\operatorname{Cap}^{\mathrm{free}}_Q(\Gamma_{\wedge}(P))}
      {\operatorname{Tr}_{\mathcal F_{\wedge}(H)}\Gamma_{\wedge}(P)}
 =\left.\frac{\mathrm d}{\mathrm d\varepsilon}\right|_{\varepsilon=0}
   \log\det(I+P+\varepsilon Q)
 =\operatorname{Tr}((I+P)^{-1}Q).
 \label{cc:eq-normalized-free-capacity}
\end{equation}
\end{corollary}

\begin{proof}
The first identity and the two inequalities are
\eqref{cc:eq-particle-hole-trace}.  To prove
\eqref{cc:eq-free-capacity-Fredholm}, first suppose that \(P\) has finite
rank and diagonalize it as \(P=\sum_i p_i e_i\otimes e_i\).  In the
occupation-number basis, \(\Gamma_{\wedge}(P)\) is diagonal.  The
off-diagonal matrix coefficients of \(Q\) therefore make no contribution
to either trace in \eqref{cc:eq-free-capacity-definition}.  For the
diagonal matrix unit \(e_i\otimes e_i\), the free capacity is the sum over
occupation sets not containing \(i\), hence
\[
 \operatorname{Cap}^{\mathrm{free}}_{e_i\otimes e_i}
   (\Gamma_{\wedge}(P))
 =\prod_{j\ne i}(1+p_j)
 =\frac{\det(I+P)}{1+p_i}.
\]
Summing against the diagonal coefficients
\(\langle Qe_i,e_i\rangle\) gives the asserted formula in finite rank.
The estimates in \eqref{cc:eq-fermionic-trace-bound}, together with
\(
 e^{\mathcal B_P}(P_0)=\Gamma_{\wedge}(P)
\), show continuity in the trace norms of \(P\) and \(Q\); finite-rank
compression therefore proves the general case.  Finally the standard
trace-class Fredholm derivative is
\[
 \left.\frac{\mathrm d}{\mathrm d\varepsilon}\right|_{\varepsilon=0}
 \log\det(I+P+\varepsilon Q)
 =\operatorname{Tr}((I+P)^{-1}Q).
\]
Combining this with
\(
 \operatorname{Tr}\Gamma_{\wedge}(P)=\det(I+P)
\)
proves \eqref{cc:eq-normalized-free-capacity}.
\end{proof}

\subsubsection{The exterior lift of chronological area}

Let \((\Omega,\mathcal F,(\mathcal F_t)_{0\leq t\leq T},\mathbb P)\)
satisfy the usual conditions, and let \(M\) be a square-integrable
c\`adl\`ag \(H\)-valued martingale with \(M_0=0\).  Its predictable
operator bracket is denoted by
\[
 R_t=\langle M\rangle_t\in\mathcal S_1(H)_+,
 \qquad V_t=\operatorname{Tr}R_t.
\]
We use the standard Hilbert-space stochastic integral, predictable-bracket
isometry, and stopping calculus in the form of
\cite{MetivierPellaumail80}.
Define
\begin{equation}
 A_t:=\operatorname{Alt}_2
       \int_{(0,t]}M_{s-}\otimes\mathrm dM_s,
 \qquad
 \operatorname{Alt}_2(u\otimes v)
 :=\frac12(u\otimes v-v\otimes u).
 \label{cc:eq-chronological-area}
\end{equation}

\begin{theorem}[Localized exterior lift and exact area realization]
\label{cc:thm-localized-exterior-lift}
Assume first that \(V_T\leq K\) almost surely for a deterministic
\(K<\infty\).  Set
\begin{equation}
 \Xi_0(t)=1,
 \qquad
 \Xi_j(t)=\int_{(0,t]}\Xi_{j-1}(s-)\wedge\mathrm dM_s,
 \qquad
 Z_t^{\lambda,N}=\bigoplus_{j=0}^N\lambda^j\Xi_j(t).
 \label{cc:eq-exterior-iterates}
\end{equation}
Then \(Z^{\lambda,N}\) converges to a c\`adl\`ag
\(\mathcal F_{\wedge}(H)\)-valued martingale \(Z^\lambda\) and
\begin{equation}
 \mathbb E\sup_{t\leq T}
 \|Z_t^{\lambda,N}-Z_t^\lambda\|^2\longrightarrow0,
 \qquad
 \mathbb E\|Z_T^\lambda\|^2\leq e^{\lambda^2K}.
 \label{cc:eq-exterior-energy}
\end{equation}
Indistinguishably in \(t\),
\begin{equation}
 Z_t^\lambda=(1+\lambda M_t)\wedge
 \exp_{\wedge}(2\lambda^2\alpha_{A_t}).
 \label{cc:eq-exact-exterior-area}
\end{equation}
In particular,
\begin{equation}
 \Xi_{2j}(t)=\frac{2^j}{j!}\alpha_{A_t}^{\wedge j},
 \qquad
 \Xi_{2j+1}(t)=\frac{2^j}{j!}
 M_t\wedge\alpha_{A_t}^{\wedge j},
 \label{cc:eq-exterior-level-formulas}
\end{equation}
and, for \(z\geq0\),
\begin{equation}
 \|P_{\mathrm{even}}Z_t^{\sqrt{z/2}}\|^2
 =\exp(\mathfrak F_{A_t}(z)).
 \label{cc:eq-even-profile-identity}
\end{equation}
For an arbitrary square-integrable \(M\), the same assertions hold locally
after a predictable scalar-clock localization.
\end{theorem}

\begin{proof}
Let \(P_{\leq N}=\sum_{j=0}^NP_j\) and
\(T_{s-}^N=Z_{s-}^{\lambda,N}\otimes Z_{s-}^{\lambda,N}\).  Predictable
bracket isometry gives
\begin{equation}
 \mathrm d\langle Z^{\lambda,N}\rangle_s
 =\lambda^2P_{\leq N}
   \mathcal B_{\mathrm dR_s}(T_{s-}^N)P_{\leq N}.
 \label{cc:eq-exterior-Fock-bracket}
\end{equation}
Write \(\mathrm dR_s=r_s\,\mathrm dV_s\), where
\(r_s\succeq0\) and \(\operatorname{Tr}r_s=1\) for
\(\mathrm dV\)-almost every \(s\).  The particle--hole estimate
\eqref{cc:eq-particle-hole-trace} shows that the predictable increasing
part of \(Y^N=\|Z^{\lambda,N}\|^2\) is bounded by
\(
 \lambda^2Y^N_{-}\,\mathrm dV
\).
With
\[
 \mathcal E(\lambda^2V)_t
 :=e^{\lambda^2V_t^c}
   \prod_{0<s\leq t}(1+\lambda^2\Delta V_s),
\]
the left-point Dol\'eans--Gronwall argument, including predictable atoms,
makes
\[
 \mathcal E(\lambda^2V)^{-1}Y^N
\]
a nonnegative local supermartingale after the usual auxiliary localization.
Since
\(
 \mathcal E(\lambda^2V)_T
 \leq e^{\lambda^2V_T}\leq e^{\lambda^2K}
\), Fatou's lemma yields
\(
 \mathbb E\|Z_T^{\lambda,N}\|^2\leq e^{\lambda^2K}
\), uniformly in \(N\).  The cutoffs are compatible, and the exterior
degrees are orthogonal.  Hence their terminal values are Cauchy in \(L^2\),
and Doob's inequality gives \eqref{cc:eq-exterior-energy}.

In finite dimension, the semimartingale It\^o formula applied to the finite
exterior polynomial gives \eqref{cc:eq-exact-exterior-area}.  Continuous
quadratic corrections are symmetric in bracket indices and antisymmetric in
exterior indices, hence vanish.  At a jump \(y=\Delta M_t\),
\[
 2\alpha_{\Delta A_t}=M_{t-}\wedge y,
 \qquad (M_{t-}\wedge y)^{\wedge2}=0,
\]
and direct multiplication gives the exact left-point update
\[
 Z_t^\lambda=Z_{t-}^\lambda\wedge(1+\lambda y).
\]
There is no finite-activity approximation in this calculation.

We spell out the infinite-dimensional passage.  Let $P_n\uparrow I_H$ be
finite-rank orthogonal projections, put $M^{(n)}=P_nM$, and denote the
finite-dimensional iterates by $\Xi_j^{(n)}$.  Naturality of stochastic
integration gives, for $m\ge n$,
\begin{equation}
 (\wedge^jP_n)\Xi_j^{(m)}=\Xi_j^{(n)}
 \quad\text{indistinguishably}.
 \label{cc:eq-exterior-projection-compatibility}
\end{equation}
The finite-dimensional Fock estimate with $\lambda=1$ implies
$\sup_n\mathbb E\|\Xi_j^{(n)}(T)\|^2\le e^K$ for every fixed $j$.
Since $\wedge^jP_n$ is an orthogonal projection,
\[
 \mathbb E\|\Xi_j^{(m)}(T)-\Xi_j^{(n)}(T)\|^2
 =\mathbb E\|\Xi_j^{(m)}(T)\|^2
  -\mathbb E\|\Xi_j^{(n)}(T)\|^2.
\]
The right-hand side tends to zero as $m,n\to\infty$; Doob's inequality
therefore makes $\Xi_j^{(n)}$ Cauchy in
$L^2(\Omega;D([0,T];\Lambda^jH))$ with the supremum norm.  Write its limit as
$\Xi_j$.  Fatou applied to the finite Fock sums yields
\[
 \sum_{j=0}^N\lambda^{2j}\mathbb E\|\Xi_j(T)\|^2
 \le e^{\lambda^2K},
\]
so $\bigoplus_{j\ge0}\lambda^j\Xi_j$ converges in the asserted Fock space
and agrees with the preceding cutoff limit.

At degree two one has
\[
 \Xi_2^{(n)}=2\alpha_{A^{P_nM}}.
\]
Hence the same argument gives convergence of $A^{P_nM}$ in
\[
 L^2\bigl(\Omega;D([0,T];\mathcal S_2(H))\bigr).
\]
Hilbert-space It\^o isometry identifies the limit with the stochastic
integral in \eqref{cc:eq-chronological-area}.  For each fixed exterior
degree, wedge multiplication is continuous on the relevant Hilbert tensor
powers, so the finite-dimensional identities pass in the supremum topology.
A diagonal subsequence gives one common full-measure set for all degrees,
and the c\`adl\`ag versions remove any dependence on a time grid.  Passing
first in each degree and then in the Fock cutoff proves
\eqref{cc:eq-exact-exterior-area}.  Comparing exterior degrees gives
\eqref{cc:eq-exterior-level-formulas}; its even part and
\eqref{cc:eq-exterior-profile} give
\eqref{cc:eq-even-profile-identity}.  Predictable localization of \(V\)
removes the deterministic bound.
\end{proof}

\begin{theorem}[Fock energy and free-capacity balance]
\label{cc:thm-Fock-free-capacity-balance}
Assume in this theorem that \(M\) is continuous and that \(V_T\leq K\)
almost surely for a deterministic \(K<\infty\).  Let \(Z^\lambda\) and
\(Z^{\lambda,N}=P_{\leq N}Z^\lambda\) be the exterior lifts above, and put
\[
 T_t^N=Z_t^{\lambda,N}\otimes Z_t^{\lambda,N},
 \qquad T_t=Z_t^\lambda\otimes Z_t^\lambda.
\]
The finite-level bracket is \eqref{cc:eq-exterior-Fock-bracket}, and its
trace-class limit is the equality of positive
\(\mathcal S_1(\mathcal F_{\wedge}(H))\)-valued measures
\begin{equation}
 \mathrm d\langle Z^\lambda\rangle_t
 =\lambda^2\mathcal B_{\mathrm dR_t}(T_t).
 \label{cc:eq-Fock-bracket-limit}
\end{equation}
More explicitly, write
\(
 \mathrm dR_t=r_t\,\mathrm dV_t
\), where \(r_t\succeq0\) and
\(\operatorname{Tr}r_t=1\) for \(\mathrm dV\)-almost every \(t\), and set
\[
 Y_t=\|Z_t^\lambda\|^2,
 \qquad
 L_t=2\int_0^t
      \langle Z_s^\lambda,\mathrm dZ_s^\lambda\rangle.
\]
Then \(L\) is a real continuous local martingale and
\begin{align}
 \mathrm dY_t
 &=\mathrm dL_t
   +\lambda^2\operatorname{Cap}^{\mathrm{free}}_{r_t}(T_t)
      \,\mathrm dV_t \notag\\
 &=\mathrm dL_t
   +\lambda^2\left(
      Y_t-\operatorname{Tr}(\mathrm d\Gamma(r_t)T_t)
     \right)\mathrm dV_t,
 \label{cc:eq-Fock-free-energy}\\
 \mathrm d[L]_t
 &\leq4\lambda^2Y_t
   \operatorname{Cap}^{\mathrm{free}}_{r_t}(T_t)
   \,\mathrm dV_t.
 \label{cc:eq-Fock-free-energy-bracket}
\end{align}
Thus the predictable growth of Fock energy is exactly the free capacity;
the occupied part is removed, rather than estimated away.
\end{theorem}

\begin{proof}
For finite \(N\), Hilbert-space It\^o isometry gives
\eqref{cc:eq-exterior-Fock-bracket}.  We justify passage to the limit as an
identity of trace-class measures.  The rank-one estimate
\[
 \|u\otimes u-v\otimes v\|_1
 \leq(\|u\|+\|v\|)\|u-v\|
\]
and \eqref{cc:eq-exterior-energy}, together with Doob's inequality, imply
\begin{equation}
 \mathbb E\sup_{t\leq T}\|T_t^N-T_t\|_1\longrightarrow0.
 \label{cc:eq-Fock-rank-one-convergence}
\end{equation}
Using \(\mathrm dR_t=r_t\,\mathrm dV_t\),
\eqref{cc:eq-fermionic-trace-bound}, and \(V_T\leq K\), we obtain
\[
 \mathbb E\int_0^T
 \|\mathcal B_{r_t}(T_t^N-T_t)\|_1\,\mathrm dV_t
 \leq K\,\mathbb E\sup_{t\leq T}\|T_t^N-T_t\|_1
 \longrightarrow0.
\]
For every trace-class \(S\),
\(
 P_{\leq N}SP_{\leq N}\to S
\)
in trace norm and
\(
 \|P_{\leq N}SP_{\leq N}-S\|_1\leq2\|S\|_1
\).
Taking \(S=\mathcal B_{r_t}(T_t)\), using
\(
 \|\mathcal B_{r_t}(T_t)\|_1\leq Y_t
\), and applying dominated convergence for the finite measure
\[
 \mu(B):=\mathbb E\int_0^T
          \mathbf1_B(\omega,t)\,\mathrm dV_t(\omega)
\]
show that the right-hand sides of the
finite bracket identities converge in expected trace-class variation to
the right-hand side of \eqref{cc:eq-Fock-bracket-limit}.  On the other
hand, \(Z^{\lambda,N}\to Z^\lambda\) in
\(L^2(\Omega;C([0,T];\mathcal F_{\wedge}(H)))\).  Passing in the scalar
martingale product identities and using uniqueness of the predictable
operator bracket identifies the limit with
\(\mathrm d\langle Z^\lambda\rangle\), proving
\eqref{cc:eq-Fock-bracket-limit}.

Continuous Hilbert-space It\^o's formula gives
\[
 \mathrm dY_t=\mathrm dL_t+
 \operatorname{Tr}(\mathrm d\langle Z^\lambda\rangle_t).
\]
Now take the Fock trace in \eqref{cc:eq-Fock-bracket-limit} and use
\eqref{cc:eq-free-capacity-definition}.  Since
\(\operatorname{Tr}T_t=Y_t\) and \(\operatorname{Tr}r_t=1\), this is
exactly \eqref{cc:eq-Fock-free-energy}.  Finally, Cauchy--Schwarz for the
stochastic-integral bracket of
\(2\int\langle Z^\lambda,\mathrm dZ^\lambda\rangle\) gives
\[
 \mathrm d[L]_t
 \leq4Y_t\operatorname{Tr}
       (\mathrm d\langle Z^\lambda\rangle_t).
\]
Substitution of \eqref{cc:eq-Fock-bracket-limit} and
\eqref{cc:eq-free-capacity-definition} proves
\eqref{cc:eq-Fock-free-energy-bracket}.
\end{proof}

\subsection{The forced weak-trace endpoint}
\label{cc:subsec-critical-split-endpoint}

Let \(H\) be a real separable Hilbert space.  Singular values are listed
nonincreasingly and with multiplicity.  We use
\[
 \|T\|_{1,\infty}:=\sup_{n\geq1}n s_n(T),
 \qquad
 \mathcal S^0_{1,\infty}(H)
 :=\{T\in\mathcal S_{1,\infty}(H):ns_n(T)\longrightarrow0\},
\]
and, for \(0<\rho<\infty\),
\[
 \|T\|_{1,\rho}^{\rho}
 :=\sum_{n\geq1}\frac{(n s_n(T))^{\rho}}{n}.
\]
The little weak ideal is precisely the
\(\|\cdot\|_{1,\infty}\)-closure of the finite-rank operators.  We shall
use only the standard completeness, ideal, and compression properties of
these symmetric ideals; see
\cite{Pietsch80,Simon05,LordSukochevZanin12}.

For \(x,y\in H\), write
\[
 x\wedge_{\mathrm{op}}y:=x\otimes y-y\otimes x.
\]
If \(E\subset H\) is finite dimensional, set
\[
 \mathscr G_{\mathrm{fin}}(E)
 :=E\times\mathcal L(E)_{\mathrm{skew}}
       \times\mathcal L(E)_{\mathrm{sa}},
 \qquad
 \mathscr G_{\mathrm{fin}}(H)
 :=\bigcup_{\dim E<\infty}\mathscr G_{\mathrm{fin}}(E),
\]
where the union is directed by isometric inclusion.  On this core define
\begin{equation}
 (x,A,Q)\circ(y,B,R)
 :=\left(x+y,A+B+\frac12x\wedge_{\mathrm{op}}y,Q+R\right).
 \label{cc:eq-critical-split-product}
\end{equation}

\begin{theorem}[Critical split completion]
\label{cc:thm-critical-split-completion}
The completion of \(\mathscr G_{\mathrm{fin}}(H)\) for the coordinate
quasi-gauge
\begin{equation}
 \|(x,A,Q)\|_{\mathrm{crit}}
 :=\|x\|_H+\|A\|_{1,\infty}+\|Q\|_1
 \label{cc:eq-critical-quasi-gauge}
\end{equation}
is canonically
\begin{equation}
 \mathscr G_{\mathrm{crit}}(H)
 =H\times\mathcal S^0_{1,\infty}(H)_{\mathrm{skew}}
       \times\mathcal S_1(H)_{\mathrm{sa}}.
 \label{cc:eq-critical-split-space}
\end{equation}
The product \eqref{cc:eq-critical-split-product} and inverse
\[
 (x,A,Q)^{-1}=(-x,-A,-Q)
\]
extend continuously and make this space a separable complete topological
group.  More precisely, there are an exponent
\(0<\vartheta\leq1\) and an equivalent quasi-norm
\(\|\cdot\|_{1,\infty,*}\) on
\(\mathcal S^0_{1,\infty}(H)\) such that
\[
 \|S+T\|_{1,\infty,*}^{\vartheta}
 \leq \|S\|_{1,\infty,*}^{\vartheta}
      +\|T\|_{1,\infty,*}^{\vartheta}.
\]
Consequently
\begin{align}
 d_{\mathrm{crit}}(g,h)
 :={}&\|x_g-x_h\|_H^{\vartheta}
 +\|A_g-A_h\|_{1,\infty,*}^{\vartheta}
 +\|Q_g-Q_h\|_1^{\vartheta}
 \label{cc:eq-critical-Aoki-metric}
\end{align}
is a compatible complete metric.  In particular,
\(\mathscr G_{\mathrm{crit}}(H)\) is Polish; no Banach-space assertion is
being made for its weak-ideal coordinate.

The physical sector
\begin{equation}
 \mathscr G^+_{\mathrm{crit}}(H)
 :=\{(x,A,Q)\in\mathscr G_{\mathrm{crit}}(H):Q\succeq0\}
 \label{cc:eq-critical-positive-sector}
\end{equation}
is a closed submonoid.  For every bounded \(L:H\to H'\),
\[
 L_{\#}(x,A,Q):=(Lx,LAL^*,LQL^*)
\]
is a continuous homomorphism, and the homogeneous dilations
\[
 \delta_a(x,A,Q):=(ax,a^2A,a^2Q),\qquad a\in\mathbb R,
\]
are continuous.
\end{theorem}

\begin{proof}
The weak Schatten quasi-triangle inequality follows from
\[
 s_{m+n-1}(S+T)\leq s_m(S)+s_n(T).
\]
It gives a universal quasi-triangle constant for
\(\|\cdot\|_{1,\infty}\).  The Aoki--Rolewicz renorming theorem
\cite{KaltonPeckRoberts84} therefore supplies \(\vartheta\) and
\(\|\cdot\|_{1,\infty,*}\) with the displayed
power triangle inequality.  This use of Aoki--Rolewicz is essential:
\(\mathcal S_{1,\infty}\) is not being treated as a Banach ideal.

The space \(\mathcal S^0_{1,\infty}\), being the weak-quasi-norm closure
of the finite-rank operators in the complete weak ideal, is complete.
Finite-rank operators with rational matrix entries relative to a fixed
countable dense subset of \(H\) form a countable dense family.  Thus this
coordinate is separable; the same is standard for \(H\) and
\(\mathcal S_1\).  The power triangle inequalities in the three
coordinates show directly that \eqref{cc:eq-critical-Aoki-metric} is a
complete metric inducing the product topology.  This also identifies the
completion of the directed finite-dimensional core with
\eqref{cc:eq-critical-split-space}.

For rank-one operators,
\begin{equation}
 \|x\wedge_{\mathrm{op}}y\|_{1,\infty}
 \leq\|x\wedge_{\mathrm{op}}y\|_1
 \leq2\|x\|_H\|y\|_H.
 \label{cc:eq-wedge-trace-bound}
\end{equation}
Hence the only nonlinear coordinate of \eqref{cc:eq-critical-split-product}
is continuous.  Indeed,
\[
 x_n\wedge_{\mathrm{op}}y_n-x\wedge_{\mathrm{op}}y
 =(x_n-x)\wedge_{\mathrm{op}}y_n
   +x\wedge_{\mathrm{op}}(y_n-y),
\]
and convergent Hilbert-space sequences are bounded.  Associativity follows
from bilinearity:
\[
 x\wedge y+(x+y)\wedge z
 =y\wedge z+x\wedge(y+z).
\]
The inverse formula is immediate because
\(x\wedge_{\mathrm{op}}x=0\).

The ideal property gives
\[
 \|LAL^*\|_{1,\infty}\leq\|L\|^2\|A\|_{1,\infty},
 \qquad
 \|LQL^*\|_1\leq\|L\|^2\|Q\|_1,
\]
which proves the pushforward statement; the dilation assertion is
coordinatewise.  Finally the positive cone is trace-norm closed, and the
third coordinate is additive under \(\circ\), so
\eqref{cc:eq-critical-positive-sector} is a closed submonoid.
\end{proof}

The critical topology is the exponent-one member of a natural family.  We
record the full family because the exterior profile detects every exponent
below the ambient Hilbert--Schmidt scale.

\begin{corollary}[Sub-Hilbert weak split completions]
\label{cc:cor-sub-Hilbert-split-completions}
For \(0<p<2\), equip \(\mathscr G_{\mathrm{fin}}(H)\) with
\[
 \|(x,A,Q)\|_{\mathrm{split},p}
 :=\|x\|_H+\|A\|_{p,\infty}+\|Q\|_1,
 \qquad
 \|A\|_{p,\infty}:=\sup_{n\geq1}n^{1/p}s_n(A).
\]
Write
\[
 \mathcal S^0_{p,\infty}(H)
 :=\{A\in\mathcal S_{p,\infty}(H):n^{1/p}s_n(A)\to0\}.
\]
Its completion is canonically
\begin{equation}
 \mathscr G^{(p)}_0(H)
 :=H\times\mathcal S^0_{p,\infty}(H)_{\mathrm{skew}}
       \times\mathcal S_1(H)_{\mathrm{sa}}.
 \label{cc:eq-sub-Hilbert-split-completion}
\end{equation}
The product \eqref{cc:eq-critical-split-product}, inverse, dilations, and
bounded pushforwards extend continuously, and
\[
 \mathscr G^{(p)}_+(H)
 :=\{(x,A,Q)\in\mathscr G^{(p)}_0(H):Q\succeq0\}
\]
is a closed submonoid.  An equivalent Aoki--Rolewicz power metric makes
each \(\mathscr G^{(p)}_0(H)\) a Polish group.  At \(p=1\) this is exactly
\(\mathscr G_{\mathrm{crit}}(H)\) with its previously defined topology.
\end{corollary}

\begin{proof}
Repeat the proof of \Cref{cc:thm-critical-split-completion} with
\(\mathcal S^0_{p,\infty}\) in place of
\(\mathcal S^0_{1,\infty}\).  The only new estimate is the rank-two bound
\[
 \|x\wedge_{\mathrm{op}}y\|_{p,\infty}
 \leq2^{1+1/p}\|x\|_H\|y\|_H,
\]
which gives continuity of the nonlinear group term.  Completeness,
separability, the ideal inequalities, Aoki--Rolewicz metrization, and
closedness of the positive cone are then identical.
\end{proof}

\begin{proposition}[Universal trace obstruction]
\label{cc:prop-critical-trace-obstruction}
Let \(\mathcal E\) be a symmetric quasi-normed ideal on the finite-rank
operators.  If the universal positive collision readout is continuous in
the sense that
\[
 \operatorname{Tr}S\leq C\|S\|_{\mathcal E}
 \qquad(S\succeq0,\ S\text{ finite rank}),
\]
then
\[
 \|S\|_1\leq C\|S\|_{\mathcal E}
 \qquad(S\succeq0,\ S\text{ finite rank}).
\]
Thus any universal topology on the physical collision cone has at least
trace-class strength.  In particular Hilbert--Schmidt control is
insufficient in infinite dimension.
\end{proposition}

\begin{proof}
For a positive finite-rank operator,
\(\|S\|_1=\operatorname{Tr}S\), which proves the first assertion.  If
\(P_N\) is an orthogonal projection of rank \(N\), then
\[
 \|P_N\|_2=N^{1/2},\qquad \operatorname{Tr}P_N=N.
\]
No dimension-free trace bound can therefore follow from the
Hilbert--Schmidt norm.
\end{proof}

\paragraph{The universal upper bound.}
Let \(M\) be a continuous square-integrable \(H\)-valued martingale with
\(M_0=0\).  Put
\[
 R_t:=\langle\!\langle M\rangle\!\rangle_t=[M]_t,
 \qquad V_t:=\operatorname{Tr}R_t,
\]
and
\begin{equation}
 A_{s,t}:=\operatorname{Alt}_2
 \int_s^t(M_u-M_s)\otimes\mathop{}\!\mathrm dM_u,
 \qquad A_t:=A_{0,t}.
 \label{cc:eq-critical-martingale-area}
\end{equation}
The integral is first an \(H^{\otimes2}\)-valued It\^o integral; under the
Hilbert tensor--operator identification \(A_{s,t}\) is skew-adjoint and
Hilbert--Schmidt.

We use the exterior spectral profile \(\mathfrak P_A=\mathfrak F_A\)
from \eqref{cc:eq-exterior-profile}; no second degree-two spectral reader is
introduced here.

For a positive trace-class operator \(Q\), also put
\begin{equation}
 \Psi_Q(x):=\log\det(I+xQ),\qquad x\geq0.
 \label{cc:eq-positive-determinant-profile}
\end{equation}

\begin{lemma}[Determinant and conditional profile decoder]
\label{cc:lem-response-profile-decoder}
The following statements hold.
\begin{enumerate}[label=\textup{(\roman*)},leftmargin=2.8em]
\item If \(0<p<1\) and \(Q\succeq0\) is trace class, then
\begin{equation}
 \sup_{x>0}x^{-p}\Psi_Q(x)
 \asymp_p\|Q\|_{\mathcal S_{p,\infty}}^p,
 \qquad
 Q\in\mathcal S^0_{p,\infty}
 \Longleftrightarrow \Psi_Q(x)=o(x^p).
 \label{cc:eq-determinant-weak-dictionary}
\end{equation}
Moreover \(Q\in\mathcal S_1\) implies \(\Psi_Q(x)=o(x)\).

\item Let \(0<p<2\), let \((A_\theta)\) be a measurable family of skew
Hilbert--Schmidt operators, and put
\[
 \Phi(z):=\sup_\theta\mathfrak P_{A_\theta}(z),
 \qquad
 W_p:=\sup_{\theta,j}j^{1/p}a_j(A_\theta).
\]
Let \(\mathcal H\) be a sigma-field.  Suppose that, for an
\(\mathcal H\)-measurable \(B\geq0\), some \(q>0\), and deterministic
\(C_0,C_1<\infty\),
\begin{equation}
 \mathbb E[e^{q\Phi(z)}\mid\mathcal H]
 \leq C_0\exp(C_1(zB)^p),\qquad z\geq0.
 \label{cc:eq-maximal-profile-hypothesis}
\end{equation}
Then, for every \(0<r<2q\),
\begin{equation}
 \left(\mathbb E[W_p^r\mid\mathcal H]\right)^{1/r}
 \leq C_{p,q,r,C_0,C_1}B.
 \label{cc:eq-maximal-profile-moment}
\end{equation}

\item Under the notation of part~\textup{(ii)}, suppose instead that
\begin{equation}
 \mathbb E[e^{\Phi(z)}\mid\mathcal H]
 \leq C_0e^{b(z)},
 \qquad b(z)=o(z^p)\quad\text{almost surely},
 \label{cc:eq-maximal-profile-little-hypothesis}
\end{equation}
where \(b(z)\) is \(\mathcal H\)-measurable.  Then
\begin{equation}
 \sup_\theta n^{1/p}s_n(A_\theta)\longrightarrow0
 \quad\text{almost surely}.
 \label{cc:eq-maximal-profile-little-conclusion}
\end{equation}
\end{enumerate}
\end{lemma}

\begin{proof}
Since
\(
 \Psi_Q(x)=2\mathfrak F_{Q^{1/2}}(\sqrt{x})
\), part~\textup{(i)} follows from
\Cref{cc:prop:weak-detection} applied to \(Q^{1/2}\) with exponent
\(2p\).  If \((q_j)\) are the eigenvalues of \(Q\in\mathcal S_1\), then
\(x^{-1}\log(1+xq_j)\leq q_j\); dominated convergence gives
\(\Psi_Q(x)=o(x)\).

For part~\textup{(ii)}, on
\[
 E_j(\lambda):=
 \{\sup_\theta j^{1/p}a_j(A_\theta)>\lambda B\}
\]
choose \(z_j=j^{1/p}/B\) on \(\{B>0\}\).  The first \(j\) positive
blocks give
\(
 \Phi(z_j)\geq j\log(1+\lambda^2)
\).
Conditional Markov inequality and
\eqref{cc:eq-maximal-profile-hypothesis} yield
\[
 \mathbb P(E_j(\lambda)\mid\mathcal H)
 \leq C_0\exp\{C_1j-qj\log(1+\lambda^2)\}.
\]
The use of the \(\mathcal H\)-measurable value \(z_j\) follows first for
simple values and then by monotone approximation and continuity in \(z\).
Summing in \(j\) gives a conditional
\(O(\lambda^{-2q})\) tail for \(W_p/B\), which integrates to
\eqref{cc:eq-maximal-profile-moment} for \(r<2q\).  The case \(B=0\)
follows by letting deterministic \(z\to\infty\).

For part~\textup{(iii)}, take \(z_j=j^{1/p}\).  For every
\(\varepsilon>0\), conditional Markov inequality bounds
\(
 \mathbb P\{\sup_\theta j^{1/p}a_j(A_\theta)>\varepsilon
 \mid\mathcal H\}
\)
by
\[
 C_0\exp\{b(j^{1/p})-j\log(1+\varepsilon^2)\},
\]
which is summable almost surely.  Conditional Borel--Cantelli, followed by
a countable sequence \(\varepsilon\downarrow0\), proves
\eqref{cc:eq-maximal-profile-little-conclusion}.
\end{proof}

\begin{proposition}[Universal weak-trace area bound]
\label{cc:prop-universal-critical-area-bound}
If \(V_T\leq K\) almost surely, then for \(q>1\) and \(z\geq0\),
\begin{equation}
 \mathbb E\exp\!\left(q\sup_{t\leq T}\mathfrak P_{A_t}(z)\right)
 \leq\left(\frac q{q-1}\right)^q
 \exp\!\left(\frac{q(2q-1)}2zK\right).
 \label{cc:eq-critical-profile-maximal}
\end{equation}
Consequently, for every finite \(r>0\),
\begin{equation}
 \left\|\sup_{t\leq T}\|A_t\|_{1,\infty}\right\|_{L^r}
 \leq C_rK,
 \label{cc:eq-universal-weak-area-bound}
\end{equation}
where \(C_r\) is independent of \(H,M,K\), and \(T\).
\end{proposition}

\begin{proof}
We first work in finite dimension, so that no Banach-space theorem is
applied to the weak ideal.  Let \(Z^\lambda\) be the exterior It\^o lift.
The exact lift theorem
\Cref{cc:thm-localized-exterior-lift}, together with
\eqref{cc:eq-even-profile-identity}, gives
\[
 \|P_{\mathrm{even}}Z_t^{\sqrt{z/2}}\|^2
 =e^{\mathfrak P_{A_t}(z)}.
\]
The full closed form and its degreewise It\^o verification belong to that
reader theorem; here we use only its even-profile identity and the
trace-class CAR bounds in \Cref{cc:lem-fermionic-calculus}.

Let \(Y_t=\|Z_t^\lambda\|^2\).  Writing
\(\mathop{}\!\mathrm dR_t=r_t\mathop{}\!\mathrm dV_t\), with
\(r_t\succeq0\) and \(\operatorname{Tr}r_t=1\) on the support of
\(\mathop{}\!\mathrm dV\), note that \(Y_t\geq1\) because the vacuum
component of \(Z^\lambda\) is one.  Define the predictable density
\[
 a_t^{\mathrm{free}}
 :=\frac{\operatorname{Tr}\mathcal B_{r_t}
        (Z_t^\lambda\otimes Z_t^\lambda)}{Y_t},
 \qquad
 \mathop{}\!\mathrm dV_t^{\mathrm{free}}
 :=a_t^{\mathrm{free}}\mathop{}\!\mathrm dV_t.
\]
The particle--hole identity gives
\(0\leq a_t^{\mathrm{free}}\leq1\).  Since \(M\), hence
\(Z^\lambda\), is continuous, its optional and predictable brackets
coincide.  Hilbert-space It\^o's formula, with
\[
 \mathop{}\!\mathrm dL_t
 :=2\operatorname{Re}
   \langle Z_t^\lambda,\mathop{}\!\mathrm dZ_t^\lambda\rangle,
\]
therefore gives
\[
 \mathop{}\!\mathrm dY_t
 =\mathop{}\!\mathrm dL_t
   +\lambda^2Y_t\mathop{}\!\mathrm dV_t^{\mathrm{free}},
 \qquad
 0\leq\mathop{}\!\mathrm dV_t^{\mathrm{free}}
 \leq\mathop{}\!\mathrm dV_t,
\]
where \(L\) is a real local martingale.  Cauchy--Schwarz for the
stochastic-integral bracket gives, pointwise as measures,
\[
 \mathop{}\!\mathrm d[L]_t
 \leq4\lambda^2Y_t^2\mathop{}\!\mathrm dV_t^{\mathrm{free}}.
\]
Thus the drift in It\^o's formula for \(Y^q\) is at most
\[
 q\lambda^2Y_t^q\mathop{}\!\mathrm dV_t^{\mathrm{free}}
 +\frac{q(q-1)}2Y_t^{q-2}\mathop{}\!\mathrm d[L]_t
 \leq q(2q-1)\lambda^2Y_t^q
       \mathop{}\!\mathrm dV_t^{\mathrm{free}}.
\]
For \(1<q<2\), this statement is obtained first for
\((Y+\epsilon)^q\) and then by the standard localized
\(\epsilon\downarrow0\) argument.
Thus
\[
 e^{-q(2q-1)\lambda^2V_t^{\mathrm{free}}}Y_t^q
\]
is a local supermartingale.  Localization and \(V_T\leq K\) give
\[
 \mathbb EY_T^q\leq e^{q(2q-1)\lambda^2K}.
\]
Since \(Y\) is a nonnegative submartingale, Doob's inequality yields
\[
 \mathbb E\sup_{t\leq T}Y_t^q
 \leq\left(\frac q{q-1}\right)^q
 e^{q(2q-1)\lambda^2K}.
\]
Taking \(\lambda^2=z/2\) and retaining the even component proves
\eqref{cc:eq-critical-profile-maximal} in finite dimension.

For general \(H\), apply the estimate to increasing finite-rank
projections \(P_nM\).  The projected areas converge in
\(L^2(\Omega;C([0,T];\mathcal S_2))\) after the usual scalar-clock
localization, by the Hilbert-space It\^o isometry and BDG inequality
\cite{MetivierPellaumail80}.  From every subsequence one may therefore
extract a further subsequence converging almost surely, uniformly in
\(t\), in \(\mathcal S_2\).  The estimate
\[
 \|A^*A-B^*B\|_1
 \leq(\|A\|_2+\|B\|_2)\|A-B\|_2
\]
shows uniform trace-norm convergence of the squared moduli.
Consequently, for each fixed \(z\), continuity of the Fredholm
determinant on \(\mathcal S_1\) gives
\[
 \sup_{t\leq T}
 |\mathfrak P_{A_t^{P_nM}}(z)-\mathfrak P_{A_t^M}(z)|
 \longrightarrow0
\]
along that further subsequence.  Fatou's lemma applied to
\(\exp(q\sup_t\mathfrak P_{A_t}(z))\) proves
\eqref{cc:eq-critical-profile-maximal}; the subsequence principle removes
the extraction.  Any auxiliary local-martingale localization is removed by
Fatou's lemma; the energy is already bounded by \(K\).  Finally apply
\Cref{cc:lem-response-profile-decoder} with \(p=1\) and choose \(q>r/2\).
This proves \eqref{cc:eq-universal-weak-area-bound} without using linear
expectation in a quasi-Banach space.
\end{proof}

\begin{corollary}[Bounded bracket-density control]
\label{cc:cor-critical-bracket-density-control}
Assume
\begin{equation}
 \mathop{}\!\mathrm dV_t\leq\Lambda\,\mathop{}\!\mathrm dt
 \qquad(0\leq t\leq T).
 \label{cc:eq-critical-bracket-density}
\end{equation}
For every finite \(q\geq2\) and \(0\leq s<t\leq T\),
\begin{align}
 \left\|\sup_{s\leq u\leq t}\|M_u-M_s\|_H\right\|_{L^q}
 &\leq C_q\Lambda^{1/2}|t-s|^{1/2},
 \label{cc:eq-critical-local-first}\\
 \left\|\sup_{s\leq u\leq t}\|A_{s,u}\|_{1,\infty}\right\|_{L^q}
 &\leq C_q\Lambda|t-s|.
 \label{cc:eq-critical-local-area}
\end{align}
For every \(\beta\in(1/3,1/2)\) there is a modification and an almost
surely finite random \(C_\beta\) such that, simultaneously for all
\(s<t\),
\begin{equation}
 \|M_t-M_s\|_H\leq C_\beta|t-s|^\beta,
 \qquad
 \|A_{s,t}\|_{1,\infty}\leq C_\beta|t-s|^{2\beta}.
 \label{cc:eq-critical-local-holder}
\end{equation}
Moreover, for every finite \(p\), the constant may be chosen so that
\begin{equation}
 \|C_\beta\|_{L^p}
 \leq C_{p,\beta,T}\bigl(\Lambda^{1/2}+\Lambda\bigr),
 \label{cc:eq-critical-holder-moments}
\end{equation}
where the constant is independent of the Hilbert-space dimension.  Thus
bounded bracket density is already a dimension-free grade-one entrance
certificate; no model-specific dyadic argument is required later.
\end{corollary}

\begin{proof}
The first estimate is Hilbert-space BDG\@.  The shifted martingale
\(u\mapsto M_{s+u}-M_s\) has scalar bracket at most
\(\Lambda(t-s)\), so
\eqref{cc:eq-critical-local-area} follows from
\eqref{cc:eq-universal-weak-area-bound}.

Fix \(p<\infty\) and \(\beta<1/2\).  Choose
\(q\geq2p\) so large that
\[
 q(1/2-\beta)>1,
 \qquad q(1-2\beta)>1.
\]
For adjacent dyadic intervals \(I\subset[0,T]\), put
\[
 D_X:=\sup_I |I|^{-\beta}\sup_{u\in I}\|M_{I^-,u}\|,
 \qquad
 D_A:=\sup_I |I|^{-2\beta}\sup_{u\in I}
       \|A_{I^-,u}\|_{1,\infty}.
\]
Using \((\sup_jY_j)^q\leq\sum_jY_j^q\), followed by
\eqref{cc:eq-critical-local-first}--\eqref{cc:eq-critical-local-area},
and summing over the \(2^k\) intervals at level \(k\), gives
\[
 \|D_X\|_{L^q}\leq C_{q,\beta,T}\Lambda^{1/2},
 \qquad
 \|D_A\|_{L^q}\leq C_{q,\beta,T}\Lambda.
\]
The canonical dyadic decomposition of an arbitrary interval uses at most
two intervals at each scale.  Additivity of the first level and Chen's
formula
\[
 A_{s,t}=\sum_\ell A_{I_\ell}
 +\frac12\sum_{\ell<j}M_{I_\ell}\wedge_{\mathrm{op}}M_{I_j}
\]
therefore imply
\[
 \sup_{s<t}\frac{\|M_{s,t}\|}{|t-s|^\beta}
 \lesssim_\beta D_X,
 \qquad
 \sup_{s<t}\frac{\|A_{s,t}\|_{1,\infty}}{|t-s|^{2\beta}}
 \lesssim_\beta D_A+D_X^2.
\]
The first sum uses the Aoki--Rolewicz power inequality from
\Cref{cc:thm-critical-split-completion}; the wedge sum is trace class by
\eqref{cc:eq-wedge-trace-bound}.  Taking \(L^p\)-norms and using
\(q\ge2p\) proves \eqref{cc:eq-critical-holder-moments}, hence also
\eqref{cc:eq-critical-local-holder}.
\end{proof}

\begin{theorem}[Critical stability and projection convergence]
\label{cc:thm-critical-projection-stability}
Let \(X,Z\) be continuous square-integrable \(H\)-valued martingales with
\(X_0=Z_0=0\), and assume, almost surely,
\[
 \operatorname{Tr}\langle\!\langle X\rangle\!\rangle_T\leq E,
 \qquad
 \operatorname{Tr}\langle\!\langle Z\rangle\!\rangle_T\leq\delta^2.
\]
For every finite \(r>0\),
\begin{align}
 \left\|\sup_{t\leq T}\|Z_t\|_H\right\|_{L^r}
 &\leq C_r\delta,
 \label{cc:eq-critical-stability-first}\\
 \left\|\sup_{t\leq T}
 \|A_t^{X+Z}-A_t^X\|_{1,\infty}\right\|_{L^r}
 &\leq C_r(\sqrt E\,\delta+\delta^2),
 \label{cc:eq-critical-stability-area}\\
 \sup_{t\leq T}
 \|\langle\!\langle X+Z\rangle\!\rangle_t
     -\langle\!\langle X\rangle\!\rangle_t\|_1
 &\leq2\sqrt E\,\delta+\delta^2.
 \label{cc:eq-critical-stability-bracket}
\end{align}
The constants are independent of the Hilbert-space dimension.

If \(P_n\to I\) strongly are increasing finite-rank orthogonal
projections and \(M\) is any continuous square-integrable Hilbert
martingale with \(M_0=0\), then its split anchors
\[
 \mathbf M_t^{(n)}
 :=\bigl(P_nM_t,A_t^{P_nM},
        \langle\!\langle P_nM\rangle\!\rangle_t\bigr)
\]
converge to a continuous
\(\mathscr G^+_{\mathrm{crit}}(H)\)-valued anchor \(\mathbf M\), uniformly
on compact time intervals in probability for \(d_{\mathrm{crit}}\).
Moreover, on one event of probability one,
\begin{equation}
 A_t^M\in\mathcal S^0_{1,\infty}(H)_{\mathrm{skew}}
 \qquad\text{for every }t\leq T.
 \label{cc:eq-individual-little-weak}
\end{equation}
Thus every individual continuous martingale realization belongs to the
little weak core; this is a samplewise assertion and does not give a
dimension-uniform decay modulus.
\end{theorem}

\begin{proof}
Equation \eqref{cc:eq-critical-stability-first} is BDG\@.  Write
\[
 A^{X+Z}-A^X=\mathcal A(X,Z)+A^Z,
 \qquad
 \mathcal A(X,Z)
 =\frac{A^{X+\lambda Z}-A^{X-\lambda Z}}{2\lambda}.
\]
Positivity of the joint operator bracket and Schatten Cauchy--Schwarz give
\[
 \operatorname{Tr}\langle\!\langle X\pm\lambda Z
 \rangle\!\rangle_T
 \leq(\sqrt E+\lambda\delta)^2.
\]
Apply \eqref{cc:eq-universal-weak-area-bound} to \(X\pm\lambda Z\), use
the weak-ideal quasi-triangle inequality, and divide by \(2\lambda\).
The resulting bound is
\[
 C_r(E/\lambda+\sqrt E\,\delta+\lambda\delta^2).
\]
Taking \(\lambda=\sqrt E/\delta\), with the cases \(E=0\) or
\(\delta=0\) obtained by a limit, bounds the mixed term by
\(C_r\sqrt E\,\delta\).  A further application of
\eqref{cc:eq-universal-weak-area-bound} to \(Z\) bounds \(A^Z\) by
\(C_r\delta^2\), proving \eqref{cc:eq-critical-stability-area}.

Put \(R^X=\langle\!\langle X\rangle\!\rangle\),
\(R^Z=\langle\!\langle Z\rangle\!\rangle\), and let \(C^{X,Z}\) be the
cross bracket.  Positivity of the \(2\times2\)
joint bracket matrix
\[
 \begin{pmatrix}
  R_t^X&C_t^{X,Z}\\
  (C_t^{X,Z})^*&R_t^Z
 \end{pmatrix}\succeq0
\]
and the Douglas factorization yields
\[
 C_t^{X,Z}=(R_t^X)^{1/2}D_t(R_t^Z)^{1/2},
 \qquad\|D_t\|\leq1.
\]
Schatten H\"older therefore gives
\[
 \|C_t^{X,Z}\|_1
 \leq
 \bigl(\operatorname{Tr}R_t^X\bigr)^{1/2}
 \bigl(\operatorname{Tr}R_t^Z\bigr)^{1/2}
 \leq\sqrt E\,\delta.
\]
Expanding the bracket of \(X+Z\) proves
\eqref{cc:eq-critical-stability-bracket}.

For the projection statement take \(X=P_nM\) and
\(Z=(I-P_n)M\).  The tail energy
\[
 \delta_n^2
 =\operatorname{Tr}\bigl((I-P_n)
 \langle\!\langle M\rangle\!\rangle_T(I-P_n)\bigr)
\]
tends to zero almost surely by trace-class monotone approximation.  Stop
at the common stopping time
\[
 \begin{aligned}
 \theta_{n,E,\delta}&:=\tau_E\wedge\sigma_{n,\delta},\\
 \tau_E&:=\inf\{t:\operatorname{Tr}R_t\geq E\}\wedge T,\\
 \sigma_{n,\delta}&:=
 \inf\{t:\operatorname{Tr}R_t^{(n)}\geq\delta^2\}\wedge T.
 \end{aligned}
\]
where
\(R_t=\langle\!\langle M\rangle\!\rangle_t\) and
\(R_t^{(n)}=(I-P_n)R_t(I-P_n)\).
These scalar bracket paths are continuous, so there is no overshoot and
the three preceding estimates apply to the stopped martingales.  For
fixed \(\delta>0\),
\[
 \mathbb P(\sigma_{n,\delta}<T)\longrightarrow0,
 \qquad
 \mathbb P(\tau_E<T)\longrightarrow0\quad(E\to\infty),
\]
because \(\operatorname{Tr}R_T^{(n)}\to0\) almost surely and
\(\operatorname{Tr}R_T<\infty\) almost surely.  Hence, in the order
\(n\to\infty\), \(\delta\downarrow0\), and \(E\uparrow\infty\), the
stopped estimates give uniform convergence in probability in all three
coordinates.

Choose a subsequence converging almost surely and uniformly for the
complete metric \eqref{cc:eq-critical-Aoki-metric}.  Its area coordinates
are finite rank.  Since
\(\mathcal S^0_{1,\infty}\) is their weak-quasi-norm closure, the limit
satisfies \eqref{cc:eq-individual-little-weak}; uniform convergence and
the completeness of the metric give a continuous path.  Standard
Hilbert-space It\^o stability identifies this limit with the original
\(H^{\otimes2}\)-valued stochastic area: indeed
\(\mathcal S_{1,\infty}\hookrightarrow\mathcal S_2\) continuously, since
\[
 \sum_{k\geq1}s_k(A)^2
 \leq\left(\sum_{k\geq1}k^{-2}\right)\|A\|_{1,\infty}^2.
\]
Uniqueness in the Hilbert tensor space, and hence uniqueness in
probability, removes dependence on the chosen subsequence.
\end{proof}

\begin{lemma}[Finite Brownian blocks]
\label{cc:lem-critical-brownian-block}
Let
\[
 \mathcal L^{(d)}
 :=\operatorname{Alt}_2\int_0^1B_t\otimes\mathop{}\!\mathrm dB_t
\]
for standard Brownian motion in \(\mathbb R^d\), with the same
half-commutator normalization as in
\eqref{cc:eq-critical-martingale-area}.  There are numerical constants
\(a,\eta,c_0>0\) such that
\begin{align}
 \mathbb P\left\{s_{\lfloor\eta d\rfloor}
       (\mathcal L^{(d)})\geq a\right\}&\longrightarrow1,
 \label{cc:eq-critical-Brownian-bulk}\\
 \mathbb P\left\{s_k(\mathcal L^{(d)})
       \geq c_0\frac d k,\ 1\leq k\leq m\right\}
 &\geq1-\varepsilon
 \label{cc:eq-critical-Brownian-edge}
\end{align}
for every fixed \(m\) and \(\varepsilon>0\), once the even dimension \(d\)
is sufficiently large.
\end{lemma}

\begin{proof}
The global singular-value law and fixed-edge asymptotics for the Brownian
L\'evy-area matrix are proved in
\cite[Sections~5.1--5.2, especially (5.1) and
(5.14)]{DelaCruzOberhauser26}.  Their convention is
\[
 \operatorname{Anti}\!\left(\int_0^1B_t\otimes\mathop{}\!\mathrm dB_t\right)_{ij}
 =\frac12\left(\int_0^1B_t^i\,\mathop{}\!\mathrm dB_t^j
              -\int_0^1B_t^j\,\mathop{}\!\mathrm dB_t^i\right),
\]
which is exactly the normalization used here.  If
\(\sigma_{(1)}^{(d)}\geq\sigma_{(2)}^{(d)}\geq\cdots\) are the distinct
positive singular values and \(d=2n\), skew-adjointness gives
\[
 s_{2j-1}(\mathcal L^{(d)})
 =s_{2j}(\mathcal L^{(d)})=\sigma_{(j)}^{(d)}.
\]
The empirical measure of the distinct positive singular values converges
almost surely to \(\lvert\operatorname{Cauchy}(0,1/2)\rvert\).  Choose a
continuity point \(a>0\) with positive upper-tail mass and then choose
\(\eta\) smaller than half that mass.  The twofold multiplicity gives
\eqref{cc:eq-critical-Brownian-bulk}.  For each fixed \(j\), (5.14) gives
\[
 \frac{\sigma_{(j)}^{(d)}}{n}
 \longrightarrow\frac{2}{(2j-1)\pi}
 \qquad\text{in probability}.
\]
Simultaneous convergence for \(1\leq j\leq\lceil m/2\rceil\), followed by
the twofold multiplicity, proves
\eqref{cc:eq-critical-Brownian-edge}; for example one may take
\(c_0=1/(4\pi)\) after increasing the dimension threshold.
\end{proof}

\begin{proposition}[Orthogonal block saturation]
\label{cc:prop-critical-block-saturation}
Let \(H=\bigoplus_{j\geq1}E_j\), where
\(d_j:=\dim E_j\uparrow\infty\).  Let \(W^Q\) be a \(Q\)-Wiener process
with
\[
 Q=\bigoplus_{j\geq1}q_jI_{E_j},
 \qquad q_j>0,\qquad\sum_jq_jd_j<\infty.
\]
If \(A\) is its area on \([0,1]\), then the compressions are independent
and
\begin{equation}
 P_{E_j}AP_{E_j}\stackrel{\mathrm d}=q_j\mathcal L^{(d_j)}.
 \label{cc:eq-critical-block-law}
\end{equation}
This distributional statement holds for every choice of the blocks and
weights above.

Conversely, given any prescribed \(m_j\uparrow\infty\), the even
dimensions \(d_j\) can be chosen recursively so that the failure
probabilities in \eqref{cc:eq-critical-Brownian-bulk} and
\eqref{cc:eq-critical-Brownian-edge} are summable.  For that choice of
dimensions, and for any positive weights satisfying
\(\sum_jq_jd_j<\infty\), the resulting almost-sure event has, eventually,
\begin{align}
 s_n(A)&\geq aq_j,\qquad
 1\leq n\leq\lfloor\eta d_j\rfloor,
 \label{cc:eq-critical-block-bulk}\\
 s_k(P_{E_j}AP_{E_j})&\geq c_0q_jd_j/k,\qquad1\leq k\leq m_j.
 \label{cc:eq-critical-block-edge}
\end{align}
\end{proposition}

\begin{proof}
The Brownian motions in the mutually orthogonal blocks are independent,
and spatial scaling by \(q_j^{1/2}\) scales area by \(q_j\), proving
\eqref{cc:eq-critical-block-law}.  Increase \(d_j\) until each failure
probability is at most \(2^{-j-1}\), and apply Borel--Cantelli.  Compression
monotonicity \(s_n(P_{E_j}AP_{E_j})\leq s_n(A)\) gives
\eqref{cc:eq-critical-block-bulk}; the edge assertion is direct.
\end{proof}

For \(K,T>0\), let \(\mathfrak C_{\mathrm{cont}}(K,T)\) be the class of
continuous square-integrable \(H\)-martingales \(M\) with \(M_0=0\)
satisfying
\[
 \operatorname{Tr}\langle\!\langle M\rangle\!\rangle_T\leq K
 \qquad\text{almost surely},
\]
and define
\[
 \begin{aligned}
 p_{\mathrm{univ}}(\mathfrak C_{\mathrm{cont}}(K,T);H)
 :=\inf\{p>0:\;&A^M_{0,T}\in\mathcal S_{p,\infty}(H)\ \text{a.s.}\\
              &\text{for every }M\in\mathfrak C_{\mathrm{cont}}(K,T)\}.
 \end{aligned}
\]

\begin{theorem}[Universal critical endpoint and its sharp boundary]
\label{cc:thm-universal-critical-endpoint}
Suppose \(H\) is infinite dimensional.  Then, for every \(K,T>0\),
\begin{equation}
 p_{\mathrm{univ}}(\mathfrak C_{\mathrm{cont}}(K,T);H)=1.
 \label{cc:eq-critical-universal-exponent}
\end{equation}
The exponent-one upper bound is dimension free and is quantified by
\eqref{cc:eq-universal-weak-area-bound}.  It is sharp in all of the
following distinct senses.
\begin{enumerate}[label=\textup{(\roman*)},leftmargin=2.8em]
\item No \(\mathcal S_{p,\infty}\) with \(p<1\) contains all terminal
areas in \(\mathfrak C_{\mathrm{cont}}(K,T)\).
\item No finite critical Lorentz ideal \(\mathcal S_{1,\rho}\),
\(0<\rho<\infty\), is universal.
\item Given any deterministic positive sequence
\(\varepsilon_n\downarrow0\), there is one
martingale in \(\mathfrak C_{\mathrm{cont}}(K,T)\) such that, almost surely,
\begin{equation}
 A_{0,T}\in\mathcal S^0_{1,\infty}(H),\qquad
 \sup_n\frac{n s_n(A_{0,T})}{\varepsilon_n}=\infty,
 \qquad
 A_{0,T}\notin\bigcup_{0<\rho<\infty}\mathcal S_{1,\rho}(H).
 \label{cc:eq-critical-no-uniform-little-modulus}
\end{equation}
\end{enumerate}
Thus there is no conflict between
\eqref{cc:eq-individual-little-weak} and universal sharpness: each fixed
martingale area has \(ns_n(A)\to0\), while dimension-varying Brownian
blocks make that convergence arbitrarily slow and destroy every fixed
finite Lorentz refinement.
\end{theorem}

\begin{proof}
It is enough first to construct a witness on \([0,1]\) with terminal trace
energy one.  Fix \(\varepsilon_n\downarrow0\), choose positive \(c_j\)
with \(\sum_jc_j=1\), and put \(q_j=c_j/d_j\).  Recursively choose
\(m_j\uparrow\infty\) and then even \(d_j\geq2m_j\), increasing the latter
so that the bulk and edge failure probabilities are at most \(2^{-j}\),
the index intervals
\([\eta d_j/2,\eta d_j]\) are disjoint, and
\begin{equation}
 c_j(\log(1+m_j))^{1/j}\geq j^2,
 \qquad
 \varepsilon_{\lfloor\eta d_j\rfloor}\leq c_j/j.
 \label{cc:eq-critical-block-selection}
\end{equation}
This is possible because \(m_j\), and then \(d_j\), may be arbitrarily
large.  The covariance
\(Q=\bigoplus_jq_jI_{E_j}\) is positive trace class with trace one.

On the almost-sure saturation event of
\Cref{cc:prop-critical-block-saturation}, put
\(n_j=\lfloor\eta d_j\rfloor\).  The bulk bound gives
\[
 n_js_{n_j}(A)\geq\frac{a\eta}{2}q_jd_j
 =\frac{a\eta}{2}c_j
\]
for all sufficiently large \(j\), because
\(\lfloor\eta d_j\rfloor\geq\eta d_j/2\),
and the second condition in \eqref{cc:eq-critical-block-selection} gives
\[
 \frac{n_js_{n_j}(A)}{\varepsilon_{n_j}}
 \geq\frac{a\eta}{2}j\longrightarrow\infty.
\]
For every orthogonal projection \(P\), approximation numbers satisfy
\(s_k(PAP)\leq s_k(A)\).  Hence, for \(0<\rho<\infty\), the edge bound
gives term by term
\begin{align*}
 \|A\|_{1,\rho}^{\rho}
 &\geq\|P_{E_j}AP_{E_j}\|_{1,\rho}^{\rho}\\
 &\geq(c_0c_j)^{\rho}\sum_{k\leq m_j}\frac1k.
\end{align*}
Because \(0<c_j\leq1\), for \(j\geq\rho\) the first condition in
\eqref{cc:eq-critical-block-selection} implies
\[
 c_j^\rho\log(1+m_j)
 \geq c_j^j\log(1+m_j)\geq j^{2j}.
\]
The harmonic sum is comparable to \(\log(1+m_j)\), so the preceding
Lorentz lower bound diverges for every fixed \(\rho\).  All these
deterministic inequalities hold on the same Borel--Cantelli event; no
intersection of \(\rho\)-dependent probability-one events is being
taken.  Thus \(A\) lies simultaneously in no finite critical Lorentz
class.  On the other hand,
\Cref{cc:thm-critical-projection-stability} gives
\(A\in\mathcal S^0_{1,\infty}\) almost surely.  This proves
\eqref{cc:eq-critical-no-uniform-little-modulus}.

To rule out \(p<1\), use the same construction with
\(\varepsilon_n=n^{1-1/p}\downarrow0\).  Membership in
\(\mathcal S_{p,\infty}\) would make
\[
 \frac{ns_n(A)}{\varepsilon_n}=n^{1/p}s_n(A)
\]
bounded, a contradiction.  Hence the universal exponent is at least one;
\Cref{cc:prop-universal-critical-area-bound} gives the opposite
inequality.

Finally set \(\widetilde M_t:=\sqrt K\,M_{t/T}\) on \([0,T]\).  Its
terminal trace energy is \(K\), its area is multiplied by \(K\), and none
of the membership or nonmembership statements changes.  This proves every
assertion for arbitrary \(K,T>0\).
\end{proof}

\begin{corollary}[Intrinsic degree-two endpoint]
\label{cc:cor-intrinsic-degree-two-endpoint}
Consider symmetric-ideal completions of the finite-rank split state with
the following four properties:
\begin{enumerate}[label=\textup{(\roman*)},leftmargin=2.8em]
\item finite-rank states are dense and every bounded pushforward
      \(L_\#\) is continuous;
\item the universal positive collision functional
      \(Q\mapsto\operatorname{Tr}Q\) is continuous;
\item the full exterior response is a faithful continuous area reader,
      and the scalar profile \(\mathfrak P_A\) is a continuous spectral
      reader for the singular-value topology;
\item the completion contains every continuous Hilbert-martingale state
      with uniformly bounded terminal trace energy.
\end{enumerate}
Then the collision coordinate must have trace-class strength.  Within the
weak and Lorentz scales, the universal area coordinate reaches exactly
the weak-trace frontier.  It cannot be replaced by any
\(\mathcal S_{p,\infty}\) with \(p<1\), by any finite
\(\mathcal S_{1,\rho}\), or by a little weak space with a prescribed
dimension-free decay modulus.  Requiring a separable finite-rank
completion therefore selects
\[
 H\times\mathcal S^0_{1,\infty}(H)_{\mathrm{skew}}
   \times\mathcal S_1(H)_{\mathrm{sa}}
\]
as the intrinsic degree-two split endpoint in this scale.
\end{corollary}

\begin{proof}
The positive collision assertion is
\Cref{cc:prop-critical-trace-obstruction}.  The profile decoder
\Cref{cc:lem-response-profile-decoder} shows that the scalar exterior
profile detects precisely the weak Schatten singular-value scale at
degree two.  The universal upper estimate is
\Cref{cc:prop-universal-critical-area-bound}, while
\Cref{cc:thm-universal-critical-endpoint} gives all three sharp
obstructions to a smaller weak/Lorentz target.  Finally,
\Cref{cc:thm-critical-split-completion} identifies the separable
finite-rank completion and proves continuity of its structural
operations.  These arguments concern scalar Hilbert and singular-value
estimates; none assumes local convexity of the weak ideal.
\end{proof}

\subsection{The canonical degree-two quotient}
\label{cc:subsec-critical-causal-quotient}

The degree-two It\^o state is not an additional stochastic decoration of a
causal character.  It is the canonical degree-two quotient of the typed
quasi-shuffle character.  We first make this statement on the algebraic core
and then pass to the critical completion.

For \(x,y\in H\), identify \(x\otimes y\) with the rank-one operator
\(z\mapsto\langle y,z\rangle x\), and put
\[
 x\wedge_{\mathrm{op}}y:=x\otimes y-y\otimes x.
\]
Let \(\mathscr I^{(2)}_{\mathrm{fin}}(H)\) be the directed union, over all
finite-dimensional subspaces \(E\subset H\), of triples
\((x,\mathbb x,Q)\) satisfying
\[
 x\in E,\qquad
 \mathbb x\in E^{\otimes 2},\qquad
 Q\in\mathcal L(E)_{\mathrm{sa}},\qquad
 \mathbb x+\mathbb x^*=x\otimes x-Q.
\]
Here the adjoint flips the two tensor legs.  The last identity is the
degree-two quasi-shuffle relation.  Define
\begin{equation}
 (x,\mathbb x,Q)\diamond(y,\mathbb y,R)
 :=(x+y,\mathbb x+\mathbb y+x\otimes y,Q+R).
 \label{cc:eq-ito-two-product}
\end{equation}

\begin{theorem}[Critical typed target and causal quotient]
\label{cc:thm-critical-typed-target}
The product \eqref{cc:eq-ito-two-product} is associative and preserves
\(\mathscr I^{(2)}_{\mathrm{fin}}(H)\).  The map
\begin{equation}
 \Pi_2(x,\mathbb x,Q)
 :=\bigl(x,\operatorname{Alt}_2\mathbb x,Q\bigr),
 \qquad
 \operatorname{Alt}_2\mathbb x:=\tfrac12(\mathbb x-\mathbb x^*),
 \label{cc:eq-split-quotient}
\end{equation}
is a group isomorphism from \(\mathscr I^{(2)}_{\mathrm{fin}}(H)\) onto
\(\mathscr G_{\mathrm{fin}}(H)\), where
\[
 (x,A,Q)\circ(y,B,R)
 =\left(x+y,A+B+\tfrac12x\wedge_{\mathrm{op}}y,Q+R\right).
\]
Its inverse is the exact reconstruction
\begin{equation}
 \mathbb x=\tfrac12x\otimes x+A-\tfrac12Q.
 \label{cc:eq-split-reconstruction}
\end{equation}
Thus, after the typed quasi-shuffle relation has been imposed, passage to
the split state loses no degree-two information.

Let \(E\subset H\) be finite dimensional and let \(X\) be an
\(E\)-valued c\`adl\`ag semimartingale.  Put
\begin{align*}
 X_{s,t}&:=X_t-X_s,\\
 \mathbb X^I_{s,t}
 &:=\int_{(s,t]}(X_{r-}-X_s)\otimes\mathop{}\!\mathrm dX_r,\\
 A_{s,t}&:=\operatorname{Alt}_2\mathbb X^I_{s,t},
 \qquad
 Q_{s,t}:=[X]_t-[X]_s.
\end{align*}
There are versions for which, on one event and for every
\(0\leq s\leq u\leq t\),
\begin{equation}
 (X_{s,t},A_{s,t},Q_{s,t})
 =(X_{s,u},A_{s,u},Q_{s,u})
  \circ(X_{u,t},A_{u,t},Q_{u,t}).
 \label{cc:eq-critical-split-Chen}
\end{equation}
Moreover,
\begin{equation}
 \mathbb X^I_{s,t}
 =\tfrac12X_{s,t}^{\otimes 2}+A_{s,t}-\tfrac12[X]_{s,t}.
 \label{cc:eq-critical-Ito-reconstruction}
\end{equation}
At every jump time \(t\), with \(z=\Delta X_t\), the relative state is
\begin{equation}
 (X_{t-,t},A_{t-,t},Q_{t-,t})=(z,0,z\otimes z).
 \label{cc:eq-critical-single-jump}
\end{equation}

For every bounded linear map \(L:H\to K\),
\[
 L_\#(x,A,Q):=(Lx,LAL^*,LQL^*)
\]
is a homomorphism on the finite-rank groups and extends continuously to the
critical completions of \Cref{cc:thm-critical-split-completion}.  On the
typed core set
\[
 L_\#(x,\mathbb x,Q):=(Lx,(L\otimes L)\mathbb x,LQL^*).
\]
Then
\(\Pi_2(L_\#(x,\mathbb x,Q))=L_\#\Pi_2(x,\mathbb x,Q)\).
\end{theorem}

\begin{proof}
The second coordinate in either association of a triple product is
\[
 \mathbb x+\mathbb y+\mathbb z+x\otimes y+x\otimes z+y\otimes z,
\]
which proves associativity.  If the two inputs satisfy the quasi-shuffle
relation, then
\begin{align*}
 &(\mathbb x+\mathbb y+x\otimes y)
  +(\mathbb x+\mathbb y+x\otimes y)^*\\
 &\qquad=(x+y)\otimes(x+y)-(Q+R),
\end{align*}
so the product is closed.  Taking antisymmetric parts gives
\[
 \operatorname{Alt}_2(\mathbb x+\mathbb y+x\otimes y)
 =\operatorname{Alt}_2\mathbb x+\operatorname{Alt}_2\mathbb y
  +\tfrac12x\wedge_{\mathrm{op}}y.
\]
Hence \(\Pi_2\) is multiplicative.  The quasi-shuffle relation fixes the
symmetric part of \(\mathbb x\) and gives
\eqref{cc:eq-split-reconstruction}, proving bijectivity.  The identity and
inverse in typed coordinates are
\[
 (0,0,0),\qquad
 (x,\mathbb x,Q)^{-1}=(-x,-\mathbb x+x\otimes x,-Q);
\]
under \(\Pi_2\) these become \((0,0,0)\) and \((-x,-A,-Q)\).

For the stochastic assertion, splitting the left-point integral at \(u\)
gives
\[
 \mathbb X^I_{s,t}
 =\mathbb X^I_{s,u}+\mathbb X^I_{u,t}
  +X_{s,u}\otimes X_{u,t}.
\]
Operator quadratic variation is additive.  The Hilbert-space It\^o product
formula gives
\[
 \mathbb X^I_{s,t}+(\mathbb X^I_{s,t})^*
 =X_{s,t}^{\otimes 2}-[X]_{s,t},
\]
which proves both the typed relation and
\eqref{cc:eq-critical-Ito-reconstruction}.  Taking antisymmetric parts
proves \eqref{cc:eq-critical-split-Chen}.  Equivalently, one first chooses
c\`adl\`ag versions of the one-parameter anchor coordinates and reconstructs
all intervals from that anchor; the identities then hold on a single event
for every time triple.  On \((t-,t]\) the left-point integrand relative to
\(X_{t-}\) is zero, whereas
\(\Delta[X]_t=z\otimes z\), proving
\eqref{cc:eq-critical-single-jump}.  Finally, deterministic bounded maps
commute with stochastic integration and operator quadratic variation, which
proves naturality.
\end{proof}

\begin{corollary}[Completed typed quotient]
\label{cc:cor-critical-typed-completion}
Set
\[
 \mathscr I^{(2)}_{\mathrm{crit}}(H)
 :=
 \left\{
 (x,\mathbb x,Q)\in
 H\times\mathcal S_2(H)\times\mathcal S_1(H)_{\mathrm{sa}}:
 \begin{array}{l}
 \mathbb x+\mathbb x^*=x\otimes x-Q,\\
 \operatorname{Alt}_2\mathbb x\in
 \mathcal S^0_{1,\infty}(H)_{\mathrm{skew}}
 \end{array}
 \right\}.
\]
Equip this space with the pullback of \(d_{\mathrm{crit}}\) under
\(\Pi_2\).  Then \(\diamond\) extends continuously,
\(\mathscr I^{(2)}_{\mathrm{crit}}(H)\) is the completion of the typed
finite-rank core, and
\[
 \Pi_2:\mathscr I^{(2)}_{\mathrm{crit}}(H)
 \longrightarrow\mathscr G_{\mathrm{crit}}(H)
\]
is a topological group isomorphism.  Its inverse is still
\eqref{cc:eq-split-reconstruction}.
In infinite dimension this pullback topology is not the ambient subspace
topology inherited from
\(H\times\mathcal S_2(H)\times\mathcal S_1(H)\): it retains the stronger
\(\mathcal S^0_{1,\infty}\)-topology of the antisymmetric component.
\end{corollary}

\begin{proof}
The inclusion
\(\mathcal S^0_{1,\infty}(H)\hookrightarrow\mathcal S_2(H)\) is continuous.
Hence \eqref{cc:eq-split-reconstruction} maps every critical split state to
the displayed typed space and is inverse to \(\Pi_2\).  Finite-rank density
in the three split coordinates and
\Cref{cc:thm-critical-split-completion} identify the completion; continuity
of the product is transported through \(\Pi_2\).
\end{proof}

Thus the typed and split presentations define the same canonical Polish
second-order entrance; the split coordinates will be used below.

\subsection{Optional and projection stability}
\label{cc:subsec-critical-optional-levy}

\begin{theorem}[Optional completion of the critical quotient]
\label{cc:thm-critical-optional-quotient}
Let \((\Omega,\mathcal F,(\mathcal F_t)_{t\geq0},\mathbb P)\) satisfy the
usual conditions, and let
\[
 \mathbf A_t=(X_t,A_t,Q_t),\qquad \mathbf A_0=(0,0,0),
\]
be an adapted c\`adl\`ag process with values in
\(\mathscr G_{\mathrm{crit}}(H)\).  For bounded stopping times
\(0\leq\sigma\leq\tau\leq T\), define
\begin{equation}
 \mathbf\Phi_{\sigma,\tau}
 :=\mathbf A_\sigma^{-1}\circ\mathbf A_\tau
 =\left(
 X_\tau-X_\sigma,
 A_\tau-A_\sigma-\tfrac12X_\sigma\wedge_{\mathrm{op}}X_\tau,
 Q_\tau-Q_\sigma
 \right).
 \label{cc:eq-critical-stopped-anchor}
\end{equation}
Then \(\mathbf\Phi_{\sigma,\tau}\) is
\(\mathcal F_\tau\)-measurable and, on one event, simultaneously for all
bounded stopping triples \(\sigma\leq\rho\leq\tau\),
\begin{equation}
 \mathbf\Phi_{\sigma,\tau}
 =\mathbf\Phi_{\sigma,\rho}\circ\mathbf\Phi_{\rho,\tau}.
 \label{cc:eq-critical-stopped-Chen}
\end{equation}
It is the unique extension of the perfected deterministic-time character
which is obtained by pointwise pasting on finite-valued stopping intervals
and is continuous in probability under nested right-dyadic approximation.

Suppose in addition that \(X\) is an \(H\)-valued c\`adl\`ag
semimartingale whose split anchor satisfies
\[
 A_t=\operatorname{Alt}_2\int_{(0,t]}X_{r-}\otimes
       \mathop{}\!\mathrm dX_r,
 \qquad
 Q_t=[X]_t.
\]
Then
\begin{align}
 X_{\sigma,\tau}&=X_\tau-X_\sigma,\notag\\
 Q_{\sigma,\tau}&=[X]_\tau-[X]_\sigma,\notag\\
 A_{\sigma,\tau}
 &=\operatorname{Alt}_2\int_{(0,T]}
  {\mathbf 1}_{\{\sigma<r\leq\tau\}}
  (X_{r-}-X_\sigma)\otimes\mathop{}\!\mathrm dX_r.
 \label{cc:eq-critical-stopped-integral}
\end{align}
Thus algebraic optional evaluation is exactly the usual stopped left-point
It\^o integral.
\end{theorem}

\begin{proof}
Apply \Cref{cc:thm-optional-completion} to the Polish group
\(\mathscr G_{\mathrm{crit}}(H)\) with trivial transport.  It gives
measurability, the simultaneous stopped Chen identity, right-dyadic
continuity, and uniqueness, while the group law gives the explicit formula
\eqref{cc:eq-critical-stopped-anchor}.  It remains only to identify this
abstract stopped increment with the classical left-point It\^o integral.
For that identification, both
\[
 r\longmapsto{\mathbf 1}_{\{\sigma<r\leq\tau\}}X_{r-}
 \quad\text{and}\quad
 r\longmapsto{\mathbf 1}_{\{\sigma<r\leq\tau\}}X_\sigma
\]
are predictable.  Indeed, \(X_-\) is predictable,
\({\mathbf 1}_{(\sigma,\tau]}\) is predictable, and the second displayed
process is adapted and left-continuous.  After localization the difference
is stochastically integrable.  The stopping identity for semimartingale
integrals, first for finite-valued stopping times and then by decreasing
right approximation, gives
\[
 \int_{(0,T]}{\mathbf 1}_{\{\sigma<r\leq\tau\}}
 (X_{r-}-X_\sigma)\otimes\mathop{}\!\mathrm dX_r
 =
 \int_{(\sigma,\tau]}(X_{r-}-X_\sigma)\otimes
 \mathop{}\!\mathrm dX_r.
\]
Taking antisymmetric parts and using
\[
 \operatorname{Alt}_2\bigl(
 X_\sigma\otimes(X_\tau-X_\sigma)\bigr)
 =\tfrac12X_\sigma\wedge_{\mathrm{op}}X_\tau
\]
identifies the middle coordinate with
\eqref{cc:eq-critical-stopped-anchor}.  The last coordinate is the standard
stopping identity for optional quadratic variation.
\end{proof}

\begin{corollary}[Projection stability on moving stopping intervals]
\label{cc:cor-critical-stopped-projection}
Let \(M\) be a continuous square-integrable \(H\)-valued martingale with
\(M_0=0\), and let
\((P_n)\) be increasing finite-rank orthogonal projections converging
strongly to the identity.  Let \(\mathbf A^{(n)}\) and \(\mathbf A\) be the
critical split anchors of \(P_nM\) and \(M\) supplied by
\Cref{cc:thm-critical-projection-stability}.  Then
\begin{equation}
 \sup_{0\leq s\leq t\leq T}
 \bigl\|\mathbf\Phi^{(n)}_{s,t}-\mathbf\Phi_{s,t}\bigr\|_{\mathrm{crit}}
 \longrightarrow0
 \qquad\text{in probability}.
 \label{cc:eq-critical-uniform-interval-projection}
\end{equation}
Here subtraction is coordinatewise and
\(\|\cdot\|_{\mathrm{crit}}\) is the coordinate quasi-gauge
\eqref{cc:eq-critical-quasi-gauge}; convergence for that gauge is
equivalent to convergence for \(d_{\mathrm{crit}}\).
Consequently, for every sequence of possibly \(n\)-dependent bounded
stopping pairs \(\sigma_n\leq\tau_n\leq T\),
\[
 d_{\mathrm{crit}}\!\left(
 \mathbf\Phi^{(n)}_{\sigma_n,\tau_n},
 \mathbf\Phi_{\sigma_n,\tau_n}
 \right)\longrightarrow0
 \qquad\text{in probability}.
\]
Thus finite-dimensional approximation commutes with optional completion,
even for intervals selected from the approximation.
\end{corollary}

\begin{proof}
Write
\(\mathbf A^{(n)}=(M^{(n)},A^{(n)},Q^{(n)})\) and
\(\mathbf A=(M,A,Q)\).  The cited stability theorem gives, after
localization,
\[
 \Delta_n:=\sup_{t\leq T}\left(
 \|M_t^{(n)}-M_t\|+
 \|A_t^{(n)}-A_t\|_{1,\infty}+
 \|Q_t^{(n)}-Q_t\|_1\right)\longrightarrow0
\]
in probability.  Formula \eqref{cc:eq-critical-stopped-anchor} and
\[
 \|x\wedge_{\mathrm{op}}y\|_{1,\infty}
 \leq2\|x\|\,\|y\|
\]
give
\[
 \sup_{0\leq s\leq t\leq T}
 \|\mathbf\Phi^{(n)}_{s,t}-\mathbf\Phi_{s,t}\|_{\mathrm{crit}}
 \leq
 C\left(1+\sup_{t\leq T}\|M_t^{(n)}\|
          +\sup_{t\leq T}\|M_t\|\right)\Delta_n.
\]
The factor in parentheses is tight, proving
\eqref{cc:eq-critical-uniform-interval-projection}.  The same pathwise
supremum bounds evaluation at every random pair
\((\sigma_n,\tau_n)\); compatibility of \(d_{\mathrm{crit}}\) with the
coordinate topology proves the last assertion.
\end{proof}

The path state below records optional quadratic variation, whereas its
conditional estimate is driven by the predictable bracket.  The two must
not be identified in the presence of jumps.



\begin{definition}[Critical second-order entrance]
\label{cc:def-critical-second-order-entrance}
A stochastic source has a \emph{critical second-order entrance} in a real
separable Hilbert space \(H\) if its perfected degree-two anchor is an
adapted c\`adl\`ag path in \(\mathscr G^+_{\mathrm{crit}}(H)\), its
finite-rank spatial projections converge ucp in the critical metric, and its
relative increments are obtained from the anchor by the stopped causal law.
A \(\beta\)-controlled entrance additionally satisfies, for a control
\(\omega\),
\[
 \|X_{s,t}\|\lesssim\omega(s,t)^\beta,
 \qquad
 \|A_{s,t}\|_{1,\infty}\lesssim\omega(s,t)^{2\beta},
 \qquad \beta\in(1/3,1/2).
\]
The source-specific analysis ends at this entrance.  Higher ancestry and
spectral conclusions are supplied by the normal and response theories.
\end{definition}
\section{Controlled critical responses and native Hall support}
\label{hall:sec-native}

The abstract response-kernel theorem reduces all lower estimates and inverse
statements to one obligation: prove full left support for a declared class of
native readers.  We first retain the source-independent upper machinery, and
then verify native support by pure-area Hall geometry.

\subsection{Critical multi-flattening spaces and sewing}
\label{cc:gen-subsection-calculus}

For \(n\geq1\), \(S\subset[n]\), and
\(K\in H^{\otimes n}\), let

\[
 \operatorname{Mat}_S(K):H^{\otimes S}\longrightarrow
 H^{\otimes S^c}
\]

be the Hilbert-space flattening obtained by placing the tensor legs in
increasing order on each side of the cut.  A different ordering changes the
operator only by unitaries.  The cases \(S=\varnothing,[n]\) are rank-one
realizations with Schatten norm \(\|K\|_{H^{\otimes n}}\).  They are retained
in the finite-\(p\) envelope below and omitted only from the critical
logarithmic maximum.

For \(1\leq p<\infty\), first endow the algebraic finite tensors with the
augmented multi-flattening norm

\begin{equation}
 \|K\|_{\mathfrak A_p^{(n)}}
 :=\max_{S\subset[n]}
 \|\operatorname{Mat}_S(K)\|_{\mathcal S_p},
 \qquad K\in H^{\otimes_{\mathrm{alg}}n},
 \label{cc:gen-augmented-schatten-multiflat}
\end{equation}

where \(H^{\otimes\varnothing}:=\mathbb R\).  Denote the completion by
\(\mathfrak A_p^{(n)}(H)\), and put
\(\mathfrak A_p^{(0)}(H):=\mathbb R\), with its absolute-value norm.  The
adjective ``augmented'' records the two trivial cuts.  If \(n\geq2\) and
\(1\leq p\leq2\), they are redundant:

\begin{equation}
 \|K\|_{\mathfrak A_p^{(n)}}
 =\max_{\varnothing\ne S\ne[n]}
 \|\operatorname{Mat}_S(K)\|_{\mathcal S_p}.
 \label{cc:gen-augmented-nontrivial-cuts}
\end{equation}

\begin{proposition}[Permutation-stable tensor algebra]
\label{cc:gen-prop-permutation-stable-tensor-algebra}
Let \(1\leq p<\infty\).  For \(m,n\geq1\) and algebraic finite tensors
\(K\in H^{\otimes_{\mathrm{alg}}m}\) and
\(L\in H^{\otimes_{\mathrm{alg}}n}\),

\begin{equation}
 \|K\otimes L\|_{\mathfrak A_p^{(m+n)}}
 =\|K\|_{\mathfrak A_p^{(m)}}
  \|L\|_{\mathfrak A_p^{(n)}}.
 \label{cc:gen-exact-schatten-tensor-crossnorm}
\end{equation}

The tensor product therefore extends to the completions and the same equality
holds there; with the scalar convention it also holds if \(m=0\) or \(n=0\).
Every tensor-leg permutation is an isometry.  More generally,
if \(Q=\sum_{\pi\in\mathfrak S_n}c_\pi U_\pi\) is a fixed element of the
permutation algebra, then

\begin{equation}
 \|QK\|_{\mathfrak A_p^{(n)}}
 \leq\left(\sum_{\pi\in\mathfrak S_n}|c_\pi|\right)
 \|K\|_{\mathfrak A_p^{(n)}}.
 \label{cc:gen-permutation-projector-bound}
\end{equation}

Consequently every fixed Young projector and every
Dynkin--Specht--Wever/free-Lie projector extends to
\(\mathfrak A_p^{(n)}(H)\), with norm independent of the dimension of \(H\).
For \(1\leq p\leq2\) and \(N\geq0\), set

\begin{equation}
 \mathfrak A_{p,\leq N}(H)
 :=\bigoplus_{n=0}^N\mathfrak A_p^{(n)}(H),
 \qquad
 \|(K_0,\ldots,K_N)\|_{\mathfrak A_{p,\leq N}}
 :=\sum_{n=0}^N\|K_n\|_{\mathfrak A_p^{(n)}}.
 \label{cc:gen-truncated-augmented-algebra}
\end{equation}

With truncated concatenation

\begin{equation}
 (K\diamond_N L)_k:=\sum_{i+j=k}K_i\otimes L_j,
 \qquad 0\leq k\leq N,
 \label{cc:gen-truncated-concatenation}
\end{equation}

this is a unital graded Banach algebra, and

\begin{equation}
 \|K\diamond_NL\|_{\mathfrak A_{p,\leq N}}
 \leq
 \|K\|_{\mathfrak A_{p,\leq N}}
 \|L\|_{\mathfrak A_{p,\leq N}}.
 \label{cc:gen-truncated-concatenation-bound}
\end{equation}
\end{proposition}

\begin{proof}
For every cut \(S\subset[m+n]\), write
\(S_1=S\cap[m]\), and let \(S_2\subset[n]\) be the translate by \(-m\)
of \(S\cap\{m+1,\ldots,m+n\}\).  Permuting tensor factors separately in
the domain and range gives unitary operators \(V_S,W_S\) such that

\begin{equation}
 \operatorname{Mat}_S(K\otimes L)
 =V_S\bigl(
   \operatorname{Mat}_{S_1}(K)\otimes
   \operatorname{Mat}_{S_2}(L)
  \bigr)W_S.
 \label{cc:gen-flattening-tensor-factorization}
\end{equation}

If \((a_j)\) and \((b_k)\) are the singular values of two finite-rank
operators, the singular values of their Hilbert tensor product are
\((a_jb_k)_{j,k}\).  Hence

\[
 \|A\otimes B\|_{\mathcal S_p}^p
 =\sum_{j,k}(a_jb_k)^p
 =\|A\|_{\mathcal S_p}^p\|B\|_{\mathcal S_p}^p.
\]

As \(S\mapsto(S_1,S_2)\) is a bijection from the cuts of \([m+n]\) to
the pairs of cuts of \([m]\) and \([n]\), taking the maximum in
\eqref{cc:gen-flattening-tensor-factorization} gives the equality
\eqref{cc:gen-exact-schatten-tensor-crossnorm}, not merely an inequality.
The trivial cut shows that \eqref{cc:gen-augmented-schatten-multiflat}
dominates the Hilbert tensor norm and is therefore a norm.  For any cut,

\[
 \|\operatorname{Mat}_S(K)\|_{\mathcal S_2}
 =\|K\|_{H^{\otimes n}}.
\]

When \(p\leq2\), monotonicity of Schatten norms gives
\(\|K\|_{H^{\otimes n}}\leq
\|\operatorname{Mat}_S(K)\|_{\mathcal S_p}\).  If \(n\geq2\), choose a
nontrivial cut; it dominates both trivial cuts and proves
\eqref{cc:gen-augmented-nontrivial-cuts}.

To pass from finite tensors to the completions, choose
\(K_a\to K\) and \(L_a\to L\) in their respective augmented norms.
The sequences are bounded, and the exact crossnorm together with the
triangle inequality gives

\begin{align*}
 \|K_a\otimes L_a-K_b\otimes L_b\|_{\mathfrak A_p^{(m+n)}}
 &\leq
 \|K_a-K_b\|_{\mathfrak A_p^{(m)}}
 \|L_a\|_{\mathfrak A_p^{(n)}}\\
 &\quad+
 \|K_b\|_{\mathfrak A_p^{(m)}}
 \|L_a-L_b\|_{\mathfrak A_p^{(n)}}.
\end{align*}

Thus \(K_a\otimes L_a\) has a limit independent of the approximating
sequences.  Passing to the limit in
\eqref{cc:gen-exact-schatten-tensor-crossnorm} proves the exact equality on
the completions.

For \(\pi\in\mathfrak S_n\), the flattening of \(U_\pi K\) across \(S\)
is, up to domain and range unitaries, the flattening of \(K\) across
\(\pi^{-1}S\).  Since \(S\mapsto\pi^{-1}S\) permutes all cuts, \(U_\pi\)
is an isometry.  The triangle inequality now proves
\eqref{cc:gen-permutation-projector-bound}.  Normalized Young projectors and
normalized Dynkin--Specht--Wever/free-Lie projectors are fixed finite linear
combinations of permutations; see \cite{Reutenauer93}.  Their algebraic
projection identities pass by continuity to the completions, and the
displayed bound depends on the fixed projector but not on \(\dim H\).

Finally, \(\mathfrak A_{p,\leq N}(H)\) is complete because it is a finite
\(\ell^1\)-sum of Banach spaces.  By the exact crossnorm and the triangle
inequality,

\begin{align*}
 \|K\diamond_NL\|_{\mathfrak A_{p,\leq N}}
 &\leq\sum_{i+j\leq N}
 \|K_i\|_{\mathfrak A_p^{(i)}}
 \|L_j\|_{\mathfrak A_p^{(j)}}\\
 &\leq
 \left(\sum_{i=0}^N\|K_i\|_{\mathfrak A_p^{(i)}}\right)
 \left(\sum_{j=0}^N\|L_j\|_{\mathfrak A_p^{(j)}}\right).
\end{align*}

Here the factors of degree zero are scalar multiplications, for which the same
crossnorm equality is immediate.  This is
\eqref{cc:gen-truncated-concatenation-bound}.  On algebraic tensors,
associativity and the unit \((1,0,\ldots,0)\) descend from the tensor algebra
after discarding all degrees above \(N\), and the component formula respects
degrees.  Continuity and density carry both identities to the completions.
This proves the graded Banach-algebra assertion.
\end{proof}

The finite-\(p\) envelope is the permutation-stable algebraic completion used
to justify tensor products and fixed projectors.  The endpoint scale needed
for critical stochastic generation is different: it records the logarithmic
loss under tensor products and is defined next.

Use the critical ideals \(\mathfrak W^{[r]}\) and
\(\mathfrak W^{[r]}_0\) from \eqref{cc:eq:critical-ideals}.  Define the big
and little multi-flattening spaces by

\begin{align}
 \mathfrak W_{\mathrm{mf}}^{[r]}(H^{\otimes n})
 &:={\left\{K\in H^{\otimes n}:
 \operatorname{Mat}_S(K)\in\mathfrak W^{[r]}
 \text{ for every }\varnothing\ne S\ne[n]\right\}},
 \label{cc:gen-multiflat-space}\\
 \|K\|_{\mathfrak W_{\mathrm{mf}}^{[r]}(n)}
 &:=\|K\|_{H^{\otimes n}}
 +\max_{\varnothing\ne S\ne[n]}
 \|\operatorname{Mat}_S(K)\|_{\mathfrak W^{[r]}},
 \label{cc:gen-multiflat-norm}\\
 \mathfrak W_{0,\mathrm{mf}}^{[r]}(H^{\otimes n})
 &:={\left\{K\in\mathfrak W_{\mathrm{mf}}^{[r]}:
 \operatorname{Mat}_S(K)\in\mathfrak W^{[r]}_0
 \text{ for every }\varnothing\ne S\ne[n]\right\}}.
 \label{cc:gen-little-multiflat}
\end{align}

The Hilbert term in \eqref{cc:gen-multiflat-norm} is included so that the
tensor represented by the compatible family of flattenings is part of the
topology.  At degree two it is controlled by the critical term, since

\[
 \|T\|_{\mathcal S_2}^2
 \leq \|T\|_{\mathfrak W^{[r]}}^2
 \sum_{N\geq1}\frac{(\log(e+N))^{2r-2}}{N^2}.
\]

\begin{lemma}[Distribution, completeness, and the little core]
\label{cc:gen-lemma-critical-spaces}
With \(\phi_r\) as in \eqref{cc:eq-critical-profile}, for a compact operator \(T\), the quantity
\(\sup_Ns_N(T)/\phi_r(N)\) is equivalent, with constants depending only
on \(r\), to the least \(A>0\) for which

\begin{equation}
 d_T(u):=\#\{N:s_N(T)>u\}
 \leq C_r\frac Au
 \left(\log\left(e+\frac Au\right)\right)^{r-1},
 \qquad u>0.
 \label{cc:gen-distribution-bound}
\end{equation}

The space \(\mathfrak W^{[r]}\) is a complete symmetric quasi-Banach
ideal.  Its little subspace is closed and is exactly the
\(\mathfrak W^{[r]}\)-closure of the finite-rank operators.  Moreover,

\[
 \mathfrak W^{[r]}_0\subset\mathfrak W^{[r+1]}_0.
\]

For every fixed \(n\),
\(\mathfrak W_{\mathrm{mf}}^{[r]}(H^{\otimes n})\) is complete, and its
little subspace is precisely the closure of the finite-dimensional tensors
in the big multi-flattening quasi-norm.  Tensor-leg permutations are
isometries of both spaces.
\end{lemma}

\begin{proof}
If \(s_N(T)\leq A\phi_r(N)\), then \(d_T(u)\geq N\) implies

\[
 \frac{N}{(\log(e+N))^{r-1}}<\frac Au.
\]

The increasing function on the left has an inverse comparable to
\(x\mapsto x(\log(e+x))^{r-1}\); this proves
\eqref{cc:gen-distribution-bound}.  Conversely, apply that bound with
\(u=s_N(T)/2\), for which \(d_T(u)\geq N\), and invert once more.

The Weyl inequality

\[
 s_{2N-1}(S+T)\leq s_N(S)+s_N(T)
\]

and the doubling relation \(\phi_r(2N)\asymp_r\phi_r(N)\) give a fixed
quasi-triangle inequality.  Symmetry and the ideal property follow from the
corresponding singular-value inequalities.  A Cauchy sequence in this
quasi-norm is Cauchy in operator norm and converges to a compact operator
\(T\).  Continuity of each \(s_N\) in operator norm and Fatou's lemma give

\[
 \|T_j-T\|_{\mathfrak W^{[r]}}
 \leq\liminf_{k\to\infty}
 \|T_j-T_k\|_{\mathfrak W^{[r]}},
\]

which proves completeness.

Let \(T^{(m)}\) be the truncation to the first \(m\) singular directions.
If

\[
 \varepsilon_N(T):=\frac{Ns_N(T)}{(\log(e+N))^{r-1}}\longrightarrow0,
\]

then

\[
 \|T-T^{(m)}\|_{\mathfrak W^{[r]}}
 \leq C_r\sup_{k\geq m+1}\varepsilon_k(T)\longrightarrow0.
\]

Conversely, if \(F\) has finite rank and
\(\|T-F\|_{\mathfrak W^{[r]}}\leq\varepsilon\), the rank-perturbation
inequality gives

\[
 \limsup_{N\to\infty}
 \frac{Ns_N(T)}{(\log(e+N))^{r-1}}
 \leq C_r\varepsilon.
\]

Approximation proves the little condition, and also its closedness.  The
inclusion into the next little grade follows from the extra logarithm in the
denominator.

For multi-flattenings, let \(K_j\) be Cauchy in
\eqref{cc:gen-multiflat-norm}.  Its Hilbert limit is a tensor \(K\).  For
each cut, completeness of \(\mathfrak W^{[r]}\) gives an operator limit,
while Hilbert convergence gives the same limit in \(\mathcal S_2\), hence
in operator norm.  Uniqueness identifies it with
\(\operatorname{Mat}_S(K)\).  There are finitely many cuts, so the whole
sequence converges in the stated quasi-norm.

If \(K\) belongs to the little space and \(P_m\uparrow I\) are finite-rank
orthogonal projections, put \(K_m=P_m^{\otimes n}K\).  For every cut,

\[
 \operatorname{Mat}_S(K_m)
 =P_m^{\otimes S^c}\operatorname{Mat}_S(K)P_m^{\otimes S}.
\]

Approximate the middle operator in \(\mathfrak W^{[r]}\) by a finite-rank
operator and use uniform convergence of \(P_m\) on its finite-dimensional
source and range.  This proves \(K_m\to K\) in every cut and in Hilbert
norm.  The converse follows because each flattening of a finite-dimensional
tensor has finite rank and the little ideal is closed.  Finally, a
tensor-leg permutation merely relabels the cuts and conjugates the
flattenings by unitaries.
\end{proof}

\begin{lemma}[Finite logarithmic addition]
\label{cc:gen-lemma-finite-addition}
For \(r\geq1\), \(m\geq2\), and compact operators \(T_1,\ldots,T_m\),

\begin{equation}
 \left\|\sum_{j=1}^mT_j\right\|_{\mathfrak W^{[r]}}
 \leq C_r(1+\log m)^r
 \sum_{j=1}^m\|T_j\|_{\mathfrak W^{[r]}}.
 \label{cc:gen-finite-addition}
\end{equation}

The same estimate, with a constant also depending on the fixed tensor
degree, holds in every multi-flattening space.
\end{lemma}

\begin{proof}
Write \(a_j=\|T_j\|_{\mathfrak W^{[r]}}\) and
\(A=\sum_ja_j\).  Fix \(u>0\).  Split the singular expansion of \(T_j\)
as \(H_j+M_j+R_j\), according as its singular values lie in

\[
 (u/4,\infty),\qquad (u/(4m),u/4],\qquad[0,u/(4m)].
\]

Then \(\|\sum_jR_j\|\leq u/4\).  By
\eqref{cc:gen-distribution-bound},

\begin{equation}
 \sum_j\operatorname{rank}H_j
 \leq C_r\frac Au
 \left(\log\left(e+\frac Au\right)\right)^{r-1}.
 \label{cc:gen-high-rank}
\end{equation}

For \(a=u/(4m)\), \(b=u/4\), scalar layer-cake and
\eqref{cc:gen-distribution-bound} give

\begin{align*}
 \sum_j\|M_j\|_1
 &\leq a\sum_jd_{T_j}(a)
 +\sum_j\int_a^bd_{T_j}(t)\,\mathrm dt\\
 &\leq C_rA(1+\log m)^r
 \left(\log\left(e+\frac Au\right)\right)^{r-1}.
\end{align*}

Here one uses \(a_j\leq A\) and

\[
 (x+y)^r-x^r\leq C_r(1+y)^r(1+x)^{r-1},
 \qquad x,y\geq0,
\]

after changing variables to logarithmic scale.  If
\(H=\sum H_j\), \(M=\sum M_j\), and \(R=\sum R_j\), then

\[
 d_{H+M+R}(u)
 \leq\operatorname{rank}H+d_M(3u/4)
 \leq\operatorname{rank}H+\frac{4}{3u}\|M\|_1.
\]

Together with \eqref{cc:gen-high-rank}, this is
\eqref{cc:gen-distribution-bound} for the sum, with parameter
\(C_r(1+\log m)^rA\).  The converse half of
\Cref{cc:gen-lemma-critical-spaces} proves
\eqref{cc:gen-finite-addition}.  Apply it to the finitely many flattenings
and use the ordinary Hilbert triangle inequality for the tensor term.
\end{proof}

\begin{theorem}[Logarithmic tensor calculus and quasi-Banach sewing]
\label{cc:gen-theorem-product-sewing}
For fixed integers \(r,s\geq1\),

\begin{align}
 \|A\otimes B\|_{\mathfrak W^{[r+s]}}
 &\leq C_{r,s}
 \|A\|_{\mathfrak W^{[r]}}
 \|B\|_{\mathfrak W^{[s]}},
 \label{cc:gen-log-product}\\
 \|A\otimes C\|_{\mathfrak W^{[r]}}
 &\leq C_r
 \|A\|_{\mathfrak W^{[r]}}\|C\|_{\mathcal S_1}.
 \label{cc:gen-trace-product}
\end{align}

If one positive-grade factor is little, the corresponding output is little.
The same assertions hold flattening by flattening for concrete tensor
products in every fixed tensor degree.  First-order Hilbert vectors and
finite-rank factors have zero logarithmic cost.

Let \(B_{s,u,t}\) be a closed three-increment with values in
\(\mathfrak W_{\mathrm{mf}}^{[r]}(H^{\otimes n})\), and suppose

\[
 \|B_{s,u,t}\|_{\mathfrak W_{\mathrm{mf}}^{[r]}(n)}
 \leq C\omega(s,t)^\gamma,
 \qquad \gamma>1,
\]

for a continuous control \(\omega\).  There is a unique two-increment
\(\Lambda B\) with \(\delta\Lambda B=B\) and

\begin{equation}
 \|\Lambda B_{s,t}\|_{\mathfrak W_{\mathrm{mf}}^{[r]}(n)}
 \leq C_{n,r,\gamma}C\omega(s,t)^\gamma.
 \label{cc:gen-sewing-bound}
\end{equation}

If all defects are little, every sewn increment is little.
\end{theorem}

\begin{proof}
After normalization, write

\[
 s_i(A)\leq\frac{(\log(e+i))^{r-1}}i,
 \qquad
 s_j(B)\leq\frac{(\log(e+j))^{s-1}}j.
\]

The singular values of \(A\otimes B\) are the decreasing rearrangement of
their pairwise products.  Dividing the \((i,j)\)-plane into the rectangles
\(2^k\leq i<2^{k+1}\), \(2^\ell\leq j<2^{\ell+1}\), and summing first over
the rectangles for which the displayed product exceeds \(u\), gives

\[
 \#\{(i,j):s_i(A)s_j(B)>u\}
 \leq C_{r,s}u^{-1}
 \bigl(\log(e+u^{-1})\bigr)^{r+s-1}.
\]

Indeed, on the \((k,\ell)\) rectangle the product is bounded by
\(C_{r,s}(1+k)^{r-1}(1+\ell)^{s-1}2^{-k-\ell}\), while the rectangle has
at most \(2^{k+\ell+2}\) points; summation over the terminal diagonal and
then over its predecessors is the discrete beta convolution

\[
 \sum_{k=0}^{L}(1+k)^{r-1}(1+L-k)^{s-1}
 \leq C_{r,s}(1+L)^{r+s-1}.
\]

The distribution dictionary proves \eqref{cc:gen-log-product}.  For a
trace-class factor, use

\[
 d_{A\otimes C}(u)
 =\sum_jd_A\!\left(\frac{u}{s_j(C)}\right)
\]

and \(\sum_js_j(C)=\|C\|_1\) in
\eqref{cc:gen-distribution-bound}; monotonicity of the logarithm gives
\eqref{cc:gen-trace-product}.  If, for example,
\(A\in\mathfrak W^{[r]}_0\), choose finite-rank \(A_m\to A\) in
\(\mathfrak W^{[r]}\).  The big product estimate gives
\(A_m\otimes B\to A\otimes B\) in
\(\mathfrak W^{[r+s]}\).  Each \(A_m\otimes B\) is a finite union of
scaled copies of the singular sequence of \(B\); the additional \(r\)
logarithmic powers in the denominator therefore put it in
\(\mathfrak W^{[r+s]}_0\).  Closedness proves the first little
assertion.  The trace-class case is identical, since a finite-rank tensor
product with a trace-class operator is trace class.  A tensor
product of flattenings represents every flattening of a concrete tensor
product, up to leg unitaries, so the multi-flattening assertions follow.

For sewing, bisect each interval successively by \(\omega\)-mass.  At
level \(k\), the difference between two consecutive approximants is a sum
of at most \(2^k\) defects, each bounded by
\(C2^{-k\gamma}\omega(s,t)^\gamma\).  By
\Cref{cc:gen-lemma-finite-addition},

\[
 \|D_k\|_{\mathfrak W_{\mathrm{mf}}^{[r]}(n)}
 \leq C_{n,r}(1+k)^r2^{-k(\gamma-1)}
 C\omega(s,t)^\gamma.
\]

By the Aoki--Rolewicz theorem \cite{KaltonPeckRoberts84}, the complete
quasi-Banach space in
\Cref{cc:gen-lemma-critical-spaces} has an equivalent \(p\)-quasi-norm,
\(0<p\leq1\), satisfying
\(\|x+y\|_*^p\leq\|x\|_*^p+\|y\|_*^p\).  The displayed geometric
decay is \(p\)-summable for every \(\gamma>1\), proving convergence and
\eqref{cc:gen-sewing-bound}.  Closedness of the little space gives the
little assertion.  The difference of two candidates is additive and has
order \(\omega^\gamma\); subdivision into \(2^k\) equal-control pieces and
\Cref{cc:gen-lemma-finite-addition} makes its norm tend to zero, proving
uniqueness.
\end{proof}

\subsection{Exact ancestry layers and controlled entrances}
\label{cc:gen-subsection-critical-class}

We now specialize the full normal filtration of
\eqref{cc:eq-exact-layer} to a fixed tensor degree over a Hilbert
alphabet.  The stable real decomposition, compatible ancestry flag, and
adapted complements are those of
\Cref{cc:thm-schur-exact-projectors}; throughout this subsection we use
the same fixed real-form-compatible splitting chosen there.  In its
notation set, explicitly,
\[
 F_rM_{\lambda,n}:=M^{\geq r}_{\lambda,n},
 \qquad
 G_{\lambda,r}:=M_{\lambda,n,r},
 \qquad
 F_rM_{\lambda,n}
 =G_{\lambda,r}\oplus F_{r+1}M_{\lambda,n}.
\]
Thus
\[
 \operatorname{gr}^{\mathrm{anc}}_rM_{\lambda,n}
 :=F_rM_{\lambda,n}/F_{r+1}M_{\lambda,n}
 \simeq G_{\lambda,r}.
\]
The theorem's \(P_{\lambda,r}\) projects onto
\({\mathbf S}_{\lambda}H_{\mathbb C}\otimes G_{\lambda,r}\), kills
\(F_{r+1}{\mathfrak L}_n\), and is a real finite permutation polynomial
natural in the following precise sense: for every bounded map
\(A:H\to K\),
\[
 A^{\otimes n}P_{\lambda,r}^{H}
 =P_{\lambda,r}^{K}A^{\otimes n}
 \quad\hbox{on }L_n(H).
\]
For \(n\geq2\), \(L_n(H)\subseteq D(H)=F_1L(H)\), so
\(F_0L_n(H)=F_1L_n(H)\) and the grade-zero multiplicity complement is
\(G_{\lambda,0}=0\).

Since every leg permutation is an isometry of the spaces in
\Cref{cc:gen-lemma-critical-spaces}, write the fixed projector as
\(P_{\lambda,r}=\sum_{j=1}^{m_P}c_jU_j\), where the \(U_j\) are leg
permutations.  Applying the fixed-\(q\), finite-sum estimate of
\Cref{cc:gen-lemma-finite-addition} (and the ordinary triangle
inequality for the Hilbert component) gives, for every
\(q\in\mathbb N_{\geq1}\),
\begin{equation}
 \|P_{\lambda,r}K\|_{\mathfrak W_{\mathrm{mf}}^{[q]}(n)}
 \leq C_{n,\lambda,r,P,q}
 \|K\|_{\mathfrak W_{\mathrm{mf}}^{[q]}(n)}.
 \label{cc:gen-projector-multiflat-bound}
\end{equation}
Here \(C_{n,\lambda,r,P,q}<\infty\) depends on the fixed finite
permutation polynomial and on \(q\), but not on the tensor being
projected or on the ambient Hilbert-space dimension; for example it may
be chosen as a fixed-degree multiple of
\[
 \left(1+C_q(1+\log m_P)^q\right)
 \sum_{j=1}^{m_P}|c_j|.
\]

\begin{definition}[Controlled critical normal entrances]
\label{cc:gen-definition-critical-class}
Fix \(\beta\in(1/3,1/2)\) and \(T>0\).  An element of
\(\mathfrak C_{\beta}^{\mathrm{crit}}(T)\) is a real Hilbert space
\(H\), a continuous control \(\omega\) on \([0,T]\), and a
multiplicative step-two geometric increment
\[
 {\mathbf X}^{(2)}_{s,t}
 =
 \left(
  1,\,
  X_{s,t},\,
  \frac12X_{s,t}^{\otimes2}+A_{s,t}
 \right)
\]
such that
\begin{equation}
 \omega(0,T)\leq1,\qquad
 \|X_{s,t}\|_H\leq\omega(s,t)^\beta,\qquad
 \|A_{s,t}\|_{\mathfrak W_{\mathrm{mf}}^{[1]}(2)}
 \leq\omega(s,t)^{2\beta}.
 \label{cc:gen-unit-critical-control}
\end{equation}
The class ranges over finite-dimensional and separable
infinite-dimensional \(H\).  A random realization is admissible on an
event \(G\) when \eqref{cc:gen-unit-critical-control} holds pathwise on
\(G\) with one deterministic control.  Spatial dilation by \(\rho\)
sends \(X,A,L_n\) to \(\rho X,\rho^2A,\rho^nL_n\); it therefore changes
constants but not logarithmic grade.
\end{definition}

\begin{proposition}[Entrance-to-Rees bridge]
\label{cc:prop-entrance-to-Rees-bridge}
Let \((X,A,Q)\) be a \(\beta\)-controlled critical second-order entrance.
Then its geometricized second level
\begin{equation}
 {\mathbf X}^{(2)}_{s,t}
 :=\left(1,X_{s,t},\tfrac12X_{s,t}^{\otimes2}+A_{s,t}\right)
 \label{cc:eq-entrance-geometricization}
\end{equation}
is multiplicative.  After one deterministic spatial rescaling and a fixed
renormalization of the control, it belongs to
\(\mathfrak C_\beta^{\mathrm{crit}}(T)\).  The collision coordinate \(Q\)
is a trace-class decoration and does not alter the positive critical cost.
The construction commutes with bounded spatial maps, optional evaluation,
and ucp projection limits.
\end{proposition}

\begin{proof}
The split Chen law gives
\[
 A_{s,t}=A_{s,u}+A_{u,t}+\tfrac12X_{s,u}\wedge_{\mathrm{op}}X_{u,t}.
\]
Substitution into \eqref{cc:eq-entrance-geometricization} is exactly the
step-two geometric Chen identity.  The entrance bounds give
\eqref{cc:gen-unit-critical-control} up to fixed constants; spatial dilation
and multiplication of the control normalize those constants and its total
mass.  The reconstruction
\(\mathbb X^I=\frac12X^{\otimes2}+A-\frac12Q\) shows that geometricization
adds only the trace-class collision correction.  The remaining assertions
follow from the naturality and optional/projection stability of the critical
split target.
\end{proof}

\begin{definition}[Universal and minimax spectral depth]
\label{cc:gen-definition-minimax-depth}
Fix the stable natural exact-layer family
\[
 G^{\mathrm{Schur}}_{\lambda,n,r}(-):
 H\longmapsto
 {\mathbf S}_{\lambda}H_{\mathbb C}\otimes
 \operatorname{gr}^{\mathrm{anc}}_rM_{\lambda,n},
 \qquad r\geq1,
\]
assume that its stable multiplicity space is nonzero, and fix the
adapted projector \(P_{\lambda,r}\).  This is a polynomial-functor
type: evaluating it on a particular Hilbert space is not part of the
spectral-depth datum.  Its ancestry depth is
\(d_{\mathrm{anc}}(G^{\mathrm{Schur}}_{\lambda,n,r})=r\).  Its
universal critical spectral depth
\(d_{\mathrm{sp}}^{(\beta)}(G^{\mathrm{Schur}}_{\lambda,n,r})\) is the
least \(q\in\mathbb N_{\geq1}\) for which a constant
\(C_{G,P,\beta}\) satisfies
\begin{equation}
 \|P_{\lambda,r}L_n(s,t)\|_{\mathfrak W_{\mathrm{mf}}^{[q]}(n)}
 \leq C_{G,P,\beta}\omega(s,t)^{n\beta}
 \label{cc:gen-universal-grade}
\end{equation}
for every \(T\), every member of
\(\mathfrak C_{\beta}^{\mathrm{crit}}(T)\), and every \(s<t\).
The constant is uniform over the natural family, in particular
independent of the ambient Hilbert-space dimension.

An integer grade \(q\in\mathbb N_{\geq1}\) is called minimax sharp if
there are \(T_G,N_G,c_G,p_G>0\) and one fixed nontrivial cut
\(S_G\subset[n]\), chosen for the natural family before the spectral
window, with the following property.  For every \(L\geq N_G\) there is
a finite-dimensional random realization and an event \(G_L\) of
probability at least \(p_G\) on
which the realization belongs to
\(\mathfrak C_{\beta}^{\mathrm{crit}}(T_G)\) and
\begin{equation}
 s_N\!\left(
 \operatorname{Mat}_{S_G}
 (P_{\lambda,r}L_n(0,T_G))
 \right)
 \geq c_G\phi_q(N),
 \qquad N_G\leq N\leq L.
 \label{cc:gen-minimax-lower}
\end{equation}
All constants and algebraic data are independent of the ambient
dimension; only the finite block dimensions may depend on \(L\).  In particular,
minimax sharpness is not the assertion that every realization, or every
flattening of one realization, is extremal.
\end{definition}

\begin{theorem}[Chronological logarithm: universal upper bound]
\label{cc:gen-theorem-chronological-upper}
Every \((H,\omega,{\mathbf X}^{(2)})\in
\mathfrak C_{\beta}^{\mathrm{crit}}(T)\) has, for every fixed
\(N_*\), a unique multiplicative geometric extension through degree
\(N_*\).  If
\[
 L_n(s,t):=
 \pi_n\log_\otimes{\mathbf X}^{(\leq N_*)}_{s,t},
\]
then every nonzero positive exact layer satisfies
\begin{equation}
 \|P_{\lambda,r}L_n(s,t)\|_
 {\mathfrak W_{\mathrm{mf}}^{[r]}(n)}
 \leq C_{n,\lambda,r,\beta,P}\,
 \omega(s,t)^{n\beta}.
 \label{cc:gen-chronological-upper}
\end{equation}
If every \(A_{s,t}\) belongs to
\(\mathfrak W_{0,\mathrm{mf}}^{[1]}\), then every displayed output
belongs to \(\mathfrak W_{0,\mathrm{mf}}^{[r]}\).
\end{theorem}

\begin{proof}
By \Cref{cc:gen-lemma-critical-spaces}, the critical degree-two norm
controls the Hilbert tensor norm.  The Hilbert tensor norms are an
admissible crossnorm system:
\[
 \|u\otimes v\|_{H^{\otimes(a+b)}}
 =\|u\|_{H^{\otimes a}}\|v\|_{H^{\otimes b}},
\]
and tensor-leg permutations are unitary.  Hence the hypotheses of the
Lyons extension theorem for a geometric \(p\)-rough path are satisfied
with \(p=1/\beta<3\); see \cite{Lyons98}.  This supplies
the unique finite-degree extension and its Hilbert estimate
\(\|L_j(s,t)\|_{H^{\otimes j}}\lesssim_j\omega(s,t)^{j\beta}\).

Multiplicativity gives
\[
 L_{s,t}=\operatorname{BCH}(L_{s,u},L_{u,t}).
\]
For \(n\geq3\), remove the two terms linear in \(L_n\).  The result is
the closed three-increment identity
\begin{equation}
 \delta L_n(s,u,t)
 =
 {\mathcal B}_n\bigl(L_{<n}(s,u),L_{<n}(u,t)\bigr),
 \label{cc:gen-bch-defect}
\end{equation}
where \({\mathcal B}_n\) is a fixed finite Lie polynomial in lower
homogeneous degrees.

Induct on \(n\).  A positive exact component of degree two is the
hypothesis.  Expand each monomial on the right of
\eqref{cc:gen-bch-defect} into exact ancestry components.  If its
positive depths are \(r_1,\ldots,r_\ell\), let
\(s=r_1+\cdots+r_\ell\).  Strong filteredness
\([F_a{\mathfrak L},F_b{\mathfrak L}]\subset F_{a+b}{\mathfrak L}\)
places the monomial in \(F_s{\mathfrak L}\).  Iterated use of
\Cref{cc:gen-theorem-product-sewing} places every flattening in
\(\mathfrak W^{[s]}\); first-level Hilbert factors are rank one in
each induced cut and add no logarithmic cost.  A monomial made only of
first-level factors has uniformly bounded finite rank.

Apply \(P_{\lambda,r}\).  Terms with \(s>r\) vanish by
\eqref{cc:eq-exact-projector-filtration}; those with \(s\leq r\) lie in
\(\mathfrak W_{\mathrm{mf}}^{[r]}\) by nesting.  If the ordinary
degrees of all factors are \(j_1,\ldots,j_k\), then
\[
 j_1+\cdots+j_k=n,
 \qquad
 \prod_i\omega(s,t)^{j_i\beta}=\omega(s,t)^{n\beta}.
\]
Since \(n\beta>1\), \Cref{cc:gen-theorem-product-sewing} sews the
projected defect to an increment
\(\widetilde L_{n,\lambda,r}\) satisfying
\eqref{cc:gen-chronological-upper}.

The actual \(P_{\lambda,r}L_n\) has the same coboundary.  Their
difference is additive and, first in Hilbert norm, is
\(O(\omega^{n\beta})\).  A subdivision into pieces of vanishing control
mass forces this additive difference to be zero because \(n\beta>1\).
This identifies the endpoint-valued sewn increment with the canonical
Lyons extension without presupposing the endpoint bound.

For the little assertion, a defect monomial containing a positive-depth
lower component is little by the inductive hypothesis and the little
product estimate.  A monomial containing only first-level vectors has
finite-rank flattenings and is therefore little as well.  The little
sewing statement of \Cref{cc:gen-theorem-product-sewing}, together with
closedness from \Cref{cc:gen-lemma-critical-spaces}, repeats the induction
in the little spaces.
\end{proof}

\subsection{The split Rees logarithm}
\label{cc:subsec-split-rees-logarithm}

For each \(n\ge2\), put
\[
 P_{n,r}:=\sum_{\lambda\vdash n}P_{\lambda,r},
 \qquad 1\le r\le\lfloor n/2\rfloor,
\]
where absent Schur sectors contribute zero.  The projectors are the fixed
real-form-compatible permutation polynomials of
\Cref{cc:thm-schur-exact-projectors}.  They split the intrinsic normal flag,
but the splitting itself is auxiliary.

\begin{definition}[Algebraic ancestry and analytic critical cost]
\label{cc:def-ancestry-response-rees-modules}
For a cutoff \(N\ge2\), the split ancestry Rees module is
\begin{equation}
 \mathscr R^{\rm anc}_{\le N}(H)
 :=H\oplus
 \bigoplus_{n=2}^N\ \bigoplus_{r=1}^{\lfloor n/2\rfloor}
 z^{-r}\operatorname{Ran}P_{n,r}.
 \label{cc:eq-ancestry-rees-module}
\end{equation}
The symbol \(z^{-r}\) records the algebraic normal order
\(F_rL_n/F_{r+1}L_n\).

For tensors set
\begin{equation}
 \mathscr K_0^{[0]}:=\mathbb R,
 \qquad \mathscr K_1^{[0]}(H):=H,
 \qquad \mathscr K_n^{[0]}(H):=\mathfrak A_1^{(n)}(H)\quad(n\ge2),
 \label{cc:eq-cost-zero-spaces}
\end{equation}
and, for \(n\ge2\), \(r\ge1\),
\begin{equation}
 \mathscr K_n^{[r]}(H):=
 \mathfrak W_{\rm mf}^{[r]}(H^{\otimes n}).
 \label{cc:eq-positive-cost-spaces}
\end{equation}
Thus cost zero means trace class in every cut, while cost \(r\) permits the
critical profile \(\phi_r(N)\).  The inclusions
\(\mathscr K_n^{[r-1]}\subseteq\mathscr K_n^{[r]}\) define the analytic
critical boundary
\begin{equation}
 \partial_{\rm cr}^0\mathscr K_n(H):=\mathscr K_n^{[0]}(H),
 \qquad
 \partial_{\rm cr}^r\mathscr K_n(H)
 :=\mathscr K_n^{[r]}(H)/\mathscr K_n^{[r-1]}(H),
 \quad r\ge1.
 \label{cc:eq-critical-associated-graded}
\end{equation}
The truncated analytic Rees space is
\begin{equation}
 \mathscr R^{\rm cr}_{\le N}(H)
 :=\bigoplus_{n=0}^N
   \bigoplus_{0\le r\le\lfloor n/2\rfloor}
   z^{-r}\mathscr K_n^{[r]}(H),
 \label{cc:eq-response-rees-module}
\end{equation}
with the degree-one summand understood as \(H\).  The direct Rees sum carries
its product quasi-topology.  By contrast, the quotients
\(\partial_{\rm cr}^r\mathscr K_n\) are used as algebraic associated-graded
spaces; no Hausdorff quotient topology is implicit.
\end{definition}

\begin{proposition}[Algebraic nature of the critical boundary]
\label{cc:prop-algebraic-critical-boundary}
For every \(r\ge1\), the inclusion
\(\mathfrak W^{[r-1]}\subset\mathfrak W^{[r]}\), with
\(\mathfrak W^{[0]}:=\mathcal S_1\), has nonclosed range in the
\(\mathfrak W^{[r]}\)-quasi-norm.  The same is true in degree two for the
multi-flattening spaces.  Consequently
\[
 \partial_{\rm cr}^r\mathscr K_n
 =\mathscr K_n^{[r]}/\mathscr K_n^{[r-1]}
\]
is deliberately an algebraic quotient.  All strictness and support statements
below depend only on membership in the filtration.  Exact or finite-window
asymptotics are asserted only after a representative-level packet certificate
has been supplied.
\end{proposition}

\begin{proof}
For \(r=1\), the diagonal operator with singular values
\[
 d_N=\frac1{N\log(e+N)}
\]
belongs to \(\mathfrak W^{[1]}_0\setminus\mathcal S_1\).  Its finite-rank
truncations lie in \(\mathcal S_1\) and converge to it in
\(\mathfrak W^{[1]}\).  For \(r\ge2\), take instead
\[
 d_N=\frac{(\log(e+N))^{r-3/2}}N.
\]
Then \(d_N/\phi_r(N)\to0\), whereas
\(d_N/\phi_{r-1}(N)\to\infty\).  Thus the corresponding diagonal operator
lies in \(\mathfrak W^{[r]}_0\setminus\mathfrak W^{[r-1]}\), while its
finite-rank truncations belong to every lower grade and converge in the
higher quasi-norm.  Degree two identifies the only nontrivial flattening
with the operator itself, proving the multi-flattening assertion.
\end{proof}

\begin{theorem}[Analytic critical Rees algebra]
\label{cc:thm-analytic-critical-rees-algebra}
Truncated tensor concatenation makes
\(\mathscr R^{\rm cr}_{\le N}(H)\) a graded topological algebra.  More
precisely, whenever the total tensor degree is at most \(N\),
\begin{equation}
 \mathscr K_m^{[a]}(H)\widehat\otimes
 \mathscr K_n^{[b]}(H)
 \longrightarrow
 \mathscr K_{m+n}^{[a+b]}(H)
 \label{cc:eq-critical-cost-addition}
\end{equation}
is continuous, with the conventions in
\eqref{cc:eq-cost-zero-spaces}.  Lowering either positive input cost by one
lowers the product cost by one.  Hence tensor product descends to the algebraic associated graded:
\begin{equation}
 \partial_{\rm cr}^a\mathscr K_m(H)\otimes
 \partial_{\rm cr}^b\mathscr K_n(H)
 \longrightarrow
 \partial_{\rm cr}^{a+b}\mathscr K_{m+n}(H),
 \qquad [A]\otimes[B]\longmapsto[A\otimes B].
 \label{cc:eq-critical-graded-product}
\end{equation}
If \(B\) is a critical atom with \(s_N(B)\sim c/N\), then
\begin{equation}
 s_N(B^{\otimes r})
 \sim\frac{c^r}{\Gamma(r)}\phi_r(N),
 \qquad
 [B]^{\otimes r}\ne0
 \quad\text{in }\partial_{\rm cr}^r.
 \label{cc:eq-exact-atom-Rees-power}
\end{equation}
Thus tensor resonance is multiplication on the analytic critical boundary,
not a separate calculation attached to each stochastic model.
\end{theorem}

\begin{proof}
For cost \(0+0\), use the exact projective tensor law in
\Cref{cc:gen-prop-permutation-stable-tensor-algebra} at \(p=1\).  A
cost-zero factor is trace class in every cut and therefore preserves every
positive cost by the trace-class decoration estimate
\eqref{cc:gen-trace-product}.  Positive costs add by
\eqref{cc:gen-log-product}.  Rank-one factors in tensor degree one obey the
same estimates.  These bounds also give
\[
 \mathscr K_m^{[a-1]}\widehat\otimes\mathscr K_n^{[b]}
 +\mathscr K_m^{[a]}\widehat\otimes\mathscr K_n^{[b-1]}
 \subseteq\mathscr K_{m+n}^{[a+b-1]},
\]
which proves that \eqref{cc:eq-critical-graded-product} is independent of
representatives.  In view of \Cref{cc:prop-algebraic-critical-boundary}, no
quotient topology is used here.  Finite quasi-Banach sums, equipped with equivalent
Aoki--Rolewicz metrics, give the asserted topological algebra.  Finally,
\Cref{cc:thm:tensor-critical-arithmetic} gives the equivalent in
\eqref{cc:eq-exact-atom-Rees-power}.  For \(r=1\), the harmonic divergence
of \(\sum_Ns_N(B)\) shows that \(B\notin\mathcal S_1\).  For \(r\ge2\),
\(\phi_r/\phi_{r-1}\to\infty\), so the tensor power does not belong to the
lower cost class.
\end{proof}

\begin{definition}[Split Rees realization]
\label{cc:def-critical-rees-compiler}
Let \((H,\omega,{\mathbf X}^{(2)})\) be a controlled critical normal
entrance, and let \(L_n(s,t)\) be the homogeneous logarithms of its geometric
extension.  Its split Rees logarithm through degree \(N\) is
\begin{equation}
 \mathbf C_{\le N}({\mathbf X})_{s,t}
 :=X_{s,t}+\sum_{n=2}^N\sum_{r=1}^{\lfloor n/2\rfloor}
 z^{-r}P_{n,r}L_n(s,t).
 \label{cc:eq-critical-rees-compiler}
\end{equation}
Changing the splitting changes the displayed lift, but not the first nonzero
normal coefficient modulo the next filtration level.
\end{definition}

\begin{theorem}[Filtered critical realization theorem]
\label{cc:thm-critical-rees-compilation}
Fix \(N\ge2\) and \(\beta\in(1/3,1/2)\).  The assignment
\eqref{cc:eq-critical-rees-compiler} has the following properties.
\begin{enumerate}[label=\textup{(\roman*)},leftmargin=2.7em]
\item It is a well-defined map from controlled critical entrances to
\(\mathscr R^{\rm cr}_{\le N}(H)\), and
\begin{equation}
 \|P_{n,r}L_n(s,t)\|_{\mathfrak W^{[r]}_{\rm mf}(n)}
 \le C_{N,\beta,P}\,\omega(s,t)^{n\beta}
 \label{cc:eq-rees-coefficient-bound}
\end{equation}
for every \(2\le n\le N\) and every exact depth \(r\).
\item If the input area is little weak trace on every interval, each Rees
coefficient belongs to the corresponding little logarithmic grade.
\item The construction is natural under bounded linear maps: the induced
map on every coefficient is the tensor pushforward, and its norm is bounded
by the appropriate power of the operator norm.
\item If \(x=L_n(s,t)\ne0\) and \(r_0=\nu_F(x)\), then
\(P_{n,r}x=0\) for \(r<r_0\), while the class of
\(P_{n,r_0}x\) in \(F_{r_0}L_n/F_{r_0+1}L_n\) is the intrinsic initial
normal symbol of \(x\).  Thus the least nonzero Rees cost is canonical even
though the full split lift is not.
\item Optional evaluation, a selected singular lift, and the split Rees realization
commute under the hypotheses of
\Cref{cc:prop-causal-normal-completion}.  In particular, no source application needs a second stopped or
gauge-covariance proof for its compiled logarithm.
\end{enumerate}
\end{theorem}

\begin{proof}
Parts~\textup{(i)} and \textup{(ii)} are the chronological upper theorem
\Cref{cc:gen-theorem-chronological-upper}, summed over the finite set of
indices.  Naturality follows from the natural permutation projectors and
functoriality of the geometric extension.  The chosen splitting satisfies
\[
 F_rL_n=\bigoplus_{q\ge r}\operatorname{Ran}P_{n,q}.
\]
Hence the first nonzero component of a vector of valuation \(r_0\) is its
class modulo \(F_{r_0+1}L_n\), proving part~\textup{(iv)}.  The final part is the functoriality statement of
\Cref{cc:prop-causal-normal-completion}, followed by naturality of the
logarithm and normal projection.
\end{proof}


\subsection{Pure-area cells and degree-exact Hall experiments}
\label{hall:subsec-pure-area}

For a Hilbert--Schmidt operator $B:E\to F$, define its skew dilation by
\begin{equation}
 \operatorname{Sk}(B):=
 \begin{pmatrix}0&-B^*\\ B&0\end{pmatrix}
 \quad\text{on }E\oplus F.
 \label{hall:eq-skew-dilation}
\end{equation}
It is skew-adjoint, its compression from $E$ to $F$ is $B$, and each nonzero
singular value of $B$ occurs twice in $\operatorname{Sk}(B)$.

\begin{lemma}[Pure-area operator cells]
\label{hall:lem-pure-area}
Let $A$ be a skew Hilbert--Schmidt operator on a real Hilbert space $H$.  For
every $p\in(2,3)$,
\begin{equation}
 \mathbf P^A_{s,t}:=\exp_\otimes((t-s)A)
 \label{hall:eq-pure-area-path}
\end{equation}
defines a weakly geometric step-two $p$-rough path with zero first level and
logarithm $(t-s)A$.  It is the $p$-variation limit of signatures of
finite-dimensional bounded-variation loops whenever $A$ is approximated in
Hilbert--Schmidt norm by finite-rank skew operators.  Its unique extension to
any fixed tensor degree is natural under bounded linear maps.
\end{lemma}

\begin{proof}
For finite-rank $A$, decompose it into finitely many planar skew blocks and
use rapidly traversed rectangular loops with the prescribed signed areas.
Trotter products converge in the step-two nilpotent group to
$\exp_\otimes(tA)$.  The homogeneous step-two distance between two pure-area
paths is bounded by a constant times $\|A-A'\|_2^{1/2}$, so finite-rank
approximation gives the general case.  The Lyons extension theorem supplies
all fixed higher levels; see \cite{FrizVictoir10,Lyons98}.
\end{proof}

A polarized area--marker Hall word is a multilinear Hall Lie word in symbols
$A_1,\ldots,A_r,x_1,\ldots,x_m$, where $\deg A_j=2$,
$\deg x_\ell=1$, and every symbol occurs once.  Its ordinary tensor degree is
$n=2r+m$.

\begin{lemma}[Area--marker spanning]
\label{hall:lem-area-marker-spanning}
For every $n,r$, the full exact layer
$G^{\mathrm{full}}_{n,r}(H)$ is spanned, after polarization, by classes of
area--marker Hall words with $r$ area symbols and $n-2r$ vector markers.
The same assertion holds after passage to every stable Schur multiplicity
quotient $\operatorname{gr}^{\mathrm{anc}}_rM_{\lambda,n}$.
\end{lemma}

\begin{proof}
By \Cref{cc:thm-normal-cone}, the exact layer is the degree-$n$ part of
$\mathfrak L_r(Q)$, where $Q=D/[D,D]$.  The Bianchi--Koszul description in
\Cref{cc:thm-bianchi-module} shows that $Q$ is generated as an $S(H)$-module
by the image of $\Lambda^2H$.  Substituting these generators into a Hall basis
of $\mathfrak L_r(Q)$ and polarizing the remaining ordinary variables proves
the claim.  Passage to a stable Schur quotient is functorial.
\end{proof}

If a geometric cell has logarithm
$g=x+A+R_{\ge3}$, append the straight closing segment $-x$ and write the
resulting group element as $g^\circ$.  The BCH formula gives
\begin{equation}
 \pi_1\log_\otimes g^\circ=0,
 \qquad
 \pi_2\log_\otimes g^\circ=A.
 \label{hall:eq-purified-cell}
\end{equation}
For group-like elements $g,h$, put
$[g,h]_{\rm grp}=ghg^{-1}h^{-1}$ and replace every bracket of a Hall word by
this group commutator.

\begin{lemma}[Degree-exact Hall realization]
\label{hall:lem-Hall-exact}
Let $w$ be a polarized area--marker Hall word.  Realize each area letter by a
pure-area cell, or by a purified geometric cell satisfying
\eqref{hall:eq-purified-cell}, and each marker by a straight segment.  The
finite path obtained by recursive group commutators satisfies
\begin{equation}
 \pi_n\log_\otimes\mathbf Z_w
 =w(A_1,\ldots,A_r;x_1,\ldots,x_m),
 \qquad n=2r+m,
 \label{hall:eq-Hall-leading}
\end{equation}
and every other logarithmic term has ordinary tensor degree strictly larger
than $n$.
\end{lemma}

\begin{proof}
If logarithms $X,Y$ begin in ordinary degrees $a,b$, the logarithm of
$e^Xe^Ye^{-X}e^{-Y}$ begins with $[X_a,Y_b]$ in degree $a+b$; every remaining
Lie monomial either repeats a positive-degree input or uses a higher-degree
term.  Induction on the Hall bracketing proves the assertion.  Formula
\eqref{hall:eq-purified-cell} handles nonzero first-level cells.
\end{proof}

\begin{lemma}[Polarized disjoint-support separation]
\label{hall:lem-polarized-label-separation}
Let
\[
 \Theta_H:(\Lambda^2H)^r\times H^m\longrightarrow H^{\otimes n},
 \qquad n=2r+m,
\]
be a nonzero multilinear natural tensor polynomial.  Polarize all slots and
let $H_r\simeq(\mathfrak S_2)^r$ act by swapping the two elementary labels in
each exterior input.  Then there is a word $b$ in pairwise distinct labels
and a representative $\tau_b\in\mathfrak S_n$ such that the signed fibre
coefficient
\begin{equation}
 \kappa_b:=\sum_{h\in H_r}c_{\tau_bh}\,\operatorname{sgn}_r(h)
 \label{hall:eq-signed-fibre-coefficient}
\end{equation}
is nonzero, where
\[
 \Theta_H=\sum_{\pi\in\mathfrak S_n}c_\pi U_\pi
\]
on the polarized elementary tensor slots.  The same coefficient is obtained
after replacing the one-dimensional labels by arbitrary mutually orthogonal
Hilbert blocks.
\end{lemma}

\begin{proof}
Naturality and multilinearity identify $\Theta$ with a finite element of the
permutation algebra on the $n$ polarized slots.  Take the universal label
space
\[
 H_{\rm lab}=\bigoplus_{j=1}^r
 (\mathbf k e_j^0\oplus\mathbf k e_j^1)
 \oplus\bigoplus_{\ell=1}^m\mathbf k f_\ell
\]
and evaluate the exterior inputs at $e_j^0\wedge e_j^1$ and the marker inputs
at $f_\ell$.  The value is nonzero: if it vanished, functoriality under
linear maps and the fact that decomposable wedges span every exterior square
would force the natural polynomial to vanish identically.

Expand the universal value in the basis of elementary label words.  The
permutations producing a fixed output word $b$ form a right coset
$\tau_bH_r$; the subgroup $H_r$ records only the internal swap in each area
pair.  Substitution of $e_j^0\wedge e_j^1$ multiplies the term indexed by
$h\in H_r$ by $\operatorname{sgn}_r(h)$.  Hence the coefficient of the basis
word $b$ is exactly \eqref{hall:eq-signed-fibre-coefficient}.  Since the
universal value is nonzero, at least one such coefficient is nonzero.
Replacing each label line by an orthogonal Hilbert block changes neither the
permutation coefficients nor the coset decomposition; coordinate projection
to the same label word recovers the identical $\kappa_b$.
\end{proof}

\begin{lemma}[Signed native corner]
\label{hall:lem-signed-corner}
Let $q$ be a nonzero labelled stable tensor probe on a polarized Hall word
$w$ of exact depth $r$.  Put $A_j=\operatorname{Sk}(B_j)$ on mutually
orthogonal copies $E_j\oplus F_j$, and place every marker on a separate
one-dimensional block.  There exist a nontrivial tensor cut $S$, norm-one
coordinate compressions $Q_-,Q_+$, tensor-leg unitaries $U,V$, and
$c_{q,w}\ne0$ such that
\begin{equation}
 Q_+\operatorname{Mat}_S\!\left(
 q_H(\pi_n\log_\otimes\mathbf Z_w)\right)Q_-
 =c_{q,w}\,U(B_1\otimes\cdots\otimes B_r)V.
 \label{hall:eq-native-corner}
\end{equation}
The reader data depend only on the finite labelled polynomial $q(w)$, not on
spectral truncation or on the dimensions of the blocks.
\end{lemma}

\begin{proof}
By the degree-exact Hall lemma,
\[
 \pi_n\log_\otimes\mathbf Z_w=w(A_1,\ldots,A_r;\text{markers});
\]
all other BCH terms have ordinary tensor degree strictly larger than $n$.
After applying $q$, polarize the $r$ area slots into their $2r$ elementary
legs.  The resulting nonzero natural multilinear polynomial satisfies
\Cref{hall:lem-polarized-label-separation}.  Choose a label word $b$ with
$\kappa_b\ne0$.

Project every elementary output leg to the block prescribed by $b$.  Terms
with another label word vanish by orthogonality.  The terms which remain form
exactly the fibre $\tau_bH_r$.  Internal swaps in an area pair change the
exterior input and the skew corner by the same sign, so the fibre collapses
to the single coefficient $\kappa_b$ rather than cancelling.  Replacing the
universal label lines by larger orthogonal blocks leaves this coefficient
unchanged.

Choose the representative of the fibre so that, for each $j$, the $E_j$ leg
lies on the input side of the cut and the $F_j$ leg lies on the output side.
Place all marker lines on either side as one-dimensional factors.  The
compression of
\[
 \operatorname{Sk}(B_j)=
 \begin{pmatrix}0&-B_j^*\\ B_j&0\end{pmatrix}
\]
from $E_j$ to $F_j$ is $B_j$; marker legs contribute fixed scalars, and the
remaining tensor reorderings are the unitaries $U,V$.  Thus the surviving
fibre is
$\kappa_bU(B_1\otimes\cdots\otimes B_r)V$.  Set
$c_{q,w}=\kappa_b$ after absorbing the fixed exterior-normalization and marker
scalars.  Hall degree exactness excludes terms from other ordinary degrees,
and the label projection has already removed every other term in degree $n$.
This proves \eqref{hall:eq-native-corner} with no same-degree or lower-cost
remainder.
\end{proof}

\subsection{The native response kernel}
\label{hall:subsec-native-kernel}

Fix $n$, a partition $\lambda\vdash n$, and the filtered stable multiplicity
space $M_{\lambda,n}$ from
\Cref{cc:gen-subsection-critical-class}.  Interpret
$q\in M_{\lambda,n}^*$ as the corresponding labelled natural tensor map.
Fix $\beta\in(1/3,1/2)$ and the universal separable Hilbert space
$\mathcal H_\infty=\ell^2$.  Let $\mathscr I_{\lambda,n,\beta}$ be the set of
tuples consisting of a unit-normalized entrance in
$\mathfrak C_\beta^{\mathrm{crit}}(1)$ on a closed subspace of
$\mathcal H_\infty$, an interval $0\le s<t\le1$, a nontrivial tensor cut,
and norm-one spatial compressions.  Fix unitary identifications of the finite
tensor powers of $\mathcal H_\infty$ with standard separable domain and range
spaces, and extend each compressed corner by zero.  For $r\ge0$ set
\begin{equation}
 K_r\mathcal Y_{\lambda,n,\beta}^{\rm nat}
 :=\ell^\infty\!\left(
 \mathscr I_{\lambda,n,\beta};\mathfrak W^{[r]}\right),
 \label{hall:eq-native-target-filtration}
\end{equation}
with the trace-class convention at cost zero, and set
\[
 \mathcal Y_{\lambda,n,\beta}^{\rm nat}
 :=K_{\lfloor n/2\rfloor}\mathcal Y_{\lambda,n,\beta}^{\rm nat}.
\]
The resulting labelled response is
\begin{equation}
 \mathcal R_{\lambda,n,\beta}^{\rm nat}:M_{\lambda,n}^*
 \longrightarrow\mathcal Y_{\lambda,n,\beta}^{\rm nat}.
 \label{hall:eq-native-response}
\end{equation}
Every separable Hilbert-space experiment is represented after an isometric
embedding into $\mathcal H_\infty$, so this reduction discards no native reader.
The normalization
\eqref{cc:gen-unit-critical-control} and
\Cref{cc:gen-theorem-chronological-upper} give the uniform bound required in
\eqref{hall:eq-native-target-filtration}; hence the response is filtered for
the dual ancestry filtration.  By
\Cref{rk:prop-robust-support}, strictness on each finite exact layer is already
certified by finitely many coordinates of this product.  We set
\begin{equation}
 d_{\rm sp}^{\rm nat}(q)
 :=\nu_K(\mathcal R_{\lambda,n,\beta}^{\rm nat}q),
 \label{hall:eq-native-spectral-depth-definition}
\end{equation}
with value $-\infty$ at $q=0$.

\begin{theorem}[Full-support native Hall kernel]
\label{hall:thm-native-kernel}
For every nonzero exact layer
$\operatorname{gr}^{\mathrm{anc}}_rM_{\lambda,n}$, finitely many pure-area
Hall experiments and signed native corners have rank-one pushed-forward Rees
kernels
\begin{equation}
 K_{r,e}=\sigma_e\otimes\tau_r,
 \qquad
 \tau_r=[D_\infty^{\otimes r}],
 \label{hall:eq-rank-one-kernels}
\end{equation}
up to nonzero scalars and fixed leg unitaries, and the symbols
$\{\sigma_e\}$ span the whole exact layer.  Hence
$\mathcal R_{\lambda,n,\beta}^{\rm nat}$ is strict and coefficient-pure.
\end{theorem}

\begin{proof}
Choose finitely many area--marker Hall words whose classes span the exact
layer, using \Cref{hall:lem-area-marker-spanning}.  In every area slot take
$B_j=D_\infty$ and apply
\Cref{hall:lem-Hall-exact,hall:lem-signed-corner}.  Each native corner is a
nonzero scalar multiple of $D_\infty^{\otimes r}$, so its pushed-forward
kernel has the form \eqref{hall:eq-rank-one-kernels}.  The left supports are
the lines $\mathbf k\sigma_e$ and their sum is the exact layer.  Apply
\Cref{rk:thm-reader-symbols,rk:thm-response-consequences}.
\end{proof}

\begin{corollary}[Native valuation exchange and flag reconstruction]
\label{hall:cor-native-valuation}
For every nonzero $q\in M_{\lambda,n}^*$,
\begin{equation}
 d_{\rm sp}^{\rm nat}(q)=a_F(q),
 \label{hall:eq-native-depth}
\end{equation}
and
\begin{equation}
 F_rM_{\lambda,n}
 =\bigcap_{\substack{q\in M_{\lambda,n}^*\\
                d_{\rm sp}^{\rm nat}(q)\le r-1}}
 \ker q.
 \label{hall:eq-native-flag}
\end{equation}
Moreover, some fixed designed geometric response of $q$ satisfies
\begin{equation}
 s_N(T_q)
 \sim C_q\frac{(\log N)^{a_F(q)-1}}{N},
 \qquad C_q>0.
 \label{hall:eq-native-exact-law}
\end{equation}
Thus forward depth and dual reconstruction require one full-support kernel,
not two separate exposure constructions.
\end{corollary}

\begin{proof}
Apply
\Cref{rk:thm-valuation-exchange,rk:thm-response-consequences} to
\Cref{hall:thm-native-kernel}.  The exact coefficient is supplied by
\eqref{rk:eq-critical-boundary-line}.
\end{proof}

\subsection{One universal Hall test path}
\label{hall:subsec-universal-path}

Fix $p\in(2,3)$.  Enumerate the countable collection of polarized
area--marker Hall patterns and their finitely many signed label corners by
$(w_k)_{k\ge1}$.  Let $H_k$ be the orthogonal Hilbert block of the $k$-th
local experiment and $\mathbf Z_k$ its pure-area Hall path.  Put
$t_k=1-2^{-k}$, $t_0=0$, and $I_k=[t_{k-1},t_k]$.  Choose scales
$\varepsilon_k\in(0,1)$ so fast that
\begin{equation}
 \sum_{k\ge1}\varepsilon_k
 \bigl(1+\|\mathbf Z_k\|_{p\text{-var}}\bigr)<\infty.
 \label{hall:eq-universal-scales}
\end{equation}
Reparametrize the homogeneous dilation
$\delta_{\varepsilon_k}\mathbf Z_k$ to $I_k$, extend it constantly outside
$I_k$, and concatenate the pieces in
$H_*:=\bigoplus_{k\ge1}H_k$.

\begin{theorem}[One universal Hall test path]
\label{hall:thm-universal-path}
The finite concatenations converge in $p$-variation to one weakly geometric
$p$-rough path $\mathbf X_*$ on $H_*$.  It is independent of tensor degree,
Schur sector, and dual probe.  For every finite stable coefficient sector
$M_{\lambda,n}$ and every nonzero $q\in M_{\lambda,n}^*$, a bounded block
projection of $\mathbf X_*$ followed by one signed native corner gives
\begin{equation}
 Q_+\operatorname{Mat}_S\!\left(
 q(\pi_n\log_\otimes(P_{k\#}\mathbf X_*)_{0,1})\right)Q_-
 =c_qUD_\infty^{\otimes a_F(q)}V,
 \qquad c_q\ne0.
 \label{hall:eq-universal-path-corner}
\end{equation}
Consequently one fixed geometric path, together with labelled bounded
readers, tests every finite primitive layer and reconstructs every finite
ancestry flag.
\end{theorem}

\begin{proof}
Let $\mathbf X^{(m)}$ be the concatenation of the first $m$ pieces, extended
constantly on $[t_m,1]$.  Reparametrization does not change variation, and
homogeneous dilation multiplies the homogeneous $p$-variation control by
$\varepsilon_k$.  The standard concatenation estimate in the step-two rough
path group therefore gives, for $n>m$,
\[
 d_{p\text{-var}}(\mathbf X^{(m)},\mathbf X^{(n)})
 \le C_p\sum_{k=m+1}^n
 \varepsilon_k\bigl(1+\|\mathbf Z_k\|_{p\text{-var}}\bigr).
\]
Thus $\mathbf X^{(m)}$ is Cauchy.  The same tail bound controls its endpoint
increments, so the limit extends continuously to $t=1$; completeness of the
weakly geometric $p$-rough path space gives a weakly geometric limit
$\mathbf X_*$.  Condition \eqref{hall:eq-universal-scales} also implies the
square summability needed for the first-level endpoint in $H_*$ and for the
Hilbert--Schmidt direct sum of the second-level pure-area pieces.

The orthogonal projection $P_k:H_*\to H_k$ kills every other block and
recovers the reparametrized dilation of the $k$-th local experiment, with
constant pieces before and after $I_k$.  For the top nonzero restriction of
$q$, choose a Hall basis word paired nontrivially with it and use
\Cref{hall:lem-signed-corner}.  Homogeneous scaling contributes only the
nonzero scalar $\varepsilon_k^n$.  The valuation and reconstruction
statements follow from \Cref{rk:thm-kernel-support}.
\end{proof}

\begin{corollary}[Finite-window little tests]
\label{hall:cor-finite-window}
Fix finitely many tensor degrees and a finite spectral window.  Replacing
$D_\infty$ in the finitely many relevant local blocks by one sufficiently
long finite-rank truncation
$D_L=\operatorname{diag}(1,1/2,\ldots,1/L)$ produces a finite-rank geometric
test path whose labelled responses agree with
\eqref{hall:eq-native-exact-law} throughout the prescribed window.  Thus
every finite part of the universal test is realizable inside the little
weak-trace class.
\end{corollary}

\begin{proof}
Only finitely many local blocks, depths, and output ranks are involved.  For
each $r$, every product involving an index larger than $L$ is at most
$1/(L+1)$, while the $L$ products
$1\cdot j^{-1}\cdot1\cdots1$, $1\le j\le L$, are at least $1/L$.
Hence the first $L$ singular values of $D_L^{\otimes r}$ agree, with
multiplicity, with those of $D_\infty^{\otimes r}$.  Taking one $L$ larger
than every requested output rank and repeating the finite Hall construction
proves the assertion, including the endpoint case $r=1$.
\end{proof}

\subsection{Natural grade-one sources and source-specific sharpness}
\label{hall:subsec-natural-sources}

\begin{definition}[Replicable grade-one saturation source]
\label{cc:def-replicable-grade-one-source}
Fix $\beta\in(1/3,1/2)$.  A sequence of one-cell random critical geometries
$\mathbf Z^{(d)}=(Z^{(d)},A^{(d)})$ is a \emph{replicable grade-one
saturation source} if:
\begin{enumerate}[label=\textup{(\roman*)},leftmargin=2.7em]
\item every finite moment of its one-cell critical H\"older control is
      bounded uniformly in $d$;
\item there is $c_*>0$ such that, for every $L$ and $\varepsilon>0$, some
      $d$ satisfies
      \begin{equation}
       \mathbb P\left\{s_k(A^{(d)}_{0,1})\ge\frac{c_*}{k},
                 \ 1\le k\le L\right\}\ge1-\varepsilon;
       \label{cc:eq-grade-one-edge-certificate}
      \end{equation}
\item finitely many independent copies may be placed on mutually orthogonal
      spatial blocks without dimension-dependent constants.
\end{enumerate}
\end{definition}



\begin{theorem}[Saturation amplification through the native kernel]
\label{cc:thm-saturation-amplification}
Assume that a replicable grade-one saturation source exists.  Then every
nonzero stable natural exact-layer family
$G^{\mathrm{Schur}}_{\lambda,n,r}$ has minimax-sharp spectral depth $r$ in
the sense of \Cref{cc:gen-definition-minimax-depth}.  The Hall template,
cut, compressions, marker increments, and algebraic coefficient are fixed by
the exact layer before the spectral window.  For each requested window only
the finite dimensions of the $r$ source cells are enlarged.
\end{theorem}

\begin{proof}
Choose one Hall word whose exact-layer projection is nonzero and fix its
signed corner by \Cref{hall:lem-area-marker-spanning,hall:lem-signed-corner}.
Place independent grade-one source cells on its area slots, close their first
levels as in \eqref{hall:eq-purified-cell}, and use deterministic straight
markers.  Closing a cell, taking inverses, and concatenating a fixed Hall
template are polynomial operations in every fixed nilpotent step.  Chen's
identity expresses the global second level as the cell areas plus finitely
many wedges of first-level increments; hence the weak-trace quasi-triangle
inequality and
$\|x\wedge_{\rm op}y\|_{1,\infty}\le2\|x\|\|y\|$ give a
polynomial bound independent of the orthogonal block dimensions.  The
canonical finite-level extension theorem then controls the fixed higher
levels required by the Hall word.  After one deterministic spatial rescaling,
the assembled path lies in a fixed critical ball on an event of positive
probability.  On the intersection of this
control event and the edge events, the degree-$n$ Hall logarithm is exact and
its native corner is an $r$-fold tensor product of the grade-one areas.  In
particular, when $r=1$ there is no residual little-weak term whose modulus
could depend on the source dimension: the native corner is already an exact
grade-one block.  Apply the window part of
\Cref{cc:thm-critical-packet-principle}.  All
probabilities, algebraic data, and rescaling constants are fixed before the
window; only the block dimensions vary.
\end{proof}

\begin{theorem}[Native critical ancestry theorem]
\label{cc:thm-source-level-compiler}
Fix $\beta\in(1/3,1/2)$ and a nonzero stable exact primitive layer
$G^{\mathrm{Schur}}_{\lambda,n,r}$.
\begin{enumerate}[label=\textup{(\roman*)},leftmargin=2.7em]
\item Every controlled critical normal entrance has analytic cost at most
      $r$, with the dimension-free bound
      \eqref{cc:eq-rees-coefficient-bound}; little grade-one input gives
      little output.  If
      $T_{S,\lambda,r}=\operatorname{Mat}_S(P_{\lambda,r}L_n(s,t))$ and
      $b=C\omega(s,t)^{n\beta}$, then
      \begin{align*}
       \frac12\log\det(I+a^2T_{S,\lambda,r}^*T_{S,\lambda,r})
       &\le C_r(ab)(\log(e+ab))^{r-1},\\
       \operatorname{Tr}|T_{S,\lambda,r}|^{1+\varepsilon}
       &\le C_rb^{1+\varepsilon}\varepsilon^{-r}.
      \end{align*}
\item The designed native Hall response has a full-support coefficient-pure
      Rees kernel.  Hence every nonzero labelled probe has spectral depth
      equal to its ancestry depth, and every finite ancestry flag is
      recovered by \eqref{hall:eq-native-flag}.
\item A replicable natural grade-one source gives minimax sharpness for the
      same exact layer.
\item Any registered exact critical packet gives the full singular, Mellin,
      and Fredholm laws of \Cref{cc:thm-critical-packet-principle}.
\end{enumerate}
An application supplies only its canonical second-order geometry, its
algebraic normal symbol, and, for source-specific sharpness, one grade-one
certificate.  It does not reprove higher tensor propagation, dual detection,
or spectral conversion.
\end{theorem}

\begin{proof}
Part~\textup{(i)} is \Cref{cc:thm-critical-rees-compilation}.
Part~\textup{(ii)} is
\Cref{rk:thm-response-consequences,hall:thm-native-kernel}.
Part~\textup{(iii)} is \Cref{cc:thm-saturation-amplification}.
Part~\textup{(iv)} is \Cref{cc:thm-critical-packet-principle}.
\end{proof}

\subsection{Finite-order classification and spectral consequences}
\label{cc:subsec-compiled-consequences}

\begin{corollary}[Low degrees, long hooks, and arbitrary depth]
\label{cc:gen-corollary-sharp-sectors}
For every controlled critical entrance the following assertions hold.
\begin{enumerate}[label=\textup{(\roman*)},leftmargin=2.8em]
\item In degree three, the unique primitive sector
$\mathbf S_{(2,1)}H_{\mathbb C}$ has exact ancestry one.  Every nontrivial
flattening satisfies the universal $N^{-1}$ upper bound, and some labelled
native Hall coordinate $T$ satisfies $s_N(T)\sim C/N$ with $C>0$.
\item In degree four,
\[
 L_4(H_{\mathbb C})
 \cong\mathbf S_{(3,1)}H_{\mathbb C}
 \oplus\mathbf S_{(2,1,1)}H_{\mathbb C},
\]
and the two sectors have exact ancestries one and two.  Their universal upper
profiles are respectively
\[
 \frac1N,
 \qquad
 \frac{\log(e+N)}N.
\]
Each profile is attained, up to a positive asymptotic coefficient, by a
labelled native Hall coordinate.
\item For every $n\ge2$, the long-hook sector
$\mathbf S_{(n-1,1)}H_{\mathbb C}$ has exact ancestry one.  Its universal
critical upper profile is $N^{-1}$, and the native Hall family attains this
profile.
\item For every $r\ge1$, some partition $\lambda_r\vdash2r$ has exact
ancestry $r$; a native coordinate then satisfies
\[
 s_N(T)\sim C\frac{(\log N)^{r-1}}N,
 \qquad C>0.
\]
Thus finite-order critical depths are unbounded.
\end{enumerate}
Under a replicable natural grade-one source the same grades are minimax sharp
on arbitrarily long finite windows.  If the input area is little weak trace,
every upper membership improves to the corresponding little logarithmic
grade.
\end{corollary}

\begin{proof}
The algebraic statements are \Cref{cc:cor-low-degree-long-hook}.  Apply
parts~\textup{(i)} and~\textup{(ii)} of
\Cref{cc:thm-source-level-compiler} for the upper and native exact
claims, and part~\textup{(iii)} for the natural-source window statement.
\end{proof}

\section{Conditional covariance and natural-source interfaces}
\label{sec:conditional-interface}

The response-kernel theorem is deterministic once a critical entrance is
available.  This section records two source mechanisms that are both useful
and genuinely independent of the Hall construction: conditional covariance
under a revealed predictable bracket, and deterministic evolution transport
for stochastic convolutions.

\subsection{Revealed brackets and conditional fermionic propagation}

Let $M$ be a square-integrable c\`adl\`ag $H$-valued martingale, let
\[
 R_t=\langle\!\langle M\rangle\!\rangle_t\in\mathcal S_1(H)_+,
 \qquad V_t=\operatorname{Tr}R_t,
\]
and let $Z^\lambda$ be the exterior lift of
\Cref{cc:thm-localized-exterior-lift}.  A sigma-field
$\mathcal G_0\subseteq\mathcal F_0$ \emph{reveals} $R$ when the complete
c\`adl\`ag path $R$ is $\mathcal G_0$-measurable.  Write
$R_t=R_t^c+\sum_{0<s\le t}\Delta R_s$ in trace-class variation and define
\begin{equation}
 \mathscr U^x_{s,t}
 :=\exp\!\left(x\mathcal B_{R_t^c-R_s^c}\right)
   \prod_{s<a\le t}(I+x\mathcal B_{\Delta R_a}).
 \label{eq:compact-fermionic-propagator}
\end{equation}
The commuting product converges in operator norm on
$\mathcal S_1(\mathcal F_\wedge(H))$ and satisfies
$\mathscr U^x_{s,t}=\mathscr U^x_{u,t}\mathscr U^x_{s,u}$.

\begin{lemma}[Conditional trace-class Stieltjes interchange]
\label{lem:conditional-stieltjes-interchange}
Let $R$ be a $\mathcal G_0$-measurable c\`adl\`ag increasing
$\mathcal S_1(H)_+$-valued path with $\mathbb EV_T<\infty$, put
$V=\operatorname{Tr}R$, and write
$\mathrm dR_s=r_s\,\mathrm dV_s$.  Let $F$ be a jointly measurable
$\mathcal S_1(\mathcal F_\wedge(H))$-valued process satisfying
\[
 \mathbb E\int_{(0,T]}\|F_s\|_1\,\mathrm dV_s<\infty.
\]
Let $\widehat F$ be its conditional expectation with respect to
$\mathcal G_0\otimes\mathcal B([0,T])$ in the Bochner space determined by
the finite measure
$\mu(\mathrm d\omega,\mathrm ds)=\mathbb P(\mathrm d\omega)\,
\mathrm dV_s(\omega)$.  Then, for every $t\le T$,
\begin{equation}
 \mathbb E\!\left[
  \int_{(0,t]}\mathcal B_{\mathrm dR_s}(F_s)
  \middle|\mathcal G_0\right]
 =\int_{(0,t]}\mathcal B_{\mathrm dR_s}(\widehat F_s)
 \quad\text{in }\mathcal S_1(\mathcal F_\wedge(H)).
 \label{eq:conditional-stieltjes-interchange}
\end{equation}
If $F_s=X_s\otimes X_s$ and
$\mathbb E\int_{(0,T]}\|X_s\|^2\,\mathrm dV_s<\infty$, then
$\widehat F_s$ may be taken to be
$\mathbb E[X_s\otimes X_s\mid\mathcal G_0]$ for $\mu$-almost every
$(\omega,s)$.
\end{lemma}

\begin{proof}
The density $r$ may be chosen
$\mathcal G_0\otimes\mathcal B([0,T])$-measurable because the complete path
$R$ is $\mathcal G_0$-measurable.  The estimate
\[
 \|\mathcal B_{r_s}(S)\|_1\le\|S\|_1
\]
makes both sides of \eqref{eq:conditional-stieltjes-interchange} Bochner
integrable.  For a simple process
$F=\sum_iX_i\mathbf1_{(a_i,b_i]}$, the identity follows from
$\mathcal G_0$-measurability of the increments of $R$ and linearity of
$\mathcal B_Q$ in both variables.  Approximate a general $F$ in
$L^1(\mu;\mathcal S_1)$ by simple processes.  The trace class is a separable
noncommutative $L^1$-space, so its Bochner conditional expectation exists and
is a contraction in that space.  The displayed trace-norm estimate passes both
Stieltjes integrals to the limit.  The last assertion follows from the same
argument after scalarization against bounded
$\mathcal G_0$-measurable functions and trace-class dual functionals.
\end{proof}

\begin{theorem}[Conditional fermionic Dol\'eans law and ideal transfer]
\label{thm:compact-revealed-doleans}
Assume that $\mathcal G_0$ reveals $R$.  Put
$E_m=\{V_T\le m\}\in\mathcal G_0$,
$M^{(m)}=\mathbf1_{E_m}M$, and $R^{(m)}=\mathbf1_{E_m}R$.
Let $\mathscr U^{x,(m)}$ denote
\eqref{eq:compact-fermionic-propagator} with $R$ replaced by $R^{(m)}$.
Then, for every $t\le T$ and $\lambda\in\mathbb R$,
\begin{equation}
 \mathbb E\!\left[
 Z_t^{\lambda,(m)}\otimes Z_t^{\lambda,(m)}\mid\mathcal G_0\right]
 =\mathscr U_{0,t}^{\lambda^2,(m)}(P_0)
 \quad\text{in }\mathcal S_1(\mathcal F_\wedge(H)),
 \label{eq:compact-doleans}
\end{equation}
with trace at most $e^{\lambda^2m}$.  The local conditional covariance is
therefore
\begin{equation}
 C_t^{\lambda,\mathrm{loc}}
 =e^{\lambda^2\mathcal B_{R_t^c}}
  \prod_{0<s\le t}(I+\lambda^2\mathcal B_{\Delta R_s})(P_0).
 \label{eq:compact-local-covariance}
\end{equation}
Its particle-two block gives the exact atomic defect
\begin{equation}
 4C_{\alpha_A,t}^{\mathrm{loc}}
 =\Lambda^2R_t-\sum_{0<s\le t}\Lambda^2\Delta R_s.
 \label{eq:compact-atomic-defect}
\end{equation}
Let $A$ be the chronological area of $M$ and $Q=R_T$.  For every $z\ge0$,
locally on the events $E_m$,
\begin{equation}
 \mathbb E\!\left[
 \exp\!\left(\sup_{t\le T}\mathfrak F_{A_t}(z)\right)
 \middle|\mathcal G_0\right]
 \le4\det\!\left(I+\frac z2Q\right).
 \label{eq:compact-revealed-profile}
\end{equation}
Consequently, if $0<p<1$ and $Q\in\mathcal S_{p,\infty}$, then for every
$0<r<2$,
\[
 \left(\mathbb E\!\left[
  \sup_{t\le T}\|A_t\|_{\mathcal S_{p,\infty}}^r
  \middle|\mathcal G_0\right]\right)^{1/r}
 \le C_{p,r}\|Q\|_{\mathcal S_{p,\infty}}.
\]
Little weak ideals transfer to little weak ideals; in particular
$Q\in\mathcal S_1$ implies
\begin{equation}
 \sup_{t\le T}n s_n(A_t)\longrightarrow0
 \qquad\text{almost surely}.
 \label{eq:compact-trace-transfer}
\end{equation}
\end{theorem}

\begin{proof}
For a particle cutoff $N$, put
$T_s^{\lambda,N}=Z_s^{\lambda,N}\otimes Z_s^{\lambda,N}$ and
$C_s^{\lambda,N}=\mathbb E[T_s^{\lambda,N}\mid\mathcal G_0]$.
Conditional It\^o isometry followed by
\Cref{lem:conditional-stieltjes-interchange}, applied to
$F_s=T_{s-}^{\lambda,N}$, gives the left-point Stieltjes equation
\[
 C_t^{\lambda,N}
 =P_0+\lambda^2\int_{(0,t]}
 P_{\le N}\mathcal B_{\mathrm dR_s}(C_{s-}^{\lambda,N})P_{\le N}.
\]
Positivity, the particle-hole trace estimate, and Stieltjes Gronwall yield
$\sup_{N,t}\operatorname{Tr}C_t^{\lambda,N}\le e^{\lambda^2m}$.
The cutoff equation is triangular in the pair of left and right particle
degrees: $\mathcal B_Q$ raises both by one, and the vacuum block is fixed.
Consequently its blocks are determined recursively, so the trace-class
left-point solution is unique.  The product in
\eqref{eq:compact-local-covariance} solves the same recursion: the continuous
part is generated by $\lambda^2\mathcal B_{\mathrm dR^c}$ and an atom
$Q=\Delta R_s$ gives the single update
$I+\lambda^2\mathcal B_Q$.

The localized exterior lifts converge in $L^2$ of the supremum norm by
\Cref{cc:thm-localized-exterior-lift}; the rank-one trace estimate permits
passage to the limit in every block and then in trace norm.  Uniqueness now
proves \eqref{eq:compact-doleans}--\eqref{eq:compact-local-covariance}.  Reading the
particle-two block gives \eqref{eq:compact-atomic-defect}.

The conditional $L^2$ maximal inequality, the even-profile identity, and
fermionic determinant domination give \eqref{eq:compact-revealed-profile}.
The determinant/profile decoder of \Cref{cc:lem-response-profile-decoder}
then gives the weak and little ideal conclusions.  At $p=1$ one uses
$\log\det(I+xQ)=o(x)$ for $Q\in\mathcal S_1$.
\end{proof}


\begin{proposition}[Revealed c\`adl\`ag critical entrance]
\label{prop:revealed-cadlag-critical-entrance}
Let $M$ be a square-integrable c\`adl\`ag $H$-valued martingale whose complete
predictable bracket path $R$ is revealed by $\mathcal G_0$ and satisfies
$\mathbb E\operatorname{Tr}R_T<\infty$.  Then its split anchor
\[
 (M_t,A_t,[M]_t)
\]
has a c\`adl\`ag version in $\mathscr G_{\mathrm{crit}}(H)$.  If
$P_k\uparrow I$ are finite-rank orthogonal projections, then
\begin{equation}
 \sup_{t\le T}d_{\mathrm{crit}}\!\left(
 (P_kM_t,P_kA_tP_k,P_k[M]_tP_k),(M_t,A_t,[M]_t)\right)
 \longrightarrow0
 \quad\text{in probability}.
 \label{eq:revealed-cadlag-projection}
\end{equation}
Consequently the anchor has the optional completion of
\Cref{cc:thm-critical-optional-quotient}.
\end{proposition}

\begin{proof}
The area is c\`adl\`ag in $\mathcal S_2$, and
\Cref{thm:compact-revealed-doleans}, after localization by
$\{\operatorname{Tr}R_T\le m\}$, gives
\begin{equation}
 \sup_{t\le T}N s_N(A_t)\longrightarrow0
 \quad\text{almost surely}.
 \label{eq:uniform-little-tail-path}
\end{equation}
We use the following elementary upgrade.  If a path $K:[0,T]\to\mathcal S_2$
is c\`adl\`ag and satisfies
$\sup_tNs_N(K_t)\to0$, then it is c\`adl\`ag in
$\mathcal S^0_{1,\infty}$.  Indeed, for $s$ near $t$, Weyl's inequality gives,
for $N$ larger than a fixed cutoff,
\[
 N s_N(K_s-K_t)
 \leq 2N\bigl(s_{\lfloor N/2\rfloor}(K_s)
                 +s_{\lfloor N/2\rfloor}(K_t)\bigr),
\]
which is uniformly small by the tail hypothesis; the finitely many remaining
indices are bounded by a fixed multiple of $\|K_s-K_t\|_2$.  This proves the
claim and applies to $A$.

The range of a c\`adl\`ag path on a compact interval is relatively compact.
Strong convergence of $P_k$ is therefore uniform on the ranges of $M$ in $H$
and of $A$ in $\mathcal S_2$, so
$P_kM\to M$ and $P_kAP_k\to A$ uniformly in those norms.  Compression does not
increase singular values; combining this convergence with
\eqref{eq:uniform-little-tail-path} and the same finite-head/tail argument gives
uniform convergence in $\mathcal S_{1,\infty}$.  Finally, $[M]$ is a positive
c\`adl\`ag trace-class path.  Positivity and trace-norm compactness give
$P_k[M]P_k\to[M]$ uniformly on $[0,T]$ in trace norm.  Removing the
localization proves \eqref{eq:revealed-cadlag-projection} in probability.
\end{proof}

\begin{corollary}[Terminal bracket versus covariance delivery]
\label{cor:terminal-bracket-delivery}
Fix a deterministic terminal bracket $Q\in\mathcal S_1(H)_+$.  Among
revealed bracket paths with $R_T=Q$, the particle-two area covariance is
\begin{equation}
 C_{\alpha_A,T}^{\mathrm{loc}}
 =\frac14\left(\Lambda^2Q-
   \sum_{0<s\le T}\Lambda^2\Delta R_s\right).
 \label{eq:terminal-bracket-delivery}
\end{equation}
Consequently
\[
 0\le C_{\alpha_A,T}^{\mathrm{loc}}\le\frac14\Lambda^2Q.
\]
Continuous delivery attains the upper bound, one centered terminal jump with
covariance $Q$ attains zero, and every scalar level
$\theta\Lambda^2Q/4$, $0\le\theta\le1$, is attainable.  Thus the terminal
predictable bracket alone does not determine area covariance; its temporal
atomic structure is essential.
\end{corollary}

\begin{proof}
The formula is \eqref{eq:compact-atomic-defect}, and each exterior square
$\Lambda^2\Delta R_s$ is positive.  For a prescribed $\theta$, put
$\rho=1-\sqrt{1-\theta}$.  Combine an independent continuous martingale with
bracket $\rho(t/T)Q$ and a centered terminal jump with covariance
$(1-\rho)Q$.  Then the unique bracket atom contributes
$(1-\rho)^2\Lambda^2Q=(1-\theta)\Lambda^2Q$ in
\eqref{eq:terminal-bracket-delivery}.
\end{proof}

\begin{proposition}[Exterior degree is not primitive ancestry]
\label{prop:exterior-free-Lie-obstruction}
Let $V$ be a complex vector space with $\dim V\ge j$.  Then
\begin{equation}
 \operatorname{Hom}_{GL(V)}(\Lambda^jV,L_j(V))=0,
 \qquad j\ge3.
 \label{eq:no-exterior-free-Lie-map}
\end{equation}
Equivalently, in stable dimension there is no nonzero homogeneous linear
natural transformation from the degree-$j$ fermionic sector to the
degree-$j$ free-Lie sector.  The coincidence $L_2(V)=\Lambda^2V$ is therefore
exceptional: exterior particle number is a covariance-reader degree, not a
second model of the Rees ancestry filtration.
\end{proposition}

\begin{proof}
The Frobenius characteristic of the multilinear Lie representation is
\[
 \operatorname{ch}\operatorname{Lie}(j)
 =\frac1j\sum_{d\mid j}\mu_{\rm ar}(d)p_d^{j/d},
\]
where $\mu_{\rm ar}$ denotes the arithmetic M\"obius function;
see \cite{Reutenauer93}.  Its sign multiplicity is
\[
 m_j=\frac1j\sum_{d\mid j}\mu_{\rm ar}(d)(-1)^{j-j/d}.
\]
It equals one for $j=1,2$.  For odd $j>1$ the ordinary M\"obius divisor sum
vanishes; for $4\mid j$ the same is true; and when $j=2m$ with odd $m>1$,
pairing the divisors $e$ and $2e$ again gives zero.  Schur--Weyl duality
identifies $m_j$ with the multiplicity of $\Lambda^jV$ in $L_j(V)$, proving
\eqref{eq:no-exterior-free-Lie-map}.
\end{proof}

\begin{corollary}[One-sided Mellin transfer]
\label{cor:compact-revealed-mellin}
If, almost surely,
\[
 \mathcal N_Q(R)\sim Ce^{pR}R^{a-1}\ell(R),
 \qquad 0<p<1,
\]
with $\ell$ slowly varying, then
\[
 \operatorname{Tr}Q^{p+\varepsilon}
 \sim pC\Gamma(a)\varepsilon^{-a}\ell(1/\varepsilon),
\]
and
\[
 \limsup_{z\to\infty}
 \frac{\log\mathbb E[\exp(\sup_{t\le T}\mathfrak F_{A_t}(z))
       \mid\mathcal G_0]}
      {z^p(\log z)^{a-1}\ell(\log z)}
 \le \frac{\pi C}{2^p\sin(\pi p)}.
\]
No converse or sample-area equivalent is asserted.
\end{corollary}

\begin{proof}
Apply the counting--Mellin--Fredholm dictionary
\Cref{cc:thm:regular-dictionary} to $Q$ and insert the determinant equivalent
into \eqref{eq:compact-revealed-profile}.
\end{proof}

\subsection{L\'evy and evolution-family entrances}

The pathwise degree-two state uses optional quadratic variation; conditional
covariance uses the predictable bracket.  They coincide on some continuous
branches but remain differently typed coordinates.


\begin{proposition}[Typed Hoffman interface and collision kernel]
\label{prop:compact-Hoffman-interface}
Let $E$ be finite dimensional and fix a stochastic weight $N\ge2$.  Enlarge
the weight-one alphabet $\mathscr V_E=E$ by the collision letters
\[
 \mathscr C_{E,N}=\bigoplus_{k=2}^N E^{\odot k}.
\]
The finite-weight It\^o iterated integrals of an $E$-valued semimartingale
form a quasi-shuffle character on this typed alphabet.  Hoffman's exponential
and logarithm give inverse coalgebra maps and turn it into a shuffle
character without deleting collision letters.  The subsequent projection
$p:\mathscr V_E\oplus\mathscr C_{E,N}\to\mathscr V_E$ induces the exact
primitive sequence
\begin{equation}
 0\longrightarrow\mathfrak I_{E,N}
 \longrightarrow\mathfrak L(\mathscr V_E\oplus\mathscr C_{E,N})
 \longrightarrow\mathfrak L(\mathscr V_E)\longrightarrow0,
 \label{eq:compact-typed-exact-sequence}
\end{equation}
where $\mathfrak I_{E,N}$ is the Lie ideal generated by the collision
letters.  Geometricization and the quotient commute with convolution,
bounded spatial maps, deterministic evolution transport, truncation, and
fixed-terminal stopping.  At degree two the weight-one geometric tensor is
\[
 \mathbb Y^H_{s,t}=\mathbb Y^I_{s,t}+\frac12[Y]_{s,t}
 =\frac12Y_{s,t}^{\otimes2}+\operatorname{Alt}_2\mathbb Y^I_{s,t}.
\]
Thus the quotient retains the geometric ancestry while the typed parent
retains optional collision information.
\end{proposition}

\begin{proof}
Stochastic integration by parts identifies products of iterated It\^o
integrals with the quasi-shuffle product.  Hoffman's substitution identity
for the series $e^z-1$ and $\log(1+z)$ gives inverse coalgebra maps and
intertwines quasi-shuffle with shuffle.  The free-Lie universal property
identifies the kernel of the letter projection with the ideal generated by
$\mathscr C_{E,N}$, proving \eqref{eq:compact-typed-exact-sequence}.
Naturality follows because all maps preserve stochastic weight and commute
with bounded deterministic pushforwards.  The displayed degree-two identity
is the It\^o product formula.
\end{proof}

\begin{example}[Compensated Poisson type check]
\label{ex:poisson-type-check}
Let $M_t=N_t-\lambda t$ for a scalar Poisson process of rate $\lambda$.
Then
\[
 [M]_t=N_t,\qquad
 \langle\!\langle M\rangle\!\rangle_t=\lambda t,
 \qquad A_t=0,
\]
while
\[
 \mathbb M_t^I=\frac12(M_t^2-N_t),
 \qquad
 \mathbb M_t^I+\frac12[M]_t=\frac12M_t^2.
\]
Thus the optional collision coordinate jumps samplewise, whereas the
predictable covariance clock is deterministic.  Replacing either by the
other would already give the wrong geometricization or the wrong
conditional covariance in dimension one.
\end{example}

\begin{theorem}[Evolution-family stochastic entrance and Galerkin stability]
\label{thm:evolution-spde-interface}
Let $U(t,s)$ be a deterministic evolution family on a separable Hilbert space
$H$.  Let $N$ be a continuous square-integrable martingale on a Hilbert noise
space $U_0$, and let $K_s:U_0\to H$ be predictable.  For every terminal time
$t$ assume that $s\mapsto U(t,s)K_s$ is stochastically integrable and put
\[
 Z_r^{s\mid t}=\int_{(s,r]}U(t,a)K_a\,\mathrm dN_a,
 \qquad s\le r\le t.
\]
Assume locally in $t$ that the terminal-frame trace energies
\[
 \mathcal E_t^{(t)}
 =\operatorname{Tr}\langle\!\langle Z^{0\mid t}\rangle\!\rangle_t
\]
are finite.  Then the split degree-two increments
$\mathbf Z_{s,t}=(Z_t^{s\mid t},A_{s,t}^{(t)},Q_{s,t}^{(t)})$ are critical
entrances.  For a bounded spatial map $B$, write
$B_\#(x,A,Q)=(Bx,BAB^*,BQB^*)$.  With this convention they satisfy the
transported Chen law
\begin{equation}
 \mathbf Z_{s,t}
 =U(t,u)_\#\mathbf Z_{s,u}\circ\mathbf Z_{u,t},
 \qquad s\le u\le t.
 \label{eq:spde-transported-Chen}
\end{equation}
Here the early increment is first constructed in the terminal $u$ frame and
then pushed to the terminal $t$ frame.

Let $K^{(m)}$ be finite-rank Galerkin approximants.  For a fixed terminal
$t$, if
\begin{equation}
 \mathbb E\operatorname{Tr}
 \langle\!\langle
  Z^{K^{(m)}-K,\,0\mid t}
 \rangle\!\rangle_t\longrightarrow0,
 \label{eq:spde-operational-convergence}
\end{equation}
then the corresponding first levels, areas, and brackets converge ucp on
$[0,t]$ in the critical split topology.  The estimates are uniform over a
family of deterministic terminal frames whenever their energy and error
bounds are uniform.  Joint ucp convergence in the terminal parameter requires
an additional terminal-frame continuity or maximal estimate.  Adding an adapted finite-variation drift whose range on each compact
interval is contained in a deterministic finite-dimensional subspace of $H$
changes the area only by trace-class terms and preserves the entrance.
\end{theorem}

\begin{proof}
For finite-rank approximants, stochastic integration by parts gives the
finite-weight quasi-shuffle character.  Hoffman geometricization commutes
with bounded maps and deterministic evolution transport; splitting the
terminal-frame integral at $u$ gives
\eqref{eq:spde-transported-Chen}.  The endpoint theorem
\Cref{cc:thm-universal-critical-endpoint} places the antisymmetric channel in
the weak-trace critical space and the collision channel in trace class.

The quantity in \eqref{eq:spde-operational-convergence} is exactly the pullback
of terminal $L^2$ convergence by the It\^o isometry.  Apply the critical
projection stability theorem \Cref{cc:thm-critical-projection-stability} to
the difference martingale; its area estimate is
$O(\sqrt E\,\delta+\delta^2)$ under energy $E$ and error energy $\delta^2$.
Localization gives ucp convergence.  The finite-variation assertion follows
by expanding the second level: all added terms have one leg in a locally
finite-dimensional range, hence are trace class.
\end{proof}

\begin{remark}[Deterministic terminal frames]
\label{rem:deterministic-terminal-frame}
The terminal parameter $t$ in \Cref{thm:evolution-spde-interface} is fixed
before stochastic integration.  For a stopping time $\tau$, the map
$a\mapsto U(\tau(\omega),a)K_a$ need not be predictable.  Stopping is
therefore applied to the already constructed deterministic-terminal
character.  Random-terminal transport requires a separate innovation or
restart theorem and is not a consequence of the transported Chen identity.
\end{remark}

\begin{corollary}[Mild SPDE interface]
\label{cor:mild-spde-interface}
Consider a mild equation
\[
 X_t=U(t,0)x_0+\int_0^tU(t,s)F_s\,\mathrm ds
      +\int_0^tU(t,s)G_s\,\mathrm dW_s.
\]
If the stochastic convolution satisfies the terminal-frame energy and
Galerkin hypotheses of \Cref{thm:evolution-spde-interface} for every fixed
terminal frame, and the drift channel satisfies the stated finite-variation
condition, then the enhanced noise part has a canonical critical degree-two
entrance compatible with stopping and deterministic evolution transport.
A joint terminal-time enhancement additionally requires the continuity or
maximal hypothesis stated in that theorem.  Every finite-order Rees
response theorem applies once a model-specific native response kernel is
shown to have full support.  By
\Cref{rk:prop-kernel-naturality,rk:prop-robust-support}, bounded terminal
observations may be reduced to finitely many scalar response coordinates.
For a fixed finite exact layer $G$, use the canonical Galerkin layer
identifications and choose a finite certificate as in
\Cref{rk:prop-robust-support}.  If the corresponding maps
$\mathsf C_G^{(m)}$ converge in operator norm to $\mathsf C_G$ and
$\kappa(\mathsf C_G)>0$, then all sufficiently large Galerkin responses are
strict.  Conversely, if $\inf_m\kappa(\mathsf C_G^{(m)})>0$, then the limiting
response is strict.  The equation's semilinear fixed-point estimates and any
state or generator observability remain separate analytic inputs.
\end{corollary}

\begin{proof}
Apply \Cref{thm:evolution-spde-interface} to the stochastic convolution and
its Galerkin approximants; the finite-variation clause treats the drift.
Kernel base change identifies bounded terminal observations with pushforwards
of the same response kernel.  The finite-reader reduction and the two
perturbative implications are exactly \Cref{rk:prop-robust-support}.
\end{proof}

\begin{corollary}[Hilbert-valued L\'evy entrance]
\label{cor:compact-levy-entrance}
Every c\`adl\`ag Hilbert-valued L\'evy process with finite second moment has a
canonical degree-two critical split lift
\[
 (L_{s,t},A^L_{s,t},[L]_{s,t})
\]
with stopped Chen identities and finite-rank projection stability.  The lift
may be chosen as a shift-covariant Borel functional on the Skorokhod path
space.  Hence, after every bounded stopping time, the lifted future is
independent of the past and has the law of the original lifted process.
At stochastic weights above two the jump alphabet contains additional
power-jump primitives, which belong to the quasi-shuffle normal geometry.
\end{corollary}

\begin{proof}
Write $L_t=tm+M_t$, where $M$ is the centered square-integrable L\'evy
martingale.  Its predictable operator bracket is deterministic:
\[
 \langle\!\langle M\rangle\!\rangle_t=tC,
 \qquad C\in\mathcal S_1(H)_+.
\]
Thus the trivial initial sigma-field reveals the complete bracket path.
Apply \Cref{thm:compact-revealed-doleans,prop:revealed-cadlag-critical-entrance}
to obtain the weak-trace area, the c\`adl\`ag critical anchor, and finite-rank
projection convergence for the martingale part.  The optional coordinate in
the split state is $[M]$, while the deterministic bracket $tC$ drives the
conditional covariance estimate.  The drift $tm$ has one-dimensional range;
expanding the second level of $tm+M_t$ shows that it adds only finite-rank,
hence trace-class, area terms.  The group law then supplies all stopped Chen
identities.

For the restart statement, fix increasing finite-rank projections
$P_k\uparrow I_H$.  On each finite-dimensional Skorokhod space use the same
Bichteler--Karandikar left-point integral map
\cite{Karandikar95,Chevyrev18,ChevyrevFriz19}.  The set on which the projected
split lifts form a Cauchy sequence in the Polish critical path space is Borel,
since its Cauchy condition is countable; define the infinite-dimensional lift
there by the projective limit and assign a fixed value off that set.  The
finite-dimensional maps commute with deterministic time shifts and with the
projections $P_k$.  Hence the convergence set is shift invariant and the
resulting Borel lift map on $D([0,\infty);H)$ is shift covariant.  Applying
this fixed map to the strong Markov shifted L\'evy path transfers stationary
independent increments to the enhanced future.
\end{proof}

\section{Brownian type separation and further spectral consequences}
\label{sec:brownian-boundaries}

\subsection{Three Brownian operators}

Let $C\in\mathcal S_1(H)_+$ have eigenpairs $(c_i,e_i)$ and let $B^C$ be
Brownian motion with covariance $C$.  Three natural operators must remain
separate.
\begin{enumerate}[label=\textup{(\roman*)},leftmargin=2.8em]
\item The raw Chen-local area $A_t^C(\omega)$ is a random skew operator on
      $H$ and belongs samplewise to $\mathcal S^0_{1,\infty}$.
\item The deterministic second-chaos coefficient root
\[
 K_{C,t}f_{ij}=\frac t2\sqrt{c_ic_j}\,e_i\wedge e_j,
 \qquad i<j,
\]
acts from the active coefficient space to $\Lambda^2H$.  Its square is the
annealed covariance of the random bivector.
\item Let $\mathsf K_T$ be the Brownian clock covariance on $L^2(0,T)$,
\[
 (\mathsf K_Tf)(s)=\int_0^T(s\wedge t)f(t)\,\mathrm dt,
\]
and let $P_{k,T}$ be its simple temporal spectral projections.  Define
\[
 \mathfrak G_T(b)=\sum_{k\ge1}|P_{k,T}b\rangle\langle P_{k,T}b|,
 \qquad
 \mathfrak D_T(b)=\mathfrak G_T(b)^{1/2}.
\]
For one Brownian sample this is a positive global auxiliary response on
$L^2(0,T)\otimes H$.
\end{enumerate}


\begin{proposition}[Second-chaos factorization and the point-evaluation gap]
\label{prop:compact-brownian-factorization}
For $i<j$, let
\[
 \Xi_{ij}(t)=t^{-1}\left(
  \int_0^tW_s^i\,\mathrm dW_s^j-
  \int_0^tW_s^j\,\mathrm dW_s^i\right).
\]
The active family $(\Xi_{ij}(t))_{c_ic_j>0}$ is orthonormal in
$L^2(\Omega)$, and second-chaos radonification defines an isometry
\[
 \operatorname{Rad}^{(2)}_{\Xi,t}:
 \mathcal S_2(\mathfrak h_{2,C},E)\longrightarrow L^2(\Omega;E).
\]
Moreover
\begin{equation}
 \alpha_{A_t^C}
 =\operatorname{Rad}^{(2)}_{\Xi,t}(K_{C,t}),
 \qquad
 \operatorname{Cov}(\alpha_{A_t^C})
 =K_{C,t}K_{C,t}^*=\frac{t^2}{4}\Lambda^2C.
 \label{eq:compact-second-chaos-factorization}
\end{equation}
If $C$ has infinite rank, then almost surely algebraic point evaluation on
the active chaos space has no bounded extension to its $L^2$ closure.
Thus \eqref{eq:compact-second-chaos-factorization} is an $L^2$
radonification identity, not a bounded samplewise contraction.
\end{proposition}

\begin{proof}
It\^o isometry and independence give
$\mathbb E[\Xi_{ij}(t)\Xi_{k\ell}(t)]
=\mathbf1_{\{i=k,j=\ell\}}$.  Expanding the area in the eigenbasis of $C$
gives
\[
 \alpha_{A_t^C}
 =\frac t2\sum_{i<j}\sqrt{c_ic_j}\,\Xi_{ij}(t)e_i\wedge e_j,
\]
which is precisely \eqref{eq:compact-second-chaos-factorization}; the
coefficient square sum equals
$t^2((\operatorname{Tr}C)^2-\operatorname{Tr}C^2)/8$.
For unboundedness, choose disjoint active pairs.  The corresponding
$\Xi$'s are independent, identically distributed, and have unit second
moment.  The strong law implies that their sample values are not in
$\ell^2$.  Riesz representation would force $\ell^2$ coefficients for any
bounded point-evaluation functional, a contradiction.
\end{proof}

\begin{proposition}[Annealed divisor spectrum and quenched zero density]
\label{prop:brownian-zero-density}
Fix $t>0$ and assume
\begin{equation}
 c_j\sim b j^{-\delta},\qquad b>0,\quad\delta>1.
 \label{eq:brownian-polynomial-covariance}
\end{equation}
Then
\begin{align}
 N_{K_{C,t}}(u)
 &\sim\frac{(tb/2)^{2/\delta}}{\delta}
 u^{-2/\delta}\log(1/u),
 \label{eq:annealed-counting}\\
 s_N(K_{C,t})
 &\sim\frac{tb}{2^{1+\delta/2}}
 N^{-\delta/2}(\log N)^{\delta/2}.
 \label{eq:annealed-singular}
\end{align}
On an event of probability one,
\begin{equation}
 A_t^C\in\mathcal S_{1/\delta,\infty}(H)
 \quad\text{and}\quad
 \frac{N_{A_t^C}(au)}{N_{K_{C,t}}(u)}\longrightarrow0
 \quad(u\downarrow0)
 \label{eq:quenched-zero-density}
\end{equation}
for every fixed $a>0$.  Thus the raw area has zero asymptotic spectral
density relative to its annealed coefficient root after every fixed threshold
rescaling.
\end{proposition}

\begin{proof}
The singular values of $K_{C,t}$ are
$(t/2)\sqrt{c_ic_j}$, $i<j$.  The distinct-pair divisor law
\[
 \#\{i<j:ij\le X\}\sim\frac12X\log X
\]
and monotone inversion give
\eqref{eq:annealed-counting}--\eqref{eq:annealed-singular}.  Since
$tC\in\mathcal S_{1/\delta,\infty}$ and $1/\delta<1$, the revealed-bracket
ideal transfer in \Cref{thm:compact-revealed-doleans} gives
$A_t^C\in\mathcal S_{1/\delta,\infty}$ almost surely.  Hence
$N_{A_t^C}(au)=O_{\omega,a}(u^{-1/\delta})$, whereas
\eqref{eq:annealed-counting} is of order
$u^{-2/\delta}\log(1/u)$, proving \eqref{eq:quenched-zero-density}.
\end{proof}

\begin{lemma}[Random harmonic rearrangement]
\label{lem:random-harmonic-rearrangement}
Let $(Y_k)_{k\ge1}$ be independent copies of a nonnegative random variable
$Y$ with $0<\mathbb EY<\infty$ and $\mathbb EY^2<\infty$.  Fix $c>0$ and
$0\le\theta<1$, and let $D_\omega$ be the positive diagonal operator whose
unordered diagonal entries are
\[
 d_k(\omega)=\frac{cY_k(\omega)}{k-\theta}.
\]
Then $D_\omega$ is compact almost surely and
\begin{equation}
 uN_{D_\omega}(u)\longrightarrow c\mathbb EY,
 \qquad
 s_N(D_\omega)\sim\frac{c\mathbb EY}{N}
 \quad\text{almost surely}.
 \label{eq:random-harmonic-rearrangement}
\end{equation}
If $Y$ has finite moments of every positive order, use the convention
$0^z=0$ for $\operatorname{Re}z>0$.  Then almost surely
$\zeta_{D_\omega}$ admits a meromorphic continuation to
$\operatorname{Re}z>1/2$ with a unique pole there at $z=1$, and
\begin{equation}
 \operatorname*{Res}_{z=1}\zeta_{D_\omega}(z)=c\mathbb EY.
 \label{eq:random-harmonic-residue}
\end{equation}
\end{lemma}

\begin{proof}
First, $d_k\to0$ almost surely because
$\sum_k\mathbb P\{Y_k>\varepsilon k\}<\infty$ for every $\varepsilon>0$;
this follows by comparing the sum with
$\varepsilon^{-1}\int_0^\infty\mathbb P\{Y>x\}\,\mathrm dx$.
For $u>0$ put $N(u)=\sum_{k\ge1}
\mathbf1_{\{cY_k>(k-\theta)u\}}$.  The tail function of $Y$ is decreasing and
integrable, so the Riemann-sum theorem gives
\[
 \mathbb E[uN(u)]
 =u\sum_{k\ge1}\mathbb P\{cY>(k-\theta)u\}
 \longrightarrow c\int_0^\infty\mathbb P\{Y>x\}\,\mathrm dx
 =c\mathbb EY.
\]
Moreover
$\operatorname{Var}(uN(u))\le u^2\sum_k\mathbb P\{cY>(k-\theta)u\}
=O(u)$.  Along every geometric sequence $u_n=\rho^n$, $0<\rho<1$,
Chebyshev and Borel--Cantelli therefore give
$u_nN(u_n)\to c\mathbb EY$ almost surely.  Monotonicity of $N$ brackets an
arbitrary $u\in[u_{n+1},u_n]$ between the two neighboring geometric values.
Taking a countable sequence $\rho\uparrow1$ proves the first limit in
\eqref{eq:random-harmonic-rearrangement}; monotone inversion proves the
second.

For the complex assertion, initially on $\operatorname{Re}z>1$,
\begin{equation}
 \zeta_{D_\omega}(z)
 =c^z\mathbb E(Y^z)\zeta_{\rm H}(z,1-\theta)+H_\omega(z),
 \qquad
 H_\omega(z)=c^z\sum_{k\ge1}
 \frac{Y_k^z-\mathbb E(Y^z)}{(k-\theta)^z},
 \label{eq:random-harmonic-dirichlet}
\end{equation}
where $\zeta_{\rm H}$ is the Hurwitz zeta function.  Let $K$ be compact in
$\{\operatorname{Re}z>1/2\}$ and choose a bounded planar domain $U$ with
$K\Subset U\Subset\{\operatorname{Re}z>1/2\}$.  The centered summands in
$H_\omega$ are independent mean-zero random variables in the Bergman Hilbert
space $A^2(U)$, and
\[
 \sum_{k\ge1}\mathbb E\left\|
 c^{(\cdot)}\frac{Y_k^{(\cdot)}-\mathbb E(Y^{(\cdot)})}
 {(k-\theta)^{(\cdot)}}\right\|_{A^2(U)}^2<\infty.
\]
Indeed the real parts on $U$ stay in a compact subset of $(1/2,\infty)$ and
all corresponding moments of $Y$ are finite.  Hilbert-space Kolmogorov
convergence gives almost-sure convergence in $A^2(U)$; the interior Bergman
estimate upgrades this to uniform convergence on $K$.  A countable exhaustion
therefore makes $H_\omega$ holomorphic on the whole half-plane.  Since
$z\mapsto\mathbb E(Y^z)$ is holomorphic there and the Hurwitz zeta function
has residue one at $z=1$, \eqref{eq:random-harmonic-residue} follows.
\end{proof}

\begin{theorem}[Auxiliary fixed-sample critical boundary]
\label{thm:compact-brownian-auxiliary-boundary}
Assume $T>0$ and $C\ne0$.  Let $G_C\sim N(0,C)$ and
$a_{C,T}=T\mathbb E\|G_C\|/\pi>0$.  On one event of probability one,
\begin{equation}
 s_N(\mathfrak D_T(B^C))\sim\frac{a_{C,T}}{N}.
 \label{eq:compact-aux-boundary}
\end{equation}
Moreover, $\zeta_{\mathfrak D_T(B^C)}$ extends meromorphically to
$\operatorname{Re}z>1/2$, has no pole there except $z=1$, and
\begin{equation}
 \operatorname*{Res}_{z=1}\zeta_{\mathfrak D_T(B^C)}(z)=a_{C,T}.
 \label{eq:compact-aux-residue}
\end{equation}
For every $r\ge1$ the tensor power satisfies
\begin{equation}
 s_N\!\left(\mathfrak D_T(B^C)^{\otimes r}\right)
 \sim\frac{a_{C,T}^r}{\Gamma(r)}
       \frac{(\log N)^{r-1}}{N},
 \qquad
 \zeta_{\mathfrak D_T(B^C)^{\otimes r}}(z)
 =\zeta_{\mathfrak D_T(B^C)}(z)^r.
 \label{eq:compact-aux-tensor-law}
\end{equation}
\end{theorem}

\begin{proof}
The clock covariance $\mathsf K_T$ has simple eigenvalues
$T^2/[\pi^2(k-\frac12)^2]$.  Hence the mutually orthogonal
Karhunen--Lo\`eve blocks of $B^C$ satisfy
\[
 \|P_{k,T}B^C\|
 \stackrel{d}=\frac{T}{\pi(k-\frac12)}\|G_k\|,
\]
where $G_k$ are independent copies of $G_C$.  Because the summands defining
$\mathfrak G_T(B^C)$ have orthogonal ranges, these block norms are precisely
the unordered singular values of $\mathfrak D_T(B^C)$.  Apply
\Cref{lem:random-harmonic-rearrangement} with
$c=T/\pi$, $\theta=1/2$, and $Y_k=\|G_k\|$.  Gaussian moments of every
positive order are finite, so the singular-value equivalent
\eqref{eq:compact-aux-boundary} and the continuation and residue
\eqref{eq:compact-aux-residue} follow on one event.  Tensor
multiplicativity of the zeta function and the exact clause of
\Cref{cc:thm-critical-packet-principle} give
\eqref{eq:compact-aux-tensor-law}.
\end{proof}

\begin{proposition}[Brownian registration trilemma]
\label{prop:compact-brownian-trilemma}
Fix $t>0$ and assume the polynomial covariance hypothesis
\eqref{eq:brownian-polynomial-covariance}.  No realization can simultaneously
\begin{enumerate}[label=\textup{(\roman*)},leftmargin=2.8em]
\item equal the raw Chen-local area $A_t^C(\omega)$;
\item arise from a bounded functional agreeing with point evaluation on all
      active second-chaos coordinates;
\item be positively spectrally faithful to $K_{C,t}$, in the sense that some
      $a,c>0$ satisfy
      \[
       N_{A_t^C}(au)\ge cN_{K_{C,t}}(u)
      \]
      for all sufficiently small $u$ on an event of positive probability.
\end{enumerate}
Thus the valid assignments are typed:
\[
 \mathcal B_{\rm ann}=K_{C,t},\qquad
 \mathcal B_{\rm que}=A_t^C(\omega),\qquad
 \mathcal B_{\rm aux}=\mathfrak D_T(B^C(\omega)).
\]
Relations between them are covariance, radonification, or nonlinear response
transforms, not operator identifications.
\end{proposition}

\begin{proof}
The zero-density statement \eqref{eq:quenched-zero-density} rules out
item~\textup{(iii)} for the raw area after every fixed threshold rescaling.
Independently, \Cref{prop:compact-brownian-factorization} shows that the
active Gaussian coordinate row does not define a bounded point-evaluation
functional on its coefficient Hilbert space almost surely.  The auxiliary
root in \Cref{thm:compact-brownian-auxiliary-boundary} drops Chen locality and
therefore does not contradict either obstruction.
\end{proof}

\begin{corollary}[Natural Brownian window certificate]
\label{cor:compact-brownian-window}
There is a numerical $c_*>0$ such that, for every $L\ge1$ and
$\varepsilon>0$, some sufficiently large even dimension $d$ satisfies
\begin{equation}
 \mathbb P\left\{
  s_k(A^{M^{(d)}}_{0,1})\ge\frac{c_*}{k}
  \text{ for }1\le k\le L\right\}\ge1-\varepsilon,
 \qquad M^{(d)}=d^{-1/2}B^{(d)}.
 \label{eq:compact-brownian-window-certificate}
\end{equation}
Through \Cref{cc:thm-saturation-amplification}, this supplies the stochastic
finite-window version of every exact primitive depth.  It remains a
family-valued minimax statement, not a full-tail theorem for one raw sample.
\end{corollary}

\begin{proof}
Area is quadratic under spatial scaling, so
$A^{M^{(d)}}_{0,1}=d^{-1}\mathcal L^{(d)}$.  Apply
\eqref{cc:eq-critical-Brownian-edge} from
\Cref{cc:lem-critical-brownian-block} with $m=L$ and take $c_*=c_0$.
The higher-depth conclusion is exactly
\Cref{cc:thm-saturation-amplification}.
\end{proof}

\subsection{Coefficient degree, Rees depth, and infinite towers}

The exact spectrum of a deterministic coefficient operator and the Rees
spectrum of a critical response answer different questions.  The following
compact statement retains the useful part of the coefficient theory.

\begin{proposition}[Pure-power tensor law and faithful transfer]
\label{prop:compact-coefficient-law}
If $T$ is positive compact and
$s_N(T)\sim cN^{-\alpha}$, then for every $k\ge1$,
\begin{equation}
 N_{T^{\otimes k}}(u)
 \sim\frac{c^{k/\alpha}}{(k-1)!\alpha^{k-1}}
 u^{-1/\alpha}(\log(1/u))^{k-1}.
 \label{eq:compact-tensor-counting}
\end{equation}
If a selected coefficient sector has a separately established positive
counting density inside $T^{\otimes k}$, and a registered operator $R_k(T)$ is
asymptotically faithful to that sector after a fixed threshold rescaling, then
the same counting, Mellin, and Fredholm laws transfer with the induced
constant.  Neither the sector density nor analytic faithfulness follows from
algebraic homogeneity alone.
\end{proposition}

\begin{proof}
The singular values of $T^{\otimes k}$ are the products
$s_{i_1}(T)\cdots s_{i_k}(T)$.  Taking logarithms reduces the counting problem
to the $k$-fold additive convolution of a measure with exponential tail;
Karamata--Delange inversion gives \eqref{eq:compact-tensor-counting}.  Once
a selected sector density is known, the transfer assertion is monotone
inversion followed by the counting--Mellin--Fredholm dictionary.
\end{proof}

\begin{remark}[Two logarithmic indices]
For coefficient spectra, every copy of $T$ is active and the logarithmic power
is governed by total operator degree $k$.  In a Rees response, transport
letters may be rank one while only $r$ critical area leaves are active; the
power is then $r-1$.  Equality of these indices is a registration theorem,
not notation.
\end{remark}

\begin{lemma}[Compact synthesis invariance]
\label{lem:compact-synthesis-invariance}
Let $W\ge0$ be compact with nonzero eigenvalues
$w_1\ge w_2\ge\cdots>0$ satisfying
$w_n=n^{-\alpha}L(n)$ for some $\alpha>0$ and slowly varying $L$.
Let $\mathscr S$ be bounded and bounded below, and suppose
$\mathscr S^*\mathscr S=I+K$ with $K$ compact.  Then the nonzero eigenvalues
of $\mathscr S W\mathscr S^*$ satisfy
\begin{equation}
 s_n(\mathscr S W\mathscr S^*)\sim w_n.
 \label{eq:compact-synthesis-invariance}
\end{equation}
The same conclusion holds whenever the counting function of $W$ is regularly
varying at zero with positive index.
\end{lemma}

\begin{proof}
The nonzero eigenvalues of $\mathscr S W\mathscr S^*$ are those of
$A=W^{1/2}(I+K)W^{1/2}$.  Fix $0<\varepsilon<1$ and choose a self-adjoint
finite-rank $F$ with $\|K-F\|<\varepsilon$.  If $m=\operatorname{rank}F$,
then
\[
 (1-\varepsilon)W+W^{1/2}FW^{1/2}
 \le A\le
 (1+\varepsilon)W+W^{1/2}FW^{1/2}.
\]
Finite-rank min--max interlacing yields, for $n>m$,
\[
 (1-\varepsilon)w_{n+m}
 \le s_n(A)\le
 (1+\varepsilon)w_{n-m}.
\]
Regular variation gives $w_{n\pm m}/w_n\to1$.  Let first $n\to\infty$ and
then $\varepsilon\downarrow0$.  The counting-function formulation is
obtained by the same argument after monotone inversion.
\end{proof}

\begin{theorem}[Primitive--PBW boundary beyond finite depth]
\label{thm:compact-infinite-tower}
Let finite-dimensional ancestry shells $G_m$ have cumulative rank
$M(x)=\sum_{m\le x}\dim G_m$, let $w_m\downarrow0$, and let
$T=\mathscr S(\bigoplus_m w_mI_{G_m})\mathscr S^*$ with
$\mathscr S^*\mathscr S=I+$compact.  If
\[
 M(x)\sim Ax^\delta(\log x)^\gamma,
 \qquad
 w_m\sim cm^{-\alpha}(\log m)^\eta,
\]
then, with $p=\delta/\alpha$ and $\theta=\gamma+p\eta$,
\[
 \mathcal N_T(R)
 \sim Ac^p\alpha^{-\theta}e^{pR}R^\theta.
\]
The corresponding real Mellin and Fredholm laws follow from
\Cref{cc:thm:regular-dictionary}.  For an exponentially attenuated free tower
on $q\ge2$ generators, the primitive zeta function has a logarithmic branch at
$p_0=(\log q)/\kappa_0$, whereas the PBW tower has a simple pole and its
$k$-fold tensor power has a pole of order $k$:
\[
 \zeta_{\rm prim}(p_0+\varepsilon)=\log(1/\varepsilon)+O(1),
 \qquad
 \zeta_{\rm PBW}(z)=\frac1{1-qe^{-\kappa_0z}}.
\]
The lattice spectrum generally produces log-periodic counting oscillations.
\end{theorem}

\begin{proof}
For the diagonal shell operator
$W=\bigoplus_m w_mI_{G_m}$, monotone inversion identifies the largest shell
visible above $e^{-R}$ and substitution into $M$ yields the displayed
counting equivalent.  Its ordered eigenvalue sequence is regularly varying;
\Cref{lem:compact-synthesis-invariance} transfers the same equivalent, with
the same leading constant, to
$T=\mathscr S W\mathscr S^*$.  For the free tower, the Witt formula gives
\[
 \zeta_{\rm prim}(z)
 =\sum_{d\ge1}\frac{\mu(d)}d
  \log\frac1{1-qe^{-\kappa_0zd}},
\]
whose $d=1$ term has the stated branch, while PBW completion is the geometric
series.  The remaining poles are the lattice translates
$p_0+2\pi i\ell/\kappa_0$.
\end{proof}

\section{Conclusion}
\label{sec:scope-conclusion}

The stochastic entrance, primitive Rees filtration, and native response kernel
form a single finite-order mechanism.  The entrance identifies the sharp
second-order operator geometry, the Rees filtration records primitive ancestry,
and full kernel support identifies ancestry depth with critical spectral depth
while reconstructing the finite labelled ancestry flag.  Coefficient-purity on
the harmonic boundary gives the exact logarithmic asymptotics.

The Hall realization supplies full support geometrically, while Brownian,
revealed-bracket, L\'evy, and evolution-family constructions provide natural
source interfaces.  For jump processes the present spectral ancestry theory is
used at degree two; higher stochastic weights carry additional quasi-shuffle
primitive letters.  Thus the finite-order continuous/geometric theory and the
second-order jump entrance fit the same causal Rees framework without changing
the response invariant.

\medskip
\noindent\textbf{Acknowledgment.}
The author used ChatGPT (GPT-5.6, OpenAI) interactively in developing proofs
and refining the presentation, and takes full responsibility for the
mathematical content.

\begin{thebibliography}{99}

\bibitem{Applebaum09}
D.~Applebaum,
\emph{L\'evy Processes and Stochastic Calculus}, second ed.,
Cambridge Studies in Advanced Mathematics 116, Cambridge University Press,
Cambridge, 2009.
\href{https://doi.org/10.1017/CBO9780511809781}
{doi:10.1017/CBO9780511809781}.

\bibitem{ApplebaumHudson84}
D.~B. Applebaum and R.~L. Hudson,
Fermion It\^o's formula and stochastic evolutions,
\emph{Comm. Math. Phys.} \textbf{96} (1984), 473--496.
\href{https://doi.org/10.1007/BF01212531}
{doi:10.1007/BF01212531}.

\bibitem{BarnettStreaterWilde82}
C.~Barnett, R.~F. Streater, and I.~F. Wilde,
The It\^o--Clifford integral,
\emph{J. Funct. Anal.} \textbf{48} (1982), no.~2, 172--212.
\href{https://doi.org/10.1016/0022-1236(82)90066-0}
{doi:10.1016/0022-1236(82)90066-0}.

\bibitem{BGT87}
N.~H. Bingham, C.~M. Goldie, and J.~L. Teugels,
\emph{Regular Variation}, Encyclopedia of Mathematics and its Applications
27, Cambridge University Press, Cambridge, 1987.
\href{https://doi.org/10.1017/CBO9780511721434}
{doi:10.1017/CBO9780511721434}.

\bibitem{Brzezniak97}
Z.~Brze\'zniak,
On stochastic convolution in Banach spaces and applications,
\emph{Stochastics Stochastics Rep.} \textbf{61} (1997), nos.~3--4,
245--295.
\href{https://doi.org/10.1080/17442509708834122}
{doi:10.1080/17442509708834122}.

\bibitem{Chevyrev18}
I.~Chevyrev,
Random walks and L\'evy processes as rough paths,
\emph{Probab. Theory Related Fields} \textbf{170} (2018), 891--932.
\href{https://doi.org/10.1007/s00440-017-0781-1}
{doi:10.1007/s00440-017-0781-1}.

\bibitem{ChevyrevFerrucci26}
I.~Chevyrev and E.~Ferrucci,
Locality of rough path lifts,
arXiv:2606.21049 [math.PR], 2026.
\href{https://doi.org/10.48550/arXiv.2606.21049}
{doi:10.48550/arXiv.2606.21049}.

\bibitem{ChevyrevFriz19}
I.~Chevyrev and P.~K. Friz,
Canonical RDEs and general semimartingales as rough paths,
\emph{Ann. Probab.} \textbf{47} (2019), no.~1, 420--463.
\href{https://doi.org/10.1214/18-AOP1264}
{doi:10.1214/18-AOP1264}.

\bibitem{ConnesKreimer00}
A.~Connes and D.~Kreimer,
Renormalization in quantum field theory and the Riemann--Hilbert problem I:
The Hopf algebra structure of graphs and the main theorem,
\emph{Comm. Math. Phys.} \textbf{210} (2000), 249--273.
\href{https://doi.org/10.1007/s002200050779}
{doi:10.1007/s002200050779}.

\bibitem{CurryEtAl14}
C.~Curry, K.~Ebrahimi-Fard, S.~J.~A. Malham, and A.~Wiese,
L\'evy processes and quasi-shuffle algebras,
\emph{Stochastics} \textbf{86} (2014), no.~4, 632--642.
\href{https://doi.org/10.1080/17442508.2013.865131}
{doi:10.1080/17442508.2013.865131}.

\bibitem{DaPratoKwapienZabczyk87}
G.~Da Prato, S.~Kwapie\'n, and J.~Zabczyk,
Regularity of solutions of linear stochastic equations in Hilbert spaces,
\emph{Stochastics} \textbf{23} (1988), no.~1, 1--23.
\href{https://doi.org/10.1080/17442508708833480}
{doi:10.1080/17442508708833480}.

\bibitem{DaPratoZabczyk14}
G.~Da Prato and J.~Zabczyk,
\emph{Stochastic Equations in Infinite Dimensions}, second ed.,
Encyclopedia of Mathematics and its Applications 152,
Cambridge University Press, Cambridge, 2014.
\href{https://doi.org/10.1017/CBO9781107295513}
{doi:10.1017/CBO9781107295513}.

\bibitem{DelaCruzOberhauser26}
Danilo Jr Dela Cruz and H.~Oberhauser,
The singular values of L\'evy's area matrix,
arXiv:2606.08928 [math.PR], 2026.
\href{https://doi.org/10.48550/arXiv.2606.08928}
{doi:10.48550/arXiv.2606.08928}.

\bibitem{Dereich10}
S.~Dereich,
Rough paths analysis of general Banach space-valued Wiener processes,
\emph{J. Funct. Anal.} \textbf{258} (2010), no.~9, 2910--2936.
\href{https://doi.org/10.1016/j.jfa.2010.01.018}
{doi:10.1016/j.jfa.2010.01.018}.

\bibitem{FeiginShoikhet07}
B.~Feigin and B.~Shoikhet,
On \([A,A]/[A,[A,A]]\) and on a \(W_n\)-action on the consecutive
commutators of free associative algebras,
\emph{Math. Res. Lett.} \textbf{14} (2007), no.~5, 781--795.
\href{https://doi.org/10.4310/MRL.2007.v14.n5.a7}
{doi:10.4310/MRL.2007.v14.n5.a7}.

\bibitem{FrizHagerTapia24}
P.~K. Friz, P.~P. Hager, and N.~Tapia,
On expected signatures and signature cumulants in semimartingale models,
in \emph{Signature Methods in Finance}, Springer Finance,
Springer, Cham, 2026, 381--424.
\href{https://doi.org/10.1007/978-3-031-97239-3_11}
{doi:10.1007/978-3-031-97239-3\_11}.

\bibitem{FrizShekhar17}
P.~K. Friz and A.~Shekhar,
General rough integration, L\'evy rough paths and a L\'evy--Kintchine-type formula,
\emph{Ann. Probab.} \textbf{45} (2017), no.~4, 2707--2765.
\href{https://doi.org/10.1214/16-AOP1123}
{doi:10.1214/16-AOP1123}.

\bibitem{FrizVictoir10}
P.~K. Friz and N.~B. Victoir,
Differential equations driven by Gaussian signals,
\emph{Ann. Inst. H. Poincar\'e Probab. Statist.} \textbf{46} (2010),
no.~2, 369--413.
\href{https://doi.org/10.1214/09-AIHP202}
{doi:10.1214/09-AIHP202}.

\bibitem{Fulton97}
W.~Fulton,
\emph{Young Tableaux}, London Mathematical Society Student Texts 35,
Cambridge University Press, Cambridge, 1997.
\href{https://doi.org/10.1017/CBO9780511626241}
{doi:10.1017/CBO9780511626241}.

\bibitem{GarridoLuSchmalfuss15}
M.~J. Garrido-Atienza, K.~Lu, and B.~Schmalfu\ss,
L\'evy-areas of Ornstein--Uhlenbeck processes in Hilbert-spaces,
in \emph{Continuous and Distributed Systems II: Theory and Applications},
Studies in Systems, Decision and Control 30, Springer, Cham, 2015,
167--188.
\href{https://doi.org/10.1007/978-3-319-19075-4_10}
{doi:10.1007/978-3-319-19075-4\_10}.

\bibitem{GrongNilssenSchmeding22}
E.~Grong, T.~Nilssen, and A.~Schmeding,
Geometric rough paths on infinite-dimensional spaces,
\emph{J. Differential Equations} \textbf{340} (2022), 151--178.
\href{https://doi.org/10.1016/j.jde.2022.08.034}
{doi:10.1016/j.jde.2022.08.034}.

\bibitem{GubinelliTindel10}
M.~Gubinelli and S.~Tindel,
Rough evolution equations,
\emph{Ann. Probab.} \textbf{38} (2010), no.~1, 1--75.
\href{https://doi.org/10.1214/08-AOP437}
{doi:10.1214/08-AOP437}.

\bibitem{Hairer14}
M.~Hairer,
A theory of regularity structures,
\emph{Invent. Math.} \textbf{198} (2014), no.~2, 269--504.
\href{https://doi.org/10.1007/s00222-014-0505-4}
{doi:10.1007/s00222-014-0505-4}.

\bibitem{Hoffman00}
M.~E. Hoffman,
Quasi-shuffle products,
\emph{J. Algebraic Combin.} \textbf{11} (2000), 49--68.
\href{https://doi.org/10.1023/A:1008791603281}
{doi:10.1023/A:1008791603281}.

\bibitem{JacodShiryaev03}
J.~Jacod and A.~N. Shiryaev,
\emph{Limit Theorems for Stochastic Processes}, second ed.,
Grundlehren der mathematischen Wissenschaften 288, Springer,
Berlin--Heidelberg, 2003.
\href{https://doi.org/10.1007/978-3-662-05265-5}
{doi:10.1007/978-3-662-05265-5}.

\bibitem{JordanOrem13}
D.~Jordan and H.~Orem,
An algebro-geometric construction of lower central series of associative
algebras,
\emph{Int. Math. Res. Not. IMRN} (2015), no.~15, 6330--6352.
\href{https://doi.org/10.1093/imrn/rnu125}
{doi:10.1093/imrn/rnu125}.

\bibitem{Kalinichenko25}
A.~A. Kalinichenko,
Rough path lifts of Banach space-valued Gaussian processes,
\emph{Stochastic Process. Appl.} \textbf{190} (2025), 104739.
\href{https://doi.org/10.1016/j.spa.2025.104739}
{doi:10.1016/j.spa.2025.104739}.

\bibitem{KaltonPeckRoberts84}
N.~J. Kalton, N.~T. Peck, and J.~W. Roberts,
\emph{An F-Space Sampler}, London Mathematical Society Lecture Note
Series 89, Cambridge University Press, Cambridge, 1984.
\href{https://doi.org/10.1017/CBO9780511662447}
{doi:10.1017/CBO9780511662447}.

\bibitem{Kapranov98}
M.~Kapranov,
Noncommutative geometry based on commutator expansions,
\emph{J. Reine Angew. Math.} \textbf{505} (1998), 73--118.
\href{https://doi.org/10.1515/crll.1998.505.73}
{doi:10.1515/crll.1998.505.73}.

\bibitem{Karandikar95}
R.~L. Karandikar,
On pathwise stochastic integration,
\emph{Stochastic Process. Appl.} \textbf{57} (1995), no.~1, 11--18.
\href{https://doi.org/10.1016/0304-4149(95)00002-O}
{doi:10.1016/0304-4149(95)00002-O}.

\bibitem{KarolNazarovNikitin08}
A.~I. Karol', A.~I. Nazarov, and Ya.~Yu. Nikitin,
Small ball probabilities for Gaussian random fields and tensor products of
compact operators,
\emph{Trans. Amer. Math. Soc.} \textbf{360} (2008), no.~3, 1443--1474.
\href{https://doi.org/10.1090/S0002-9947-07-04233-X}
{doi:10.1090/S0002-9947-07-04233-X}.

\bibitem{LedouxLyonsQian02}
M.~Ledoux, T.~Lyons, and Z.~Qian,
L\'evy area of Wiener processes in Banach spaces,
\emph{Ann. Probab.} \textbf{30} (2002), no.~2, 546--578.
\href{https://doi.org/10.1214/aop/1023481002}
{doi:10.1214/aop/1023481002}.

\bibitem{LordSukochevZanin12}
S.~Lord, F.~Sukochev, and D.~Zanin,
\emph{Singular Traces: Theory and Applications}, De Gruyter Studies in
Mathematics 46, De Gruyter, Berlin--Boston, 2012.
\href{https://doi.org/10.1515/9783110262551}
{doi:10.1515/9783110262551}.

\bibitem{Lyons98}
T.~J. Lyons,
Differential equations driven by rough signals,
\emph{Rev. Mat. Iberoam.} \textbf{14} (1998), no.~2, 215--310.
\href{https://doi.org/10.4171/RMI/240}
{doi:10.4171/RMI/240}.

\bibitem{MetivierPellaumail80}
M.~M\'etivier and J.~Pellaumail,
\emph{Stochastic Integration}, Probability and Mathematical Statistics,
Academic Press, New York, 1980.

\bibitem{PapadimaSuciu04}
S.~Papadima and A.~I. Suciu,
Chen Lie algebras,
\emph{Int. Math. Res. Not.} (2004), no.~21, 1057--1086.
\par\noindent
\href{https://doi.org/10.1155/S1073792804132017}
{doi:10.1155/S1073792804132017}.

\bibitem{Pietsch80}
A.~Pietsch,
\emph{Operator Ideals}, North-Holland Mathematical Library 20,
North-Holland, Amsterdam, 1980.

\bibitem{Quillen68}
D.~G. Quillen,
On the associated graded ring of a group ring,
\emph{J. Algebra} \textbf{10} (1968), 411--418.
\href{https://doi.org/10.1016/0021-8693(68)90069-0}
{doi:10.1016/0021-8693(68)90069-0}.

\bibitem{Rastegaev18}
N.~V. Rastegaev,
On spectral asymptotics of the tensor product of operators with almost
regular asymptotics,
\emph{St. Petersburg Math. J.} \textbf{29} (2018), no.~6, 1007--1029.
\href{https://doi.org/10.1090/spmj/1525}
{doi:10.1090/spmj/1525}.

\bibitem{Reutenauer93}
C.~Reutenauer,
\emph{Free Lie Algebras}, London Mathematical Society Monographs, New
Series 7, Clarendon Press, Oxford, 1993.
\href{https://doi.org/10.1093/oso/9780198536796.001.0001}
{doi:10.1093/oso/9780198536796.001.0001}.

\bibitem{Simon05}
B.~Simon,
\emph{Trace Ideals and Their Applications}, second ed.,
Mathematical Surveys and Monographs 120, American Mathematical Society,
Providence, RI, 2005.
\href{https://doi.org/10.1090/surv/120}
{doi:10.1090/surv/120}.

\end{thebibliography}

\end{document}
